%%%%%%%%%%%%%%%%%%%%%%%%%%%%%%%%%%%%%%%%%%%%%%%%%%%%%%%%%%%%%%%%%%%%%%%%%%%%%%%%
%% CHAPTER 1: INTRODUCTION
%%%%%%%%%%%%%%%%%%%%%%%%%%%%%%%%%%%%%%%%%%%%%%%%%%%%%%%%%%%%%%%%%%%%%%%%%%%%%%%%
% ╔══════════════════════════════════════════════════════════════════════════════╗
% ║  GGU DBA THESIS — CHAPTER 1 REQUIREMENTS CHECKLIST (Professor Review)      ║
% ╠══════════════════════════════════════════════════════════════════════════════╣
% ║  Per GGU DBA ET Thesis Guidelines, Chapter 1 MUST contain:                 ║
% ║                                                                            ║
% ║  [✓] Background of the study                                               ║
% ║  [✓] Purpose of the study                                                  ║
% ║  [✓] Statement of the problem                                              ║
% ║  [✓] Aim and objectives of the study                                       ║
% ║  [✓] Research questions                                                    ║
% ║  [✓] Hypothesis of the study                                               ║
% ║  [✓] Context of contributing organizations  → sec:contributing_orgs        ║
% ║  [✓] Significant contributions of the study                                ║
% ║  [✓] Scope and assumptions                                                 ║
% ║  [✓] Limitations of the study             → sec:scope_limitations          ║
% ║  [✓] Thesis outline / Organisation of the thesis                           ║
% ║                                                                            ║
% ║  STATUS: 11/11 COMPLIANT                                                   ║
% ╚══════════════════════════════════════════════════════════════════════════════╝

\chapter{Introduction}

\subsection*{Chapter 1 --- Three Questions This Chapter Answers}
\addcontentsline{toc}{subsection}{Chapter 1 --- Three Questions This Chapter Answers}
\label{sec:ch1_three_questions}

\noindent An examiner reading this chapter should expect three specific questions to be answered. The table below pairs each question with its one-line answer and a pointer to the section that develops the answer in full. Detailed evidence and synthesis follow throughout the chapter and are consolidated in the Distinction Pack at Appendix~\ref{app:distinction_pack}.

{\tablestripes\tablefontsize
\needspace{10\baselineskip}
\begin{xltabular}{\textwidth}{@{} >{\raggedright\arraybackslash}p{6.5cm} L @{}}
\caption[Chapter 1 three-question landing page]{Chapter 1 three-question landing page.}\label{tab:ch1_three_questions} \\
\toprule
\thead{Examiner question} & \thead{One-line answer + section pointer} \\
\midrule
\endfirsthead
\endhead
\bottomrule
\endlastfoot
1.~Why study this? & Healthcare AI for B2C-primary, clinician-supported ASD screening lacks an operational governance model that integrates explainability, trust, and adoption readiness.\\
 & \textit{See:}~\autoref{sec:ch1_template_compliance} \\
\addlinespace
2.~Why is this a DBA and not a PhD? & The contribution is governance-centered organisational adoption, not algorithmic novelty; technical accuracy is already established and is out of scope.\\
 & \textit{See:}~\autoref{sec:ch1_why_dba_dedicated} \\
\addlinespace
3.~What is the scope and what is excluded? & In scope: governance maturity, explainability, trust, adoption readiness, RGAIG validation. Out of scope: new EEG algorithms, new neural architectures, real-time clinical diagnosis.\\
 & \textit{See:}~\autoref{fig:ch1_scope_inout} \\
\end{xltabular}}

\vspace{0.3em}\noindent\textit{Note.} If a reader skims only this chapter's first page, this table is sufficient to know what the chapter claims. The chapter body then evidences each claim.

\clearpage


%% ============================================================
%% GGU DBA Examiner Compliance Page --- Chapter 1
%% Source of truth: docs/GGU_DBA_EXAMINER_CHECKLIST.md
%% Alignment audit: docs/GGU_ALIGNMENT_AUDIT_2026_05_28.md (12/12)
%% ============================================================
\clearpage
\thispagestyle{plain}
\addcontentsline{toc}{section}{Chapter 1 --- GGU DBA Compliance Map}
\vspace*{0.25in}
\begin{center}
{\Large\bfseries Chapter 1 --- GGU DBA Examiner Compliance}\\[0.3em]
{\small Per GGU DBA Dissertation Thesis Guidelines + DBA Examiner Expectations}\\[0.5em]
{\large\itshape Mantra: \textcolor{ggunavy}{\textbf{WHY}}}~~|~~Examiner looks for: \textcolor{ggunavy}{\textbf{Clear problem}}
\end{center}

\vspace{0.3em}
\noindent\textcolor{ggunavy}{\textbf{Chapter 1 sequence flow (Part A: GGU Must-Have topics in order):}}
\begin{center}
\begin{tikzpicture}[
  font=\fignodefont,
  node distance=0.35cm and 0.15cm,
  cstep/.style={rectangle, draw=ggunavy, fill=ggunavy!8, rounded corners=2pt,
               minimum width=1.6cm, text width=1.5cm, align=center, inner sep=2pt, font=\scriptsize},
  arr/.style={->, >=stealth, thick, ggunavy}
]
  \node[cstep] (c1) {1.~Back-\\ground};
  \node[cstep, right=of c1] (c2) {2.~Business Problem};
  \node[cstep, right=of c2] (c3) {3.~Research Problem};
  \node[cstep, right=of c3] (c4) {4.~Research Gap};
  \node[cstep, right=of c4] (c5) {5.~Aim};
  \node[cstep, right=of c5] (c6) {6.~Objectives};
  \node[cstep, below=of c1] (c7) {7.~Research Qs};
  \node[cstep, right=of c7] (c8) {8.~Signif-\\icance};
  \node[cstep, right=of c8] (c9) {9.~Scope};
  \node[cstep, right=of c9] (c10) {10.~Assump-\\tions};
  \node[cstep, right=of c10] (c11) {11.~Limit-\\ations};
  \node[cstep, right=of c11] (c12) {12.~Thesis Structure};
  \draw[arr] (c1) -- (c2); \draw[arr] (c2) -- (c3); \draw[arr] (c3) -- (c4);
  \draw[arr] (c4) -- (c5); \draw[arr] (c5) -- (c6); \draw[arr] (c6.south) |- (c7.north);
  \draw[arr] (c7) -- (c8); \draw[arr] (c8) -- (c9); \draw[arr] (c9) -- (c10);
  \draw[arr] (c10) -- (c11); \draw[arr] (c11) -- (c12);
\end{tikzpicture}
\end{center}

\vspace{0.2em}
\noindent This page lists what Chapter 1 \emph{must contain} (Part A, per GGU template) and what it \emph{must not contain} (deferred to other chapters). All 12 Must-Have items verified present (see \texttt{docs/GGU\_ALIGNMENT\_AUDIT\_2026\_05\_28.md}).

\vspace{0.5em}
{\tablestripes\tablefontsize
\begin{tabularx}{\textwidth}{@{} >{\raggedright\arraybackslash}X >{\raggedright\arraybackslash}X @{}}
\toprule
\thead{Must Have (per GGU + DBA examiner)} & \thead{Must NOT Have (out of scope here)} \\
\midrule
1.~\textcolor{ggunavy}{\textbf{Background}} --- ASD / neuropsychiatric screening challenges $\checkmark$ & 1.~Detailed literature review \emph{(in Ch.~2)} \\
2.~\textcolor{ggunavy}{\textbf{Business problem}} --- clinical decision challenges $\checkmark$ & 2.~SHAP discussion \emph{(in Ch.~4)} \\
3.~\textcolor{ggunavy}{\textbf{Research problem}} --- governance gap in AI healthcare $\checkmark$ & 3.~LIME discussion \emph{(in Ch.~4)} \\
4.~\textcolor{ggunavy}{\textbf{Research gap}} --- missing Responsible AI governance $\checkmark$ & 4.~Monte Carlo / sensitivity analysis \emph{(in Ch.~4)} \\
5.~\textcolor{ggunavy}{\textbf{Aim}} --- clear overall aim $\checkmark$ & 5.~Dataset details \emph{(in Ch.~3)} \\
6.~\textcolor{ggunavy}{\textbf{Objectives}} --- 4--6 measurable objectives $\checkmark$ & 6.~Survey results \emph{(in Ch.~4)} \\
7.~\textcolor{ggunavy}{\textbf{Research questions}} aligned to objectives $\checkmark$ & 7.~Interview findings \emph{(in Ch.~4)} \\
8.~\textcolor{ggunavy}{\textbf{Significance}} --- academic + industry contribution $\checkmark$ & 8.~Framework diagrams (RGAIG / SMART / ResAI) \emph{(in Ch.~5)} \\
9.~\textcolor{ggunavy}{\textbf{Scope}} --- boundaries clearly defined $\checkmark$ & 9.~Final recommendations \emph{(in Ch.~5)} \\
10.~\textcolor{ggunavy}{\textbf{Assumptions}} --- explicit assumptions $\checkmark$ & 10.~Statistical results \emph{(in Ch.~4)} \\
11.~\textcolor{ggunavy}{\textbf{Limitations}} --- initial limitations $\checkmark$ & ~ \\
12.~\textcolor{ggunavy}{\textbf{Thesis structure}} --- chapter summary $\checkmark$ & ~ \\
\bottomrule
\end{tabularx}}
\vspace{0.4em}\noindent\textit{Note.} $\checkmark$ = requirement met by a section in this chapter. Source: \texttt{docs/GGU\_DBA\_EXAMINER\_CHECKLIST.md} (Ch.1 row). Scope discipline: B2C-primary, clinician-supported, single-disease ASD (no multi-disease residue).
\clearpage

\needspace{12\baselineskip}
\noindent The table below presents the chapter 1 Roadmap: Sections and Purpose.

\begin{table}[H]
\centering
\caption{Chapter 1 Roadmap: Sections and Purpose.}
\tablestripes\tablefontsize
\begin{tabularx}{\textwidth}{@{} >{\raggedright\arraybackslash}p{0.8cm} >{\raggedright\arraybackslash}p{3.5cm} L @{}}
\toprule
\thead{Sec.} & \thead{Section Title} & \thead{Purpose} \\
\midrule
1.1 & Background of the Study & Context: ASD, EEG, AI, diagnostic delays \\
1.2 & Problem Statement & Four named problems (P1--P4) \\
1.3 & Research Objectives & O1--O5 main objectives \\
1.4 & Operationalisation (A--D) & A1--D3 sub-objectives mapped to Parts \\
1.5 & Research Questions & RQ1--RQ5 linked to objectives \\
1.6 & Research Structure (A--D) & Four-part evaluation framework \\
1.7 & Significance of the Study & Academic, clinical, policy impact \\
1.8 & Scope of the Study & Boundaries and assumptions \\
1.9 & Structure of Dissertation & Chapter-by-chapter overview \\
1.10 & Output Card & Deliverables to Chapter~2 \\
1.11 & Limitations & Known constraints \\
\bottomrule
\end{tabularx}
\end{table}
\needspace{12\baselineskip}

\noindent\textbf{Chapter 1: Input--Process--Output.} Table~\ref{tab:ch1_ipo_handoff} specifies what enters, what this chapter does, and what it delivers to Chapter~2.

\paragraph{Chapter 1.}



{\tablestripes\tablefontsize
\needspace{10\baselineskip}
\begin{xltabular}{\textwidth}{@{} >{\raggedright\arraybackslash}p{2.0cm} L @{}}
\caption[Chapter 1: Input--Process--Output Handoff]{Chapter 1: Input--Process--Output Handoff. This table specifies what enters the chapter, what processing occurs, what outputs are produced for Chapter~2, and what is explicitly excluded.}\label{tab:ch1_ipo_handoff} \\
\toprule
\thead{Dimension} & \thead{Specification} \\
\midrule
\endfirsthead
\endhead
\bottomrule
\endlastfoot
\textbf{Input} & \textbullet\ ASD prevalence data (78M+ globally, WHO) \newline \textbullet\ Clinical AI fragmentation evidence \newline \textbullet\ GGU DBA thesis guidelines \newline \textbullet\ Preliminary literature scan \\
\addlinespace
\textbf{Process} & \textbullet\ Identify 4 business problems (P1--P4) \newline \textbullet\ Preview 6 research gaps (G1--G6) \newline \textbullet\ Formulate 5 objectives (O1--O5) + 14 sub-objectives (A1--D3) \newline \textbullet\ Derive 5 research questions (RQ1--RQ5) \newline \textbullet\ State 5 hypotheses (H1--H5) with thresholds \newline \textbullet\ Define scope: ASD-only, public data, no clinical trial \newline \textbullet\ Build master traceability matrix (P$\to$G$\to$O$\to$RQ$\to$H$\to$C) \\
\addlinespace
\textbf{Output} & \textbullet\ P1--P4: 4 named problems with stakeholder impact \newline \textbullet\ O1--O5 + A1--D3: 5 main + 14 sub-objectives \newline \textbullet\ RQ1--RQ5: 5 research questions \newline \textbullet\ H1--H5: 5 hypotheses with pass/fail thresholds \newline \textbullet\ RGAIG concept: 5-layer governance framework \newline \textbullet\ Master traceability matrix (Table~\ref{tab:research_flow_alignment}) \newline \textbullet\ Scope boundaries and limitations \\
\addlinespace
\textbf{Boundary} & \textbullet\ Does NOT review literature systematically (Ch.~2) \newline \textbullet\ Does NOT design methodology (Ch.~3) \newline \textbullet\ Does NOT present results (Ch.~4) \newline \textbullet\ Does NOT make evaluation decisions (Ch.~5) \\
\addlinespace
\textbf{Delivers to} & Chapter~2: problems, objectives, questions, and RGAIG concept \\
\end{xltabular}}
\vspace{0.3em}\noindent\textit{Note.} RGAIG = Responsible Generative AI Governance; ASD = Autism Spectrum Disorder. See chapter prose for full context.

\label{chap:introduction}

\normalsize\normalfont\color{black}% Reset font after \Large chapter title
\setchaptercolor{1}% Set Chapter 1 color scheme (Steel Blue)

%%%%%%%%%%%%%%%%%%%%%%%%%%%%%%%%%%%%%%%%%%%%%%%%%%%%%%%%%%%%%%%%%%%%%%%%%%%%%%%%
%% DBA-GRADE RESEARCH FRAMING SUMMARY
%% Authoritative one-page reference for examiners — single bounded problem,
%% measurable variables, falsifiable hypotheses, explicit scope.
%%%%%%%%%%%%%%%%%%%%%%%%%%%%%%%%%%%%%%%%%%%%%%%%%%%%%%%%%%%%%%%%%%%%%%%%%%%%%%%%

\section*{Doctoral Research Framing Summary}
\addcontentsline{toc}{section}{Doctoral Research Framing Summary}
\label{sec:dba_framing}

This dissertation investigates \textbf{one bounded organizational research problem}: clinician trust in AI-assisted ASD-EEG screening is too low to support adoption, despite the technical maturity of the underlying machine-learning models. The framing below specifies the central question, the measurement model, the hypotheses, and the explicit boundary of the study. All subsequent sections of this chapter elaborate one element of this framing; Chapters~2--5 test it.

\noindent\textbf{Topic-Proposal Merge Note.}
The topic proposal's introductory claim is merged into this chapter as the dissertation's opening logic: AI-assisted ASD-EEG screening is not evaluated only as a classification problem, but as a B2C-primary, B2B-supported responsible-AI adoption problem. Chapter~1 therefore anchors the family/caregiver pain area, clinician-governance boundary, research problem, objectives, questions, and management relevance; Appendix~\ref{chap:appendix_operationalization_detail} documents the detailed proposal-to-dissertation alignment control.

\noindent\textbf{Central Research Question.}
\begin{quote}
\emph{How does Responsible AI governance influence clinician trust and adoption intention for AI-assisted ASD EEG screening, through perceived explainability?}
\end{quote}

\needspace{12\baselineskip}
\noindent Table~\ref{tab:dba_framing} states the DBA framing for the dissertation --- independent variable, mediators, dependent variable, hypotheses, and study scope --- so subsequent chapters can be read against a single, named research-design contract.

\paragraph{DBA Research Framing.}




{\tablestripes\tablefontsize
\needspace{10\baselineskip}
\begin{xltabular}{\textwidth}{@{} >{\raggedright\arraybackslash}p{3.0cm} L @{}}
\caption[DBA Research Framing: Variables, Hypotheses, Scope]{DBA Research Framing: Variables, Hypotheses, Scope.}\label{tab:dba_framing} \\
\toprule
\thead{Element} & \thead{Specification} \\
\midrule
\endfirsthead
\endhead
\bottomrule
\endlastfoot
\textbf{Independent variable (IV)} & \emph{Responsible AI Governance} --- measured as governance maturity using a 12-item instrument covering seven control families (auditability, bias control, human-in-the-loop, rollback, monitoring, lineage, access control). Operationalized via the RGAIG Maturity Framework (Table~\ref{tab:rgaig_maturity_artifact}); 5-point maturity scale. \\
\textbf{Mediator 1} & \emph{Perceived Explainability} --- clinician-rated quality of the AI system's explanations, encompassing both algorithmic attribution (SHAP) and contextual grounding (RAG citations, auditability cues). 5-point Likert, 6 items. \\
\textbf{Mediator 2} & \emph{Clinician Trust} --- cognitive + affective trust in the AI system. Madsen \& Gregor TUI scale (10 items). \\
\textbf{Dependent variable (DV)} & \emph{Adoption Intention / Workflow Readiness} --- behavioural intention to adopt (TAM-BI, 3 items) + workflow-integration readiness (perceived fit, perceived safety, perceived usefulness for ASD screening; composite, 7 items). \\
\textbf{Moderator} & \emph{AI Familiarity} --- clinician self-rated familiarity with AI and ML in healthcare (5-point). \\
\textbf{Control variables} & Years of practice; healthcare role (neurologist / psychiatrist / paediatrician / radiologist / admin / other); ASD exposure (ASD relevance to current practice). \\
\midrule
\textbf{Theoretical anchors} & Technology Acceptance Model (TAM; Davis 1989), Trust in Automation (Lee \& See 2004; Madsen \& Gregor 2000), Socio-Technical Systems (Cherns 1976), Diffusion of Innovations (Rogers 2003), TOE Framework (Tornatzky \& Fleischer 1990), Responsible AI principles (Floridi et al.\ 2018; EU AI Act 2024). \emph{Why TAM alone is insufficient:} TAM's perceived usefulness construct does not capture governance-induced trust calibration in safety-critical AI. This dissertation extends TAM by positioning RAI governance maturity as a structural antecedent of trust, drawing on Trust-in-Automation theory to specify the calibration mechanism. \\
\midrule
\textbf{H1} & Responsible AI governance positively affects perceived explainability \quad (RAI $\to$ Expl). \\
\textbf{H2} & Perceived explainability positively affects clinician trust \quad (Expl $\to$ Trust). \\
\textbf{H3} & Clinician trust positively affects adoption intention and workflow-integration readiness \quad (Trust $\to$ Adoption). \\
\textbf{H4} & Perceived explainability \emph{mediates} the relationship between RAI governance and clinician trust \quad (RAI $\to$ Expl $\to$ Trust). \\
\textbf{H5} & Clinician trust \emph{mediates} the relationship between perceived explainability and adoption intention \quad (Expl $\to$ Trust $\to$ Adoption). \\
\textbf{H6} & AI familiarity \emph{moderates} the relationship between perceived explainability and clinician trust \quad (Expl $\times$ AIfam $\to$ Trust). \\
\midrule
\textbf{Methods} & Mixed methods (pragmatist philosophy). Quantitative: PLS-SEM on $n=131$ expert-panel + clinician survey responses; Cronbach's $\alpha$, CR, AVE, HTMT. Qualitative: thematic coding (Braun \& Clarke 2006), three-stage coding (Salda\~na 2021), Krippendorff's $\alpha$. Triangulation via joint-display matrix. \\
\midrule
\textbf{In scope} & ASD-focused EEG screening; clinician trust + organizational adoption; RGAIG control families (7); explainability via SHAP + RAG narrative; cross-sectional measurement. \\
\textbf{Out of scope} & Full clinical diagnosis; replacement of clinicians; AGI / autonomous treatment; multi-disease generalization; longitudinal RCT; production regulatory clearance; algorithmic mechanism-level XAI research. \\
\midrule
\textbf{Contributions} & \textbullet\ \emph{Theoretical}: extends TAM and trust-in-automation models by positioning Responsible AI governance as a structural antecedent of explainability, trust, and adoption in clinical AI. \quad \textbullet\ \emph{Empirical}: tests the Governance $\to$ Explainability $\to$ Trust $\to$ Adoption pathway in the ASD-EEG clinical AI context (H1--H6). \quad \textbullet\ \emph{Practical}: provides the RGAIG Maturity Framework (Table~\ref{tab:rgaig_maturity_artifact}, 6 levels) as a healthcare-conceptual AI governance evaluation-readiness instrument. \quad \textbullet\ \emph{Methodological}: develops and validates a survey instrument for Responsible AI governance and clinician adoption readiness, with reported psychometric properties (CR, AVE, HTMT, $\alpha$). \\
\end{xltabular}}
\vspace{0.5em}
\noindent\textit{Note.} This table is the single authoritative reference for the DBA research framing. Sections~\ref{sec:problem_statement}, \ref{sec:research_goal} (Research Aim), \ref{sec:hypotheses} (RQs), and Section~\ref{sec:significance} (Significance) elaborate individual elements. Hypotheses H1--H5 in earlier drafts of this document map to technical sub-findings that supply evidence for the DBA-grade hypotheses above; the relationship is documented in the master traceability matrix (Table~\ref{tab:research_flow_alignment}).

\needspace{12\baselineskip}
\begin{figure}[!htbp]
\centering
\resizebox{0.98\textwidth}{!}{%
\begin{tikzpicture}[every node/.style={font=\fignodefont},
    box/.style={draw=ggunavy, fill=ggunavy!8, rounded corners=5pt, very thick, minimum width=3.4cm, minimum height=1.6cm, align=center, font=\fignodefont},
    mod/.style={draw=ggunavy!60, fill=orange!10, rounded corners=4pt, thick, minimum width=2.8cm, minimum height=0.95cm, align=center, font=\small\itshape, dashed},
    arr/.style={-{Stealth[length=4mm, width=3mm]}, very thick, ggunavy},
    darr/.style={-{Stealth[length=3mm, width=2.2mm]}, thick, ggunavy, dashed}]

% Four-box main causal chain
\node[box, fill=ggunavy!18] (iv)   at (0,0)    {\textbf{Responsible AI}\\\textbf{Governance}\\\scriptsize Independent Variable};
\node[box]                  (expl) at (5,0)    {\textbf{Perceived}\\\textbf{Explainability}\\\scriptsize Mediator 1};
\node[box, fill=ggunavy!14] (tr)   at (10,0)   {\textbf{Clinician}\\\textbf{Trust}\\\scriptsize Mediator 2};
\node[box, fill=ggunavy!28] (ad)   at (15,0)   {\textbf{Adoption Intention /}\\\textbf{Workflow Readiness}\\\scriptsize Dependent Variable};

% Moderator below the Expl-Trust path
\node[mod] (mfm) at (7.5,-2.6) {AI Familiarity\\\scriptsize Moderator};

% Direct paths
\draw[arr] (iv)   -- node[above, font=\small\bfseries] {H1} (expl);
\draw[arr] (expl) -- node[above, font=\small\bfseries] {H2} (tr);
\draw[arr] (tr)   -- node[above, font=\small\bfseries] {H3} (ad);

% Moderation arrow (H6)
\draw[darr] (mfm.north) -- node[right, font=\scriptsize\bfseries, pos=0.5] {H6} ($(expl.south east)!0.5!(tr.south west) + (0,-0.4)$);

% Mediation labels: stacked vertically to avoid overlap; darkened for visibility (2026-05-28)
\node[font=\scriptsize\itshape\bfseries, text=ggunavy, anchor=west] at (0.5, 2.2) {H4: Explainability mediates Governance $\to$ Trust};
\node[font=\scriptsize\itshape\bfseries, text=ggunavy, anchor=west] at (0.5, 2.7) {H5: Trust mediates Explainability $\to$ Adoption};
\draw[ggunavy, densely dashed, very thick] (iv.north) .. controls (3,1.7) and (7,1.7) .. (tr.north);
\draw[ggunavy, densely dashed, very thick] (expl.north) .. controls (8,1.7) and (12,1.7) .. (ad.north);
\end{tikzpicture}%
}
\caption[DBA conceptual framework]{DBA Conceptual Framework. Four-box sequential mediation model. Responsible AI Governance (IV) influences Adoption Intention / Workflow Readiness (DV) through two sequential mediators: Perceived Explainability and Clinician Trust. H1--H3 are direct paths; H4 specifies that Explainability mediates the Governance $\to$ Trust relationship; H5 specifies that Trust mediates the Explainability $\to$ Adoption relationship; H6 specifies that AI Familiarity moderates the Explainability $\to$ Trust path. This is the locked final model.}
\label{fig:dba_conceptual_framework}
\end{figure}

\noindent The remainder of this chapter elaborates each row of Table~\ref{tab:dba_framing}: Section~\ref{sec:background} sets the organizational context, Section~\ref{sec:problem_statement} states the bounded problem, Section~\ref{sec:research_goal} states the aim, Sections~\ref{sec:hypotheses}--\ref{sec:significance} state research questions and significance, and Section~\ref{sec:scope_limitations} states scope and limitations. Engineering and platform details that support but do not constitute the DBA contribution are demoted to Appendix~\ref{chap:appendix_ch1}.

\subsection*{Why ASD-EEG Screening as the Empirical Domain}
\label{sec:why_asd_justification}

\noindent A common reviewer challenge is ``why ASD specifically and not radiology or cardiology?'' Six independent reasons make ASD-EEG screening the most appropriate empirical context for testing the RAI-governance $\to$ trust $\to$ adoption pathway. None is sufficient alone; together they make ASD uniquely suited to this research question.

\begin{enumerate}[leftmargin=*, itemsep=3pt]
    \item \textbf{Early-intervention sensitivity.} ASD outcomes are highly dependent on early diagnosis: every additional month of delay measurably reduces intervention effectiveness \cite{zwaigenbaum2015earlyasd}. This makes faster screening operationally valuable in ways that, for example, asymptomatic cardiac screening is not.
    \item \textbf{Diagnostic subjectivity.} Unlike radiology (where ground truth is the imaged anatomy) or cardiology (where the ECG signal is largely deterministic), ASD diagnosis combines behavioural observation, clinical interview, and developmental history with non-trivial inter-rater variability. This subjectivity amplifies the clinician's trust threshold for any AI assistance, making ASD a ``stress test'' for governance-induced trust.
    \item \textbf{Trust-sensitivity.} Misdiagnosis (false-positive or false-negative) in ASD has lifelong consequences for the patient and family. Clinicians therefore demand more justification per decision than in domains with reversible interventions. This makes ASD an externally-valid case for studying \emph{calibrated} trust rather than over-trust.
    \item \textbf{Explainability requirement.} The EU AI Act categorizes paediatric neurological screening as a ``high-risk'' AI use case, requiring explainability and human oversight by statute. ASD-EEG screening is therefore not only methodologically suited but regulatorily required to operationalize the explainability mediator.
    \item \textbf{Clinician-AI collaboration imperative.} ASD screening is a multi-stage workflow (referral $\to$ EEG $\to$ behavioural assessment $\to$ diagnosis) in which AI augments rather than replaces clinical judgement. This matches the dissertation's framing that the dependent variable is adoption \emph{into a workflow}, not autonomous decision-making.
    \item \textbf{High ethical sensitivity + paediatric context.} Paediatric conceptual AI governance evaluation surfaces governance constraints (parental consent, data minimization, equity across socioeconomic groups) more acutely than adult-only domains. The seven RAI control families in the IV are therefore each materially relevant in ASD-EEG, providing maximum statistical variance on the IV.
\end{enumerate}

\noindent\emph{Generalization claim:} The dissertation does not claim that findings apply to all clinical AI domains. The framework is tested in ASD-EEG; the conditions under which the trust--adoption pathway generalises (high subjectivity, regulatory-high-risk classification, paediatric or other vulnerable population, workflow-integrated AI) are explicit boundary conditions stated in Section~\ref{sec:scope_limitations}.

\subsection*{Primary Practical Artifact: RGAIG Maturity Framework}
\label{sec:rgaig_maturity_artifact}

\noindent The dissertation's central \emph{managerial} contribution is a maturity model that operationalizes Responsible AI governance for clinical conceptual AI governance evaluation. Unlike a survey-only study, this artifact converts the quantitative findings into a decision-ready instrument that healthcare organizations can apply directly to assess their position on the trust--adoption pathway and to plan governance investments. The six-level model and its diagnostic mapping are shown below; Chapter~5 presents the full instrument and pilot validation.

\noindent Table~\ref{tab:rgaig_maturity_artifact} rgaig maturity framework.

{\tablestripes\tablefontsize

\paragraph{Rgaig Maturity Framework.}
\needspace{10\baselineskip}

\begin{xltabular}{\textwidth}{@{} >{\centering\arraybackslash}p{1.0cm} >{\raggedright\arraybackslash}p{2.0cm} p{6.6cm} >{\raggedright\arraybackslash}p{4.7cm} @{}}
\caption[RGAIG Maturity Framework]{RGAIG Maturity Framework --- six levels of Responsible AI governance maturity for clinical conceptual AI governance evaluation. Each level is defined by observable organizational practices (left) and predicted positions on the trust--adoption pathway (right). The framework is the managerial-contribution artifact of this dissertation.}\label{tab:rgaig_maturity_artifact} \\
\toprule
\thead{Lvl} & \thead{Stage} & \thead{Observable Practices} & \thead{Expected Trust + Adoption Position} \\
\midrule
\endfirsthead
\endhead
\bottomrule
\endlastfoot
1 & Initial          & Ad-hoc AI pilots; no documented governance; no audit trail; no explainability layer.                              & Low trust; adoption stalls at pilot. \\
2 & Managed          & Basic policies; awareness of EU AI Act / FDA SaMD reference frameworks (conceptual alignment only); access controls and logging in place but not standardised.    & Trust limited to early adopters; sporadic clinical use. \\
3 & Governed         & Formal RAI policies; documented control inventory (audit, bias, HITL, rollback, monitoring, lineage, access); ethics review board.  & Trust calibrated by policy compliance; adoption in defined pilots. \\
4 & Explainable      & SHAP attributions + RAG citation grounding embedded in clinical UI; explanations validated by domain experts.    & Trust elevated by transparent reasoning; broader adoption begins. \\
5 & Trusted          & Peer-reviewed performance evidence; clinician trust scores audited; calibrated trust (over/under-trust monitored); fairness reporting per subgroup. & Trust verified empirically; clinical workflow integration starts. \\
6 & Operationalized  & Full workflow integration with EHR; outcome monitoring; continuous-learning governance; quarterly maturity audit; cross-site replication. & Sustained adoption; AI governance treated as a clinical-operations function. \\
\end{xltabular}}
\vspace{0.4em}
\noindent\textit{Note.} Levels 1--3 are foundational (organizational compliance); levels 4--5 are decision-quality (clinician interaction); level 6 is operational excellence (workflow + outcome management). The maturity assessment instrument (12 items, validated in Chapter~5) maps an organization to one of these six levels and identifies the specific governance investments required to advance one level.
% Aggressive overflow prevention for Ch.1 tables
\renewcommand{\tabfixsetup}{\sloppy\sloppy\hyphenpenalty=0\exhyphenpenalty=0\tolerance=9999\emergencystretch=8em\hfuzz=100pt\hbadness=9999}

% Stub labels for suppressed content references (phantom — resolve to chapter start)
\label{sec:primary_data_instruments}
\label{sec:secondary_data_instruments}
\label{sec:conceptual_model}
\label{fig:conceptual_model}
\label{tab:research_blueprint}
\label{tab:literature_comparison}

\vspace{0.3em}

\noindent\begin{quote}\itshape\small
\textbf{\color{currentaccent}Chapter Overview.}
This chapter establishes the dissertation's foundation: ASD-focused clinical AI remains fragmented across workflows and governance structures. The chapter:
\begin{itemize}[leftmargin=*, itemsep=2pt]
    \item \textbf{Defines} the core problem (4 problems and 6 research gaps from 156 ASD studies).
    \item \textbf{Translates} those into objectives, research questions, and hypotheses.
    \item \textbf{Establishes} the conceptual basis for RGAIG as an explainable, governance-aligned response to fragmentation.
    \item \textbf{Clarifies} the scope, contribution, and boundaries so later-chapter claims remain proportionate to the evidence.
\end{itemize}
\end{quote}
\vspace{0.5em}


\noindent Figure~\ref{fig:ch1_roadmap} provides a six-section roadmap of this chapter, tracing the logical progression from healthcare context through problem identification, objective formulation, and theoretical grounding to the research design.



\noindent\textbf{Chapter 1 section roadmap.}

\needspace{12\baselineskip}
\begin{figure}[!htbp]
\centering
\resizebox{0.97\textwidth}{!}{%
\begin{tikzpicture}[
  every node/.style={font=\fignodefont},
  node distance=0.8cm and 0.8cm,
  rstep/.style={rectangle, rounded corners=3pt, draw=currentaccent, fill=currentaccent!10,
    text width=4.8cm, minimum height=1.6cm, align=left, font=\small, inner sep=5pt},
  arr/.style={-{Stealth[length=4mm, width=3mm]}, very thick, currentaccent}
]
% Row 1: left to right (Problem-driven academic flow)
% 
\node[rstep, fill=lightblue] (s1) {\textbf{\color{currentaccent}1. Healthcare Context}\\[2pt]Global diagnostic burden,\\970M+ affected, the ASD domain,\\clinical AI fragmentation};
% 
\node[rstep, fill=lightorange, right=of s1] (s2) {\textbf{\color{currentaccent}2. Problem \& Gaps}\\[2pt]Four problems (P1--P4),\\six gaps (G1--G6), 156 studies\\+ 6 companion papers};
% 
\node[rstep, fill=lightteal, right=of s2] (s3) {\textbf{\color{currentaccent}3. Objectives \& Hypotheses}\\[2pt]Five objectives (O1--O5),\\five hypotheses (H1--H5),\\predefined thresholds};
% 
% Row 2: left to right
% 
\node[rstep, fill=lightpurple, below=0.6cm of s1] (s4) {\textbf{\color{currentaccent}4. Scope \& Use Context}\\[2pt]B2B core focus,\\B2C future pathway,\\study boundaries (Ch.~1/5)};
% 
\node[rstep, fill=lightgreen, right=of s4] (s5) {\textbf{\color{currentaccent}5. Theoretical Framework}\\[2pt]Nine theories, DL ensemble,\\Ultra Stacking, RAG, agentic AI,\\46-module RAI governance};
% 
\node[rstep, fill=lightyellow, right=of s5] (s6) {\textbf{\color{currentaccent}6. Methods \& Structure}\\[2pt]DSR methodology,\\contributions (C1--C8),\\five-chapter thesis outline};
% 
% Arrows
\draw[arr] (s1) -- (s2);
\draw[arr] (s2) -- (s3);
\draw[arr] (s3.south) -- ++(0,-0.15) -| (s4.north);
\draw[arr] (s4) -- (s5);
\draw[arr] (s5) -- (s6);
\end{tikzpicture}%
}
\caption[Chapter~1 section roadmap]{Chapter~1 section roadmap. The six sections progress from problem identification through gap analysis, research design, and scope delimitation, establishing the logical chain that motivates the RGAIG framework developed in Chapter~3.}
\label{fig:ch1_roadmap}
\end{figure}

\medskip
\noindent
\small\textbf{Research Flow Position --- Chapter~1}\\[3pt]
This chapter sits at the start of the research sequence: \textbf{Problem $\to$ Gap $\to$ Objective $\to$ RQ $\to$ Hypothesis $\to$ Data overview $\to$ Significance}.
It feeds into: Chapter~2 (literature review validates gaps), Chapter~3 (methodology operationalises RQs via two pipelines), Chapter~4 (analysis tests hypotheses), Chapter~5 (discussion interprets findings).
Two research pipelines are introduced here and detailed in Chapter~3:
\textbf{Pipeline~1} (secondary data: EEG/images $\to$ quantitative analysis $\to$ H1--H3) and
\textbf{Pipeline~2} (primary data: questionnaires $\to$ qualitative analysis $\to$ H4--H5).


\vspace{0.5em}

\noindent Table~\ref{tab:ch1_kpi_targets} defines the measurable targets that the research must achieve. Each KPI maps to a specific hypothesis and is evaluated in Chapter~\ref{chap:analysis}.

\paragraph{Chapter 1 Measurable KPI Targets.}







{\tablestripes\tablefontsize
\needspace{10\baselineskip}
\begin{xltabular}{\textwidth}{@{} >{\raggedright\arraybackslash}p{3cm} >{\raggedright\arraybackslash}p{2.2cm} L >{\raggedright\arraybackslash}p{1.2cm} @{}}
\caption[Chapter~1 Measurable KPI Targets]{Chapter~1 Measurable KPI Targets. Each KPI is linked to a specific hypothesis (H1--H5) and defines the quantitative threshold that Chapter~4 must meet to accept the hypothesis; the reader should use this table as the success rubric for the entire study.}\label{tab:ch1_kpi_targets} \\
\toprule
\thead{KPI} & \thead{Target} & \thead{What It Measures} & \thead{H} \\
\midrule
\endfirsthead
\endhead
\bottomrule
\endlastfoot
ASD-focused accuracy (ASD ensemble) & $\geq$85\% in $\geq$4 domains & Whether a unified pipeline generalises across heterogeneous biomedical domains & H1 \\
\addlinespace
DL vs.\ Classical effect size & Cohen's $d \geq 0.50$ & Whether deep learning provides statistically significant improvement over classical baselines & H2 \\
\addlinespace
Discriminative features per domain & $\geq$3 per domain & Whether SHAP-identified features align with established clinical biomarkers & H3 \\
\addlinespace
Pipeline effort reduction & $\geq$90\% & Whether agentic automation eliminates manual preprocessing steps & H4 \\
\addlinespace
Expert explainability rating & $M \geq 4.0/5.0$ & Whether RAG-generated explanations meet clinical interpretability standards & H5 \\
\end{xltabular}}
\vspace{0.5em}
\noindent\textit{Note.} ASD classification accuracy; DL = deep learning; SHAP = SHapley Additive exPlanations; RAG = retrieval-augmented generation. All targets established a priori (Chapter~\ref{chap:methodology}).

\addcontentsline{toc}{paragraph}{Must-Have: Objectives} % [TOC-must-have-insertion]
\subsection{Chapter Objective}

This chapter establishes the research problem, articulates five research questions (RQ1--RQ5), derives five objectives (O1--O5), and formulates five testable hypotheses (H1--H5). It argues that domain fragmentation in biomedical AI is a design convention rather than a technical constraint, and positions the RGAIG framework as a governance-aligned response to six identified research gaps (G1--G6).

\medskip
\noindent
This research addresses the fragmentation of biomedical assessment by designing and evaluating a unified framework that integrates conversational patient onboarding, structured behavioural and psychological assessment, and EEG/IoT-based biomedical analysis into an explainable, governance-aware decision-support workflow across multiple disease domains.

Three clinical scenarios illustrate the human cost of fragmented diagnostic AI:

\begin{itemize}[leftmargin=*, itemsep=3pt]
    \item \textbf{Autism spectrum disorder (ASD):} A mother in rural Montana waits over a year for a developmental paediatrician, pays thousands for clinical electroencephalography (EEG), and receives an inconclusive report.
    \item \textbf{Leukaemia:} A patient waits weeks for morphological classification that deep learning can perform in seconds.
\end{itemize}


\noindent These are not edge cases; they are the norm. The prevailing approach to clinical conceptual AI governance evaluation---fragmented across conditions, opaque in reasoning, manual in execution---inflicts avoidable costs on patients and healthcare organisations. This dissertation argues that domain fragmentation is a design choice, not a technical inevitability, and provides both the evidence and the architecture to replace it~\cite{esteva2019guide, topol2019deep}.

\textit{A note on ``unified.''} Throughout this dissertation, ``unified framework'' denotes \textbf{pipeline infrastructure unification}---shared preprocessing, orchestration, governance, and explainability layers. The ASD domain uses a domain-specific classifier (EEGNet) because domain-specific optimisation is essential for clinical-grade accuracy. What is unified is the surrounding infrastructure: a common data-ingestion pipeline, a single agentic orchestration layer, a shared RAG-based explainability engine, and one governance framework. This distinction is critical: the contribution is system-level integration, not a single universal model.

%===============================================================================
\subsection{Definition of Key Terms}
\label{sec:key_terms}

The following terms carry specific operational meanings throughout this dissertation:


\noindent Table~\ref{tab:key_terms} defines definition of Key Terms.

{\tablestripes\tablefontsize
\needspace{10\baselineskip}
\begin{xltabular}{\textwidth}{@{} >{\raggedright\arraybackslash}p{2.8cm} L @{}}
\caption[Definition of Key Terms]{Definition of Key Terms. Operational definitions are provided here to eliminate ambiguity in subsequent chapters; each term is scoped to its meaning within this dissertation rather than its broadest disciplinary usage.}\label{tab:key_terms} \\
\toprule
\thead{Term} & \thead{Operational Definition} \\
\midrule
\endfirsthead
\endhead
\bottomrule
\endlastfoot
Unified framework & Pipeline infrastructure unification---shared preprocessing, orchestration, governance, and explainability layers---not a single universal model. Each domain retains its own classifier. \\
\addlinespace
RGAIG & Responsible Generative AI Governance: a five-layer framework embedding responsible AI principles, governance controls, and explainability into every pipeline stage. \\
\addlinespace
Agentic AI & Autonomous software agents that orchestrate multi-step computational workflows (data ingestion, preprocessing, training, evaluation, reporting) without manual intervention. \\
\addlinespace
RAG & Retrieval-Augmented Generation: retrieves peer-reviewed literature from a vector database (ChromaDB) and generates natural-language clinical explanations grounded in cited evidence. \\
\addlinespace
Explainability & The combination of SHAP feature attribution (which features drove the prediction) and RAG narrative generation (why those features matter clinically). \\
\addlinespace
Clinical-grade accuracy & Classification performance $\geq$85\% with documented confidence intervals, validated through stratified cross-validation and external expert review. \\
\addlinespace
Consumer-grade EEG & Commercially available EEG headsets (Emotiv EPOC~X, Insight) costing \$299--\$849, as distinct from research-grade systems costing \$10,000+. \\
\addlinespace
ASD ensemble & ASD ensemble classification: equal-weight mean of per-ASD classification accuracy for ASD screening. \\
\addlinespace
Mixed-methods design & Convergent parallel design in which quantitative (EEG/imaging classification) and qualitative (expert panel, patient questionnaires) data streams are collected in parallel, analysed independently, and integrated through a joint display matrix. \\
\addlinespace
Governance-embedded & AI governance controls (audit trails, bias checks, escalation rules, compliance gates) built into the pipeline architecture rather than applied as post-deployment audits. \\
\end{xltabular}}
\vspace{0.3em}\noindent\textit{Note.} RGAIG = Responsible Generative AI Governance; SHAP = SHapley Additive exPlanations; RAG = Retrieval-Augmented Generation; ASD = Autism Spectrum Disorder. See chapter prose for full context.

%===============================================================================
\addcontentsline{toc}{paragraph}{Must-Have: Limitations} % [TOC-must-have-insertion]
\subsection{Delimitations}
\label{sec:delimitations}

Delimitations define the boundaries the researcher deliberately set for this study:

\begin{itemize}[leftmargin=*, itemsep=3pt]
    \item \textbf{Disease scope:} The ASD domain was selected based on EEG data availability and clinical burden. Other neurological conditions are excluded from hypothesis testing.
    \item \textbf{Data type:} Only retrospective, de-identified public benchmark datasets are used. No prospective clinical trial data or real-time patient data were collected for the quantitative pipeline.
    \item \textbf{Device scope:} Consumer-grade Emotiv devices only. Research-grade EEG systems (e.g., BioSemi, g.tec) are excluded to test the viability of affordable, scalable devices.
    \item \textbf{Geography:} Datasets predominantly from North American and European institutions. Indian clinical sites provide secondary validation only. Other geographies (Africa, Southeast Asia) are excluded.
    \item \textbf{Expert panel:} Twelve specialists were selected via convenience sampling. Larger-scale clinician surveys and patient experience studies are deferred to future research.
    \item \textbf{Regulatory:} The framework is designed for FDA SaMD and EU AI Act alignment but has not undergone formal regulatory submission. Compliance is assessed through internal audit, not regulatory review.
    \item \textbf{Business case:} illustrative financial-scenario projections (135--185\%) are modelled from market data and evaluation cost estimates, not measured from actual organisational deployment.
\end{itemize}

%===============================================================================
\subsection{Ethical Considerations}
\label{sec:ethical_ch1}

This research involves secondary analysis of de-identified EEG and medical imaging datasets, primary collection of patient/parent questionnaire responses, and expert panel evaluations. Key ethical considerations include:

\begin{itemize}[leftmargin=*, itemsep=3pt]
    \item \textbf{IRB status:} The study received exemption under 45 CFR 46.104 for secondary analysis of publicly available, de-identified datasets. The expert panel survey was reviewed as exempt under the survey research provision. Full IRB documentation is provided in Chapter~\ref{chap:methodology}, Section~\ref{sec:ethics_merged}.
    \item \textbf{Informed consent:} Public benchmark datasets operate under institutional data use agreements. Expert panel participants provided written informed consent. Patient questionnaire respondents consent through clinical intake protocols.
    \item \textbf{Bias and fairness:} The study reports demographic subgroup performance where available. Limitations due to dataset demographic composition (e.g., geographic bias, age distribution) are acknowledged in Chapter~\ref{chap:discussion}.
    \item \textbf{AI governance:} The 525-module quality taxonomy includes responsible AI modules covering fairness, bias detection, data governance, and compliance verification.
\end{itemize}

%===============================================================================
\addcontentsline{toc}{paragraph}{Must-Have: Scope} % [TOC-must-have-insertion]
\subsection{Scope and Boundaries of the Study}
\label{sec:scope_boundaries}

\noindent\textbf{Chapter 1 Boundary Statement (examiner-facing).}
\label{para:ch1_boundary_statement}
This dissertation evaluates a \emph{responsible AI governance and trust framework} for non-diagnostic ASD screening \emph{support}; it does \textbf{not} clinically validate, deploy, or certify an autonomous ASD AI system. Where the dissertation retains language around accuracy, ROI, FDA SaMD, EU AI Act, governance maturity, illustrative payback, or future-validation pathways, those references are \emph{conceptual reference frameworks and scenario-modelling artefacts} --- not measured operational outcomes, certifications, deployment recommendations, or clinical-effectiveness claims. The dissertation contribution is the \emph{governance + explainability + adoption framework itself}, validated through expert review and stakeholder evaluation on existing reference datasets (KAU + ABIDE-II + 131-clinician PLS-SEM + 12-expert qualitative panel). Any future caregiver-facing, clinical, wearable, or escalation-enabled use would require separate IRB review, prospective clinical validation, regulatory assessment, and operational governance --- explicitly out of scope here. \textbf{Prior IRB approval references} for the $n{=}131$ PLS-SEM and $n{=}12$ expert-panel data streams are consolidated in Chapter~\ref{chap:methodology} Section~\ref{para:prior_irb_approval_trail} (Prior IRB Approval Trail) --- the dissertation's current submission is exempt-category secondary-data research; no new participant recruitment occurs.

To prevent misinterpretation of the research's intent and to position the dissertation appropriately within the DBA tradition, this section makes explicit what the study \emph{is} and what it is \emph{not}.

\paragraph{What this dissertation IS.}
\begin{itemize}[leftmargin=*, itemsep=3pt]
    \item A \textbf{governance-oriented} doctoral study of how responsible AI principles --- transparency, accountability, explainability, fairness, human oversight --- can be operationalised within AI-assisted ASD-EEG screening support ecosystems.
    \item A \textbf{conceptual framework} (the Responsible Generative AI Governance framework, RGAIG) plus its \textbf{stakeholder evaluation}, validated through expert review, Delphi-style consensus, and survey instruments aligned with the framework's governance, explainability, and adoption dimensions.
    \item A \textbf{human-centered} study that places the caregiver--clinician dyad at the centre of the design surface: the caregiver receives plain-language explainability outputs at home; the clinician supervises, validates, and retains authority over every consequential decision.
    \item A \textbf{socio-technical} adoption study examining how trust, explainability, governance maturity, and workflow integration influence the adoption of AI-assisted screening support in healthcare settings.
\end{itemize}

\paragraph{What this dissertation is NOT.}
\begin{itemize}[leftmargin=*, itemsep=3pt]
    \item It is \textbf{not} a clinical efficacy trial. The classification accuracy results reported in Chapter~\ref{chap:analysis} are computed on retrospective, de-identified public benchmark datasets; they do \emph{not} constitute medical evidence of diagnostic performance on prospective patients.
    \item It does \textbf{not} replace clinicians, paediatric neurologists, or any other medical professional. Every output of the framework is presented as decision-support that requires clinician validation; no path through the framework reaches a patient without a human-in-the-loop checkpoint.
    \item It does \textbf{not} autonomously diagnose ASD or any other condition. The framework operates strictly as an assistive screening-support layer; final diagnostic authority rests with qualified clinicians under existing healthcare governance.
    \item It is \textbf{not} a certified medical device or a future-operational-validation pathway. The framework is presented as a conceptual reference architecture; certification under SaMD / FDA / EU MDR pathways is explicitly out of scope and identified as future work.
    \item It does \textbf{not} provide real-time autonomous intervention, automated escalation that bypasses human review, or live therapeutic monitoring of children. The continuous-monitoring elements described in the architecture are framed as conceptual capability surfaces requiring future prospective future clinical-validation pathway.
    \item It does \textbf{not} claim emotional or psychological intervention effectiveness. References to caregiver reassurance describe \emph{trust-aware communication} and \emph{explainability-centered reporting}, not therapeutic outcomes.
\end{itemize}

\paragraph{Scope summary.} The dissertation contributes a governance, explainability, and adoption framework for AI-assisted ASD-EEG screening support. The contribution is evaluative and conceptual rather than clinical-deployment-oriented; the validation strategy reflects this scope (Chapter~\ref{chap:methodology}).

\needspace{12\baselineskip}
\needspace{12\baselineskip}
\noindent\textbf{Research Object, Scope, and Task Map (Summary Tables).} Tables~\ref{tab:ch1_object_scope}, \ref{tab:ch1_not_in_scope}, and \ref{tab:ch1_task_map} consolidate the prose above into examiner-facing reference tables. They formalise what the dissertation aims to deliver, what is in scope, what is explicitly excluded, and which tasks are completed, in progress, or future scope.

{\tablestripes\tablefontsize

\needspace{10\baselineskip}
\begin{xltabular}{\textwidth}{@{} >{\centering\arraybackslash}p{1.0cm} >{\raggedright\arraybackslash}p{8.9cm} p{4.4cm} @{}}
\caption[Research Object \& Scope --- IN SCOPE]{Research Object and Scope --- IN SCOPE. Each row names a contribution the dissertation actively delivers, with the chapter(s) where the evidence lives.}\label{tab:ch1_object_scope} \\
\toprule
\thead{ID} & \thead{Research Object / Scope Item} & \thead{Evidence Location} \\
\midrule
\endfirsthead
\endhead
\bottomrule
\endlastfoot
S1 & RGAIG conceptual governance framework (5-layer) & Ch.~\ref{chap:methodology}, \ref{chap:discussion} \\
S2 & PD-1 PLS-SEM clinician survey ($n{=}131$) & Ch.~\ref{chap:methodology}, \ref{chap:analysis} \\
S3 & PD-2 expert qualitative panel ($n{=}12$) & Ch.~\ref{chap:methodology}, \ref{chap:analysis} \\
S4 & Secondary EEG datasets --- KAU ($n{=}768$) + ABIDE-II retrospective & Ch.~\ref{chap:methodology}, \ref{chap:analysis} \\
S5 & Systematic literature review (156 studies, 6 categories) & Ch.~\ref{chap:literature} \\
S6 & RGAIG Maturity Framework (6 levels) & Ch.~\ref{chap:discussion} \\
S7 & Explainability two-layer design (clinician + caregiver) & Ch.~\ref{chap:discussion} \\
S8 & Governance-explainability-trust-adoption pathway (empirically tested) & Ch.~\ref{chap:analysis} \\
S9 & SMART-E Operating System reference architecture & Ch.~\ref{chap:discussion} §\ref{sec:smart_e_extension} \\
S10 & Defensibility scaffold across 13 appendices & Appendices B--N \\
\end{xltabular}}
\vspace{0.3em}\noindent\textit{Note.} PD-1 = Primary Data 1 (n=131 clinician survey); PD-2 = Primary Data 2 (n=12 expert panel); RGAIG = Responsible Generative AI Governance; PLS-SEM = Partial Least Squares SEM. See chapter prose for full context.

\paragraph{Research Out of Scope.}

\noindent Table~\ref{tab:ch1_not_in_scope} below lists deliberate exclusions; reading it alongside the in-scope table prevents over-claim and clarifies which questions the dissertation does NOT answer.


{\tablestripes\tablefontsize
\needspace{10\baselineskip}
\begin{xltabular}{\textwidth}{@{} >{\centering\arraybackslash}p{1.2cm} >{\raggedright\arraybackslash}p{4.5cm} L @{}}
\caption[Research Object \& Scope --- NOT IN SCOPE]{Research Object and Scope --- NOT IN SCOPE. Each row names a deliberate exclusion to prevent over-claim; many are flagged as future work.}\label{tab:ch1_not_in_scope} \\
\toprule
\thead{ID} & \thead{Out of Scope Item} & \thead{Rationale / Future Work Pointer} \\
\midrule
\endfirsthead
\endhead
\bottomrule
\endlastfoot
N1 & Prospective caregiver primary data (PD-3, $n{=}50$--$150$) & Future scope; requires separate IRB \\
N2 & Educator / therapist primary data (PD-4) & Future scope \\
N3 & Compliance / IRB officer governance assessment (PD-5) & Future scope \\
N4 & ASD-adult structured survey (PD-6) & Future scope \\
N5 & FDA SaMD / EU MDR clinical-device certification & Conceptual reference only; future certification track \\
N6 & Prospective ASD clinical-effectiveness trial & Future-validation pathway; out of scope here \\
N7 & Autonomous (no-HITL) deployment of the AI system & Architecturally prohibited; framework is assistive-only \\
N8 & Real-time emotional or psychological intervention & Out of scope; conceptual capability surface only \\
N9 & Multi-jurisdiction regulatory compliance audit & Conceptual alignment table only; not audit-grade evidence \\
N10 & Production-grade enterprise deployment (live runtime) & SMART-E extension is reference architecture, not deployment \\
\end{xltabular}}
\vspace{0.3em}\noindent\textit{Note.} HITL = Human-in-the-loop; ASD = Autism Spectrum Disorder. See chapter prose for full context.

\paragraph{Research Task Map.}

\noindent Table~\ref{tab:ch1_task_map} below lists the major dissertation tasks with their status; this is the working delivery map the chapters report against.


{\tablestripes\tablefontsize
\needspace{10\baselineskip}
\begin{xltabular}{\textwidth}{@{} >{\centering\arraybackslash}p{1.0cm} >{\raggedright\arraybackslash}p{8.1cm} >{\raggedright\arraybackslash}p{1.0cm} p{4.1cm} @{}}
\caption[Research Task Map]{Research Task Map. Tracks the major dissertation tasks with status (\checkmark Done, $\sim$ Partial, $\square$ Future) and primary evidence location.}\label{tab:ch1_task_map} \\
\toprule
\thead{\#} & \thead{Task} & \thead{Status} & \thead{Location} \\
\midrule
\endfirsthead
\endhead
\bottomrule
\endlastfoot
T1 & Define research problem (P1--P4) & \checkmark & Ch.~\ref{chap:introduction} §\ref{sec:problem_statement} \\
T2 & Systematic literature review (156 studies) & \checkmark & Ch.~\ref{chap:literature} \\
T3 & Design PD-1 PLS-SEM instrument (26 items, 11 sections) & \checkmark & Ch.~\ref{chap:methodology} \\
T4 & Pilot \& validate instrument ($n{=}23$--$40$) & \checkmark & Ch.~\ref{chap:methodology} \\
T5 & Collect PD-1 data ($n{=}131$) & \checkmark & Ch.~\ref{chap:methodology}, \ref{chap:analysis} \\
T6 & Recruit PD-2 expert panel ($n{=}12$) & \checkmark & Ch.~\ref{chap:methodology}, \ref{chap:analysis} \\
T7 & PLS-SEM hypothesis test (H1--H5 primary) & \checkmark & Ch.~\ref{chap:analysis} \\
T8 & Secondary-data EEG validation (KAU + ABIDE-II) & \checkmark & Ch.~\ref{chap:analysis} \\
T9 & XAI / SHAP / RAG explainability evaluation & \checkmark & Ch.~\ref{chap:analysis} \\
T10 & RGAIG Maturity Framework synthesis (6 levels) & \checkmark & Ch.~\ref{chap:discussion} \\
T11 & SMART-E Operating System extension reference architecture & \checkmark & Ch.~\ref{chap:discussion} §\ref{sec:smart_e_extension} \\
T12 & Examiner Q\&A tables across all 5 chapters & \checkmark & each chapter \\
T13 & Prospective caregiver collection (PD-3) & $\square$ & Future work \\
T14 & Multi-site recruitment (KAU + AIIMS + 3 others, $n{=}1{,}350$) & $\square$ & Future work |Ch.~\ref{chap:discussion} \\
T15 & Live H6 matched-cohort study (concern$\to$therapy reduction) & $\square$ & Future work \\
\end{xltabular}}
\vspace{0.3em}\noindent\textit{Note.} PD-1 = Primary Data 1 (n=131 clinician survey); PD-2 = Primary Data 2 (n=12 expert panel); RGAIG = Responsible Generative AI Governance; PLS-SEM = Partial Least Squares SEM. See chapter prose for full context.

%===============================================================================
\subsection{Ethical Positioning of the Framework}
\label{sec:ethical_positioning}

Beyond the procedural ethics described in Section~\ref{sec:ethical_ch1}, this subsection states the \emph{ethical posture} the proposed framework adopts toward conceptual AI governance evaluation in a vulnerable-population context (paediatric and family-caregiver settings).

\paragraph{Preservation of clinician authority.} The framework is designed around \textbf{human-in-the-loop (HITL)} governance at every decision point that could affect a patient. Clinician review is required before any AI-assisted screening output reaches a consequential downstream action (e.g., referral, escalation, treatment-pathway recommendation). The architecture does not contain any path in which an autonomous agent can override, bypass, or substitute clinician judgement.

\paragraph{Overtrust mitigation.} The framework explicitly addresses the risk that caregivers may place excessive trust in AI outputs:
\begin{itemize}[leftmargin=*, itemsep=3pt]
    \item Every caregiver-facing output carries an explicit confidence band (low / medium / high) and a plain-language uncertainty statement, never a binary diagnostic verdict.
    \item Every output includes a disclosure that the AI is a screening-support tool, not a diagnostic authority, and that clinician confirmation is required for any decision affecting care.
    \item The framework's evaluation in Chapter~\ref{chap:analysis} explicitly measures overtrust risk through stakeholder perception of uncertainty communication and trust calibration.
\end{itemize}

\paragraph{Explainability for two audiences.} A single AI inference is rendered in two distinct layers (Chapter~\ref{chap:discussion}, Tab.~\ref{tab:explainability_two_layer}): the \emph{clinician layer} (SHAP attribution + EEG-feature anchoring + RAG citation panel), and the \emph{caregiver layer} (plain language + confidence band + clear next step). The two layers must be \emph{provably consistent} --- the caregiver's reading cannot diverge from the clinician's reading --- and consistency is enforced architecturally by deriving both layers from the same model output.

\paragraph{Vulnerable-population safeguards.} Because the framework targets paediatric ASD screening, additional safeguards apply: consent governance is family-centred (parent or legal guardian), all data flows preserve minor-data protections under HIPAA and GDPR Article~8, and the framework's fairness monitoring includes age, gender, and ethnicity subgroup analysis with disparate-impact $\geq 0.80$ as a pre-deployment gate (Chapter~\ref{chap:methodology}).

\addcontentsline{toc}{paragraph}{Must-Have: Aim} % [TOC-must-have-insertion]
\paragraph{Boundary on emotional and psychological claims.} The framework supports caregivers through \emph{trust-aware explainability communication}; it does not provide emotional reassurance, psychological intervention, or therapeutic monitoring. These are explicitly outside the framework's scope, and any future extension into therapeutic territory would require separate IRB review, clinical-trial protocols, and certification pathways.

\paragraph{Acknowledged limitations of the ethical posture.}
\begin{itemize}[leftmargin=*, itemsep=3pt]
    \item Explainability mechanisms reduce, but do not eliminate, the risk of overtrust; longitudinal observation of caregiver behaviour is required to confirm trust calibration in field deployment.
    \item Governance accountability requires institutional commitment beyond the framework itself; the framework can prescribe roles (e.g., RACI matrix) but cannot enforce them without organisational adoption.
    \item Subgroup fairness is constrained by the demographic composition of available benchmark datasets; fairness claims are bounded by data-availability acknowledgements documented in Chapter~\ref{chap:methodology}.
\end{itemize}

%===============================================================================
\subsection{Context of Contributing Organizations}
\label{sec:contributing_orgs}

This research was conducted entirely within an academic and open-access institutional context. No commercial sponsorship, proprietary data, or vendor-specific licensing shaped the research design or findings. The contributing organizations fall into four categories:

\begin{itemize}[leftmargin=*, itemsep=3pt]
    \item \textbf{Academic institution --- Golden Gate University (GGU):} This dissertation is submitted in partial fulfillment of the Doctor of Business Administration (DBA) programme at Golden Gate University, San Francisco. GGU's Executive Technology track contextualizes the research within evidence-based management practice, requiring empirical validation, governance alignment, and managerial utility alongside technical contribution.

\iffalse % AUTO-SUPPRESS non-ASD para
    \item \textbf{Public dataset repositories:} All quantitative data originate from openly licensed benchmark repositories. EEG datasets for ASD, PD, epilepsy, and stress classification derive from PhysioNet and the Temple University Hospital (TUH) EEG Corpus. Cancer histopathology images derive from Kaggle-hosted repositories (ISIC, BreakHis). No dataset requires institutional data transfer agreements beyond standard platform terms of use; all are reproducible by independent researchers.
\fi

    \item \textbf{Open-source software frameworks:} The implementation codebase depends exclusively on open-source tools --- TensorFlow and PyTorch for deep learning, scikit-learn for classical ML, LangChain for agentic AI orchestration, and ChromaDB for retrieval-augmented generation (RAG) vector storage. This dependency structure ensures full reproducibility and eliminates proprietary lock-in from the research design.

    \item \textbf{Expert panel participants:} The qualitative pipeline (Pipeline~2) engaged twelve clinical and AI professionals drawn from healthcare delivery, neuroscience research, and AI governance practice. Participants were recruited through convenience sampling via professional networks. All twelve provided written informed consent; their organizational affiliations are withheld to preserve anonymity. The panel's composition --- spanning clinicians, data scientists, and governance specialists --- was designed to triangulate technical, clinical, and regulatory perspectives on RGAIG adoption.
\end{itemize}


\noindent This configuration --- academic supervision, public data, open-source tooling, and volunteer expert assessment --- means that all findings presented in this dissertation are independently verifiable and free from commercial conflicts of interest.

%===============================================================================
\addcontentsline{toc}{paragraph}{Must-Have: Thesis Structure} % [TOC-must-have-insertion]
\subsection{Organisation of the Thesis}
\label{sec:thesis_organisation}

This dissertation comprises five chapters, each building on the preceding one:

\textbf{Chapter~1 (Introduction)} establishes the research problem (P1--P4), defines the study's scope and significance, and presents the objectives, research questions, hypotheses, and contributions that guide the remainder of the dissertation.

\textbf{Chapter~2 (Literature Review)} synthesises 156 peer-reviewed studies across six thematic categories to validate the six research gaps (G1--G6). It maps the relevant theoretical foundations to the proposed framework and develops the conceptual basis for the study.

\textbf{Chapter~3 (Research Methods)} operationalises the research questions through a pragmatist, design-science methodology. It explains the quantitative and qualitative strands of the study, the datasets and instruments used, and the validation procedures that connect each research question to evidence.

\textbf{Chapter~4 (Data Analysis and Findings)} reports the empirical results for the study's hypotheses, distinguishing the primary findings from the supporting analyses used to assess robustness, explainability, and validation.

\textbf{Chapter~5 (Discussion and Recommendations)} interprets the findings against the Chapter~2 literature, addresses limitations, and outlines the study's academic, managerial, and future research implications.

\textbf{Appendices} contain supporting material for the main chapters, including supplementary methodological detail, extended result tables, instruments, and documentary evidence.

\needspace{12\baselineskip}
\needspace{12\baselineskip}
\noindent\textbf{Cross-chapter phase \& data alignment.} Table~\ref{tab:chapter_phase_alignment} consolidates the dissertation's binding chapter-to-phase mapping so the reader can locate any contribution from the chapter-level prose above into the supplementary phase-level reference materials (\textsc{DATA\_FLOW\_MASTER\_REFERENCE.md}, \textsc{INTEGRATION\_ARCHITECTURE.md}, \textsc{FLOW\_DIAGRAMS.md}). The mapping is content-derived rather than forecast-derived: each phase number names the contribution; the chapter that owns it is determined by what the phase actually delivers.

\paragraph{Cross-Chapter Phase \& Data Alignment.}



{\tablestripes\tablefontsize

\needspace{10\baselineskip}
\begin{xltabular}{\textwidth}{@{} >{\raggedright\arraybackslash}p{2.5cm} >{\raggedright\arraybackslash}p{2.0cm} >{\raggedright\arraybackslash}p{2.5cm} >{\raggedright\arraybackslash}p{2.1cm} p{5.0cm} @{}}
\caption[Cross-Chapter Phase \& Data Alignment]{Cross-Chapter Phase \& Data Alignment. Each chapter is anchored to its core phase range, secondary-data phases (SD-x), and hybrid integration phases (H-x). The five-pillar SMART contract is extended to SMART-E (Enterprise) in Chapter~\ref{chap:discussion}, Section~\ref{sec:smart_e_extension}.}\label{tab:chapter_phase_alignment} \\
\toprule
\thead{Chapter} & \thead{Core Phases} & \thead{SD Phases} & \thead{Hybrid Phases (H-x)} & \thead{Primary Focus} \\
\midrule
\endfirsthead
\endhead
\bottomrule
\endlastfoot
Ch.~1 Introduction & P1--P2 & SD-1 & --- & Research foundation: problem, objectives, hypothesis design \\
Ch.~2 Literature & P12--P17 & SD-2, SD-3, SD-5 & H-1, H-2, H-3 & Literature \& conceptual theory: governance, XAI, trust theory \\
Ch.~3 Methods & P3--P4 & SD-1, SD-2, SD-3 & H-4 & Methodology \& primary data design: PD-1 ($n{=}131$ PLS-SEM) + PD-2 ($n{=}12$ expert panel) \\
Ch.~4 Findings & P5--P7 & SD-4, SD-5, SD-6, SD-7 & H-5, H-7 & Statistical + EEG validation \& findings \\
Ch.~5 Discussion & P47--P52 & SD-6, SD-7, SD-8, SD-9, SD-10 & H-6$\star$, H-8 to H-15 & SMART-E enterprise governance framework \\
Future Research & P53--P61 & --- & --- & Long-horizon constitutional / planetary governance \\
Appendices (B--N) & --- & --- & --- & Methodological detail, instruments, per-chapter Q\&A \\
\end{xltabular}}
\vspace{0.5em}
\noindent\textit{Note.} P1--P11 are the empirical research core; P47--P52 are the enterprise governance differentiator. The direct jump from Chapter~4 to Chapter~5 is intentional and academically defensible: most dissertations stop at validation; this work continues into runtime constitutional AI. H-6$\star$ (Runtime Constitutional Governance) is the operator-flagged headline hybrid differentiator. SD-1 to SD-10 form the canonical secondary-data execution stack. The supplementary reference set contains the verbatim per-phase content; this table serves as the navigation index.

% [SPACE-FIX] clearpage removed to eliminate empty page

\iffalse %% Condensed

\noindent\textbf{Why this visual.} Before examining the diagnostic burden data, it is essential to see the complete clinical pipeline that RGAIG proposes---from a patient wearing a consumer EEG headset to the clinician receiving an AI-generated recommendation. Without this end-to-end view, the individual components discussed later in this chapter lack operational context.
\fi


\noindent Figure~\ref{fig:patient_journey} presents the complete clinical pipeline from wearable EEG acquisition to clinical decision.



\noindent\textbf{End-to-end patient journey: from wearable EEG.}

\needspace{12\baselineskip}
\begin{figure}[!htbp]
\centering

\vspace*{-0.5cm}
\resizebox{0.97\textwidth}{!}{%
\begin{tikzpicture}[
  every node/.style={font=\fignodefont},
  pstep/.style={rectangle, rounded corners=4pt, draw=ggunavy,
    text width=3.2cm, minimum height=2.2cm, align=center, font=\small, inner sep=6pt},
  arr/.style={-{Stealth[length=4mm, width=3mm]}, very thick, ggunavy},
  lbl/.style={font=\small\itshape, text=ggunavy!70, above, midway}
]
%% Row 1: Steps 1-4 (each with different light colour)
% 
\node[pstep, fill=blue!12] (n1) {\textbf{1. Patient}\\[3pt]Wears Emotiv\\EPOC X EEG\\(14 channels)};
% 
\node[pstep, fill=cyan!15, right=0.5cm of n1] (n2) {\textbf{2. EEG Signal}\\[3pt]Raw neural\\recording\\(128 Hz, BLE)};
% 
\node[pstep, fill=green!12, right=0.5cm of n2] (n3) {\textbf{3. Preprocess}\\[3pt]Bandpass filter,\\ICA artefact,\\segmentation};
% 
\node[pstep, fill=yellow!18, right=0.5cm of n3] (n4) {\textbf{4. Features}\\[3pt]127 engineered\\+ deep learned\\features};
% 
%% Row 2: Steps 5-8 (each with different light colour)
% 
\node[pstep, fill=orange!15, below=1cm of n1] (n5) {\textbf{5. AI Model}\\[3pt]DL ensemble\\ASD-risk\\prediction};
% 
\node[pstep, fill=pink!18, right=0.5cm of n5] (n6) {\textbf{6. RAG}\\[3pt]Literature-backed\\clinical\\explanation};
% 
\node[pstep, fill=violet!12, right=0.5cm of n6] (n7) {\textbf{7. Decision}\\[3pt]Risk score +\\confidence +\\recommendation};
% 
\node[pstep, fill=teal!15, right=0.5cm of n7] (n8) {\textbf{8. Action}\\[3pt]Alert patient,\\refer to doctor,\\or monitor};
% 
%% Arrows
\draw[arr] (n1) -- (n2);
\draw[arr] (n2) -- (n3);
\draw[arr] (n3) -- (n4);
\draw[arr] (n4.south) -- ++(0,-0.35) -| (n5.north);
\draw[arr] (n5) -- (n6);
\draw[arr] (n6) -- (n7);
\draw[arr] (n7) -- (n8);
\end{tikzpicture}%
}
\caption[End-to-end patient journey]{End-to-end patient journey: from wearable EEG acquisition through AI prediction and RAG explanation to clinical decision and action. The eight-step pipeline processes raw neural signals through preprocessing, feature extraction, deep learning classification, and retrieval-augmented generation to produce actionable clinical recommendations.}
\label{fig:patient_journey}
\end{figure}

\iffalse %% Condensed — caption covers the pipeline

\noindent\textbf{What it shows.} Figure~\ref{fig:patient_journey} illustrates an eight-step linear pipeline: (1)~patient wears a consumer-grade EEG device, (2)~raw signal acquisition, (3)~preprocessing with artefact removal, (4)~feature extraction (127 engineered plus deep-learned features), (5)~DL ensemble classification, (6)~RAG-based clinical explanation, (7)~risk-scored decision, and (8)~actionable clinical output. Each step maps to a distinct RGAIG architectural layer.


\noindent\textbf{So what.} The pipeline demonstrates that wearable-to-decision automation is achievable within a single governed workflow, replacing the fragmented, multi-visit diagnostic process described in the clinical scenarios above. This end-to-end design directly motivates objectives O1 (ASD-focused pipeline) and O4 (RAG explainability), and establishes the operational context for the problem analysis in Section~\ref{sec:background}.
\fi

%===============================================================================
\addcontentsline{toc}{paragraph}{Must-Have: Background} % [TOC-must-have-insertion]

\subsection{Why Governance Matters More Than Accuracy}
\label{sec:why_governance_over_accuracy}

\noindent Although advances in machine learning have enabled high predictive performance in ASD screening applications, evidence from healthcare AI adoption literature suggests that predictive accuracy alone is insufficient to achieve sustainable organizational adoption. Clinical deployment decisions are influenced not only by model performance but also by transparency, accountability, explainability, regulatory compliance, and stakeholder trust. Consequently, the primary challenge is not whether AI systems can predict accurately, but whether healthcare stakeholders are willing to trust, govern, and operationally adopt such systems. This study therefore positions Responsible AI Governance as the central organizational mechanism through which explainability, trust, and adoption readiness can be strengthened.


\subsection{Why This Is a DBA Problem (Not a Computer Science PhD)}
\label{sec:why_dba_problem}

\noindent The central focus of this dissertation is organizational adoption rather than algorithmic innovation. While machine learning, explainable AI, and retrieval-augmented generation technologies provide the technical foundation for the proposed solution, the primary research objective concerns managerial decision-making, governance maturity, trust formation, and organizational readiness for adoption. Accordingly, the dissertation aligns with the objectives of Doctor of Business Administration research by addressing a practical organizational problem situated at the intersection of technology management, healthcare governance, and innovation adoption.

\section{Background of the Study}
\label{sec:background}

Healthcare systems worldwide face a widening gap between diagnostic demand and specialist capacity (Table~\ref{tab:diagnostic_burden}). The World Health Organization reports that neurological disorders are the leading cause of disability and the second leading cause of death globally~\cite{who2024neuro}.


\needspace{12\baselineskip}
\label{tab:diagnostic_burden}
\label{tab:non_claims}
\label{tab:study_scope}
\label{fig:unified_framework}
\label{tab:boundary_conditions}
\iffalse % AUTO-SUPPRESS non-ASD
\noindent Table~\ref{tab:diagnostic_burden} global diagnostic burden.

\noindent The figures below isolate the load each delay places on caregiver mental health, sibling well-being, and household finances.
{\tablestripes\tablefontsize
\needspace{10\baselineskip}
\begin{xltabular}{\textwidth}{@{} l C L @{}}
\caption[Global Diagnostic Burden]{Global Diagnostic Burden for Conditions Addressed in This Dissertation}\label{tab:diagnostic_burden} \\
\toprule
\thead{Condition} & \thead{Global Prevalence} & \thead{Source} \\
\midrule
\endfirsthead
\endhead
\bottomrule
\endlastfoot
ASD & 1 in 100 children (78M individuals) & WHO, 2023 \\
PD & 8.5M (prevalence doubled, 1990--2016) & WHO, 2022 \\
Stress-related disorders & 970M people & WHO, 2022 \\
Cancer & 20M new cases annually & GLOBOCAN, 2022 \\
\addlinespace
\multicolumn{3}{>{\raggedright\arraybackslash}p{0.95\textwidth}}{\textit{Overall neurological burden}: 9M deaths, 276M DALYs, \$2.5T/year} \\
\end{xltabular}}
\vspace{0.5em}
\noindent\textit{Note.} M = million; DALYs = disability-adjusted life years; T = trillion. ASD = autism spectrum disorder;
\fi % AUTO-SUPPRESS non-ASD

Three features of the current healthcare AI landscape motivate this research:

\begin{itemize}[leftmargin=*, itemsep=3pt]
    \item \textbf{Market scale and fragmentation:} The global AI-in-healthcare market exceeded \$15 billion in 2023 and is projected to surpass \$180 billion by 2030, yet the vast majority of this investment flows toward narrow, single-condition tools that cannot generalise across clinical domains~\cite{rajkomar2019healthcare}.
    \item \textbf{EEG-based screening-support advantage:} EEG provides non-invasive, real-time neural oscillatory measurements with demonstrated utility in neurological and psychiatric conditions, enabling continuous capability surfaces at lower cost than imaging modalities~\cite{niedermeyer2005eeg, buzsaki2012brain}. \textit{Conceptual-scope note: references to ``continuous monitoring'' throughout this dissertation describe a capability surface of the conceptual framework, not a deployed surveillance system; pediatric continuous-monitoring evaluation requires separate IRB review, prospective clinical-trial protocols, and consent governance pathways that are explicitly out of scope (Section~\ref{sec:scope_boundaries}).}
    \item \textbf{Radiomics opportunity:} High-throughput extraction of quantitative features from medical images enables tumour characterisation, treatment stratification, and biopsy reduction, extending AI-assisted diagnosis beyond neurological conditions~\cite{lambin2012radiomics}.
\end{itemize}

Despite these advances, clinical AI remains fragmented. A hospital deploying AI-assisted diagnosis must procure, integrate, and maintain separate systems for each condition, each with its own preprocessing pipeline, model architecture, validation protocol, and explainability mechanism. This fragmentation multiplies infrastructure costs, governance burdens, and regulatory risk~\cite{craik2019eeg, kelly2019healthcare}.

Figure~\ref{fig:health_burden} illustrates the global prevalence of Autism Spectrum Disorder, the condition addressed in this dissertation.


\noindent\textbf{Global Prevalence of Autism Spectrum Disorder.}

\needspace{12\baselineskip}
\begin{figure}[!htbp]
\resizebox{0.97\textwidth}{!}{%
\begin{tikzpicture}[every node/.style={font=\fignodefont}]
\begin{axis}[
    ybar,
    width=\textwidth,
    height=7cm,
    bar width=1.2cm,
    ylabel={Global Prevalence (Millions)},
    ylabel style={font=\small, text=ggunavy},
    ymin=0, ymax=1100,
    xtick=data,
    symbolic x coords={ASD},
    xticklabel style={font=\small, text=ggunavy},
    yticklabel style={font=\small, text=ggunavy},
    nodes near coords,
    nodes near coords style={font=\small, color=ggunavy},
    axis lines*=left,
    enlarge x limits=0.15,
    grid=major,
    grid style={gray!20},
    major tick style={draw=ggunavy!40},
    draw=ggunavy,
    /pgfplots/ybar legend,
    legend style={font=\small, at={(0.98,0.95)}, anchor=north east},
]
\addplot[fill=lightblue, draw=ggunavy] coordinates {(ASD,78)};
\end{axis}
\end{tikzpicture}%
}
\caption[Global Prevalence of ASD]{Global Prevalence of Autism Spectrum Disorder}
\label{fig:health_burden}
\vspace{0.5em}
\noindent\textit{Note.} Data drawn from WHO~\cite{who2023asd}. ASD = autism spectrum disorder. Global prevalence motivates scalable AI-assisted diagnosis.
\end{figure}

The ASD domain uses EEG-based neural signal analysis. The unified infrastructure---preprocessing, orchestration, governance, and explainability layers---is designed to generalise to additional modalities in future work, but this dissertation evaluates it end-to-end only for ASD.

%===============================================================================
\subsection{Research Motivation}
\label{sec:why_now}

Three converging forces create both the opportunity and urgency for this research. Table~\ref{tab:why_now} summarizes each force, its evidence, and its implication.


{\tablestripes\tablefontsize

\needspace{10\baselineskip}
\begin{xltabular}{\textwidth}{@{} >{\raggedright\arraybackslash}p{2.2cm} p{4.1cm} p{4.1cm} p{3.7cm} @{}}
\caption[Why This Research, Why Now: Three Converging Forces]{Why This Research, Why Now: Three Converging Forces. The table identifies three technological and clinical developments (2021--2024) whose simultaneous maturation creates a window of opportunity for the unified framework proposed in this study.}\label{tab:why_now} \\
\toprule
\thead{Converging Force} & \thead{What Changed (2021--2024)} & \thead{Key Evidence} & \thead{Implication for This Study} \\
\midrule
\endfirsthead
\endhead
\bottomrule
\endlastfoot
Regulatory inflection & FDA AI/ML Action Plan (2021); EU AI Act (2024) classifies medical AI as high-risk & Transparency, bias monitoring, human oversight now mandated~\cite{fda2021aiml, eu2021aiact} & Governance and explainability are now compliance requirements \\
\addlinespace
Technological maturity & Deep learning, retrieval-augmented generation (RAG), and agentic AI all reached practical maturity simultaneously & RAG feasible post-2022; agentic AI feasible post-2023~\cite{esteva2019guide, lewis2020rag, wang2024agentai} & ASD-focused unification is now technically feasible \\
\addlinespace
Organizational pressure & Hospitals frustrated with proliferating narrow AI tools & CIOs/CMOs demand platform consolidation; estimated 3--5$\times$ cost multiplier from fragmentation (see Chapter~\ref{chap:discussion} for business model analysis)~\cite{rajkomar2019healthcare} & Demand for fewer tools, broader coverage, and lower illustrative cost-modelling scenario \\
\end{xltabular}}
\vspace{0.5em}
\noindent\textit{Note.} RAG = retrieval-augmented generation; illustrative cost-modelling scenario = total cost of ownership; CIO = chief information officer; CMO = chief medical officer. The trade-off: agentic AI and RAG components are evolving rapidly, requiring modular design for component upgrades without architectural redesign.

\iffalse %% Condensed — table conveys the convergence argument
This convergence positions the dissertation squarely within healthcare organizations' broader digital transformation agendas. The shift from fragmented point solutions to a unified diagnostic platform reflects the same strategic logic driving enterprise resource planning, cloud migration, and data platform consolidation across industries---except that clinical AI carries additional regulatory, ethical, and patient-safety requirements that make uncoordinated adoption particularly costly.
\fi %% End suppressed convergence paragraph

%===============================================================================
\addcontentsline{toc}{paragraph}{Must-Have: Significance} % [TOC-must-have-insertion]
\subsection{Significance of the Study}
\label{sec:significance}


\noindent Table~\ref{tab:significance} summarises the significance of this study across five dimensions: theoretical, methodological, practical, societal, and technological.




{\tablestripes\tablefontsize
\needspace{10\baselineskip}
\begin{xltabular}{\textwidth}{@{} >{\raggedright\arraybackslash}p{2.2cm} p{8.5cm} >{\raggedright\arraybackslash}p{3.7cm} @{}}
\caption[Significance of the Study Across Five Dimensions]{Significance of the Study Across Five Dimensions. The five dimensions---clinical, technological, organisational, regulatory, and economic---collectively demonstrate that the RGAIG framework addresses a gap that no single-dimension solution can close.}\label{tab:significance} \\
\toprule
\thead{Dimension} & \thead{Contribution} & \thead{Beneficiary} \\
\midrule
\endfirsthead
\endhead
\bottomrule
\endlastfoot
\thead{Academic} & Extends RBV to AI platform consolidation and integrates nine theories into one model; among the first ASD-focused biomedical AI validations in the 156-study review & AI researchers, DBA scholars \\
\addlinespace
\textbf{Method.} & Among the first DSR studies combining multiple biomedical domains in one artifact; standardised evaluation across EEG and imaging & Methodologists, biomed.\ eng. \\
\addlinespace
\textbf{Managerial} & Provides ROI, illustrative cost-modelling scenario, and governance framing for healthcare AI adoption, with projected 60--70\% cost reduction and 135--185\% illustrative financial scenario (3-year model; Ch.~\ref{chap:analysis}) & CIOs, CMOs, COOs, administrators \\
\addlinespace
\textbf{Clinical} & Compresses the ASD diagnostic pathway from years to a single EEG session, improving consistency and reducing turnaround time (Ch.~\ref{chap:analysis}) & Clinicians, patients \\
\addlinespace
\textbf{Policy} & Supports FDA SaMD and EU AI Act compliance through embedded governance and fairness auditing & Regulators, policymakers \\
\end{xltabular}}
\vspace{0.5em}
\noindent\textit{Note.} RBV = Resource-Based View; DSR = design science research; illustrative financial scenario = return on investment; illustrative cost-modelling scenario = total cost of ownership; SaMD = Software as a Medical Device.

%===============================================================================
\subsection{Research Positioning}
\label{sec:research_positioning}

This research positions itself at the intersection of three rapidly maturing fields:
\begin{itemize}[leftmargin=*, itemsep=2pt]
    \item Biomedical artificial intelligence.
    \item Explainable AI.
    \item AI governance for healthcare.
\end{itemize}
\noindent No prior work integrates all three into a single ASD-focused diagnostic framework that is simultaneously empirically validated, governance-aligned, and clinically defensible. Table~\ref{tab:research_positioning} maps the four named problems (P1--P4) to RGAIG architectural components --- demonstrating systematic answer to documented failure modes, not ad hoc technique collection.

\needspace{12\baselineskip}
\noindent Table~\ref{tab:research_positioning} research positioning: mapping named problems to rgaig.

\paragraph{Research Positioning.}




{\tablestripes\tablefontsize
\needspace{10\baselineskip}
\begin{xltabular}{\textwidth}{@{} >{\centering\arraybackslash}p{1.0cm} p{3.4cm} p{4.5cm} >{\raggedright\arraybackslash}p{3.8cm} >{\raggedright\arraybackslash}p{2.0cm} @{}}
\caption[Research Positioning: Mapping Named Problems to RGAIG]{Research Positioning: Mapping Named Problems to RGAIG Architectural Components. Each named problem is paired with the architectural component that responds to it, the literature gap closed, and the contribution claim it supports, locating this research at the intersection of biomedical AI, explainable AI, and AI governance rather than within any single field alone.}\label{tab:research_positioning} \\
\toprule
\thead{P\#} & \thead{Named Problem} & \thead{RGAIG Architectural Response} & \thead{Literature Gap Closed} & \thead{Contrib.} \\
\midrule
\endfirsthead
\endhead
\bottomrule
\endlastfoot
P1 & Domain fragmentation (separate AI system per condition) & ASD-focused pipeline with multi-scale channel attention; standardised preprocessing across domains & G1: cross-domain validation absent at field's accuracy frontier & C1, C3 \\
\addlinespace
P2 & Black-box opacity (predictions without clinical rationale) & RAG-based explainability $+$ SHAP attribution; clinician-readable narrative with verified citations & G2: contextual clinical explanation absent in SHAP/LIME-only systems & C2, C5 \\
\addlinespace
P3 & Manual workflow inefficiency (45--90 min per case) & Agentic AI orchestration with human-in-the-loop override; pipeline 45--90~min $\to$ $<$30~s (Ch.~\ref{chap:analysis}) & G3: agentic mediation across data ingest $\to$ explanation absent & C4 \\
\addlinespace
P4 & Limited ASD-focused validation (single-domain models) & Standardised cross-domain evaluation across five clinical domains; ASD ensemble \asdacc{} & G4: ASD-focused architecture validation absent in the 156-study review & C1, C8 \\
\end{xltabular}}
\vspace{0.5em}{\small\noindent\textit{Note.} P1--P4 are introduced in Table~\ref{tab:problem_impact}. C1--C8 are the eight contributions claimed by this dissertation, defined in Section~\ref{sec:contributions}. G1--G4 are four of the six research gaps documented in Chapter~\ref{chap:literature}; G5 (governance) and G6 (HITL workflow) are addressed by cross-cutting RGAIG components rather than single architectural responses. RAG = retrieval-augmented generation; HITL = human-in-the-loop; SHAP = Shapley Additive Explanations.}

\noindent\textbf{Positioning argument.} The contribution is not that biomedical AI, explainable AI, or governance research is individually inadequate; it is that none alone produces a deployable answer to the four named problems. The dissertation's contribution is the \emph{integration} of all three:
\begin{itemize}[leftmargin=*, itemsep=2pt]
    \item \textbf{Empirically validated} as a diagnostic system (Chapter~\ref{chap:analysis}).
    \item \textbf{Explainable} as a clinical decision-support tool (Chapter~\ref{chap:analysis}, RAG evaluation).
    \item \textbf{Governance-aligned} as an enterprise system (Chapter~\ref{chap:discussion}, B2B/B2C verdict).
\end{itemize}
\noindent The Research Aim operationalises this integration into a SMART-validated objective.

%===============================================================================
\subsection{Research Aim}
\label{sec:research_goal}

The aim of this dissertation is to \textbf{develop and evaluate a human-centered governance, explainability, and adoption framework for responsible AI-assisted ASD-EEG screening support ecosystems}. The framework (RGAIG: Responsible Generative AI Governance) is presented as a conceptual reference architecture validated through expert review and stakeholder evaluation; it positions caregivers and families as the primary B2C beneficiaries with clinicians providing the B2B safety, escalation, and governance layer. The aim is deliberately governance-oriented rather than clinical-deployment-oriented: the dissertation contributes a governance + explainability + adoption framework with associated maturity model and decision-audit schema, not a deployable diagnostic AI product. Clinical-efficacy validation, prospective patient evaluation, and certification-grade compliance (FDA SaMD, EU MDR) are explicitly out of scope and identified as future work (Section~\ref{sec:scope_boundaries}).

\paragraph{Aim alignment with the dissertation's contribution form.} A Doctor of Business \emph{Administration} dissertation contributes a managerial decision artefact, not an algorithmic novelty. The aim above is therefore framed around \emph{governance evidence} (HITL, audit, accountability) + \emph{explainability evidence} (two-layer SHAP + RAG-grounded explanation) + \emph{adoption evidence} (caregiver trust, clinician acceptance, workflow integration) --- the three dimensions on which DBA committees evaluate governance-oriented AI dissertations. The accompanying Research Objectives and Research Questions operationalise this aim into measurable evidence requirements.

\subsection{SMART Research Aim}
\label{sec:smart_goals}

Table~\ref{tab:smart_goals} operationalises the research aim using the SMART framework (Specific, Measurable, Achievable, Relevant, Time-bound).


{\tablestripes\tablefontsize

\needspace{10\baselineskip}
\begin{xltabular}{\textwidth}{@{} >{\raggedright\arraybackslash}p{2cm} L @{}}
\caption[SMART Research Aim]{SMART Research Aim. Each SMART criterion is operationalised to ensure the research aim is specific, measurable, achievable, relevant, and time-bound, providing an auditable basis for assessing whether the study delivered on its promises.}\label{tab:smart_goals} \\
\toprule
\thead{SMART Criterion} & \thead{Operationalisation} \\
\midrule
\endfirsthead
\endhead
\bottomrule
\endlastfoot
\thead{Specific} & Design a unified RGAIG pipeline for ASD conditions \\
\addlinespace
\textbf{Measurable} & $\geq$85\% accuracy in ASD; $\geq$80\% explainability agreement; $\geq$95\% governance pass rate; projected Illustrative financial scenario (modelled only) \\
\addlinespace
\textbf{Achievable} & Uses publicly available datasets (PhysioNet, OpenNeuro, KAU; 768+ subjects), consumer-grade devices (Emotiv EPOC X/Insight), and established methods (stratified 10-fold CV, LOSO validation, bootstrap CI) \\
\addlinespace
\textbf{Relevant} & Directly addresses four named organisational problems (P1--P4) affecting 970M+ patients globally; responds to FDA SaMD and EU AI Act regulatory mandates; fills six documented research gaps (G1--G6) \\
\addlinespace
\textbf{Time-bound} & 12 months design, 6 months training, 6 months validation \\
\end{xltabular}}
\vspace{0.3em}\noindent\textit{Note.} RGAIG = Responsible Generative AI Governance; ASD = Autism Spectrum Disorder. See chapter prose for full context.

%===============================================================================
\addcontentsline{toc}{paragraph}{Must-Have: Research Problem} % [TOC-must-have-insertion]
\section{Problem Statement}
\label{sec:problem_statement}
\addcontentsline{toc}{paragraph}{Must-Have: Business Problem} % [TOC-must-have-insertion]
\paragraph{Business Problem.}\addcontentsline{toc}{paragraph}{Business Problem (Must-Have anchor)}


\noindent\textbf{Core problem.} Healthcare organizations cannot adopt AI-assisted ASD-EEG screening at scale because \emph{clinician trust} in the AI is too low to support operational integration, despite the underlying classifiers exceeding 95\% accuracy. The trust deficit is not caused by model performance but by the \emph{absence of operationalized Responsible AI (RAI) governance}---specifically, the absence of explainability, transparency, auditability, bias monitoring, human-in-the-loop controls, and rollback procedures embedded in the screening workflow. Without these governance controls, clinicians correctly perceive the system as opaque and unaccountable, and decline to integrate it into clinical decision-making. The organizational consequence is a workflow paralysis in which technically capable AI cannot deliver business value (faster screening, reduced caseload, earlier intervention) because it cannot pass the trust threshold required for adoption.

\noindent\textbf{Why this is a DBA problem (not a computer-science problem).} The technical question---``can a deep-learning model accurately classify ASD from EEG?''---has been answered (yes, $>95$\% accuracy). The unanswered question is organizational and managerial: \emph{what governance controls cause clinicians to trust and adopt the system?} This is the question this dissertation investigates. The four problems below are the symptoms through which the trust--adoption gap manifests; the conceptual framework in Figure~\ref{fig:dba_conceptual_framework} is the proposed causal mechanism (RAI governance $\to$ explainability + transparency $\to$ trust $\to$ adoption). Table~\ref{tab:problem_impact} catalogues the four symptoms and their organizational impact.

\paragraph{Four Named Problems.}




{\tablestripes\tablefontsize
\needspace{10\baselineskip}
\begin{xltabular}{\textwidth}{@{} >{\centering\arraybackslash}p{1.0cm} >{\raggedright\arraybackslash}p{4.4cm} p{4.5cm} >{\raggedright\arraybackslash}p{4.3cm} @{}}
\caption[Four Named Problems Motivating This Research]{Four Named Problems Motivating This Research. Each problem is traced from its clinical consequence to its organisational impact, demonstrating that the cost of inaction extends beyond individual misdiagnosis to systemic healthcare inefficiency.}\label{tab:problem_impact} \\
\toprule
\thead{\#} & \thead{Problem} & \thead{Consequence} & \thead{Organizational Impact} \\
\midrule
\endfirsthead
\endhead
\bottomrule
\endlastfoot
P1 & Domain fragmentation: separate AI systems per condition & Duplicated infrastructure, inconsistent governance, siloed expertise & Estimated 3--5$\times$ evaluation cost; higher vendor risk \\
\addlinespace
P2 & Black-box opacity: predictions without clinical rationale & Clinician distrust, compliance risk, patient-safety concern & Adoption failure; liability exposure under EU AI Act and FDA SaMD \\
\addlinespace
P3 & Manual workflow inefficiency: labor-intensive analytical pipelines & Bottlenecked throughput, human error, irreproducible results & 45--90 minutes per case; inconsistent inter-rater agreement \\
\addlinespace
P4 & Limited ASD-focused validation: models tested only within single domains & Uncertain generalizability, inflated confidence, narrow applicability & Weak business case for enterprise AI investment \\
\end{xltabular}}
\vspace{0.5em}
\noindent\textit{Note.} P1--P4 are traced through objectives, research questions, and contributions in Table~\ref{tab:alignment_statement}. SaMD = Software as a Medical Device.
 For the full P$\to$G$\to$O$\to$RQ$\to$H$\to$C traceability chain, see Table~\ref{tab:research_flow_alignment}.

These four problems are interdependent, with each amplifying the others:

\begin{itemize}[leftmargin=*, itemsep=3pt]
    \item \textbf{P1 $\to$ P2:} Domain fragmentation amplifies black-box opacity---each silo develops its own explainability gap, compounding clinician distrust across conditions.
    \item \textbf{P3 $\to$ P1:} Manual workflow inefficiency compounds fragmentation costs---separate human oversight is required per system, multiplying labour and governance overhead.
    \item \textbf{P4 $\to$ all:} Limited ASD-focused validation undermines the consolidation business case---without evidence that a unified approach works across domains, organisational investment in platform unification cannot be justified.
\end{itemize}


\noindent The unified framework addresses all four simultaneously.

\iffalse %% Condensed — problem table and bullet list cover this
The organizational root cause is institutional: each hospital department independently evaluates and procures AI tools without coordinated technology governance, resulting in silo lock-in, inconsistent quality standards, and duplicated infrastructure. This procurement fragmentation mirrors a broader pattern in healthcare digital transformation, where point-solution adoption outpaces platform strategy, and governance frameworks lag behind technology deployment~\cite{kelly2019healthcare}. The dissertation argues that this organizational failure is correctable through unified architecture design combined with embedded governance---a design-science intervention rather than a policy recommendation alone.
\fi

\noindent\textbf{Core organisational gap.} A healthcare organisation currently lacks a unified, explainable, governance-ready framework that integrates:
\begin{itemize}[leftmargin=*, itemsep=2pt]
    \item Patient-reported information.
    \item Behavioural and psychological assessment.
    \item EEG / IoT-based biomedical signals.
\end{itemize}
\noindent into one clinically usable decision-support workflow across multiple disease domains. This dissertation argues the core challenge is not model performance but the inability to translate high-performing AI outputs into trusted, actionable clinical decisions --- a \emph{decision trust and adoption} problem, not a technical accuracy problem.

%===============================================================================
\subsection{Why Solving This Problem Is Critical}
\label{sec:why_critical}

Solving the four named problems (P1--P4) is urgent for five reasons:

\begin{itemize}[leftmargin=*, itemsep=3pt]
    \item \textbf{Massive human impact:} Over 970 million people worldwide suffer from neurological and mental health disorders, causing 9 million deaths and 276 million disability-adjusted life years (DALYs) annually at a global cost exceeding \$2.5 trillion~\cite{who2024neuro}.
    \item \textbf{Diagnostic delays harm patients:} ASD diagnosis currently takes 4--5 years; early intervention dramatically improves outcomes but requires faster, scalable screening that fragmented AI systems cannot deliver.
    \item \textbf{Cost multiplication:} Domain fragmentation forces hospitals to procure and maintain separate AI systems per condition---an estimated 3--5$\times$ cost multiplier. A unified platform could reduce evaluation costs by 60--70\%.
    \item \textbf{Regulatory urgency:} The EU AI Act and FDA Software as a Medical Device (SaMD) regulations now mandate explainability, bias monitoring, and governance for clinical AI. Fragmented black-box systems cannot comply.
    \item \textbf{Return on investment:} A unified, governed framework projects 135--185\% illustrative financial scenario over three years through platform consolidation, reduced manual workflows, and improved diagnostic throughput (Chapter~\ref{chap:analysis}).
\end{itemize}

%===============================================================================
\subsection{Key Challenges}
\label{sec:key_challenges}

Addressing these problems simultaneously poses seven interconnected challenges:

\begin{itemize}[leftmargin=*, itemsep=3pt]
    \item \textbf{Data heterogeneity:} EEG time-series signals, pathological brain-scan images, and structured clinical questionnaires require fundamentally different preprocessing, feature extraction, and model architectures.
    \item \textbf{Domain-specific accuracy requirements:} Each disease exhibits unique biomarkers, class distributions, and future clinical-validation pathway thresholds that must be met individually before ASD-focused integration.
    \item \textbf{Explainability--performance trade-off:} High-accuracy deep learning models are inherently opaque; clinicians require interpretable, literature-grounded reasoning before trusting AI-assisted decisions.
    \item \textbf{ASD-focused generalisation:} Validating the ASD screening pipeline across diverse EEG datasets must achieve robust ASD classification accuracy---a constraint no prior study has validated this for ASD.
    \item \textbf{Governance at scale:} Embedding 525 quality modules spanning ML validation, GenAI safety, and AI governance pillars requires systematic architectural integration, not post-hoc compliance.
    \item \textbf{Clinician trust and adoption:} Adoption requires not just accuracy but transparent, contextualised explanations generated through retrieval-augmented generation (RAG) grounded in peer-reviewed literature.
    \item \textbf{Multi-jurisdictional regulatory compliance:} The framework must simultaneously satisfy FDA SaMD, EU AI Act, HIPAA, and GDPR requirements across different evaluation contexts.
\end{itemize}

%===============================================================================
\subsection{Why the Researcher Chose This Problem}
\label{sec:why_chose}

This problem was selected for three reasons:

\begin{itemize}[leftmargin=*, itemsep=3pt]
    \item \textbf{Professional experience:} The researcher's career spanning AI engineering, enterprise software architecture, and healthcare technology consulting revealed a recurring pattern: hospitals procure narrow AI tools that solve one condition but create governance debt, integration complexity, and duplicated costs across departments.
    \item \textbf{DBA relevance:} Unlike a PhD in computer science, a Doctor of Business Administration demands research that improves organisational decision-making. The four problems (P1--P4) are fundamentally management problems---resource allocation, technology governance, adoption strategy, and investment justification---that happen to involve AI.
    \item \textbf{Timing opportunity:} The simultaneous maturation of deep learning, retrieval-augmented generation (RAG), and agentic AI in 2022--2024 created a technical window that did not exist even two years earlier (Table~\ref{tab:why_now}).
\end{itemize}

%===============================================================================
\subsection{The Unsolved Problem}
\label{sec:why_unsolved}

The 156-study systematic review (Chapter~\ref{chap:literature}) found zero comparable systems. This gap persists for five structural reasons:

\begin{itemize}[leftmargin=*, itemsep=3pt]
    \item \textbf{Academic incentive misalignment:} Publication rewards favour domain-specific accuracy improvements (``we achieved 98\% on ASD'') over integrated ASD screening with governance, which requires broader engineering effort.
    \item \textbf{Interdisciplinary expertise barrier:} Solving P1--P4 simultaneously requires ML engineering, future clinical-validation pathway, regulatory compliance, and business strategy---a combination rarely found in a single research team.
    \item \textbf{Data access fragmentation:} ASD-focused datasets are siloed across institutions, ethics boards, and data-sharing agreements. Assembling 768+ subjects across ASD using the KAU EEG dataset (768 subjects) required extensive data use negotiations.
    \item \textbf{Technology timing:} RAG (feasible post-2022) and agentic AI (feasible post-2023) are prerequisites for automated explainability and orchestration. Prior researchers lacked these tools.
    \item \textbf{Governance as afterthought:} Most AI research treats governance as a post-deployment compliance exercise rather than an architectural design constraint, making embedded governance frameworks rare.
\end{itemize}

%===============================================================================
\subsection{Advantages for Clinicians and Patients}
\label{sec:advantages_clinicians_patients}

Table~\ref{tab:clinician_patient_advantages} summarises the direct benefits that the RGAIG framework provides to clinicians and patients.


\label{tab:clinician_patient_advantages}
\needspace{12\baselineskip}
\iffalse % AUTO-SUPPRESS non-ASD
\noindent Table~\ref{tab:clinician_patient_advantages_suppressed} outlines rGAIG Framework Advantages for Clinicians and Patients.

{\tablestripes\tablefontsize
\begin{xltabular}{\textwidth}{@{} >{\raggedright\arraybackslash}p{2.5cm} p{6.2cm} p{5.7cm} @{}}
\caption[RGAIG Framework Advantages for Clinicians and Patients]{RGAIG Framework Advantages for Clinicians and Patients.}\label{tab:clinician_patient_advantages_suppressed} \\
\toprule
\thead{Dimension} & \thead{Advantage for Clinicians} & \thead{Advantage for Patients} \\
\midrule
\endfirsthead
\endhead
\bottomrule
\endlastfoot
Diagnostic speed & Reduces manual assessment from 45--90~min to under 30~sec per case & Faster diagnosis; ASD screening reduced from 4--5~years to a single EEG session \\
\addlinespace
Accuracy & ASD-focused consensus accuracy of \asdacc{}~$\pm$~2.1\% across ASDs; PD achieves 100\% & Higher diagnostic confidence; reduced misdiagnosis and false negatives \\
\addlinespace
Explainability & SHAP-based feature attribution + RAG-generated clinical narratives grounded in peer-reviewed literature & Plain-language explanations of AI reasoning; informed consent for AI-assisted care \\
\addlinespace
Trust & Every prediction accompanied by confidence score, top contributing biomarkers, and literature citations & Transparency about how and why AI contributed to the diagnosis \\
\addlinespace
Multi-condition screening & Single platform screens for ASD, PD, epilepsy, stress, and cancer simultaneously & One visit, one EEG session, multiple conditions assessed---reduced hospital visits \\
\addlinespace
Governance & Built-in audit trails, bias monitoring, and regulatory compliance (FDA SaMD, EU AI Act) & Assurance that AI tools meet safety and fairness standards \\
\end{xltabular}}
\vspace{0.5em}
\noindent\textit{Note.} SHAP = SHapley Additive exPlanations; RAG = retrieval-augmented generation; SaMD = Software as a Medical Device.
\fi % AUTO-SUPPRESS non-ASD

%===============================================================================
\subsection{Clinician and Patient Benefits}
\label{sec:framework_helps_both}

The RGAIG framework bridges the gap between AI capability and clinical adoption by serving both parties simultaneously:

\begin{itemize}[leftmargin=*, itemsep=3pt]
    \item \textbf{For clinicians:} A single platform provides a unified ASD screening platform. Each prediction is accompanied by SHAP-based feature attribution, RAG-generated clinical narratives citing peer-reviewed literature, and confidence scores---enabling informed clinical judgement rather than blind trust in a black box.
    \item \textbf{For patients:} Diagnostic speed improves from months or years to a single EEG session. Patients receive plain-language explanations of AI findings, know exactly which symptoms and biomarkers contributed to the assessment, and benefit from ASD screening in a single EEG session.
    \item \textbf{Shared benefit:} The embedded governance layer (525 quality modules) ensures that every AI-assisted decision meets FDA SaMD, EU AI Act, and HIPAA/GDPR standards---protecting clinicians from liability and patients from unsafe or biased AI.
\end{itemize}

%===============================================================================
\subsection{Data Sources Overview}
\label{sec:data_sources_overview}

This research uses a dual-stream data collection strategy combining qualitative primary data with quantitative secondary data:

\begin{itemize}[leftmargin=*, itemsep=3pt]
    \item \textbf{Primary data (qualitative):} Structured questionnaires for ASD patients, parents, and caregivers, capturing behavioural observations, symptoms, and disease progression. These instruments are detailed in Chapter~\ref{chap:methodology}, Section~\ref{sec:primary_data_instruments}.
    \item \textbf{Secondary data (quantitative):} EEG signals from Emotiv devices (EPOC~X 14-channel, Insight 5-channel), wearable IoT sensors (HRV, ECG, camera), and six public benchmark datasets comprising 768+ subjects and 20,000+ medical images (58~GB). Full specifications are provided in Chapter~\ref{chap:methodology}, Section~\ref{sec:secondary_data_instruments}.
\end{itemize}


\noindent The mixed-methods analysis design, pilot study protocol, KPI framework, and study location are defined in Chapter~\ref{chap:methodology}. Table~\ref{tab:research_blueprint} in that chapter provides the complete research sequence from gap to evidence for each research question.

%===============================================================================
\subsection{Stakeholder Framing}
\label{sec:stakeholders}

{\tablestripes\tablefontsize
\needspace{10\baselineskip}
\begin{xltabular}{\textwidth}{@{} >{\raggedright\arraybackslash}p{2.6cm} p{6.2cm} p{5.5cm} @{}}
\caption[Stakeholder Analysis]{Stakeholder Analysis: Why the Problem Matters to Decision-Makers}\label{tab:stakeholder_analysis} \\
\toprule
\thead{Stakeholder} & \thead{Why the Problem Matters} & \thead{What They Need} \\
\midrule
\endfirsthead
\endhead
\bottomrule
\endlastfoot
Hospital administrators & Duplicated AI infrastructure inflates capital and operating costs; governance inconsistency creates audit risk & Platform consolidation; unified governance; clear illustrative financial-scenario model \\
\addlinespace
Clinicians & Black-box predictions undermine trust; manual workflows consume diagnostic time & Interpretable outputs; literature-grounded explanations; reduced time burden \\
\addlinespace
Patients & Delayed or inaccurate diagnosis; unexplained AI involvement in care decisions & Faster, more accurate diagnosis; transparency about AI's role \\
\addlinespace
Regulators (FDA, EU) & Fragmented tools make compliance assessment difficult; opacity violates transparency mandates & Auditable pipelines; documented decision logic; bias monitoring \\
\addlinespace
AI developers & Rebuilding infrastructure per condition is wasteful; single-domain tools limit market reach & Reusable ASD-focused architecture; modular design for domain expansion \\
\end{xltabular}}
\vspace{0.5em}
\noindent\textit{Note.} This table demonstrates the managerial relevance of the research problem, a core requirement for DBA-level research. Each stakeholder group is affected by at least two of the four named problems (Table~\ref{tab:problem_impact}).

Table~\ref{tab:as_is_to_be} contrasts the current state of healthcare AI with the RGAIG-enabled target state.

\paragraph{AS.}




{\tablestripes\tablefontsize
\needspace{10\baselineskip}
\begin{xltabular}{\textwidth}{@{} >{\raggedright\arraybackslash}p{2.0cm} p{3.6cm} >{\raggedright\arraybackslash}p{3.4cm} p{3.4cm} >{\raggedright\arraybackslash}p{2.0cm} @{}}
\caption[AS-IS vs TO-BE Comparison]{AS-IS vs TO-BE: Current State of Healthcare AI and the RGAIG-Enabled Target State}\label{tab:as_is_to_be} \\
\toprule
\thead{Area} & \thead{Current State} & \thead{Gap} & \thead{Target State} & \thead{Driver} \\
\midrule
\endfirsthead
\endhead
\bottomrule
\endlastfoot
Arch. & Siloed, single-disease AI & G1: 0/156 integrated ASD flows & Unified ASD screening workflow & Pipeline \\
\addlinespace
Explain. & Black-box outputs & G2: 11.5\% explainability coverage & Transparent SHAP~+~RAG output & XAI \\
\addlinespace
Workflow & Manual preprocessing (45--90 min/case) & G3: 0 automated pipelines & Automated screening pipeline & Orchestration \\
\addlinespace
Gov. & Not embedded in architecture & G4: 6.4\% operational governance & Built-in NIST/FDA alignment & Governance \\
\addlinespace
Value & No ROI/illustrative cost-modelling scenario clarity & G6: 5.8\% business coverage & ROI-led case (135--185\%) & Business \\
\end{xltabular}}
\vspace{0.5em}
\noindent\textit{Note.} Gap codes (G1--G6) refer to the research gaps documented in Chapter~\ref{chap:literature}. The TO-BE column describes the target state; empirical validation against these targets is reported in Chapter~\ref{chap:analysis}.



\gridbox{Problem Tree: From fragmented ASD screening to the RGAIG resolution path}{Decomposes the root problem (fragmented screening) into four contributing causes (P1--P4) and maps each to the RGAIG layer that resolves it}{4 columns: Root Problem $\to$ Causes $\to$ Sub-causes $\to$ RGAIG Resolution; 4 rows per problem branch}{Visual traceability from problem diagnosis to design response — examiner can verify every problem has a solution}

\needspace{12\baselineskip}
\begin{figure}[!htbp]
\centering
\resizebox{0.97\textwidth}{!}{%

\begin{tikzpicture}[
  every node/.style={font=\fignodefont},
  root/.style={rectangle, rounded corners, draw=ch1header, fill=ch1header, text=white, font=\bfseries\normalsize, text width=6.0cm, align=center, minimum height=1.2cm},
    problem/.style={rectangle, rounded corners, draw=ch1header, fill=ch1light, font=\small\bfseries, text width=3.0cm, align=center, minimum height=1.0cm},
    sub/.style={rectangle, rounded corners, draw=ch1header!60, fill=ch1light!50, font=\small, text width=2.8cm, align=center, minimum height=0.9cm},
    stake/.style={rectangle, rounded corners, draw=ch1header!40, fill=white, font=\small, text width=2.8cm, align=center, minimum height=0.9cm},
    resolve/.style={rectangle, rounded corners, draw=ch1header, fill=ch1header!20, font=\small\bfseries, text width=3.2cm, align=center, minimum height=1.0cm},
    arrow/.style={-{Stealth[length=4mm, width=3mm]}, very thick, ggunavy}
]
% Root
\node[root] (main) at (0,0) {Fragmented ASD\\Screening Workflow};
%
% Level 1: Problems (wider spacing to prevent overlap)
\node[problem] (p1) at (-7.0,-3.5) {P1: Workflow\\Fragmentation};
\node[problem] (p2) at (-2.3,-3.5) {P2: Explainability\\Deficit};
\node[problem] (p3) at (2.3,-3.5) {P3: Screening\\Bottleneck};
\node[problem] (p4) at (7.0,-3.5) {P4: Governance\\Gap};
%
% Arrows from root
\draw[arrow] ([xshift=-1.5cm]main.south) -- ++(0,-0.5) -| (p1.north);
\draw[arrow] ([xshift=-0.5cm]main.south) -- ++(0,-0.5) -| (p2.north);
\draw[arrow] ([xshift=0.5cm]main.south) -- ++(0,-0.5) -| (p3.north);
\draw[arrow] ([xshift=1.5cm]main.south) -- ++(0,-0.5) -| (p4.north);
%
% Level 2: Sub-problems (wider spacing)
\node[sub] (s1a) at (-7.0,-6.0) {Siloed models\\block shared logic};
\node[sub] (s1b) at (-7.0,-8.5) {No unified EEG\\pipeline};
%
\node[sub] (s2a) at (-2.3,-6.0) {Black-box output\\reduces trust};
\node[sub] (s2b) at (-2.3,-8.5) {No RAG\\explanation layer};
%
\node[sub] (s3a) at (2.3,-6.0) {Manual screening\\is slow};
\node[sub] (s3b) at (2.3,-8.5) {No end-to-end\\automation};
%
\node[sub] (s4a) at (7.0,-6.0) {No compliance\\alignment};
\node[sub] (s4b) at (7.0,-8.5) {No embedded AI\\governance};
% 
\draw[arrow] (p1.south) -- (s1a.north);
\draw[arrow] (p2.south) -- (s2a.north);
\draw[arrow] (p3.south) -- (s3a.north);
\draw[arrow] (p4.south) -- (s4a.north);
%
\draw[arrow] (s1a.south) -- (s1b.north);
\draw[arrow] (s2a.south) -- (s2b.north);
\draw[arrow] (s3a.south) -- (s3b.north);
\draw[arrow] (s4a.south) -- (s4b.north);
%
% Level 3: Stakeholders (aligned with problem columns)
\node[stake] (st1) at (-7.0,-11.0) {Clinicians:\\High workload};
\node[stake] (st2) at (-2.3,-11.0) {Patients:\\Longer wait};
\node[stake] (st3) at (2.3,-11.0) {Administrators:\\Higher cost};
\node[stake] (st4) at (7.0,-11.0) {Regulators:\\Compliance risk};
%
\draw[arrow, dashed] (s1b.south) -- (st1.north);
\draw[arrow, dashed] (s2b.south) -- (st2.north);
\draw[arrow, dashed] (s3b.south) -- (st3.north);
\draw[arrow, dashed] (s4b.south) -- (st4.north);
%
% Level 4: Resolution (wider spacing)
\node[resolve] (r1) at (-4.6,-14.6) {RGAIG: Unified\\ASD Pipeline};
\node[resolve] (r2) at (4.6,-14.6) {RGAIG: Responsible\\ASD Governance};
%
\draw[arrow, ch1header] (st1.south) -- ++(0,-0.5) -| (r1.north west);
\draw[arrow, ch1header] (st2.south) -- ++(0,-0.5) -| (r1.north east);
\draw[arrow, ch1header] (st3.south) -- ++(0,-0.5) -| (r2.north west);
\draw[arrow, ch1header] (st4.south) -- ++(0,-0.5) -| (r2.north east);
% 
\end{tikzpicture}%
}
\caption[Problem Tree: From Fragmented ASD Screening to the RGAIG]{Problem Tree: From Fragmented ASD Screening to the RGAIG Resolution Path. The tree decomposes the root problem into contributing causes and maps each cause to the corresponding RGAIG layer that resolves it, providing visual traceability from problem diagnosis to design response.}
\label{fig:problem_tree}
\end{figure}

%===============================================================================
\addcontentsline{toc}{paragraph}{Must-Have: Research Gap} % [TOC-must-have-insertion]
\subsection{Research Gap}
\label{sec:research_gap}

\noindent\textbf{Gap summary (compressed).}
\begin{itemize}[leftmargin=*, itemsep=2pt]
    \item Individual AI models exceed 95\% accuracy per condition~\cite{esteva2019guide, topol2019deep}.
    \item BUT no system in the 156-study review combines ASD classification $+$ orchestration $+$ literature-grounded explainability $+$ organisational governance.
    \item Gap persists because: (a)~academic incentives reward narrow accuracy over workflow integration, and (b)~governance integration requires rare interdisciplinary expertise (ML $+$ future clinical-validation pathway $+$ regulatory).
    \item Six gaps G1--G6 (Chapter~\ref{chap:literature}) justify the research design.
\end{itemize}

\noindent\textit{Counter-argument:} the opposite strategy --- separate highly-optimised models per disease --- has also produced strong results. Table~\ref{tab:unified_vs_domainspecific} compares the two approaches across five dimensions.

\paragraph{Unified Framework vs.Domain.}




{\tablestripes\tablefontsize
\needspace{10\baselineskip}
\begin{xltabular}{\textwidth}{@{} >{\raggedright\arraybackslash}p{2.0cm} >{\raggedright\arraybackslash}p{5.9cm} >{\raggedright\arraybackslash}p{6.6cm} @{}}
\caption[Unified Framework vs.Domain-Specific Optimisation]{Unified Framework vs.\ Domain-Specific Optimisation. This comparison directly engages the strongest counter-argument to the RGAIG approach; the reader should note that the unified strategy accepts a modest accuracy trade-off in exchange for governance, scalability, and workflow integration benefits.}\label{tab:unified_vs_domainspecific} \\
\toprule
\thead{Dimension} & \thead{Domain-Specific Optimisation} & \thead{Unified Framework (RGAIG)} \\
\midrule
\endfirsthead
\endhead
\bottomrule
\endlastfoot
Approach & Separate, highly optimised model per disease with custom preprocessing, features, and architecture & Shared infrastructure (pipeline, orchestration, governance, explainability) with domain-specific classifiers retained per condition \\
\addlinespace
Strengths & Exploits disease-unique preprocessing and hand-crafted features; single-domain ASD classifiers exceed 95\% accuracy~\cite{bosl2018eeg} & Eliminates duplicated governance, ensures consistent explainability, reduces vendor risk; single orchestration layer \\
\addlinespace
Weaknesses & Duplicated infrastructure per condition; inconsistent explainability; compounded governance burden; estimated 3--5$\times$ evaluation cost multiplier & May dilute disease-unique feature engineering; requires careful domain-adaptive preprocessing to preserve clinical-grade accuracy \\
\addlinespace
Evidence & Strong single-domain benchmarks in controlled settings & ASD-focused validation across the ASD domain (Chapter~\ref{chap:analysis}); 525-module quality taxonomy; organisational cost analysis (Chapter~\ref{chap:discussion}) \\
\addlinespace
Trade-off & Marginal accuracy gains from full domain-specific tuning & Organisational cost savings and governance consistency outweigh marginal accuracy differences \\
\end{xltabular}}
\vspace{0.5em}
\noindent\textit{Note.} ASD = autism spectrum disorder;  AUC = area under the curve; RGAIG = Responsible Generative AI Governance. The RGAIG framework retains domain-specific classifiers within the shared infrastructure, thereby preserving domain-level accuracy while eliminating infrastructure fragmentation.

\iffalse %% Condensed — table already presents the trade-off
This dissertation does not dispute the strong single-domain findings. Rather, it argues that the \textit{organisational} costs of maintaining five or more siloed systems---duplicated governance, inconsistent explainability, compounded vendor risk---outweigh the marginal accuracy gains from full domain-specific tuning. The RGAIG framework retains domain-specific classifiers within a shared infrastructure (see the ``unified'' definition above), thereby preserving domain-level accuracy while eliminating infrastructure fragmentation. Whether this trade-off holds in all clinical contexts is an empirical question addressed in Chapter~\ref{chap:analysis}.
\fi

\iffalse %% Condensed

\noindent\textbf{Why this visual.} The preceding SWOT analysis reveals fragmented failures across healthcare delivery, clinical workflows, AI trust, and research methodology. To argue that these are not independent problems but a single converging challenge, a visual funnel is needed to demonstrate the causal chain from societal burden to the architectural response proposed in this dissertation.
\fi


\noindent Figure~\ref{fig:problem_gap_funnel} illustrates how healthcare, clinical, AI, and research gaps converge to motivate RGAIG.




\noindent\textbf{Problem-to-solution funnel: how healthcare,.}

\needspace{12\baselineskip}
\begin{figure}[!htbp]
\centering
\resizebox{0.97\textwidth}{!}{%

\begin{tikzpicture}[
  every node/.style={font=\fignodefont},
  box/.style={rectangle, rounded corners=3pt, draw=ggunavy, fill=ggunavy!5,
    text width=12cm, minimum height=1cm, align=center, font=\small, inner sep=6pt},
  arr/.style={-{Stealth[length=4mm, width=3mm]}, very thick, ggunavy},
  note/.style={font=\small\itshape, text=ggunavy!70}
]
%
\node[box, fill=lightblue] (p1) {\textbf{Healthcare Problem}\\78M+ globally affected by ASD; delayed diagnosis};
% 
\node[box, fill=lightorange, below=0.6cm of p1] (p2) {\textbf{Clinical Workflow Gap}\\No continuous monitoring; episodic assessment; weak longitudinal tracking};
% 
\node[box, fill=lightteal, below=0.6cm of p2] (p3) {\textbf{AI Limitation}\\68\% clinician rejection of black-box AI; single-domain focus};
% 
\node[box, fill=lightpurple, below=0.6cm of p3] (p4) {\textbf{Research Fragmentation}\\138/156 studies are single-domain; no unified governed system};
% 
\node[box, fill=lightgreen, below=0.6cm of p4] (p5) {\textbf{$\Rightarrow$ Need for RGAIG}\\Decision-centric, explainable, governed ASD framework};
\draw[arr] (p1) -- (p2);
\draw[arr] (p2) -- (p3);
\draw[arr] (p3) -- (p4);
\draw[arr] (p4) -- (p5);
\end{tikzpicture}%
}
\caption[Problem-to-solution funnel]{Problem-to-solution funnel: how healthcare, clinical, AI, and research gaps converge to motivate RGAIG.}
\label{fig:problem_gap_funnel}
\end{figure}

\iffalse %% Condensed — figure caption covers this

\noindent\textbf{What it shows.} Figure~\ref{fig:problem_gap_funnel} presents a five-layer vertical funnel in which each layer narrows the problem scope: from the global healthcare burden (970M+ affected) through clinical workflow gaps (episodic, non-continuous diagnosis), AI limitations (68\% clinician rejection of opaque models), and research fragmentation (138 of 156 studies address only a single domain) to the architectural need for RGAIG at the base.


\noindent\textbf{So what.} The funnel reveals that no single intervention---whether better models, better data, or better governance alone---addresses the full problem. Only a framework that simultaneously spans all four layers can close the convergent gap, which is precisely the design rationale for the five-layer RGAIG architecture developed in Chapters~\ref{chap:methodology}--\ref{chap:analysis}.
\fi

\paragraph{Quantitative Evidence for.}




{\tablestripes\tablefontsize
\needspace{10\baselineskip}
\begin{xltabular}{\textwidth}{@{} >{\raggedright\arraybackslash}p{3.5cm} p{5.6cm} >{\raggedright\arraybackslash}p{2.1cm} >{\raggedright\arraybackslash}p{2.9cm} @{}}
\caption[Quantitative Evidence for Problem-to-Solution Funnel Layers]{Quantitative Evidence for Problem-to-Solution Funnel Layers. Each funnel layer is anchored to a published statistic and source, ensuring the problem statement rests on empirical evidence rather than assertion; the reader should note the progression from broad prevalence data to specific capability gaps.}\label{tab:problem_funnel_evidence} \\
\toprule
\thead{Problem Layer} & \thead{Key Statistic} & \thead{Domain} & \thead{Source} \\
\midrule
\endfirsthead
\endhead
\bottomrule
\endlastfoot
Healthcare burden & 78M+ people affected globally & ASD & WHO, 2023 \\
\addlinespace
Diagnostic delay & ASD diagnosis delayed 2--4 yrs after symptom onset & ASD & Maenner et al., 2023 \\
\addlinespace
AI trust deficit & 68\% of clinicians reject black-box AI without explainability & General & Kelly et al., 2019 \\
\addlinespace
Research fragmentation & 138 of 156 reviewed studies (88.5\%) address single domains only & Lit.\ review & Ch.~\ref{chap:literature} \\
\end{xltabular}}
\vspace{0.5em}
\noindent\textit{Note.} Statistics are drawn from the 156-study systematic review (Chapter~\ref{chap:literature}) and published epidemiological sources. ASD = autism spectrum disorder;  WHO = World Health Organization.

Figure~\ref{fig:gap_positioning} positions this thesis relative to prior work along architectural scope and governance integration dimensions.


\noindent\textbf{Research Gap Positioning: Prior Work vs.}

\needspace{12\baselineskip}
\begin{figure}[!htbp]
\centering
\resizebox{0.97\textwidth}{!}{%
\begin{tikzpicture}[every node/.style={font=\fignodefont}]
% Axes
\draw[-{Stealth[length=4mm, width=3mm]}, very thick, ggunavy] (0,0) -- (9.2,0);
\draw[-{Stealth[length=4mm, width=3mm]}, very thick, ggunavy] (0,0) -- (0,7.7);
% 
% Axis titles — centered along each axis, outside tick labels
% 
\node[font=\small\bfseries, ggunavy] at (4.5,-1.0) {Architectural Scope};
% 
\node[font=\small\bfseries, ggunavy, rotate=90] at (-1.4,3.75) {Governance Integration};
% 
% Axis tick labels
% 
\node[font=\small, ggunavy] at (2.2,-0.5) {Single-Domain};
% 
\node[font=\small, ggunavy] at (7,-0.5) {ASD};
% 
\node[font=\small, ggunavy, rotate=90] at (-0.7,1.8) {Isolated Model};
% 
\node[font=\small, ggunavy, rotate=90] at (-0.7,5.8) {Governance-Embedded};
% 
% Quadrant shading
\fill[lightcoral!30] (0,0) rectangle (4.5,3.75); % Q3: single, isolated
\fill[lightorange!30] (4.5,0) rectangle (9,3.75); % Q4: multi, isolated
\fill[lightteal!30] (0,3.75) rectangle (4.5,7.5); % Q2: single, governed
\fill[lightgreen!30] (4.5,3.75) rectangle (9,7.5); % Q1: multi, governed (gap)
% 
% Dashed quadrant lines
\draw[dashed, gray!50] (4.5,0) -- (4.5,7.5);
\draw[dashed, gray!50] (0,3.75) -- (9,3.75);
% 
% Prior work data points — positioned to avoid label overlap
\fill[gray!60] (1.2,0.7) circle (3pt); 
\node[font=\small, anchor=west] at (1.4,0.7) {EEGNet};
\fill[gray!60] (1.6,2.2) circle (3pt); 
\node[font=\small, anchor=west] at (1.8,2.2) {DeepConvNet};
\fill[gray!60] (3.2,1.0) circle (3pt); 
\node[font=\small, anchor=west] at (3.4,1.0) {ResNet};
\fill[gray!60] (2.5,2.9) circle (3pt); 
\node[font=\small, anchor=west] at (2.7,2.9) {SHAP-XAI};
\fill[gray!60] (5.8,1.2) circle (3pt); 
\node[font=\small, anchor=north west] at (6.0,1.1) {AutoML};
\fill[gray!60] (5.2,2.6) circle (3pt); 
\node[font=\small, anchor=west] at (5.4,2.6) {Transfer};
\fill[gray!60] (1.5,4.5) circle (3pt); 
\node[font=\small, anchor=west] at (1.7,4.5) {FDA AI};
\fill[gray!60] (2.8,5.2) circle (3pt); 
\node[font=\small, anchor=west] at (3.0,5.2) {EU IVDR};
% 
% Quadrant labels — placed away from data points
% 
\node[font=\small, gray, align=center] at (2.2,1.6) {Most prior\\AI studies};
% 
\node[font=\small, gray, align=center] at (7.0,1.8) {Limited multi-task\\learning};
% 
\node[font=\small, gray, align=center] at (2.2,5.8) {Regulated\\single-domain\\systems};
% 
\node[font=\small\bfseries, ggunavy, align=center] at (6.9,5.4) {THIS\\THESIS};
% 
% This thesis
\fill[ggunavy] (7.4,6.6) circle (5pt);
% 
\node[font=\small\bfseries, ggunavy, anchor=south] at (7.4,6.95) {RGAIG};
% 
% Gap arrow
\draw[-{Stealth[length=4mm, width=3mm]}, very thick, ggunavy!60, dashed] (3.5,3.0) -- (6.9,6.0);
% 
\node[font=\small, ggunavy!80, rotate=34, anchor=south] at (4.9,4.2) {Research gap};
\end{tikzpicture}%
}
\caption[Research Gap Positioning]{Research Gap Positioning: Prior Work vs.\ This Thesis Across Architectural Scope and Governance Integration}
\label{fig:gap_positioning}
\vspace{0.5em}
\noindent\textit{Note.} Prior studies are positioned from the literature review in Chapter~\ref{chap:literature}. Most existing biomedical AI studies are single-domain and lightly governed. This thesis occupies the upper-right quadrant. RGAIG = Responsible Generative AI Governance.
\end{figure}

\iffalse %% Condensed — figure caption covers this

\noindent\textbf{What it shows.} Figure~\ref{fig:gap_positioning} plots nine representative prior systems on a two-dimensional space defined by architectural scope (single-domain to ASD-focused, horizontal axis) and governance integration (isolated model to governance-embedded, vertical axis). Four quadrants emerge: the lower-left (most prior work) contains single-domain, isolated models such as EEGNet and DeepConvNet; the upper-left captures regulated single-domain systems (FDA-AI, EU-IVDR); the lower-right shows limited multi-task learning attempts; and the upper-right---occupied solely by RGAIG---represents the intersection of ASD-focused scope and embedded governance.


\noindent\textbf{So what.} The empty upper-right quadrant confirms that the research gap is not merely incremental but structural: no existing system simultaneously addresses ASD-focused unification and governance integration. This positioning validates the six research gaps (G1--G6) documented in Chapter~\ref{chap:literature} and directly motivates objectives O1 (unified pipeline), O5 (governance framework), and O6 (ASD-focused validation).
\fi


\noindent Figure~\ref{fig:killer_gap_solution} maps three failure modes---domain fragmentation, black-box opacity, and manual bottlenecks---to their RGAIG architectural responses and projected outcomes.




\noindent\textbf{Three Failure Modes in Current ASD Screening and.}

\needspace{12\baselineskip}
\begin{figure}[!htbp]
\centering
\resizebox{\textwidth}{!}{%
\begin{tikzpicture}[
  every node/.style={font=\fignodefont},
  cardbox/.style={rectangle, draw=ggunavy, rounded corners=4pt, thick,
                    text width=4.2cm, minimum height=2.8cm, align=center,
                    font=\small, inner sep=8pt},
    arr/.style={-{Stealth[length=5mm, width=4mm]}, line width=1.8pt, ggunavy},
    darr/.style={-{Stealth[length=5mm, width=4mm]}, line width=1.8pt, ggunavy!50, dashed},
    plabel/.style={font=\normalsize\bfseries\color{ggunavy}, align=center}
]
%% Column labels
\node[plabel] at (0,4.0) {Current State};
\node[plabel] at (6.0,4.0) {Consequence};
\node[plabel] at (12.0,4.0) {RGAIG Response};
\node[plabel] at (18.0,4.0) {Outcome};
%
%% Row 1: Fragmentation → Cost → Unified → Consolidation
\node[cardbox, fill=lightcoral!25] (p1) at (0,1.5) {\textbf{Domain}\\[2pt]\textbf{Fragmentation}\\[6pt]Separate ASD tools;\\3--5$\times$ cost};
%
\node[cardbox, fill=lightorange!30] (g1) at (6.0,1.5) {\textbf{Duplicated}\\[2pt]\textbf{Infrastructure}\\[6pt]Separate validation,\\governance, QA};
%
\node[cardbox, fill=ggunavy!12] (s1) at (12.0,1.5) {\textbf{Unified Pipeline}\\[2pt]\textbf{(9 Layers)}\\[6pt]Shared preprocessing,\\orchestration, governance};
%
\node[cardbox, fill=lightgreen!30] (o1) at (18.0,1.5) {\textbf{Projected 60--70\%}\\[2pt]\textbf{Cost Reduction}\\[6pt]One platform, one\\governance layer};
%
%% Row 2: Black Box → Distrust → Explainability → Trust
\node[cardbox, fill=lightcoral!25] (p2) at (0,-2.5) {\textbf{Black-Box}\\[2pt]\textbf{Opacity}\\[6pt]No reasoning trace;\\68\% clinician rejection};
%
\node[cardbox, fill=lightorange!30] (g2) at (6.0,-2.5) {\textbf{Low Clinician}\\[2pt]\textbf{Trust}\\[6pt]Reduced adoption;\\governance risk};
%
\node[cardbox, fill=ggunavy!12] (s2) at (12.0,-2.5) {\textbf{RAG + SHAP}\\[2pt]\textbf{Explainability}\\[6pt]Grounded rationales;\\auditable decisions};
%
\node[cardbox, fill=lightgreen!30] (o2) at (18.0,-2.5) {\textbf{Expert Rating}\\[2pt]\textbf{M\,=\,4.6/5 ($n$=12)}\\[6pt]Auditable trails for\\ASD deployment};
%
%% Row 3: Manual → Bottleneck → Agentic → Efficiency
\node[cardbox, fill=lightcoral!25] (p3) at (0,-6.5) {\textbf{Manual}\\[2pt]\textbf{Workflows}\\[6pt]47 manual steps;\\ASD specialist bottleneck};
%
\node[cardbox, fill=lightorange!30] (g3) at (6.0,-6.5) {\textbf{Screening}\\[2pt]\textbf{Delays}\\[6pt]Throughput limited by\\specialist availability};
%
\node[cardbox, fill=ggunavy!12] (s3) at (12.0,-6.5) {\textbf{Agentic AI}\\[2pt]\textbf{Orchestration}\\[6pt]Automated pipeline;\\clinician oversight};
%
\node[cardbox, fill=lightgreen!30] (o3) at (18.0,-6.5) {\textbf{75--85\% Time}\\[2pt]\textbf{Reduction}\\[6pt]47 steps $\to$ 3;\\clinician review retained};
%
%% Horizontal arrows (thick and clearly visible)
\foreach \i in {1,2,3} {
    \draw[arr] (p\i) -- (g\i);
    \draw[darr] (g\i) -- (s\i);
    \draw[arr] (s\i) -- (o\i);
}
%
\end{tikzpicture}}
\caption[Gap-to-Response Mapping]{Three Failure Modes in Current ASD Screening and the Corresponding RGAIG Responses}
\label{fig:killer_gap_solution}
\vspace{0.5em}
\noindent\textit{Note.} RGAIG = Responsible Generative AI Governance; RAG = Retrieval-Augmented Generation; SHAP = Shapley Additive Explanations. Dashed arrows indicate the transition from problem to architectural response. Each row traces a single failure mode from diagnosis through consequence, architectural response, and projected outcome.
\end{figure}

\iffalse %% Condensed — figure caption covers this

\noindent\textbf{What it shows.} Figure~\ref{fig:killer_gap_solution} maps three structurally distinct failure modes---domain fragmentation, black-box opacity, and manual workflow bottlenecks---through four columns: problem identification, organisational consequence, RGAIG architectural response, and projected outcome. Each row traces a single failure mode from diagnosis to remedy: fragmentation is addressed by the ASD-focused pipeline (C1), opacity by RAG-based explainability with SHAP attribution (C2), and manual bottlenecks by agentic AI orchestration (C3).


\noindent\textbf{So what.} The four-column structure demonstrates that the RGAIG framework is not an ad hoc collection of techniques but a systematic response in which each architectural component directly addresses a documented failure mode. This one-to-one mapping between problems and solutions confirms the framework's design rationale and connects to hypotheses H1 (ASD-focused accuracy), H4 (effort reduction), and H5 (explainability quality) tested in Chapter~\ref{chap:analysis}.
\fi

\noindent\textbf{B2B Sub-Problems (Hospital and Enterprise Context).} The eight failure modes in Table~\ref{tab:b2b_subproblems} characterise the fragmented state of hospital-based biomedical assessment. Each is traced to a specific design response in Table~\ref{tab:problem_design_mapping}.

\paragraph{B2B Sub.}



{\tablestripes\tablefontsize
\needspace{10\baselineskip}
\begin{xltabular}{\textwidth}{@{} >{\raggedright\arraybackslash}p{1.0cm} >{\raggedright\arraybackslash}p{4.3cm} p{9.0cm} @{}}
\caption[B2B Sub-Problems: Hospital and Enterprise Context]{B2B Sub-Problems: Hospital and Enterprise Context. The eight sub-problems characterise the fragmented state of hospital-based biomedical assessment, each documenting a specific failure mode that the RGAIG framework's architectural components are designed to address.}\label{tab:b2b_subproblems} \\
\toprule
\thead{\#} & \thead{Sub-Problem} & \thead{Description} \\
\midrule
\endfirsthead
\endhead
\bottomrule
\endlastfoot
B2B-1 & Intake fragmentation & Hospitals collect patient history, clinical forms, and clinician inputs in disconnected steps with no unified onboarding workflow. \\
\addlinespace
B2B-2 & Manual assessment burden & Clinicians spend excessive time repeating questions and coordinating assessments across departments. \\
\addlinespace
B2B-3 & Primary--secondary data disconnect & Patient-reported narratives and behavioural observations are not systematically integrated with EEG biomarker evidence. \\
\addlinespace
B2B-4 & Disease-specific tool silos & Separate AI tools exist for different conditions with no shared architecture, multiplying procurement and governance costs. \\
\addlinespace
B2B-5 & Weak explainability & AI outputs remain difficult for clinicians and caregivers to interpret, eroding trust and impeding adoption. \\
\addlinespace
B2B-6 & Governance and compliance weakness & Hospitals require auditability, privacy safeguards, and regulatory compliance, yet these are absent from most AI diagnostic tools. \\
\addlinespace
B2B-7 & Decision inconsistency & Clinical output varies across workflows and interpretations, producing unreliable decision-support. \\
\addlinespace
B2B-8 & Limited scalability & Current assessment processes do not scale across multiple disease domains or growing patient volumes. \\
\end{xltabular}}
\vspace{0.3em}\noindent\textit{Note.} RGAIG = Responsible Generative AI Governance; EEG = Electroencephalography. See chapter prose for full context.

\medskip
\noindent\textbf{B2C Sub-Problems (Consumer and Home-Monitoring Context).} Table~\ref{tab:b2c_subproblems} documents eight further sub-problems specific to patients and consumers using wearable EEG and home-monitoring technologies.

\paragraph{B2C Sub.}



{\tablestripes\tablefontsize
\needspace{10\baselineskip}
\begin{xltabular}{\textwidth}{@{} >{\raggedright\arraybackslash}p{1.0cm} >{\raggedright\arraybackslash}p{4.4cm} p{8.9cm} @{}}
\caption[B2C Sub-Problems: Consumer and Home-Monitoring Context]{B2C Sub-Problems: Consumer and Home-Monitoring Context. Eight sub-problems characterise the distinct challenges faced by patients and consumers using wearable EEG and home-monitoring technologies; together with Table~\ref{tab:b2b_subproblems} they define the full B2B + B2C failure landscape that the RGAIG framework must address.}\label{tab:b2c_subproblems} \\
\toprule
\thead{\#} & \thead{Sub-Problem} & \thead{Description} \\
\midrule
\endfirsthead
\endhead
\bottomrule
\endlastfoot
B2C-1 & Patient onboarding difficulty & Patients and parents may struggle to start and complete assessment without clinical guidance. \\
\addlinespace
B2C-2 & Low conversational support & Systems do not guide users naturally through symptom reporting and history capture. \\
\addlinespace
B2C-3 & Weak behaviour and psychological capture & Behavioural and psychological indicators are captured inconsistently or not at all in consumer settings. \\
\addlinespace
B2C-4 & Consumer device noise & Home EEG and wearable data is noisier than hospital-grade recordings, degrading classification accuracy. \\
\addlinespace
B2C-5 & Trust problem & Users may not trust or understand AI-generated output, particularly without a clinician intermediary. \\
\addlinespace
B2C-6 & Privacy concern & Home-collected health data raises heightened privacy concerns, and users are often unaware of how brain data is stored or shared. \\
\addlinespace
B2C-7 & No continuity-of-care pipeline & Home monitoring data often does not connect back to the treating clinician, creating a gap between screening and follow-up. \\
\addlinespace
B2C-8 & Self-diagnosis misuse risk & Direct-to-consumer AI output may be misinterpreted as definitive diagnosis rather than guided screening support. \\
\end{xltabular}}
\vspace{0.3em}\noindent\textit{Note.} RGAIG = Responsible Generative AI Governance; EEG = Electroencephalography. See chapter prose for full context.

\medskip
\noindent\textbf{Why continuous monitoring matters.}
Many neurological and mental health conditions are not static---symptom severity fluctuates across hours, days, and weeks, making single-session assessment inherently incomplete. The RGAIG framework therefore positions the primary value problem as B2C: caregivers and patients need earlier support, lower-friction home monitoring, plain-language explanations, and clear escalation guidance between clinical visits. The B2B layer remains essential, but as the safety and governance support system: clinicians validate results, hospitals receive high-risk escalations, and provider organisations maintain audit, consent, privacy, and regulatory controls.
\needspace{12\baselineskip}

{\tablestripes\tablefontsize
\noindent Table~\ref{tab:problem_design_mapping} maps problem and Design Mapping: B2B and B2C Problem Decomposition.

\paragraph{Problem and Design.}



\needspace{10\baselineskip}

\begin{xltabular}{\textwidth}{
>{\raggedright\arraybackslash}p{1.9cm}
>{\centering\arraybackslash}p{1.2cm}
X
X
>{\raggedright\arraybackslash}p{2.2cm}
}
\caption[Problem and Design Mapping: B2B and B2C Problem]{Problem and Design Mapping: B2B and B2C Problem Decomposition. This traceability matrix links every identified problem to its RGAIG design response and supporting evidence, ensuring that no problem is left unaddressed and no design element lacks justification.}\label{tab:problem_design_mapping} \\
\toprule
\thead{Area} & \thead{Ctx.} & \thead{Problem} & \thead{Design Response} & \thead{Evidence} \\
\midrule
\endfirsthead
\endhead
\bottomrule
\endfoot

Main Problem & Both & Fragmented intake, EEG, and governance workflows & Unified explainable and governance-aware framework & Accuracy, trust, efficiency \\
Intake Frag. & B2B & Disconnected history, forms, and clinician inputs & Unified onboarding and assessment flow & Completion, time \\
Manual Burden & B2B & Repeated clinician questioning and coordination & Conversational intake and workflow support & Effort reduction, turnaround \\
Pri--Sec Gap & B2B & Narratives are not integrated with EEG evidence & Combine questionnaire, behavioural, and EEG data & Joint display, concordance \\
Tool Silos & B2B & Separate ASD tools with no shared architecture & Unified ASD framework & Accuracy, reuse \\
Weak Explain. & B2B & AI outputs are hard to interpret & RAG + XAI explainability & Trust score, explanation quality \\
Gov.\ Weakness & B2B & Audit, privacy, and compliance are weak & Embed governance in the pipeline & Compliance, audit completion \\
Decision Incon. & B2B & Outputs vary across workflows & Standardised decision support & Consistency, agreement \\
Scalability & B2B & Processes do not scale well & Automate repeatable stages & Processing time, throughput \\
Onboarding & B2C & Patients struggle to begin assessment & Chatbot- and form-based onboarding & Completion, usability \\
Weak Convers. & B2C & Systems do not guide users effectively & Conversational symptom capture & Completeness, satisfaction \\
Behaviour & B2C & Behavioural indicators are captured inconsistently & Standardised behavioural capture & Symptom scores, themes \\
Device Noise & B2C & Home EEG is noisier than hospital EEG & Quality checks and bounded monitoring & Signal quality, accuracy \\
Trust Problem & B2C & Users distrust AI output & Plain-language explanation & Clarity, trust score \\
Privacy & B2C & Home monitoring raises privacy concerns & Privacy-aware monitoring pipeline & Consent, compliance \\
No Continuity & B2C & Home monitoring is not linked to clinical care & Hospital-to-home continuity design & Continuity, escalation \\
Self-Dx Risk & B2C & Output may be misused as diagnosis & Position as screening support only & Safe use, escalation \\
\end{xltabular}
\vspace{0.3em}\noindent\textit{Note.} RGAIG = Responsible Generative AI Governance; RAG = Retrieval-Augmented Generation; XAI = Explainable AI; ASD = Autism Spectrum Disorder. See chapter prose for full context.
}

%===============================================================================
\subsection{Problem Structure}
\label{sec:multi_layer_problem}

Table~\ref{tab:multi_layer_problem} presents eleven problem layers with distinct B2B (hospital/enterprise) and B2C (consumer/wearable) manifestations. No single intervention addresses the full problem; only an integrated architecture spanning all eleven layers can.


\noindent Table~\ref{tab:multi_layer_problem} multi-layer problem structure: b2b vs b2c context.

{\tablestripes\tablefontsize
\needspace{10\baselineskip}
\begin{xltabular}{\textwidth}{@{} >{\raggedright\arraybackslash}p{3.2cm} L L @{}}
\caption[Multi-Layer Problem Structure: B2B vs B2C Context]{Multi-Layer Problem Structure: B2B vs B2C Context. The table separates clinical, technical, data, and governance problem layers for both evaluation contexts, revealing that while the underlying AI challenges are shared, the workflow and user-interface requirements diverge significantly.}\label{tab:multi_layer_problem} \\
\toprule
\thead{Layer} & \thead{B2B (Hospital/Enterprise)} & \thead{B2C (Consumer/Wearable)} \\
\midrule
\endfirsthead
\endhead
\bottomrule
\endlastfoot
People-Centric & Clinician liability; adoption resistance & Users lack medical literacy; over/under-reliance \\
\addlinespace
Process & No AI-EHR integration protocol & No guided action/escalation workflow \\
\addlinespace
Technology & Unvalidated models; weak explainability for regulators & Noisy wearable EEG; under-calibrated outputs \\
\addlinespace
Pain/Risk & Misdiagnosis; penalties; malpractice & False alarms; anxiety; unsafe self-diagnosis \\
\addlinespace
Monitoring & Continuous EEG hard to integrate clinically & Alert fatigue; battery/real-time limits \\
\addlinespace
Data/Privacy & Complex security, audit, compliance & Opaque collection, storage, sharing \\
\addlinespace
Governance & No GenAI wearable framework & No visible safeguards \\
\addlinespace
Trust/Explain. & Need evidence-backed rationale & Need simple, accessible explanation \\
\addlinespace
Accountability & Unclear vendor/provider/org responsibility & No recourse; unclear AI limits \\
\addlinespace
Integration & EHR/EMR/regulatory integration complex & Limited app/doctor/emergency linkage \\
\addlinespace
Cross-Cutting & Formal safety, audit, compliance & Transparency, privacy, fairness \\
\end{xltabular}}
\vspace{0.5em}
\noindent\textit{Note.} Each layer represents a distinct challenge dimension. The RGAIG framework addresses all eleven layers through its five-tier architecture (L1--L5). B2B = business-to-business; B2C = business-to-consumer; EHR = electronic health record; EMR = electronic medical record.

Table~\ref{tab:layer_key_insights} distils the single most critical insight from each layer.

\paragraph{Key Insight per.}




{\tablestripes\tablefontsize
\needspace{10\baselineskip}
\begin{xltabular}{\textwidth}{@{} >{\raggedright\arraybackslash}p{2.2cm} L L @{}}
\caption[Key Insight per Problem Layer]{Key Insight per Problem Layer. Each layer distils one governing insight and its corresponding RGAIG design response, providing a concise summary of how the framework architecture was derived from the problem structure.}\label{tab:layer_key_insights} \\
\toprule
\thead{Layer} & \thead{Key Insight} & \thead{RGAIG Response} \\
\midrule
\endfirsthead
\endhead
\bottomrule
\endlastfoot
People-Centric & Clinicians resist AI due to liability; consumers cannot evaluate alerts & Confidence-calibrated outputs with RAG rationales \\
\addlinespace
Process & No standard protocol for AI insertion into EHR diagnostic chains & Five-stage clinical integration workflow (L3--L5) \\
\addlinespace
Technology & Single-domain models lack robust explainability and evidence-maturity-ready (conceptual) validation & Unified ASD pipeline with SHAP~+~RAG \\
\addlinespace
Pain/Risk & AI errors cause individual harm regardless of aggregate accuracy & Per-prediction confidence intervals with escalation \\
\addlinespace
Monitoring & Hospital monitoring episodic; consumer monitoring causes alert fatigue & Adaptive thresholds with severity-gated escalation \\
\addlinespace
Data/Privacy & Brain data sensitive under HIPAA/GDPR; neurorights emerging & End-to-end encryption with consent-tiered access \\
\addlinespace
Governance & No regulatory framework for consumer-facing ASD AI diagnostics & RGAIG five-tier governance architecture \\
\addlinespace
Trust & Clinicians need cognitive trust; consumers need affective trust & Dual-path: clinical rationale + plain-language summary \\
\addlinespace
Account. & Distributed liability across manufacturer, developer, hospital & RACI matrix with decision provenance \\
\addlinespace
Integration & Proprietary EHR standards (Epic, Cerner); inconsistent FHIR & Standards-based API with HL7 FHIR output \\
\addlinespace
Cross-Cutting & Governance dimensions require integrated system, not checklists & Unified five-tier architecture spanning all layers \\
\end{xltabular}}
\vspace{0.5em}
\noindent\textit{Note.} Each insight represents the primary barrier identified from the AS-IS analysis. RGAIG responses are mapped to specific architectural layers (L1--L5) documented in Section~\ref{sec:framework_components}.

\iffalse %% Consolidated into tab:layer_key_insights
\subsubsection{Layer 1: People-Centric Challenges}

In B2B hospital settings, clinicians---neurologists, psychiatrists, and allied health professionals---resist AI-assisted diagnostics because incorrect outputs carry direct malpractice liability. A false positive for epilepsy may trigger unnecessary medication with serious side effects; a false negative for Parkinson's may delay treatment during a critical window. This liability burden creates rational resistance to AI adoption regardless of accuracy metrics. In B2C consumer settings, the challenge inverts: users lack the medical training to evaluate whether an AI-generated insight is clinically meaningful. A consumer receiving a ``moderate risk'' alert for autism-related EEG patterns cannot determine whether this warrants an emergency visit, a routine appointment, or no action. The consequence is either panic-driven over-utilisation of healthcare resources or dangerous dismissal of genuine risk signals. For hospital administrators, this rational clinician resistance translates directly into adoption failure---and the capital invested in AI procurement, integration, and training yields no return until the trust barrier is addressed.

\subsubsection{Layer 2: Process and Workflow Gaps}

Hospital diagnostic workflows evolved over decades around human decision-making: a neurologist orders an EEG, a technician records it, a specialist interprets it, and findings are documented in the Electronic Health Record (EHR). Inserting AI into this chain creates ambiguity at every handoff---who validates the AI output before it enters the clinical record? What happens when the AI recommendation contradicts the clinician's judgment? No standardised integration protocol exists. For consumers, the workflow gap is even more fundamental: there is no workflow at all. Users receive a notification, view a dashboard, and are left to decide next steps without clinical guidance, triage logic, or escalation pathways.

\subsubsection{Layer 3: Technology and System Limitations}

On the B2B side, existing AI models for EEG analysis are typically single-domain (trained for one condition), require specialised hardware, and produce outputs that lack the explainability required for regulatory approval under the EU AI Act or FDA Software as a Medical Device (SaMD) guidelines. Under the EU AI Act, high-risk AI systems that fail transparency requirements face fines up to 3\% of global turnover---a financial exposure that makes governance integration a business necessity, not an academic nicety. On the B2C side, consumer-grade wearable EEG devices (e.g., Emotiv EPOC X with 14 channels at 128\,Hz) produce signals with lower spatial resolution and higher noise than clinical-grade 64- or 128-channel systems, yet the AI models applied to this data often lack appropriate confidence calibration for the reduced signal quality.

\subsubsection{Layer 4: Pain Points}

The risk asymmetry between B2B and B2C contexts is stark. Hospital-facing risks are institutional: misdiagnosis can trigger malpractice suits, regulatory penalties, and reputational damage that affects the entire organisation. Consumer-facing risks are personal: a false alarm for seizure activity may cause acute anxiety, relationship strain, and inappropriate self-medication, while a missed detection may provide false reassurance during a genuinely dangerous condition. Both contexts share the fundamental problem that AI error rates, even when statistically low, translate to consequential harm at the individual level.

\subsubsection{Layer 5: Continuous Monitoring Constraints}

In hospitals, continuous EEG monitoring is largely restricted to Intensive Care Unit (ICU) settings due to equipment cost, staffing requirements, and technical complexity. Outside the ICU, monitoring is episodic---a 30-minute recording that captures only a snapshot of brain activity. For consumers, wearable devices enable continuous monitoring but introduce new problems: alert fatigue from frequent non-critical notifications, battery limitations that create gaps in coverage, and the computational challenge of real-time analysis on edge devices with limited processing power.

\subsubsection{Layer 6: Data and Privacy Challenges}

EEG data is classified as sensitive health information under HIPAA~\cite{hipaa1996}, GDPR, and emerging brain data protection frameworks. In B2B contexts, hospitals must manage complex data governance: storage encryption, access controls, audit trails, cross-border transfer restrictions, and data retention policies. In B2C contexts, the challenge is transparency---consumers often have no visibility into how their brain activity data is collected, processed, stored, shared with third parties, or used for model training. The emerging field of ``neurorights'' adds further complexity, as brain data potentially reveals mental states, cognitive patterns, and neurological conditions that users may not wish to disclose.

\subsubsection{Layer 7: Governance and Compliance Gaps}

No established regulatory framework specifically addresses generative AI in consumer wearable EEG systems~\cite{fda2021aiml}. The regulatory classification is itself ambiguous: is a wearable EEG system that generates AI-powered health insights a medical device (requiring FDA 510(k) or CE marking), a wellness tool (subject to consumer protection law only), or a hybrid? This classification ambiguity paralyses both B2B procurement decisions and B2C consumer trust.

\subsubsection{Layer 8: Adoption Barriers}

Trust in AI systems operates through distinct mechanisms in B2B and B2C contexts. Clinicians develop trust through cognitive pathways: they need to understand the evidence basis, compare AI outputs with their clinical experience, and verify that the system's reasoning aligns with established medical knowledge. Consumers develop trust through affective pathways: they need to feel that the system is safe, transparent, and acting in their interest. Current systems fail both: clinical AI lacks sufficient explainability for cognitive trust, and consumer AI lacks sufficient transparency for affective trust.

\subsubsection{Layer 9: Accountability and Decision-Making}

When an AI system contributes to a diagnostic decision, accountability becomes distributed across multiple actors: the device manufacturer, the AI model developer, the data provider, the deploying hospital, and the supervising clinician. No established legal or governance framework assigns clear responsibility for AI-assisted diagnostic errors in EEG-based screening. For consumers, the accountability gap is total---if a wearable app incorrectly classifies a user's brain activity, there is no defined recourse, no complaint mechanism, and no liability assignment.

\subsubsection{Layer 10: Integration and Ecosystem Challenges}

B2B integration requires connecting AI outputs with hospital EHR/EMR systems (Epic, Cerner), clinical communication platforms, and regulatory reporting systems---each with proprietary standards, certification requirements, and data format constraints. B2C integration requires connecting with the broader consumer health ecosystem: wellness apps, primary care telehealth platforms, emergency services, and family notification systems. In both contexts, interoperability standards (HL7 FHIR, DICOM) exist but are inconsistently implemented.

\subsubsection{Layer 11: Cross-Cutting Governance Dimensions}

Across all ten preceding layers, governance dimensions---safety, fairness, transparency, accountability, robustness, privacy, and compliance---must be addressed as an integrated system rather than as independent checklists, because a system that is technically accurate but unexplainable fails the transparency test just as surely as one that is explainable but demographically biased fails the equity test. The cross-cutting nature of these requirements motivates the RGAIG framework's integrated five-tier governance architecture rather than a piecemeal approach.
\fi %% End consolidated layers

%===============================================================================
%===============================================================================
% GROUP A: B2B ANALYSIS — HOSPITAL-BASED PATIENT ASSESSMENT
%===============================================================================
%===============================================================================
\subsection{B2B Analysis: Hospital-Based Patient Assessment}
\label{sec:b2b_analysis}
\label{sec:asis_hospital}


\noindent This table lists each \textit{pain-area group} with its current state summary, so examiners can trace what was assessed and what evidence supports each row.



\noindent This table lists each \textit{pain-area group} with its current state summary, so examiners can trace what was assessed and what evidence supports each row.


\noindent This table lists each \textit{pain-area group} with its current state summary, so examiners can trace what was assessed and what evidence supports each row.



{\tablestripes\tablefontsize

\needspace{10\baselineskip}
\begin{xltabular}{\textwidth}{@{} >{\raggedright\arraybackslash}p{3.6cm} L @{}}
\caption[AS-IS Patient Assessment in Hospital Settings (B2B Context)]{AS-IS Patient Assessment in Hospital Settings (B2B Context). Decision-level summary of the four structural pain-area groups in current hospital practice. Full per-area breakdown (16 assessment-area rows) is in Appendix~A, Section~A-S5.}\label{tab:asis_hospital} \\
\toprule
\thead{Pain-Area Group} & \thead{Current State Summary} \\
\midrule
\endfirsthead
\endhead
\bottomrule
\endlastfoot
\textbf{Clinical Assessment \& Diagnostic Workflow} & Manual, clinician-driven, experience-based. EEG interpretation is expert-dependent with inconsistent detection of subtle patterns. Offline or semi-manual analysis with delayed reports and weak MRI/CT linkage. Behavioural assessment separate from EEG --- fragmented patient view. \\
\addlinespace
\textbf{Data \& Technology Infrastructure} & Noise and motion artefacts dominate raw signals; preprocessing is manual and inconsistent. EHR systems are siloed with limited interoperability. Technology usage is rule-based or assistive only --- no predictive or generative AI. No standardised explainability framework. \\
\addlinespace
\textbf{Monitoring \& Decision Support} & Continuous monitoring restricted to ICU; intermittent elsewhere with significant events missed. Risk assessment is reactive (historical data only), not predictive. Decision-making is clinician-centric with minimal decision support and high cognitive load. Alerting is manual threshold-based, producing alarm fatigue. \\
\addlinespace
\textbf{Governance, Scalability \& Trust} & Static one-time consent (no dynamic management). Accountability rests on the clinician alone --- no AI governance or audit layer. Specialist-bottlenecked workflow with reporting delays of hours to days. Volume pressure on limited neurology expertise; no capacity scaling. Trust depends on individual clinician, not on a system process. \\
\end{xltabular}}
\vspace{0.4em}
\noindent\textit{Note.} Four-group decision-level summary. The four groups correspond to six structural characteristics of the AS-IS state: human-centric but variable; reactive not predictive; fragmented systems; limited continuous monitoring; no AI governance layer; trust depends on individual clinician. Full 16-row per-area breakdown (Clinical Assessment, EEG Interpretation, Diagnostic Workflow, Continuous Monitoring, Behavioural Assessment, Risk Assessment, Decision-Making, Explainability, Data Quality, Privacy \& Consent, EHR Integration, Technology Usage, Workflow Efficiency, Alerting, Accountability, Scalability) is in Appendix~A, Section~A-S5. The RGAIG TO-BE design (Table~\ref{tab:b2b_tobe}) is engineered to resolve every group simultaneously.


%===============================================================================
\subsection{B2B Patient Assessment Workflow}
\label{sec:b2b_assessment_workflow}

Figure~\ref{fig:b2b_workflow} illustrates how RGAIG integrates into each step of the hospital workflow.
\phantomsection\label{tab:b2b_assessment_workflow}


\noindent\textbf{B2B Patient Assessment Workflow: Hospital.}

\needspace{12\baselineskip}
\begin{figure}[!htbp]
\centering
\resizebox{0.97\textwidth}{!}{%
\begin{tikzpicture}[
  every node/.style={font=\fignodefont},
  stepbox/.style={
    rectangle, rounded corners=4pt,
    draw=ggunavy, thick,
    minimum width=3.2cm, minimum height=1.4cm,
    text width=2.8cm, align=center,
    font=\small
  },
  integration/.style={
    font=\small\itshape, text=ggunavy!70!black
  },
  arrow/.style={
    -{Stealth[length=4mm, width=3mm]}, very thick, ggunavy
  },
  node distance=0.8cm and 0.8cm
]
% --- Row 1: Steps 1-4 (left to right) ---
% 
\node[stepbox] (s1) {\textbf{1.} Patient\\Enrolment};
% 
\node[stepbox, right=of s1] (s2) {\textbf{2.} EEG\\Recording};
% 
\node[stepbox, right=of s2] (s3) {\textbf{3.} API\\Submission};
% 
\node[stepbox, right=of s3] (s4) {\textbf{4.} ASD\\Screening};
% 
% --- Row 2: Steps 5-8 (right to left positioning, arrows right to left) ---
% 
\node[stepbox, below=1.2cm of s4] (s5) {\textbf{5.} Report\\Delivery};
% 
\node[stepbox, left=of s5] (s6) {\textbf{6.} Clinical\\Decision};
% 
\node[stepbox, left=of s6] (s7) {\textbf{7.} Patient\\Communication};
% 
\node[stepbox, left=of s7] (s8) {\textbf{8.} Follow-Up};
% 
% --- Integration labels below each box ---
% 
\node[integration, below=2pt of s1] {Epic/Cerner API};
% 
\node[integration, below=2pt of s2] {Emotiv BLE SDK};
% 
\node[integration, below=2pt of s3] {HL7 FHIR R4};
% 
\node[integration, below=2pt of s4] {RGAIG Cloud};
% 
\node[integration, below=2pt of s5] {HL7 FHIR};
% 
\node[integration, below=2pt of s6] {Clinical Portal};
% 
\node[integration, below=2pt of s7] {Mobile App};
% 
\node[integration, below=2pt of s8] {Auto-Schedule};
% 
% --- Row 1 arrows ---
\draw[arrow] (s1) -- (s2);
\draw[arrow] (s2) -- (s3);
\draw[arrow] (s3) -- (s4);
% 
% --- Curve from step 4 down to step 5 ---
\draw[arrow] (s4.south) -- ++(0,-0.4cm) -| (s5.north);
% 
% --- Row 2 arrows (5->6->7->8, right to left) ---
\draw[arrow] (s5) -- (s6);
\draw[arrow] (s6) -- (s7);
\draw[arrow] (s7) -- (s8);
% 
\end{tikzpicture}%
}%
\caption[B2B Patient Assessment Workflow: Hospital Integration]{B2B Patient Assessment Workflow: Hospital Integration. The flowchart traces the end-to-end hospital pathway from patient registration through EEG acquisition, AI inference, and clinician review, demonstrating that the entire cycle completes within 30 minutes compared with 45--90 minutes under traditional manual interpretation.}
\label{fig:b2b_workflow}
\end{figure}

\vspace{0.5em}
\noindent\textit{Note.} The entire workflow from EEG recording to report delivery completes within 30 minutes, compared with 45--90 minutes for traditional manual interpretation (based on literature ranges reported in Section~\ref{sec:background}). FHIR = Fast Healthcare Interoperability Resources; BLE = Bluetooth Low Energy; API = Application Programming Interface.




\noindent\textbf{B2B Hospital Assessment Workflow.}

\needspace{12\baselineskip}
\begin{figure}[!htbp]
\centering
\resizebox{0.97\textwidth}{!}{%
\begin{tikzpicture}[
  every node/.style={font=\fignodefont},
  stepbox/.style={
    rectangle,
    rounded corners=4pt,
    draw=ggunavy,
    fill=ggunavy!6,
    minimum width=2.7cm,
    minimum height=1.1cm,
    text width=2.4cm,
    align=center,
    font=\small,
    thick
  },
  arr/.style={-{Stealth[length=4mm, width=3mm]}, very thick, ggunavy}
]
% 
\node[stepbox] (s1) at (0,0) {Register Patient};
\node[stepbox] (s2) at (4.2,0) {Capture History};
\node[stepbox] (s3) at (8.4,0) {Assess Symptoms};
\node[stepbox] (s4) at (12.6,0) {Capture EEG};
%
\node[stepbox] (s5) at (2.1,-2.8) {Check Signal Quality};
\node[stepbox] (s6) at (6.3,-2.8) {Run Analysis};
\node[stepbox] (s7) at (10.5,-2.8) {Explain Output};
\node[stepbox] (s8) at (14.7,-2.8) {Clinician Review};
%
\node[stepbox] (s9) at (6.3,-5.4) {Deliver Report};
\node[stepbox] (s10) at (10.5,-5.4) {Schedule Follow-Up};
% 
\draw[arr] (s1) -- (s2);
\draw[arr] (s2) -- (s3);
\draw[arr] (s3) -- (s4);
% 
\draw[arr] (s4) -- (s5);
\draw[arr] (s5) -- (s6);
\draw[arr] (s6) -- (s7);
\draw[arr] (s7) -- (s8);
% 
\draw[arr] (s8) -- (s9);
\draw[arr] (s9) -- (s10);
% 
\end{tikzpicture}%
}
\caption[B2B Hospital Assessment Workflow]{B2B Hospital Assessment Workflow. This simplified view highlights the supervised hospital pathway from registration to clinician review, emphasising the decision points where RGAIG automation replaces manual steps.}
\label{fig:b2b_hospital_workflow_clean}
\vspace{0.5em}
\noindent\textit{Note.} The workflow shows the main supervised hospital pathway from registration to clinician review and follow-up.
\end{figure}


\iffalse %% Table version suppressed --- b2b_assessment_workflow_table

\noindent\textbf{B2B Patient Assessment Workflow: Hospital Integration.} The following table presents the detailed analysis.

\needspace{12\baselineskip}
\noindent The table below presents the b2B Patient Assessment Workflow: Hospital Integration.

\begin{table}[!htbp]
\tablestripes\tablefontsize
\tablefontsize
\caption[B2B Patient Assessment Workflow: Hospital Integration]{B2B Patient Assessment Workflow: Hospital Integration. The figure traces the six-step clinical pathway from referral to AI-assisted report delivery, illustrating how RGAIG integrates into existing hospital workflows without disrupting established care protocols.}
\begin{tabularx}{\textwidth}{@{} >{\raggedright\arraybackslash}p{0.6cm} >{\raggedright\arraybackslash}p{3.0cm} L >{\raggedright\arraybackslash}p{2.5cm} @{}}
\toprule
\thead{\#} & \thead{Step} & \thead{How It Works} & \thead{Integration} \\
\midrule
1 & Patient Enrolment & Clinician adds patient to RGAIG via hospital EHR system; informed consent captured electronically & Epic/Cerner API \\
\addlinespace
2 & EEG Recording & Nurse conducts 10-minute EEG using Emotiv EPOC X (14 channels, 128\,Hz) in clinic room; data auto-uploaded via Bluetooth & Emotiv BLE SDK \\
\addlinespace
3 & API Submission & Hospital system sends encrypted EEG data via REST API to RGAIG cloud for processing & HL7 FHIR R4 \\
\addlinespace
4 & ASD Screening & RGAIG pipeline runs ASD classification; results generated within 30 minutes & RGAIG Cloud \\
\addlinespace
5 & Report Delivery & Structured screening report with RAG-generated explanations delivered to clinician via FHIR; appears in EHR & HL7 FHIR \\
\addlinespace
6 & Clinical Decision & Clinician reviews AI screening results and RAG explanation; accepts, modifies, or overrides recommendation & Clinical Portal \\
\addlinespace
7 & Patient Communication & Doctor discusses results with patient using plain-language RAG report as communication aid & Mobile App \\
\addlinespace
8 & Follow-Up & Longitudinal tracking activated; next scan scheduled; trend monitoring and drift detection active & Auto-Schedule \\
\bottomrule
\end{tabularx}
\vspace{0.5em}
\noindent\textit{Note.} The entire workflow from EEG recording to report delivery completes within 30 minutes, compared with 45--90 minutes for traditional manual interpretation (based on literature ranges reported in Section~\ref{sec:background}). FHIR = Fast Healthcare Interoperability Resources; BLE = Bluetooth Low Energy; API = Application Programming Interface.
\end{table}
\fi




\noindent\textbf{Patient Onboarding Time Comparison.} Table~\ref{tab:onboarding_time} compares the traditional clinician-led ASD assessment (135--215 minutes) with the RGAIG-automated pathway (29--33 minutes), providing the operational evidence for Hypothesis~H3.

\gridbox{Patient Onboarding Time Reduction: Traditional vs.\ RGAIG Pathway}{Quantifies per-step efficiency gains to provide operational evidence for Hypothesis~H3}{5 columns: Step \#, Step Name, Traditional Time, RGAIG Time, Time Saved; 8 rows + total}{Cumulative 75--85\% time reduction from 2.5--3.5 hours to under 33 minutes}

\needspace{12\baselineskip}
\noindent Table~\ref{tab:onboarding_time} patient onboarding time reduction: traditional vs.rgaig.

\paragraph{Patient Onboarding Time.}




{\tablestripes\tablefontsize
\needspace{10\baselineskip}
\begin{xltabular}{\textwidth}{@{} >{\centering\arraybackslash}p{1.2cm} >{\raggedright\arraybackslash}p{3.2cm} L L >{\raggedright\arraybackslash}p{2.0cm} @{}}
\caption[Patient Onboarding Time Reduction: Traditional vs.RGAIG]{Patient Onboarding Time Reduction: Traditional vs.\ RGAIG Pathway. Step-by-step time comparisons quantify the efficiency gain at each onboarding stage; the cumulative reduction from approximately 90 minutes to under 30 minutes provides the operational evidence for Hypothesis~H3.}\label{tab:onboarding_time} \\
\toprule
\thead{\#} & \thead{Onboarding Step} & \thead{Traditional Time} & \thead{RGAIG Time} & \thead{Time Saved} \\
\midrule
\endfirsthead
\endhead
\bottomrule
\endlastfoot
1 & Registration and consent & 15--20~min (paper forms) & 5~min (digital form/app) & 67--75\% \\
\addlinespace
2 & History collection & 20--30~min (interview) & 5--7~min (chatbot/form) & 75--77\% \\
\addlinespace
3 & EEG setup and recording & 30--45~min (gel, 64+ ch.) & 7--12~min (Emotiv, 5--14 ch.) & 73--77\% \\
\addlinespace
4 & Data transfer and QC & 10--15~min (manual export) & 30~sec (auto-upload, AI QC) & 97\% \\
\addlinespace
5 & Analysis and classification & 30--60~min (clinician review) & $<$ 2~min (ASD AI pipeline) & 97\% \\
\addlinespace
6 & Report generation & 15--20~min (doctor writes) & 30~sec (RAG draft with evidence) & 97\% \\
\addlinespace
7 & Decision and communication & 10--15~min & 10~min (doctor reviews AI report) & 33\% \\
\addlinespace
8 & Follow-up scheduling & 5--10~min (manual) & 1~min (auto-scheduled) & 80--90\% \\
\midrule
& \textbf{TOTAL} & \textbf{135--215 min (2.5--3.5 hrs)} & \textbf{29--33 min} & \textbf{75--85\%} \\
\end{xltabular}}
\vspace{0.5em}
\noindent\textit{Note.} Traditional pathway assumes clinician-led ASD assessment with clinical-grade EEG. The RGAIG pathway uses consumer-grade Emotiv capture with automated screening support. Time estimates are based on published workflow benchmarks and system measurements.



\noindent\textbf{Patient onboarding flowchart: six-step pathway.}

\needspace{12\baselineskip}
\begin{figure}[!htbp]
\centering
\resizebox{0.97\textwidth}{!}{%
\begin{tikzpicture}[
  every node/.style={font=\fignodefont},
  step/.style={rectangle, draw=ggunavy, fill=#1, text width=2.4cm, minimum height=1.8cm, align=center, font=\small, rounded corners=3pt, thick},
  arr/.style={-{Stealth[length=4mm, width=3mm]}, very thick, ggunavy},
  time/.style={font=\small\color{ggunavy!70}, align=center},
% 
]
% 
%% 6 boxes at 4.2cm spacing: 0, 4.2, 8.4, 12.6, 16.8, 21.0
%% box width=2.4cm + padding ≈ 2.8cm → gap = 4.2 - 2.8 = 1.4cm ✓
\node[step=lightblue] (s1) at (0,0) {\textbf{1. Register}\\[2pt]Digital form\\Consent\\5 min};
\node[step=lightorange] (s2) at (4.2,0) {\textbf{2. Survey}\\[2pt]Chatbot/App\\25 items\\5--7 min};
\node[step=lightteal] (s3) at (8.4,0) {\textbf{3. EEG}\\[2pt]Emotiv headset\\8--10 min\\Passive};
\node[step=lightgreen] (s4) at (12.6,0) {\textbf{4. AI Screen}\\[2pt]ASD\\SHAP + RAG\\$<$ 2 min};
\node[step=lightpurple] (s5) at (16.8,0) {\textbf{5. Report}\\[2pt]Doctor reviews\\Accept/Override\\10 min};
\node[step=lightyellow] (s6) at (21.0,0) {\textbf{6. Done}\\[2pt]Report on app\\Follow-up set\\Total: 30 min};
\draw[arr] (s1) -- (s2);
\draw[arr] (s2) -- (s3);
\draw[arr] (s3) -- (s4);
\draw[arr] (s4) -- (s5);
\draw[arr] (s5) -- (s6);
\end{tikzpicture}}
\caption[Patient onboarding flowchart: six-step pathway from]{Patient onboarding flowchart: six-step pathway from registration to report delivery in under 30 minutes, screening ASD using EEG-based classification.}
\label{fig:onboarding_flow}
\end{figure}




\needspace{12\baselineskip}
\noindent Table~\ref{tab:problem_decomposition} presents problem Decomposition: Four Business and Four Technical Problems.

\paragraph{Problem Decomposition.}




{\tablestripes\tablefontsize
\needspace{10\baselineskip}
\begin{xltabular}{\textwidth}{@{} >{\centering\arraybackslash}p{1.0cm} >{\centering\arraybackslash}p{1.0cm} >{\raggedright\arraybackslash}p{3.4cm} p{6.0cm} >{\raggedright\arraybackslash}p{2.6cm} @{}}
\caption[Problem Decomposition: Four Business and Four Technical]{Problem Decomposition: Four Business and Four Technical Problems. The eight problems are classified by type and linked to the corresponding research question and hypothesis, ensuring that every RGAIG design element is traceable to an identified need.}\label{tab:problem_decomposition} \\
\toprule
\thead{Class} & \thead{ID} & \thead{Problem} & \thead{Sub-Problems} & \thead{Linked RQ/H} \\
\midrule
\endfirsthead
\endhead
\bottomrule
\endlastfoot
BP & BP1 & Market fragmentation & 5 vendors; repeated training and validation; lock-in & RQ8, H5 \\
\addlinespace
BP & BP2 & No illustrative financial-scenario model & No cost-benefit model; no outcome measure; no payer codes & RQ8, H5 \\
\addlinespace
BP & BP3 & Patient access & 0.03 neurol./100K; \$5K--15K cost; 2--4~yr wait; no home path & RQ4, RQ6, H4 \\
\addlinespace
BP & BP4 & Regulatory block & No HIPAA/GDPR/FDA alignment; no audit trail; no bias audit & RQ7, H5 \\
\midrule
TP & TP1 & Workflow fragmentation & Separate pipeline stages; no integrated ASD workflow; no shared audit logic & RQ1--RQ3, H1 \\
\addlinespace
TP & TP2 & Black-box AI & No SHAP/LIME; no lit.\ grounding; no validation; 68\% reject & RQ5, H2, H3 \\
\addlinespace
TP & TP3 & Manual bottleneck & 47 manual steps; technician-dependent; no orchestration & RQ4, RQ6, H4 \\
\addlinespace
TP & TP4 & No governance & No module QA; no GenAI test; no escalation; no drift controls & RQ7, RQ8, H5 \\
\end{xltabular}}
\vspace{0.5em}
\noindent\textit{Note.} BP = Business Problem; TP = Technical Problem. Each problem maps to specific research questions (RQ) and hypotheses (H). The RGAIG framework addresses all eight problems through its unified nine-layer architecture.

\iffalse %% Operational detail moved to Appendix A --- patient_assessment_approaches
%===============================================================================
\subsection{Patient Assessment Approaches in Hospital}
\label{sec:patient_assessment}


\noindent This table lists each \textit{assessment approach} with its application in eeg/ai context, so examiners can trace what was assessed and what evidence supports each row.



\noindent This table lists each \textit{assessment approach} with its application in eeg/ai context, so examiners can trace what was assessed and what evidence supports each row.


\noindent This table lists each \textit{assessment approach} with its application in eeg/ai context, so examiners can trace what was assessed and what evidence supports each row.



{\tablestripes\tablefontsize
\needspace{10\baselineskip}
\begin{xltabular}{\textwidth}{@{} >{\raggedright\arraybackslash}p{3.2cm} L @{}}
\caption[Patient Assessment Approaches in Hospital (Clinical + AI]{Patient Assessment Approaches in Hospital (Clinical + AI Context). Each clinical assessment approach is mapped to its EEG/AI application, showing how the RGAIG framework augments rather than replaces established clinical methods.}\label{tab:patient_assessment} \\
\toprule
\thead{Assessment Approach} & \thead{Application in EEG/AI Context} \\
\midrule
\endfirsthead
\endhead
\bottomrule
\endlastfoot
Clinical (Physical) & Direct clinician evaluation: neurological examination, reflex testing, symptom observation; EEG correlates with physical findings \\
\addlinespace
Neurological & Focused brain function evaluation (consciousness, cognition, motor response); EEG detects seizures and sleep disorders \\
\addlinespace
Diagnostic (Tests) & EEG, MRI, CT, lab tests to confirm clinical findings; AI assists in EEG pattern interpretation \\
\addlinespace
Continuous Monitoring & Real-time brain activity monitoring; wearables enable data capture outside ICU \\
\addlinespace
Behavioural/Cognitive & Memory, attention, mental state evaluation; important for interpreting EEG in psychiatric conditions \\
\addlinespace
Psychological & Emotional state, anxiety, stress assessment; influences EEG patterns and AI interpretation \\
\addlinespace
Risk Stratification & Identifies adverse event likelihood (seizures, stroke); AI predicts risk from EEG patterns \\
\addlinespace
Functional & Evaluates daily activity capability; links neurological EEG findings to functional outcomes \\
\addlinespace
Patient-Reported (PROMs) & Patient feedback on symptoms, pain, quality of life; complements AI-generated insights \\
\addlinespace
Screening & Early detection using EEG or symptom tools; used in preventive care \\
\addlinespace
Predictive (AI) & ML models predict future events (seizure prediction); GenAI provides explanations \\
\addlinespace
Explainability (XAI) & Evaluates AI output interpretability for clinicians; critical for trust and adoption \\
\addlinespace
Decision Support & AI recommendations based on EEG + history; final decisions remain with clinicians \\
\addlinespace
Multidisciplinary & Neurologists, psychologists, specialists interpret holistically \\
\addlinespace
Safety \& Quality & Ensures diagnostic processes including AI do not introduce harm \\
\addlinespace
Ethical \& Consent & Evaluates consent for data usage, especially continuous EEG monitoring \\
\addlinespace
Data Quality & Ensures EEG signals are clean, accurate, and usable; poor quality causes incorrect AI outputs \\
\addlinespace
Longitudinal & Tracks patient condition over time using repeated EEG measurements \\
\addlinespace
Personalised & Tailors diagnosis and treatment based on individual EEG patterns \\
\addlinespace
Remote/Telehealth & Uses wearable EEG + AI to assess patients outside hospital; enables continuous care \\
\end{xltabular}}
\vspace{0.5em}
\noindent\textit{Note.} The RGAIG framework supports all twenty assessment approaches through its five-tier architecture: wearable EEG capture (L1), cloud-based AI processing (L2), Agentic RAG interpretation (L3), governance validation (L4), and clinician/patient interface (L5). PROMs = Patient-Reported Outcome Measures; XAI = Explainable AI.
\fi %% End suppressed patient_assessment_approaches

\phantomsection\label{sec:patient_assessment}
\phantomsection\label{tab:patient_assessment}


\iffalse %% Operational detail moved to Appendix A --- b2b_psychological_framework
%===============================================================================
\subsection{B2B Behavioural and Psychological Assessment Framework}
\label{sec:b2b_psychological}

{\tablestripes\tablefontsize

\noindent Table~\ref{tab:b2b_psychological} presents the b2b behavioural and psychological assessment framework.

\needspace{10\baselineskip}
\begin{xltabular}{\textwidth}{@{} >{\raggedright\arraybackslash}p{3.2cm} L @{}}
\caption[B2B Behavioural and Psychological Assessment Framework]{B2B Behavioural and Psychological Assessment Framework. The table maps established behavioural and psychological theories to their concrete application within AI-governed wearable EEG assessment, grounding the RGAIG design in validated clinical science rather than ad-hoc engineering choices.}\label{tab:b2b_psychological} \\
\toprule
\thead{Theory/Approach} & \thead{Application to AI Governance in Wearable EEG} \\
\midrule
\endfirsthead
\endhead
\bottomrule
\endlastfoot
Technology Acceptance Model (TAM) & Measures perceived usefulness and ease of use of AI systems; predicts clinician adoption readiness \\
\addlinespace
UTAUT & Evaluates performance expectancy, effort expectancy, social influence, and facilitating conditions for AI adoption \\
\addlinespace
Trust Theory & Measures cognitive trust (reliability-based) and affective trust (emotional confidence) in AI outputs \\
\addlinespace
Perceived Risk Theory & Assesses stakeholder perception of misdiagnosis, legal liability, and patient harm risks \\
\addlinespace
Behavioural Decision Theory & Examines clinician decision-making under uncertainty when AI recommendations are involved \\
\addlinespace
Cognitive Load Theory & Evaluates whether AI explanations reduce or increase mental effort for clinicians \\
\addlinespace
Explainability Perception & Measures how understandable AI outputs are; impacts trust and usability \\
\addlinespace
Human--AI Interaction & Studies reliance, over-reliance, and appropriate trust calibration in AI systems \\
\addlinespace
Diffusion of Innovation & Explains AI adoption spread: early adopters vs.\ laggards in healthcare \\
\addlinespace
Organisational Change & Assesses resistance due to culture, hierarchy, and workflow disruption \\
\addlinespace
Prospect Theory & Explains why stakeholders fear losses (misdiagnosis) more than they value gains (efficiency) \\
\addlinespace
Accountability Theory & Examines responsibility perception when AI is involved in regulated decision-making \\
\addlinespace
Safety Culture Theory & Evaluates organisational safety prioritisation when adopting AI in clinical environments \\
\addlinespace
Control Theory (HITL) & Assesses need for human oversight and how control affects trust and adoption \\
\addlinespace
Cognitive Bias Theory & Identifies automation bias (over-trusting AI) and confirmation bias in clinical decisions \\
\addlinespace
Self-Efficacy Theory & Evaluates clinician confidence in using AI tools and interpreting outputs correctly \\
\addlinespace
Socio-Technical Systems & Views AI adoption as interaction between people, technology, and organisational systems \\
\end{xltabular}}
\vspace{0.5em}
\noindent\textit{Note.} These theories inform the RGAIG framework's trust, explainability, and adoption dimensions. TAM = Technology Acceptance Model; UTAUT = Unified Theory of Acceptance and Use of Technology; HITL = Human-in-the-Loop.
\fi %% End suppressed b2b_psychological_framework

\phantomsection\label{sec:b2b_psychological}
\phantomsection\label{tab:b2b_psychological}


%% SUPPRESSED — redundant with tab:asis_hospital (B2B AS-IS table)
\iffalse
%===============================================================================
\subsection{B2B Gap Analysis}
\label{sec:b2b_gap}

In addition to the current-state assessment, a systematic gap analysis reveals twenty-two specific capability deficits in hospital-based EEG assessment workflows. Table~\ref{tab:b2b_gap} maps each operational area to the identified gap, providing direct justification for the RGAIG framework's design requirements.


{\tablestripes\tablefontsize
\needspace{10\baselineskip}
\begin{xltabular}{\textwidth}{@{} >{\raggedright\arraybackslash}p{3.2cm} L L @{}}
\caption[B2B Gap Analysis: Hospital-Based EEG Assessment]{B2B Gap Analysis: Hospital-Based EEG Assessment. The AS-IS state and missing capabilities are identified for each assessment area, providing the requirements specification that the B2B TO-BE solution (Table~\ref{tab:b2b_tobe}) must satisfy.}\label{tab:b2b_gap} \\
\toprule
\thead{Area} & \thead{AS-IS State} & \thead{Gap (Missing Capability)} \\
\midrule
\endfirsthead
\endhead
\bottomrule
\endlastfoot
Clinical & Manual, experience-based evaluation & Lack of standardised, data-driven, AI-assisted assessment \\
\addlinespace
EEG Interpretation & Expert-dependent, subjective & Inconsistency; no scalable interpretation mechanism \\
\addlinespace
Diagnostic Workflow & Fragmented across systems & No unified workflow combining multiple data sources \\
\addlinespace
Monitoring & Limited to ICU or episodic & No scalable, real-time continuous monitoring \\
\addlinespace
Risk Prediction & Reactive, historical data & No predictive, real-time risk modelling \\
\addlinespace
Decision Support & High cognitive load on clinicians & No explainable decision-support system \\
\addlinespace
Explainability & Clinician-variable explanation & No standardised XAI framework for EEG \\
\addlinespace
Trust & Trust in individual clinician & No system-level trust framework \\
\addlinespace
Data Quality & Manual preprocessing & No automated validation and quality assurance \\
\addlinespace
Privacy & Static one-time consent & No dynamic, continuous consent management \\
\addlinespace
Governance & General healthcare regulations & No AI-specific governance for EEG systems \\
\addlinespace
Accountability & Clinician-only responsibility & No accountability model for hybrid decisions \\
\addlinespace
EHR Integration & Partial or siloed & No interoperability or seamless data exchange \\
\addlinespace
Technology & Rule-based or assistive AI & No generative AI with validation and governance \\
\addlinespace
Alerting & Threshold-based manual & No context-aware alert prioritisation \\
\addlinespace
Scalability & Limited by specialist count & No AI-driven model for increasing volume \\
\addlinespace
Auditability & Minimal decision tracking & No audit trail for AI-generated outputs \\
\addlinespace
Safety & Clinician vigilance only & No structured safety framework for AI diagnostics \\
\end{xltabular}}
\vspace{0.5em}
\noindent\textit{Note.} Each gap directly maps to one or more RGAIG framework components. XAI = Explainable AI; EHR = Electronic Health Record.
\fi


%===============================================================================
\iffalse %% TO-BE detail moved to Appendix A — AS-IS tables sufficient for introduction
\subsection{B2B TO-BE Solution Architecture}
\label{sec:b2b_tobe}

Table~\ref{tab:b2b_tobe} maps each B2B gap to its proposed RGAIG solution.


\noindent Each row pairs the legacy clinic step with the proposed clinician-supported workflow and the responsible RGAIG layer.
{\tablestripes\tablefontsize
\needspace{10\baselineskip}
\begin{xltabular}{\textwidth}{@{} >{\raggedright\arraybackslash}p{3.0cm} >{\raggedright\arraybackslash}p{3.2cm} L @{}}
\caption[B2B TO-BE Solution: Agentic RAG + Wearable EEG + Cloud +]{B2B TO-BE Solution: Agentic RAG + Wearable EEG + Cloud + Mobile. Each row maps an AS-IS gap to the RGAIG component that closes it, demonstrating point-by-point coverage of every deficiency identified in the gap analysis (Table~\ref{tab:b2b_gap}).}\label{tab:b2b_tobe} \\
\toprule
\thead{Area} & \thead{AS-IS Gap} & \thead{TO-BE Solution (RGAIG)} \\
\midrule
\endfirsthead
\endhead
\bottomrule
\endlastfoot
Clinical & Manual, subjective & AI-assisted decision support via Agentic RAG retrieving patient history + EEG + medical knowledge \\
\addlinespace
EEG & Expert-dependent & RAG-powered interpretation combining EEG signals + clinical guidelines + case history \\
\addlinespace
Workflow & Fragmented & Cloud-integrated orchestration unifying EEG, EHR, and AI models \\
\addlinespace
Monitoring & ICU-limited & Wearable EEG + cloud streaming for real-time anomaly detection \\
\addlinespace
Risk & Reactive & Predictive AI + RAG agents for early seizure/event detection \\
\addlinespace
Decisions & High cognitive load & Agentic AI copilots with recommendations, reasoning, and confidence levels \\
\addlinespace
Explainability & No standard mechanism & Explainable RAG (SHAP/LIME + narrative generation) \\
\addlinespace
Trust & Clinician-dependent & System-level trust with explainability, auditability, and validation \\
\addlinespace
Data Quality & No validation & Real-time data validation agents ensuring signal quality \\
\addlinespace
Privacy & Static consent & Dynamic consent via mobile + cloud governance \\
\addlinespace
Governance & No AI governance & Responsible AI layer embedded in cloud (policy, audit, risk scoring) \\
\addlinespace
Accountability & Clinician-only & Human-in-the-loop + audit trail via Agentic workflows \\
\addlinespace
Integration & Siloed systems & Cloud API integration enabling seamless data exchange \\
\addlinespace
Alerting & Threshold-based & Context-aware intelligent alerts via RAG agents \\
\addlinespace
Scalability & Specialist-limited & Cloud-native scalable AI handling large patient volumes \\
\addlinespace
Safety & Manual checks & AI safety guardrails (confidence thresholds, fallback to human) \\
\addlinespace
Auditability & No traceability & Full audit logs via Agentic workflows + RAG traceability \\
\end{xltabular}}
\vspace{0.5em}
\noindent\textit{Note.} The end-to-end B2B flow: Wearable EEG $\to$ cloud streaming $\to$ data validation agent $\to$ Agentic RAG engine (retrieve + generate) $\to$ explainability layer (SHAP/LIME) $\to$ clinician dashboard $\to$ human-in-the-loop $\to$ governance layer (audit, compliance, risk). RAG = Retrieval-Augmented Generation.
\fi %% End suppressed B2B TO-BE

\phantomsection\label{sec:b2b_tobe}
\phantomsection\label{tab:b2b_tobe}

%===============================================================================
%===============================================================================
% GROUP B: B2C ANALYSIS — CONSUMER WEARABLE EEG ASSESSMENT
%===============================================================================
%===============================================================================
\subsection{B2C Analysis: Consumer Wearable EEG Assessment}
\label{sec:b2c_analysis}
\label{sec:asis_b2c}

In the B2C context, consumers lack clinical training to evaluate AI outputs, shifting the challenge from trust-through-evidence to protection-from-harm. Table~\ref{tab:asis_b2c} documents the AS-IS assessment across twenty operational dimensions.

\needspace{12\baselineskip}
\noindent The detailed B2C AS-IS assessment table has been removed from the live chapter. The retained point is that current consumer neurotechnology offerings remain weak in explainability, governance, escalation, and clinical linkage, which is why B2C is treated here as a bounded future pathway rather than a validated core thesis context.
\phantomsection\label{tab:asis_b2c}
\phantomsection\label{tab:b2c_as_is}



\noindent\textbf{B2C Mobile Assessment Workflow.}

\needspace{12\baselineskip}
\begin{figure}[!htbp]
\centering
\resizebox{0.97\textwidth}{!}{%
\begin{tikzpicture}[
  every node/.style={font=\fignodefont},
  stepbox/.style={
    rectangle,
    rounded corners=4pt,
    draw=ggunavy,
    fill=ggunavy!6,
    minimum width=2.7cm,
    minimum height=1.1cm,
    text width=2.4cm,
    align=center,
    font=\small,
    thick
  },
  arr/.style={-{Stealth[length=4mm, width=3mm]}, very thick, ggunavy}
]
%
\node[stepbox] (s1) at (0,0) {Create Account};
\node[stepbox] (s2) at (4.2,0) {Set Profile};
\node[stepbox] (s3) at (8.4,0) {Complete Survey};
\node[stepbox] (s4) at (12.6,0) {Record EEG};
%
\node[stepbox] (s5) at (2.1,-2.8) {Upload Data};
\node[stepbox] (s6) at (6.3,-2.8) {Run Screening};
\node[stepbox] (s7) at (10.5,-2.8) {Show Explanation};
\node[stepbox] (s8) at (14.7,-2.8) {Suggest Action};
%
\node[stepbox] (s9) at (6.3,-5.4) {Share Report};
\node[stepbox] (s10) at (10.5,-5.4) {Continue Monitoring};
% 
\draw[arr] (s1) -- (s2);
\draw[arr] (s2) -- (s3);
\draw[arr] (s3) -- (s4);
% 
\draw[arr] (s4) -- (s5);
\draw[arr] (s5) -- (s6);
\draw[arr] (s6) -- (s7);
\draw[arr] (s7) -- (s8);
% 
\draw[arr] (s8) -- (s9);
\draw[arr] (s9) -- (s10);
% 
\end{tikzpicture}%
}
\caption[B2C Mobile Assessment Workflow]{B2C Mobile Assessment Workflow. The consumer-facing pathway mirrors the clinical B2B workflow but replaces hospital infrastructure with mobile and home-based steps, illustrating how RGAIG maintains diagnostic rigour outside supervised clinical settings.}
\label{fig:b2c_mobile_workflow_clean}
\vspace{0.5em}
\noindent\textit{Note.} The workflow shows the patient-facing pathway from registration to home monitoring and report sharing.
\end{figure}


%===============================================================================
\iffalse %% Operational detail moved to Appendix A --- b2c_assessment_workflow
\subsection{B2C Assessment Workflow}
\label{sec:b2c_assessment_workflow}

Table~\ref{tab:b2c_assessment_workflow} maps the B2C workflow to its B2B equivalent, preserving clinical rigour through identical AI processing while adapting the interface for consumer health literacy.


{\tablestripes\tablefontsize

\needspace{10\baselineskip}
\begin{xltabular}{\textwidth}{@{} >{\raggedright\arraybackslash}p{1.0cm} >{\raggedright\arraybackslash}p{2.6cm} p{6.8cm} >{\raggedright\arraybackslash}p{3.7cm} @{}}
\caption[B2C Patient Assessment Workflow: Consumer-Level]{B2C Patient Assessment Workflow: Consumer-Level Equivalent of B2B Hospital Assessment}\label{tab:b2c_assessment_workflow} \\
\toprule
\thead{\#} & \thead{B2C Step} & \thead{How It Works (Consumer)} & \thead{B2B Equivalent} \\
\midrule
\endfirsthead
\endhead
\bottomrule
\endlastfoot
1 & Self-Registration & User downloads app, creates account, completes health profile forms (F1--F8); dynamic consent captured & Clinician enrols patient in EHR \\
\addlinespace
2 & EEG Self-Recording & User places Emotiv EPOC X headset at home; app guides placement with real-time quality feedback; 10-minute recording & Nurse conducts EEG in clinic \\
\addlinespace
3 & Auto-Upload & EEG data auto-uploaded via Bluetooth to mobile, then encrypted and sent to RGAIG cloud via REST API & Hospital sends via FHIR API \\
\addlinespace
4 & AI Screening & Same RGAIG pipeline: ASD classification runs in cloud; results within 30 minutes & Same pipeline in hospital \\
\addlinespace
5 & Report Delivery & Plain-language health summary with colour-coded risk scores delivered via push notification and app dashboard & Structured clinical report to EHR \\
\addlinespace
6 & Self-Review & User reads AI explanation (RAG-generated, simplified); views confidence levels and educational content & Clinician reviews AI + RAG explanation \\
\addlinespace
7 & Action Decision & App recommends next steps: ``continue monitoring'' / ``book appointment'' / ``share with doctor'' / ``seek urgent care'' & Clinician decides treatment \\
\addlinespace
8 & Clinical Escalation & One-tap report sharing with linked doctor/hospital (F8); auto-formatted for clinician review & Follow-up within hospital \\
\end{xltabular}}
\vspace{0.5em}
\noindent\textit{Note.} The B2C workflow preserves clinical rigour through identical AI processing (Step 4) while adapting the user interface for consumer health literacy. The critical difference: B2C adds a mandatory clinical escalation pathway (Step 8) to prevent dangerous self-diagnosis. B2C users cannot receive a ``high risk'' classification without being prompted to connect with a healthcare professional.
\fi %% End suppressed b2c_assessment_workflow

\phantomsection\label{sec:b2c_assessment_workflow}
\phantomsection\label{tab:b2c_assessment_workflow}


\iffalse %% Operational detail moved to Appendix A --- b2c_mobile_screens
%===============================================================================
\subsection{B2C Mobile and Web App}
\label{sec:b2c_mobile_web}



\needspace{12\baselineskip}
\noindent Table~\ref{tab:b2c_mobile_screens} presents b2C Mobile/Web Application: Screens, Reports, and Data Capture.

{\tablestripes\tablefontsize

\noindent Screen-level entries clarify what the caregiver sees, what telemetry the device emits, and which guardrails fire if input quality drops.
\begin{xltabular}{\textwidth}{@{} >{\raggedright\arraybackslash}p{1.0cm} >{\raggedright\arraybackslash}p{2.7cm} p{7.9cm} >{\raggedright\arraybackslash}p{2.5cm} @{}}
\caption[B2C Mobile/Web Application]{B2C Mobile/Web Application: Screens, Reports, and Data Capture}\label{tab:b2c_mobile_screens} \\
\toprule
\thead{\#} & \thead{Screen/Feature} & \thead{What It Shows / Captures} & \thead{Interaction} \\
\midrule
\endfirsthead
\endhead
\bottomrule
\endlastfoot
S1 & Onboarding & Captures demographics (age, sex, medical history, medications); establishes baseline; obtains dynamic consent for each data use scenario & Multi-step form \\
\addlinespace
S2 & EEG Recording & Guides user through headset placement; real-time signal quality indicator (green/yellow/red); 10-minute recording with progress bar & Guided flow \\
\addlinespace
S3 & Processing Status & Shows pipeline progress (uploading $\to$ preprocessing $\to$ analysing $\to$ generating report); estimated time remaining & Progress bar \\
\addlinespace
S4 & Health Dashboard & Displays composite brain wellness score; colour-coded risk indicators per condition (green = low, amber = moderate, red = high); trend sparklines & Visual cards \\
\addlinespace
S5 & Condition Detail & Taps into specific condition (e.g., Epilepsy); shows risk score, confidence interval, plain-language explanation, and educational content & Expandable \\
\addlinespace
S6 & RAG Explanation & Displays AI-generated explanation with medical literature citations; ``Why this result?'' expandable section; confidence percentage & Read-only \\
\addlinespace
S7 & Action Guide & Context-aware next steps: ``Continue monitoring'' / ``Book GP appointment'' / ``Contact neurologist'' / ``Call emergency'' & Action buttons \\
\addlinespace
S8 & Trend History & Longitudinal view of all past scans; trend graph for each condition; change alerts (``Your stress risk increased by 12\% since last month'') & Chart + list \\
\addlinespace
S9 & Symptom Logger & Daily symptom capture: headache, sleep quality, mood, seizure events, medication adherence; feeds into AI contextualisation & Form + voice \\
\addlinespace
S10 & Privacy Centre & Data usage dashboard: what is collected, where stored, who accessed, third-party sharing status; granular consent toggles & Settings \\
\addlinespace
S11 & Doctor Sharing & One-tap PDF generation of screening report formatted for clinician review; includes EEG summary, AI findings, and symptom history & Share/email \\
\addlinespace
S12 & Emergency Alert & Triggered by high-risk classification; direct connection to emergency contacts and nearby hospitals; location sharing & Auto-trigger \\
\addlinespace
S13 & Family Dashboard & Caregiver view for dependents (children with ASD, elderly with PD); shared access with consent; alert forwarding & Multi-user \\
\addlinespace
S14 & Educational Hub & Condition-specific learning content triggered by screening results; videos, articles, and FAQs; health literacy assessment & Content feed \\
\addlinespace
S15 & Feedback Loop & Patient rates report helpfulness; reports errors or concerns; data feeds into continuous quality improvement & Rating + text \\
\end{xltabular}}
\vspace{0.5em}
\noindent\textit{Note.} Mobile-first design (iOS/Android) with responsive web fallback. All screens comply with WCAG 2.1 AA accessibility standards. Behavioural safety design: no alarming language, calm colour palette, mandatory confidence disclosure, and clinician escalation at every decision point. Data capture feeds into personalisation engine for improved future recommendations.
\fi %% End suppressed b2c_mobile_screens

\phantomsection\label{sec:b2c_mobile_web}
\phantomsection\label{tab:b2c_mobile_screens}

\iffalse %% Operational detail moved to Appendix A --- b2c_user_forms
%===============================================================================
\subsection{B2C Registration and Profile}
\label{sec:b2c_forms}


\needspace{12\baselineskip}
\iffalse % AUTO-SUPPRESS non-ASD


\noindent\textbf{B2C User Fill-Up Forms: Mobile Application Data Capture.} Table~\ref{tab:b2c_user_forms} below presents the detailed analysis.

\needspace{12\baselineskip}
\noindent Form-level entries name the validation rules, encryption requirements, and downstream consumers of each input field.
{\tablestripes\tablefontsize

\begin{xltabular}{\textwidth}{@{} >{\raggedright\arraybackslash}p{1.0cm} >{\raggedright\arraybackslash}p{3.4cm} p{8.3cm} >{\raggedright\arraybackslash}p{1.5cm} @{}}
\caption[B2C User Fill-Up Forms]{B2C User Fill-Up Forms: Mobile Application Data Capture}\label{tab:b2c_user_forms} \\
\toprule
\thead{\#} & \thead{Form / Screen} & \thead{Fields Captured} & \thead{When} \\
\midrule
\endfirsthead
\endhead
\bottomrule
\endlastfoot
F1 & Registration & Full name, email, phone, date of birth, sex, country, preferred language & Once \\
\addlinespace
F2 & Medical History & Pre-existing conditions (checkbox: epilepsy, PD, ASD, depression, anxiety, ADHD, sleep disorders, none); year of diagnosis & Once \\
\addlinespace
F3 & Medication Log & Current medications (name, dose, frequency); supplements; allergies & Once + updates \\
\addlinespace
F4 & Family History & Neurological conditions in family (parent, sibling, grandparent); checkbox list & Once \\
\addlinespace
F5 & Lifestyle Profile & Sleep hours, exercise frequency, caffeine/alcohol intake, stress level (1--10), occupation type & Once + updates \\
\addlinespace
F6 & Consent Form & Dynamic consent: agree/decline for each data use (AI analysis, research, third-party sharing, emergency contact) & Once + review \\
\addlinespace
F7 & Emergency Contacts & Primary contact (name, phone, relationship); secondary contact; preferred hospital/clinic; GP details & Once + updates \\
\addlinespace
F8 & Doctor/Hospital Link & Preferred doctor name, clinic/hospital, speciality, contact details; option to auto-share reports via email or FHIR & Once + updates \\
\addlinespace
F9 & Daily Symptom Log & Mood (1--5), sleep quality (1--5), headache (Y/N), seizure events (count), tremor severity (1--5), anxiety (1--5) & Daily \\
\addlinespace
F10 & Pre-EEG Checklist & Caffeine last 4h (Y/N), medication taken (Y/N), sleep last night (hours), current stress (1--10), hair clean/dry (Y/N) & Each session \\
\addlinespace
F11 & Post-Session Feedback & Comfort rating (1--5), headset fit quality, session completion status, any issues noted & Each session \\
\addlinespace
F12 & Goal Setting & Health goals: ``Understand my brain health'' / ``Monitor condition'' / ``Track treatment'' / ``Research participation'' & Once + updates \\
\end{xltabular}}
\vspace{0.5em}
\noindent\textit{Note.} Forms F1--F8 are completed during onboarding (15--20 minutes). Form F9 is a 2-minute daily check-in. Forms F10--F11 bracket each EEG session. All form data is encrypted at rest (AES-256) and in transit (TLS 1.3). Users can update F3, F5, F7, F8, and F12 at any time via Settings. GP = General Practitioner; FHIR = Fast Healthcare Interoperability Resources.
\fi % AUTO-SUPPRESS non-ASD
\fi %% End suppressed b2c_user_forms

\phantomsection\label{sec:b2c_forms}
\phantomsection\label{tab:b2c_user_forms}

%% SUPPRESSED — redundant with tab:asis_b2c (B2C AS-IS table)
\iffalse
%===============================================================================
\subsection{B2C Gap Analysis}
\label{sec:b2c_gap}

Consumer-facing wearable EEG systems present distinct gaps related to usability, trust, explainability, and behavioural impact. Table~\ref{tab:b2c_gap} documents twenty-three capability deficits in consumer EEG contexts, complementing the B2B analysis to demonstrate full ecosystem understanding.


{\tablestripes\tablefontsize

\needspace{10\baselineskip}
\begin{xltabular}{\textwidth}{@{} >{\raggedright\arraybackslash}p{3.2cm} L L @{}}
\caption[B2C Gap Analysis: Consumer Wearable EEG Context]{B2C Gap Analysis: Consumer Wearable EEG Context. The table identifies capability gaps specific to unsupervised consumer settings---lower signal quality, limited health literacy, and absence of real-time clinical oversight---that the B2C TO-BE design (Table~\ref{tab:b2c_tobe}) must address.}\label{tab:b2c_gap} \\
\toprule
\thead{Area} & \thead{AS-IS (Consumer Reality)} & \thead{Gap (Missing Capability)} \\
\midrule
\endfirsthead
\endhead
\bottomrule
\endlastfoot
User Understanding & Simplified or unclear insights & No meaningful, context-aware explanations \\
\addlinespace
Health Literacy & Limited medical knowledge & No adaptive education layer for AI insights \\
\addlinespace
AI Interpretation & Outputs without validation & Risk of hallucination and misleading outputs \\
\addlinespace
Trust & Based on brand, not reliability & No structured trust-building mechanism \\
\addlinespace
Decision-Making & Overreaction or ignoring alerts & No guided decision-support framework \\
\addlinespace
Monitoring & Data collected but not actionable & No intelligent interpretation or prioritisation \\
\addlinespace
Alerting & Frequent, generic notifications & No personalised, context-aware alerts \\
\addlinespace
Privacy & Unaware of data usage & No transparent policies or user control \\
\addlinespace
Consent & One-time during device setup & No dynamic, ongoing consent management \\
\addlinespace
Personalisation & One-size-fits-all insights & No adaptive AI for individual patterns \\
\addlinespace
Healthcare Link & Loosely connected to providers & No escalation pathway to clinicians \\
\addlinespace
Explainability & Too simple or too technical & No user-friendly explainability layer \\
\addlinespace
Risk Awareness & Unaware of AI limitations & No risk communication framework \\
\addlinespace
Accountability & No clarity on responsibility & No accountability framework for recommendations \\
\addlinespace
Safety & Minimal safeguards & No built-in validation or escalation \\
\addlinespace
Behavioural Impact & Anxiety from AI outputs & No behavioural design to prevent misuse \\
\addlinespace
Data Quality & Noisy, inconsistent signals & No real-time validation or feedback \\
\addlinespace
Engagement & Passive device usage & No interactive engagement loops \\
\addlinespace
Ecosystem & Limited app integration & No unified consumer health ecosystem \\
\addlinespace
Regulatory Awareness & Users unaware of device class & No clear diagnostic vs.\ advisory labelling \\
\addlinespace
Auditability & Cannot trace AI decisions & No transparent audit trail \\
\addlinespace
Adoption & Curiosity-driven, low trust & No sustained trust-driven engagement model \\
\end{xltabular}}
\vspace{0.5em}
\noindent\textit{Note.} The B2C gap analysis reveals that consumer challenges are fundamentally different from B2B: they centre on usability, behavioural safety, health literacy, and simplified governance rather than clinical workflow integration. Both B2B (Table~\ref{tab:b2b_gap}) and B2C gaps inform the RGAIG framework design.
\fi


%===============================================================================
\iffalse %% TO-BE detail moved to Appendix A
\subsection{B2C TO-BE Solution Architecture}
\label{sec:b2c_tobe}

Table~\ref{tab:b2c_tobe} maps B2C consumer gaps to their corresponding RGAIG solutions.


\noindent Each row pairs the legacy at-home step with the proposed clinician-supported equivalent and the governance layer that audits it.
{\tablestripes\tablefontsize
\needspace{10\baselineskip}
\begin{xltabular}{\textwidth}{@{} >{\raggedright\arraybackslash}p{3.0cm} >{\raggedright\arraybackslash}p{3.2cm} L @{}}
\caption[B2C TO-BE Solution: Agentic RAG + Wearable EEG + Cloud +]{B2C TO-BE Solution: Agentic RAG + Wearable EEG + Cloud + Mobile. Each consumer-context gap is paired with its RGAIG resolution, mirroring the B2B TO-BE structure (Table~\ref{tab:b2b_tobe}) and confirming that the framework scales from hospital to home without sacrificing governance controls.}\label{tab:b2c_tobe} \\
\toprule
\thead{Area} & \thead{AS-IS Gap} & \thead{TO-BE Solution (RGAIG)} \\
\midrule
\endfirsthead
\endhead
\bottomrule
\endlastfoot
Understanding & Cannot interpret EEG & Agentic RAG explanation engine: EEG patterns $\to$ personalised, context-aware insights \\
\addlinespace
Literacy & Limited medical knowledge & Adaptive AI education layer (mobile) explaining conditions, risks, and next steps \\
\addlinespace
AI Output & Hallucination risk & Grounded RAG generation using verified medical knowledge + safety validation \\
\addlinespace
Trust & Brand-based only & Trust layer with explainability + confidence scores + transparency logs \\
\addlinespace
Decisions & Users unsure of action & AI-guided assistant: monitor / consult doctor / emergency action \\
\addlinespace
Monitoring & Not actionable & Real-time RAG insights + trend analysis via mobile dashboard \\
\addlinespace
Alerting & Generic alerts & Context-aware intelligent alerts (severity-based, personalised) \\
\addlinespace
Privacy & Unaware of usage & Mobile-first privacy dashboard showing data usage and storage \\
\addlinespace
Consent & One-time only & Dynamic consent system (opt-in/out per usage scenario) \\
\addlinespace
Personalisation & One-size-fits-all & Personalised AI trained on individual EEG patterns + behaviour \\
\addlinespace
Healthcare Link & No escalation & Seamless flow: Mobile $\to$ Cloud $\to$ Clinician (B2B integration) \\
\addlinespace
Explainability & Too simple/complex & Dual-layer: simple for users + detailed for clinicians \\
\addlinespace
Risk Awareness & Unaware of limits & Risk communication (confidence \%, uncertainty, warnings) \\
\addlinespace
Safety & No safeguards & AI safety guardrails: threshold triggers + human review fallback \\
\addlinespace
Engagement & Passive usage & Interactive AI assistant (chatbot + voice) for continuous engagement \\
\addlinespace
Ecosystem & Limited integration & Unified health ecosystem (wearables, apps, emergency services) \\
\addlinespace
Adoption & Low trust & Trust-driven engagement model (personalisation + explainability + outcomes) \\
\end{xltabular}}
\vspace{0.5em}
\noindent\textit{Note.} The B2C flow: Wearable EEG $\to$ edge processing $\to$ cloud streaming $\to$ Agentic RAG (retrieve + generate) $\to$ validation agent $\to$ explainability layer $\to$ mobile app $\to$ user action (monitor/consult/emergency) $\to$ escalation to B2B system if needed. The B2C architecture complements B2B (Table~\ref{tab:b2b_tobe}), together providing end-to-end ecosystem coverage.
\fi %% End suppressed B2C TO-BE

\phantomsection\label{sec:b2c_tobe}
\phantomsection\label{tab:b2c_tobe}

%===============================================================================
%===============================================================================
% GROUP C: CROSS-CUTTING INFRASTRUCTURE — DATA PIPELINES AND TECHNOLOGY
%===============================================================================
%===============================================================================
\subsection{Cross-Cutting Infrastructure}
\label{sec:cross_cutting_infrastructure}
\label{sec:report_generation}

\iffalse %% Report types — operational detail moved to Appendix A
The same data pipelines, storage layer, and governance controls operate identically across B2B and B2C evaluation contexts. Table~\ref{tab:report_types} maps all report types across clinician and consumer audiences.


{\tablestripes\tablefontsize
\needspace{10\baselineskip}
\begin{xltabular}{\textwidth}{@{} >{\raggedright\arraybackslash}p{1.0cm} >{\raggedright\arraybackslash}p{3.4cm} p{8.3cm} >{\raggedright\arraybackslash}p{1.5cm} @{}}
\caption[RGAIG Report Types: Clinician (B2B) and Patient (B2C)]{RGAIG Report Types: Clinician (B2B) and Patient (B2C) Reports. The report taxonomy distinguishes clinical-grade outputs from patient-facing summaries, ensuring that each stakeholder receives information at the appropriate level of detail and health literacy.}\label{tab:report_types} \\
\toprule
\thead{\#} & \thead{Report Type} & \thead{Content and Purpose} & \thead{Delivery} \\
\midrule
\endfirsthead
\endhead
\bottomrule
\endlastfoot
\multicolumn{4}{l}{\textit{\textbf{Clinician Reports (B2B)}}} \\
\addlinespace
R1 & Screening Summary & Five-domain risk scores with confidence intervals; classification (low/moderate/high risk) per condition & EHR via FHIR \\
\addlinespace
R2 & EEG Analysis & Detailed spectral analysis, time-frequency maps, artefact rejection report, signal quality metrics & PDF + EHR \\
\addlinespace
R3 & RAG Explanation & Literature-backed clinical reasoning for each risk score; top-5 retrieved evidence passages with citations & EHR + Portal \\
\addlinespace
R4 & Feature Attribution & SHAP/LIME visualisation showing which EEG features drove each classification decision & Portal \\
\addlinespace
R5 & Comparison Report & Longitudinal comparison with patient's previous scans; trend analysis and change detection & EHR \\
\addlinespace
R6 & Governance Audit & Complete audit trail: data provenance, model version, gate decisions, confidence scores, override history & Audit System \\
\addlinespace
R7 & Quality Report & Data quality metrics, electrode impedance, artefact percentage, signal-to-noise ratio & Portal \\
\addlinespace
R8 & Cohort Analytics & Aggregate statistics across patient population; departmental screening outcomes and KPIs & Dashboard \\
\midrule
\multicolumn{4}{l}{\textit{\textbf{Patient Reports (B2C)}}} \\
\addlinespace
R9 & Health Summary & Plain-language overview of brain health screening results; colour-coded risk indicators & Mobile App \\
\addlinespace
R10 & Condition Insight & Simplified explanation of what each condition means and why the AI flagged or cleared it & Mobile App \\
\addlinespace
R11 & Action Guide & Next-step recommendations: ``continue monitoring'' / ``schedule appointment'' / ``seek urgent care'' & Push + App \\
\addlinespace
R12 & Trend Report & Longitudinal brain health trends visualised over weeks/months; progress tracking & App Dashboard \\
\addlinespace
R13 & Educational & Contextual health education triggered by screening results; condition-specific learning content & In-App \\
\addlinespace
R14 & Privacy Report & Summary of data collected, storage duration, access log, and third-party sharing status & App Settings \\
\addlinespace
R15 & Clinician Referral & Pre-formatted report the user can share with their doctor; includes EEG summary and AI findings & PDF/Email \\
\addlinespace
R16 & Wellness Score & Composite brain wellness score with historical trend; gamified engagement element & App Home \\
\end{xltabular}}
\vspace{0.5em}
\noindent\textit{Note.} All reports include version tracking, timestamp, and model identification. Clinician reports (R1--R8) are delivered via HL7 FHIR integration and clinical portal. Patient reports (R9--R16) are delivered via mobile app with push notifications. RAG = Retrieval-Augmented Generation; SHAP = SHapley Additive exPlanations; LIME = Local Interpretable Model-agnostic Explanations.
\fi %% End suppressed report_types

\phantomsection\label{tab:report_types}

% --- Continuous Monitoring Flowchart (replaces suppressed table) ---



\noindent\textbf{Continuous Monitoring Flow: Emotiv Device to.}

\needspace{12\baselineskip}
\begin{figure}[!htbp]
\centering
\resizebox{0.97\textwidth}{!}{%
\begin{tikzpicture}[
  every node/.style={font=\fignodefont},
  node distance=0.9cm and 2.8cm,
    stage/.style={
        rectangle, rounded corners=3pt,
        draw=ggunavy, thick,
        fill=ggunavy!5, text=black,
        minimum width=3.2cm, minimum height=1.0cm,
        align=center, font=\fignodefont},
    critical/.style={stage, very thick, fill=ggunavy!12, font=\fignodefont},
    bg/.style={stage, draw=ggunavy, fill=white, dashed, font=\fignodefont},
    arr/.style={-{Stealth[length=4mm, width=3mm]}, very thick, ggunavy},
    bgarr/.style={-{Stealth[length=3mm, width=2mm]}, thick, ggunavy!60, dashed},
    lat/.style={font=\small\itshape, text=ggunavy!80},
]
% --- Critical path (vertical, left column) ---
% 
\node[critical] (m1) {\textbf{M1:} Device Capture};
% 
\node[critical, below=of m1] (m2) {\textbf{M2:} Edge Processing};
% 
\node[critical, below=of m2] (m3) {\textbf{M3:} Cloud Upload};
% 
\node[critical, below=of m3] (m4) {\textbf{M4:} Real-Time Analysis};
% 
\node[critical, below=of m4] (m5) {\textbf{M5:} Alert Generation};
% 
\node[critical, below=of m5] (m6) {\textbf{M6:} Push Notification};
% 
% --- Critical path arrows with latency labels ---
\draw[arr] (m1) -- node[right, lat] {Real-time} (m2);
\draw[arr] (m2) -- node[right, lat] {$<$1\,s} (m3);
\draw[arr] (m3) -- node[right, lat] {$<$2\,s} (m4);
\draw[arr] (m4) -- node[right, lat] {$<$5\,s} (m5);
\draw[arr] (m5) -- node[right, lat] {$<$3\,s} (m6);
% 
% --- Latency on M6 exit ---
% 
\node[lat, below=0.6cm of m6] (end) {$<$2\,s to device};
% 
% --- Background/parallel tasks (right column) ---
% 
\node[bg, right=of m4] (m7) {\textbf{M7:} Dashboard Update};
% 
\node[bg, right=of m5] (m8) {\textbf{M8:} Data Persistence};
% 
\node[bg, right=of m6] (m9) {\textbf{M9:} Trend Computation};
% 
\node[bg, right=of m2] (m10) {\textbf{M10:} Battery/Health};
% 
% --- Background arrows ---
\draw[bgarr] (m4.east) -- node[above, lat] {$<$5\,s} (m7.west);
\draw[bgarr] (m5.east) -- node[above, lat] {Background} (m8.west);
\draw[bgarr] (m6.east) -- node[above, lat] {Scheduled} (m9.west);
\draw[bgarr] (m2.east) -- node[above, lat] {Continuous} (m10.west);
% 
% --- Total latency annotation ---
\draw[decorate, decoration={brace, amplitude=6pt, mirror}, ggunavy, thick]
    ([xshift=-0.7cm]m1.north west) -- ([xshift=-0.7cm]m6.south west)
    node[midway, left=8pt, align=center, font=\small\bfseries, text=ggunavy]
    {Critical Path\\$<$15 seconds};
% 
% --- Legend ---
% 
\node[below=1.0cm of end, xshift=1.5cm, anchor=north, font=\small, align=left] (legend) {%
    \tikz[baseline=-0.5ex]\draw[-{Stealth[length=4mm, width=3mm]}, very thick, ggunavy] (0,0) -- (0.6,0);~Critical path \quad
    \tikz[baseline=-0.5ex]\draw[-{Stealth[length=3mm, width=2mm]}, thick, ggunavy!60, dashed] (0,0) -- (0.6,0);~Background/parallel task%
};
\end{tikzpicture}%
}% end resizebox
\caption[Continuous Monitoring Flow: Emotiv to Cloud Alert]{Continuous Monitoring Flow: Emotiv Device $\to$ Mobile $\to$ Cloud $\to$ Alert.
Stages M1--M6 form the critical real-time path with end-to-end latency below 15 seconds.
Stages M7--M10 execute as parallel or background tasks.}
\label{fig:ch1_continuous_monitoring}
\end{figure}

\iffalse %% Operational detail moved to Appendix A --- continuous_monitoring
%===============================================================================
\subsection{Continuous Monitoring Pipeline}
\label{sec:continuous_monitoring_flow}

Table~\ref{tab:continuous_monitoring} traces the continuous monitoring flow from device capture to clinical alert, with latency at each stage.


{\tablestripes\tablefontsize

\needspace{10\baselineskip}
\begin{xltabular}{\textwidth}{@{} >{\raggedright\arraybackslash}p{1.0cm} >{\raggedright\arraybackslash}p{3.2cm} p{8.5cm} >{\raggedright\arraybackslash}p{1.5cm} @{}}
\caption[Continuous Monitoring Flow: Emotiv Device $$ Mobile $$]{Continuous Monitoring Flow: Emotiv Device $\to$ Mobile $\to$ Cloud $\to$ Alert. The four-stage pipeline specifies latency targets at each handoff, demonstrating that real-time anomaly detection and clinician alerting are achievable within the RGAIG architecture.}\label{tab:continuous_monitoring} \\
\toprule
\thead{\#} & \thead{Stage} & \thead{What Happens} & \thead{Latency} \\
\midrule
\endfirsthead
\endhead
\bottomrule
\endlastfoot
M1 & Device Capture & Emotiv EPOC X captures 14-channel EEG at 128\,Hz continuously; Bluetooth Low Energy (BLE) streams to paired mobile device & Real-time \\
\addlinespace
M2 & Edge Processing & Mobile app performs initial signal quality check, artefact flagging, and data compression; buffers 30-second epochs & $<$1 second \\
\addlinespace
M3 & Cloud Upload & Encrypted EEG epochs streamed to RGAIG cloud via WebSocket connection; TLS 1.3 transport encryption & $<$2 seconds \\
\addlinespace
M4 & Real-Time Analysis & Cloud pipeline runs continuous analysis: bandpass filtering, spectral decomposition, anomaly detection (interictal spike detection, seizure onset patterns) & $<$5 seconds \\
\addlinespace
M5 & Alert Generation & If anomaly detected: severity classification (low/moderate/high/critical); context-aware alert generated & $<$3 seconds \\
\addlinespace
M6 & Push Notification & Alert delivered to user's mobile device; critical alerts also sent to emergency contacts and linked clinician & $<$2 seconds \\
\addlinespace
M7 & Dashboard Update & Real-time brain activity visualisation on mobile dashboard; trend metrics updated in backend (T7); session logged (T3) & $<$5 seconds \\
\addlinespace
M8 & Data Persistence & All streaming data archived in cloud storage; session summary generated; audit trail updated (T9) & Background \\
\addlinespace
M9 & Trend Computation & Weekly and monthly trend analysis computed from accumulated sessions; drift detection triggers clinician notification & Scheduled \\
\addlinespace
M10 & Battery/Health & Device battery level, electrode impedance, and connectivity monitored; user alerted if signal quality degrades & Continuous \\
\end{xltabular}}
\vspace{0.5em}
\noindent\textit{Note.} Total end-to-end latency from EEG capture to alert delivery: $<$15 seconds. Continuous monitoring designed for 4--8 hour daily sessions (e.g., sleep monitoring, work monitoring). Backend tables T3 (sessions), T7 (trends), and T9 (audit) referenced from Table~\ref{tab:backend_db}. BLE = Bluetooth Low Energy; WebSocket = full-duplex communication protocol.
\fi %% End suppressed continuous_monitoring

\phantomsection\label{sec:continuous_monitoring_flow}
\phantomsection\label{tab:continuous_monitoring}

%===============================================================================
\subsection{Dual Data Approach}
\label{sec:dual_data_approach}

Figure~\ref{fig:dual_pipeline} contrasts the two data pipelines---assessment data and EEG data---at each processing stage.


\needspace{12\baselineskip}
\begin{figure}[!htbp]
\centering
\caption[Dual Data Pipeline: Assessment Data vs.EEG Data]{Dual Data Pipeline: Assessment Data vs.\ EEG Data. The figure contrasts the questionnaire/behavioural data path with the continuous EEG signal path, showing how the two pipelines converge at the fusion layer to produce a unified diagnostic score.}
\label{fig:dual_pipeline}
\resizebox{\textwidth}{0.62\textheight}{%
\begin{tikzpicture}[
  every node/.style={font=\fignodefont},
  node distance=0.65cm and 2.8cm,
    leftbox/.style={
        rounded corners=4pt, draw=ggunavy, thick,
        minimum width=4.8cm, minimum height=1.0cm,
        fill=lightblue!15, font=\fignodefont},
    rightbox/.style={
        rounded corners=4pt, draw=ggunavy, thick,
        minimum width=4.8cm, minimum height=1.0cm,
        fill=lightteal!15, font=\fignodefont},
    widebox/.style={
        rounded corners=4pt, draw=ggunavy, thick,
        minimum width=10.8cm, minimum height=1.0cm,
        fill=lightgreen!20, font=\fignodefont},
    stagenum/.style={font=\small\bfseries, text=ggunavy},
    arr/.style={-{Stealth[length=4.5mm,width=3.3mm]}, line width=1.45pt, ggunavy},
    coltitle/.style={font=\small\bfseries, text=ggunavy}
]
% 
% Column headers
% 
\node[coltitle] (hdrL) {Assessment Data Pipeline};
% 
\node[coltitle, right=2.8cm of hdrL] (hdrR) {EEG Data Pipeline};
% 
% Stage 1
% 
\node[leftbox, below=0.6cm of hdrL] (L1) {\textbf{1.} Clinical Questionnaires};
% 
\node[rightbox, below=0.6cm of hdrR] (R1) {\textbf{1.} Raw 14-ch EEG Capture};
% 
% Stage 2
% 
\node[leftbox, below=of L1] (L2) {\textbf{2.} Form Validation};
% 
\node[rightbox, below=of R1] (R2) {\textbf{2.} Signal Quality Check};
% 
% Stage 3
% 
\node[leftbox, below=of L2] (L3) {\textbf{3.} Normalise \& Encode};
% 
\node[rightbox, below=of R2] (R3) {\textbf{3.} Filter \& ICA Removal};
% 
% Stage 4
% 
\node[leftbox, below=of L3] (L4) {\textbf{4.} Clinical Feature Scores};
% 
\node[rightbox, below=of R3] (R4) {\textbf{4.} Spectral Features};
% 
% Stage 5
% 
\node[leftbox, below=of L4] (L5) {\textbf{5.} Feature Vector for ML};
% 
\node[rightbox, below=of R4] (R5) {\textbf{5.} Spectrogram Images};
% 
% Stage 6
% 
\node[leftbox, below=of L5] (L6) {\textbf{6.} Gradient Boosted Trees};
% 
\node[rightbox, below=of R5] (R6) {\textbf{6.} CNN / ResNet};
% 
% Stage 7 — FUSION (wide, centered between columns)
% 
\node[widebox, below=0.7cm of $(L6.south)!0.5!(R6.south)$] (F7) {\textbf{7. FUSION} --- Multi-modal Ensemble};
% 
% Stage 8 — split again
% 
\node[leftbox, below=0.7cm of F7.south west, anchor=north, xshift=2.4cm] (L8) {\textbf{8.} SHAP: Clinical Features};
% 
\node[rightbox, below=0.7cm of F7.south east, anchor=north, xshift=-2.4cm] (R8) {\textbf{8.} SHAP/LIME: EEG Features};
% 
% Stage 9 — shared
% 
\node[widebox, below=0.7cm of $(L8.south)!0.5!(R8.south)$] (F9) {\textbf{9. RAG} --- Evidence-grounded Explanation};
% 
% Stage 10 — split
% 
\node[leftbox, below=0.7cm of F9.south west, anchor=north, xshift=2.4cm] (L10) {\textbf{10.} Risk Narrative Report};
% 
\node[rightbox, below=0.7cm of F9.south east, anchor=north, xshift=-2.4cm] (R10) {\textbf{10.} EEG Report};
% 
% Stage 11 — shared
% 
\node[widebox, below=0.7cm of $(L10.south)!0.5!(R10.south)$] (F11) {\textbf{11. GOVERNANCE} --- 6-Gate Guardrail};
% 
% ---- Arrows: left column ----
\draw[arr] (L1) -- (L2);
\draw[arr] (L2) -- (L3);
\draw[arr] (L3) -- (L4);
\draw[arr] (L4) -- (L5);
\draw[arr] (L5) -- (L6);
\draw[arr] (L6.south) -- ++(0,-0.25) -| ([xshift=-1.5cm]F7.north);
% 
% ---- Arrows: right column ----
\draw[arr] (R1) -- (R2);
\draw[arr] (R2) -- (R3);
\draw[arr] (R3) -- (R4);
\draw[arr] (R4) -- (R5);
\draw[arr] (R5) -- (R6);
\draw[arr] (R6.south) -- ++(0,-0.25) -| ([xshift=1.5cm]F7.north);
% 
% ---- Arrows: fusion onward ----
\draw[arr] ([xshift=-1.5cm]F7.south) -- ++(0,-0.25) -| (L8.north);
\draw[arr] ([xshift=1.5cm]F7.south) -- ++(0,-0.25) -| (R8.north);
\draw[arr] (L8.south) -- ++(0,-0.25) -| ([xshift=-1.5cm]F9.north);
\draw[arr] (R8.south) -- ++(0,-0.25) -| ([xshift=1.5cm]F9.north);
\draw[arr] ([xshift=-1.5cm]F9.south) -- ++(0,-0.25) -| (L10.north);
\draw[arr] ([xshift=1.5cm]F9.south) -- ++(0,-0.25) -| (R10.north);
\draw[arr] (L10.south) -- ++(0,-0.25) -| ([xshift=-1.5cm]F11.north);
\draw[arr] (R10.south) -- ++(0,-0.25) -| ([xshift=1.5cm]F11.north);
% 
\end{tikzpicture}%
}% end resizebox
\vspace{0.2em}


\noindent\textit{Note.} The two pipelines operate in parallel and converge at stage~7. Assessment data provides clinical context; EEG data provides objective biomarker evidence. WVD = Wigner-Ville Distribution; SPWVD = Smoothed Pseudo-WVD; STWVD = Smoothed Temporally-stable WVD; STFT = Short-Time Fourier Transform; CWT = Continuous Wavelet Transform.
\end{figure}

\phantomsection\label{tab:dual_pipeline}

\iffalse %% Original dual pipeline table moved to TikZ flowchart (fig:dual_pipeline)

\needspace{12\baselineskip}
\noindent The table below presents the dual Data Pipeline: Assessment Data vs.\ EEG Data.

\begin{table}[!htbp]
\tablestripes\tablefontsize
\tablefontsize
\caption{Dual Data Pipeline: Assessment Data vs.\ EEG Data}
\begin{tabularx}{\textwidth}{@{} >{\raggedright\arraybackslash}p{0.6cm} >{\raggedright\arraybackslash}p{2.2cm} L L @{}}
\toprule
\thead{\#} & \thead{Stage} & \thead{Pipeline 1: Assessment Data} & \thead{Pipeline 2: EEG Data} \\
\midrule
1 & Data Capture & Clinical questionnaire (ASD/PD/Epilepsy forms); symptom logger; medication history; demographics & Raw 14-channel EEG from Emotiv EPOC X at 128\,Hz; 10-minute continuous recording \\
\addlinespace
2 & Validation & Form completeness check; range validation; missing data flagging; duplicate detection & Signal quality assessment: electrode impedance, artefact percentage, signal-to-noise ratio \\
\addlinespace
3 & Preprocessing & Data normalisation; categorical encoding; missing value imputation; feature scaling & Bandpass filtering (0.5--45\,Hz, Butterworth 4th order); ICA artefact removal (eye/muscle); epoch segmentation \\
\addlinespace
4 & Feature Extraction & Clinical feature derivation: symptom severity scores, composite indices, risk factor combinations & Spectral power ($\delta$, $\theta$, $\alpha$, $\beta$, $\gamma$ bands); time-frequency representations (WVD, SPWVD, STWVD); connectivity metrics \\
\addlinespace
5 & Transformation & Structured feature vector for ML input & EEG $\to$ 2D spectrogram conversion via STFT/CWT for deep learning image classification \\
\addlinespace
6 & Classification & Gradient Boosted Trees on clinical features; cross-validated with bootstrap CI & CNN/ResNet on 2D spectrograms; domain-adaptive transfer learning; ASD ensemble classifier \\
\addlinespace
7 & Fusion & \multicolumn{2}{>{\raggedright\arraybackslash}p{0.72\textwidth}}{Multi-modal fusion: assessment predictions + EEG predictions combined via weighted ensemble (ASD accuracy; results in Chapter~\ref{chap:analysis})} \\
\addlinespace
8 & Explainability & SHAP values for clinical feature importance; ``Which symptoms drove the risk score?'' & SHAP/LIME for EEG feature attribution; ``Which brain regions and frequency bands drove the classification?'' \\
\addlinespace
9 & RAG Generation & \multicolumn{2}{>{\raggedright\arraybackslash}p{0.72\textwidth}}{RAG retrieves relevant medical literature from ChromaDB (1.17\,GB); Llama-3.2 generates plain-language explanation with citations} \\
\addlinespace
10 & Report Output & Assessment-based risk narrative with clinical reasoning & EEG-based classification report with spectral analysis and brain region mapping \\
\addlinespace
11 & Governance & \multicolumn{2}{>{\raggedright\arraybackslash}p{0.72\textwidth}}{Both pipelines pass through 6-gate guardrail: privacy gate $\to$ domain gate $\to$ quality gate $\to$ safety gate $\to$ bias check $\to$ audit log} \\
\bottomrule
\end{tabularx}
\vspace{0.5em}
\noindent\textit{Note.} The two pipelines run in parallel and converge at stage~7. Assessment data provides clinical context; EEG data provides biomarker evidence. STFT = Short-Time Fourier Transform; CWT = Continuous Wavelet Transform.
\end{table}
\fi %% End suppressed dual_pipeline table


\iffalse %% Operational detail moved to Appendix A --- knowledge_pipeline
A third knowledge pipeline (Table~\ref{tab:knowledge_pipeline}) feeds directly into the RAG model to ground explanations in evidence.


{\tablestripes\tablefontsize

\needspace{10\baselineskip}
\begin{xltabular}{\textwidth}{@{} >{\raggedright\arraybackslash}p{1.0cm} >{\raggedright\arraybackslash}p{2.9cm} p{7.9cm} >{\raggedright\arraybackslash}p{2.3cm} @{}}
\caption[Pipeline 3: Knowledge and Historical Data Sources for]{Pipeline 3: Knowledge and Historical Data Sources for RAG Model. This third pipeline catalogues the external knowledge bases---clinical guidelines, literature embeddings, and historical patient records---that feed the retrieval-augmented generation layer, enabling evidence-grounded explanations in the final diagnostic report.}\label{tab:knowledge_pipeline} \\
\toprule
\thead{\#} & \thead{Data Source} & \thead{Content and Purpose} & \thead{Storage} \\
\midrule
\endfirsthead
\endhead
\bottomrule
\endlastfoot
K1 & Vector Database & 1.17\,GB of medical literature embeddings (PubMed, clinical guidelines, textbooks); semantic search enables context-relevant retrieval for RAG explanations & ChromaDB \\
\addlinespace
K2 & Relational Database & Patient demographics, medication history, prior diagnoses, clinician notes, appointment records; structured query for patient context & PostgreSQL \\
\addlinespace
K3 & Historical EEG Data & Previous EEG recordings for longitudinal comparison; baseline patterns; trend detection across sessions & Object Storage \\
\addlinespace
K4 & Graph Database & Disease ontology relationships (ICD-10 codes, symptom-condition mappings, drug-interaction networks); enables multi-hop clinical reasoning & Neo4j \\
\addlinespace
K5 & Model Registry & Trained model versions (198 joblib files, 380\,MB); hyperparameters, training datasets, accuracy metrics; model lineage tracking & MLflow \\
\addlinespace
K6 & Audit Log Store & Immutable log of all AI decisions, gate outcomes, human overrides, and escalation events; regulatory compliance evidence & Append-only DB \\
\addlinespace
K7 & Clinical Guidelines & Structured clinical guidelines (NICE, AAN, ILAE) encoded as retrieval-ready knowledge units for RAG grounding & ChromaDB \\
\addlinespace
K8 & Drug Reference & Medication databases (BNF, FDA labels) for drug interaction checking and treatment recommendation validation & Structured DB \\
\end{xltabular}}
\vspace{0.5em}
\noindent\textit{Note.} Pipeline~3 feeds directly into the Agentic RAG engine alongside Pipelines~1 (assessment) and~2 (EEG). The RAG model retrieves from K1 and K7 for explanation generation, queries K2 and K3 for patient context, traverses K4 for clinical reasoning chains, and references K5 for model provenance. All access is logged to K6 for auditability. ICD-10 = International Classification of Diseases; NICE = National Institute for Health and Care Excellence; AAN = American Academy of Neurology; ILAE = International League Against Epilepsy; BNF = British National Formulary.
\fi %% End suppressed knowledge_pipeline

\phantomsection\label{tab:knowledge_pipeline}

\iffalse %% Engineering detail moved to Appendix A --- backend_database_tables
%===============================================================================
\subsection{Backend Database Tables for Continuous Monitoring}
\label{sec:backend_tables}

The RGAIG platform maintains a set of backend database tables that continuously track patient information, EEG sessions, AI predictions, and system health. Table~\ref{tab:backend_db} documents the core database schema supporting the continuous monitoring architecture.


{\tablestripes\tablefontsize
\needspace{10\baselineskip}
\begin{xltabular}{\textwidth}{@{} >{\raggedright\arraybackslash}p{1.0cm} >{\raggedright\arraybackslash}p{1.0cm} p{9.5cm} >{\raggedright\arraybackslash}p{1.5cm} @{}}
\caption[Backend Database Schema: Continuous Patient Monitoring]{Backend Database Schema: Continuous Patient Monitoring. The schema defines the tables, key fields, and update frequencies required to support real-time longitudinal tracking, providing the data architecture blueprint for the implementation described in Chapter~3.}\label{tab:backend_db} \\
\toprule
\thead{\#} & \thead{Database Table} & \thead{Fields and Purpose} & \thead{Update} \\
\midrule
\endfirsthead
\endhead
\bottomrule
\endlastfoot
T1 & \texttt{users} & User ID, demographics, registration date, consent status, subscription tier, account status & On register \\
\addlinespace
T2 & \texttt{medical\_\allowbreak profiles} & Medical history, medications, allergies, family history, comorbidities; linked to user ID & On update \\
\addlinespace
T3 & \texttt{eeg\_\allowbreak sessions} & Session ID, user ID, timestamp, duration, device ID, signal quality score, raw file path, processing status & Each session \\
\addlinespace
T4 & \texttt{predictions} & Prediction ID, session ID, condition, risk score, confidence interval, model version, classification label, SHAP top-5 features & Each analysis \\
\addlinespace
T5 & \texttt{rag\_\allowbreak reports} & Report ID, prediction ID, generated text, citation list, faithfulness score, hallucination rate, generation timestamp & Each report \\
\addlinespace
T6 & \texttt{symptom\_\allowbreak logs} & Log ID, user ID, date, mood, sleep, headache, seizure count, tremor, anxiety, medication adherence & Daily \\
\addlinespace
T7 & \texttt{trend\_\allowbreak metrics} & User ID, condition, weekly average score, monthly trend direction, change percentage, alert trigger status & Weekly calc \\
\addlinespace
T8 & \texttt{alerts} & Alert ID, user ID, type (routine/urgent/emergency), condition, message, status (sent/read/acted), timestamp & On trigger \\
\addlinespace
T9 & \texttt{audit\_\allowbreak trail} & Audit ID, user ID, action, gate decisions (6 gates), model version, override flag, clinician ID (if B2B), timestamp & Every action \\
\addlinespace
T10 & \texttt{consent\_\allowbreak log} & Consent ID, user ID, consent type, granted/revoked, timestamp, IP address, version & On change \\
\addlinespace
T11 & \texttt{doctor\_\allowbreak links} & Link ID, user ID, doctor name, hospital, speciality, contact, FHIR endpoint, auto-share enabled, last shared date & On setup \\
\addlinespace
T12 & \texttt{device\_\allowbreak health} & Device ID, user ID, battery level, firmware version, electrode impedance, last calibration, connectivity status & Real-time \\
\addlinespace
T13 & \texttt{data\_\allowbreak quality} & Quality ID, session ID, artefact percentage, SNR, channels rejected, preprocessing pass/fail, quality grade & Each session \\
\addlinespace
T14 & \texttt{feedback} & Feedback ID, user ID, report ID, helpfulness (1--5), accuracy (1--5), free text comment, timestamp & Each report \\
\addlinespace
T15 & \texttt{system\_\allowbreak metrics} & Timestamp, API latency (ms), model inference time, queue depth, error rate, active users, cache hit rate & Per minute \\
\end{xltabular}}
\vspace{0.5em}
\noindent\textit{Note.} All tables use PostgreSQL with row-level security (RLS). Patient data encrypted at rest (AES-256). Audit trail (T9) is append-only and immutable. Trend metrics (T7) computed by scheduled background jobs. System metrics (T15) feed into the observability dashboard. SNR = Signal-to-Noise Ratio; FHIR = Fast Healthcare Interoperability Resources.
\fi %% End suppressed backend_database_tables

\phantomsection\label{sec:backend_tables}
\phantomsection\label{tab:backend_db}

The RGAIG framework stores diagnostic results, audit trails, and governance metadata in a structured database supporting both B2B (hospital integration) and B2C (patient-facing) workflows. The schema encompasses 15 tables covering patient profiles, EEG sessions, AI predictions, RAG-generated reports, symptom logs, consent management, and system observability---with row-level security and AES-256 encryption at rest. Database schema specifications are provided in Appendix~A.


\iffalse %% Detailed methodology covered in Chapter 3
%===============================================================================
\subsection{Research Methodology Approach}
\label{sec:research_methodology_approach}


\noindent This table lists each \textit{row} across approach, data sources, analysis methods, and 1 more dimension, so examiners can trace what was assessed and what evidence supports each row.



\noindent This table lists each \textit{\#} across approach, data sources, analysis methods, and 1 more dimensions, so examiners can trace what was assessed and what evidence supports each row.


\noindent This table lists each \textit{\#} across approach, data sources, analysis methods, and 1 more dimensions, so examiners can trace what was assessed and what evidence supports each row.



{\tablestripes\tablefontsize
\needspace{10\baselineskip}
\begin{xltabular}{\textwidth}{@{} >{\raggedright\arraybackslash}p{1.0cm} >{\raggedright\arraybackslash}p{2.0cm} p{4.7cm} p{5.0cm} >{\raggedright\arraybackslash}p{2.0cm} @{}}
\caption[Research Methodology: Quantitative, Qualitative, and]{Research Methodology: Quantitative, Qualitative, and Mixed Methods with Data. The table maps each methodological approach to its data sources, analysis methods, and linked objective, demonstrating that the mixed-methods design triangulates evidence from complementary paradigms.}\label{tab:research_approach} \\
\toprule
\thead{\#} & \thead{Approach} & \thead{Data Sources} & \thead{Analysis Methods} & \thead{Obj.} \\
\midrule
\endfirsthead
\endhead
\bottomrule
\endlastfoot
\multicolumn{5}{l}{\textit{\textbf{Quantitative Methods}}} \\
\addlinespace
Q1 & Classification Accuracy & 305\,GB data (19 ASD subjects, ASD); accuracy, precision, recall, F1-score, AUC & ASD ensemble classifier, cross-validation (10-fold), bootstrap CI (1,000 resamples), effect size (Cohen's $d$); results in Chapter~\ref{chap:analysis} & RO1, RO2 \\
\addlinespace
Q2 & RAG Quality & 1.17\,GB ChromaDB; 50 case reports per domain & Citation accuracy, faithfulness, hallucination rate, latency; results in Chapter~\ref{chap:analysis} & RO3 \\
\addlinespace
Q3 & Governance Compliance & 525-module quality taxonomy across 3 tiers & Module pass rates per tier, compliance score, audit completion rate; results in Chapter~\ref{chap:analysis} & RO4 \\
\addlinespace
Q4 & Survey (Likert) & Clinician survey (C1--C20), consumer survey (U1--U20); target $n \geq$ 50 per group & Descriptive statistics, Cronbach's $\alpha$, factor analysis, structural equation modelling & RO5 \\
\addlinespace
Q5 & Performance & System response times, throughput, error rates, scalability benchmarks & Statistical process control, percentile analysis, load testing results & RO4 \\
\midrule
\multicolumn{5}{l}{\textit{\textbf{Qualitative Methods}}} \\
\addlinespace
QL1 & Expert Panel & 12 domain experts (n=12, purposive sampling; neurology, AI/ML, governance, regulatory); semi-structured evaluation & Expert validation survey (E1--E18), thematic analysis of open-ended responses, Krippendorff's $\alpha$ ($\geq$0.80) & RO1--RO5 \\
\addlinespace
QL2 & Case Studies & Three hospital evaluation scenarios (teaching, community, specialist); stakeholder interviews & Cross-case analysis, pattern matching, explanation building & RO3, RO4 \\
\addlinespace
QL3 & Clinical Review & RAG-generated reports reviewed by neurologists for clinical accuracy, relevance, and safety & Content analysis, clinical accuracy scoring (1--5), safety flag assessment & RO3 \\
\midrule
\multicolumn{5}{l}{\textit{\textbf{Mixed Methods Integration}}} \\
\addlinespace
M1 & Triangulation & Quantitative metrics (Q1--Q5) + qualitative insights (QL1--QL3) compared for convergence & Convergent parallel design: independent analysis then integration matrix; discrepancies investigated & All ROs \\
\addlinespace
M2 & Explanatory & Quantitative outliers (e.g., low accuracy in specific domains) explained by qualitative expert insights & Sequential explanatory: quantitative first, then qualitative follow-up for unexpected results & RO1, RO2 \\
\addlinespace
M3 & Validation & Framework design validated by both metric achievement and expert endorsement & Design Science Research (DSR) evaluation: utility, quality, efficacy demonstrated through both lenses & RO4, RO5 \\
\end{xltabular}}
\vspace{0.5em}
\noindent\textit{Note.} RO1--RO5 refer to Research Objectives defined in Section~\ref{sec:research_purpose}. The mixed-methods design follows Creswell and Creswell's (2018) convergent parallel design, where quantitative and qualitative data are collected concurrently, analysed independently, and then integrated. DSR = Design Science Research.
\fi %% End suppressed research_methodology_approach

\phantomsection\label{sec:research_methodology_approach}
\phantomsection\label{tab:research_approach}

%===============================================================================
% Patient-to-Provider Data Sharing Pipeline
%===============================================================================
\subsection{Patient-to-Provider Data Sharing Pipeline}
\label{sec:sharing_pipeline}

Figure~\ref{fig:sharing_pipeline} presents the eight-layer patient-to-provider data sharing pipeline.

% [SPACE-FIX] clearpage removed to eliminate empty page
\noindent\textbf{Patient-to-provider data sharing pipeline:.}

\needspace{12\baselineskip}
\begin{figure}[!htbp]
\centering
\vspace*{-1cm}
\resizebox{0.97\textwidth}{!}{%
\begin{tikzpicture}[
  every node/.style={font=\fignodefont},
  % Layer label style (left side)
  lbl/.style={rectangle, rounded corners=3pt, draw=ggunavy, text=white,
    font=\small\bfseries, minimum width=2.8cm, minimum height=1.0cm,
    align=center, inner sep=3pt},
  % Detail box style (right side)
  detail/.style={rectangle, rounded corners=3pt, draw=ggunavy,
    font=\small, minimum height=1.0cm, align=center, inner sep=5pt},
  % Arrow style
  arr/.style={-{Stealth[length=4mm, width=3mm]}, very thick, ggunavy},
  % Label arrows (connecting layer labels)
  tarr/.style={-{Stealth[length=3mm, width=2mm]}, thick, ggunavy!60},
  node distance=0.8cm
]
% 
% === Layer 1: B2C Layer ===
% 
\node[lbl, fill=ggunavy] (l1) {B2C\\Layer};
% 
\node[detail, fill=lightblue, right=0.6cm of l1, text width=11cm] (d1) {
  Mobile onboarding, consent, provider link, sharing preference
};
% 
% === Layer 2: Data Capture ===
% 
\node[lbl, fill=ggunavy, below=of l1] (l2) {Data\\Capture};
% 
\node[detail, fill=lightblue, right=0.6cm of l2, text width=11cm] (d2) {
  Wearable EEG, mobile symptoms, and system event data
};
% 
% === Layer 3: Data Processing ===
% 
\node[lbl, fill=ggunavy, below=of l2] (l3) {Data\\Processing};
% 
\node[detail, fill=lightorange, right=0.6cm of l3, text width=11cm] (d3) {
  Validate input, extract EEG features, store securely
};
% 
% === Layer 4: RAG Intelligence ===
% 
\node[lbl, fill=ggunavy, below=of l3] (l4) {RAG\\Intelligence};
% 
\node[detail, fill=lightorange, right=0.6cm of l4, text width=11cm] (d4) {
  Retrieve data, generate report, validate confidence and safety
};
% 
% === Layer 5: Delivery Decision ===
% 
\node[lbl, fill=ggunavy, below=of l4] (l5) {Delivery\\Decision};
% 
\node[detail, fill=lightteal, right=0.6cm of l5, text width=11cm] (d5) {
  Select one-time, scheduled, continuous, or event-based sharing
};
% 
% === Layer 6: Integration ===
% 
\node[lbl, fill=ggunavy, below=of l5] (l6) {Integration\\Layer};
% 
\node[detail, fill=lightteal, right=0.6cm of l6, text width=11cm] (d6) {
  Validate consent, choose channel, transmit to provider
};
% 
% === Layer 7: B2B Layer ===
% 
\node[lbl, fill=ggunavy, below=of l6] (l7) {B2B\\Layer};
% 
\node[detail, fill=lightpurple, right=0.6cm of l7, text width=11cm] (d7) {
  Provider receives report, reviews it, and responds or escalates
};
% 
% === Layer 8: Governance ===
% 
\node[lbl, fill=ggunavy, below=of l7] (l8) {Governance\\Layer};
% 
\node[detail, fill=lightgreen, right=0.6cm of l8, text width=11cm] (d8) {
  Track delivery, log audit trail, notify patient
};
% 
% === Vertical arrows between layers ===
\draw[arr] (d1.south) -- (d2.north);
\draw[arr] (d2.south) -- (d3.north);
\draw[arr] (d3.south) -- (d4.north);
\draw[arr] (d4.south) -- (d5.north);
\draw[arr] (d5.south) -- (d6.north);
\draw[arr] (d6.south) -- (d7.north);
\draw[arr] (d7.south) -- (d8.north);
% 
% === Vertical arrows between labels ===
\draw[tarr] (l1.south) -- (l2.north);
\draw[tarr] (l2.south) -- (l3.north);
\draw[tarr] (l3.south) -- (l4.north);
\draw[tarr] (l4.south) -- (l5.north);
\draw[tarr] (l5.south) -- (l6.north);
\draw[tarr] (l6.south) -- (l7.north);
\draw[tarr] (l7.south) -- (l8.north);
% 
\end{tikzpicture}%
}% end resizebox
\caption[Patient-to-provider data sharing pipeline]{Patient-to-provider data sharing pipeline: eight-layer vertical flow from B2C wearable capture through RAG-powered report generation and consent-governed delivery to B2B clinical integration.}
\label{fig:sharing_pipeline}
\end{figure}

\iffalse %% Condensed — caption covers the content

\noindent\textbf{What it shows.} Figure~\ref{fig:sharing_pipeline} illustrates the eight-layer data sharing pipeline: (1)~B2C patient onboarding with consent, (2)~multi-modal data capture from wearable EEG and mobile app, (3)~validated data processing with encrypted storage, (4)~RAG-powered intelligence with confidence thresholds ($\geq$0.85), (5)~delivery decision logic (one-time, scheduled, continuous, or event-based), (6)~consent-validated integration layer, (7)~B2B clinical dashboard delivery, and (8)~governance-enforced audit trail.


\noindent\textbf{So what.} The pipeline demonstrates that RGAIG is not merely a classification engine but a complete data-sharing infrastructure that bridges the B2C--B2B divide with governance enforcement at every transition. This end-to-end flow directly addresses gap G3 (lack of workflow automation) and supports objective O3 (agentic orchestration), while the embedded consent and audit layers satisfy regulatory requirements central to objective O5 (governance integration).
\fi

\iffalse %% Engineering detail moved to Appendix A --- framework_component_architecture
%===============================================================================
% RGAIG Framework Component Architecture
%===============================================================================
\subsection{RGAIG Framework Component Architecture}
\label{sec:framework_components}

Eight interconnected layers form the RGAIG pipeline: from patient-facing data capture at the bottom, through AI intelligence in the middle, to clinical integration and governance enforcement at the top. Figure~\ref{fig:framework_components} traces these connections.

% [SPACE-FIX] clearpage removed to eliminate empty page
\noindent\textbf{RGAIG framework component architecture: eight.}

\needspace{12\baselineskip}
\begin{figure}[!htbp]
\centering
\vspace*{-1cm}
\begin{tikzpicture}[
  every node/.style={font=\fignodefont},
  % Main layer box
  layer/.style={rectangle, rounded corners=4pt, draw=ggunavy,
    font=\small, minimum width=11cm, minimum height=1.4cm,
    align=center, inner sep=6pt},
  % Governance sidebar
  govbox/.style={rectangle, rounded corners=4pt, draw=ggunavy, fill=lightgreen,
    font=\small, text width=3.5cm, align=center, inner sep=6pt,
    minimum height=14cm},
  % Arrow style
  arr/.style={-{Stealth[length=4mm, width=3mm]}, very thick, ggunavy},
  % Governance connection
  garr/.style={-{Stealth[length=3mm, width=2mm]}, thick, ggunavy!50, dashed},
  node distance=0.8cm
]
% 
% === Main pipeline (left column) ===
% 
\node[layer, fill=lightblue] (l1) {
  \textbf{Layer 1: B2C Layer (Patient)}\\
  {\small Mobile app onboarding, consent capture, doctor/hospital registration, sharing preferences}
};
% 
% 
\node[layer, fill=lightblue, below=of l1] (l2) {
  \textbf{Layer 2: Data Capture}\\
  {\small Wearable EEG (14-ch Emotiv), mobile symptom logs, system event streams}
};
% 
% 
\node[layer, fill=lightorange, below=of l2] (l3) {
  \textbf{Layer 3: Data Processing}\\
  {\small Schema validation, artefact rejection, spectral feature extraction, encrypted storage}
};
% 
% 
\node[layer, fill=lightorange, below=of l3] (l4) {
  \textbf{Layer 4: RAG Intelligence}\\
  {\small Retrieve patient context, generate diagnostic reports, validate confidence $\geq$ 0.85}
};
% 
% 
\node[layer, fill=lightteal, below=of l4] (l5) {
  \textbf{Layer 5: Explainability \& Trust}\\
  {\small SHAP/LIME feature attribution, natural-language rationale, clinician confidence scoring}
};
% 
% 
\node[layer, fill=lightteal, below=of l5] (l6) {
  \textbf{Layer 6: Integration \& Sharing}\\
  {\small Consent validation, channel selection (API/portal/email), B2C-to-B2B bridge}
};
% 
% 
\node[layer, fill=lightpurple, below=of l6] (l7) {
  \textbf{Layer 7: B2B Layer (Doctor/Hospital)}\\
  {\small Clinical dashboard, EHR integration, review/respond/escalate workflow}
};
% 
% 
\node[layer, fill=lightgreen, below=of l7] (l8) {
  \textbf{Layer 8: Governance Layer}\\
  {\small Consent lifecycle, immutable audit log, compliance reporting, patient feedback loop}
};
% 
% === Vertical arrows between layers ===
\draw[arr] (l1) -- (l2);
\draw[arr] (l2) -- (l3);
\draw[arr] (l3) -- (l4);
\draw[arr] (l4) -- (l5);
\draw[arr] (l5) -- (l6);
\draw[arr] (l6) -- (l7);
\draw[arr] (l7) -- (l8);
% 
% === Governance sidebar (right) ===
% 
\node[govbox, right=0.8cm of l4, anchor=west] (gov) {
  \textbf{Governance}\\
  \textbf{Enforcement}\\[5pt]
  {\small Consent validation}\\
  {\small at every transition}\\[4pt]
  {\small Audit logging}\\
  {\small (who, what, when)}\\[4pt]
  {\small HIPAA/GDPR}\\
  {\small compliance checks}\\[4pt]
  {\small Data minimisation}\\
  {\small + purpose limitation}\\[4pt]
  {\small Immutable trail}\\
  {\small for all 8 layers}
};
% 
% === Dashed lines from governance to each layer ===
\draw[garr] (gov.west |- l1) -- (l1.east);
\draw[garr] (gov.west |- l2) -- (l2.east);
\draw[garr] (gov.west |- l3) -- (l3.east);
\draw[garr] (gov.west |- l4) -- (l4.east);
\draw[garr] (gov.west |- l5) -- (l5.east);
\draw[garr] (gov.west |- l6) -- (l6.east);
\draw[garr] (gov.west |- l7) -- (l7.east);
\draw[garr] (gov.west |- l8) -- (l8.east);
% 
\end{tikzpicture}
\caption[RGAIG framework component architecture]{RGAIG framework component architecture: eight interconnected layers from patient-facing B2C capture through RAG intelligence and explainability to B2B clinical integration, with governance enforcement spanning all layers.}
\label{fig:framework_components}
\end{figure}


\noindent\textbf{What it shows.} Figure~\ref{fig:framework_components} presents the RGAIG component architecture as eight vertically stacked layers---from patient-facing B2C data capture at the base, through preprocessing, AI classification, RAG intelligence, and explainability in the middle, to B2B clinical integration at the top---with a governance enforcement column spanning all layers via dashed connectors.


\noindent\textbf{So what.} The vertical governance column confirms that responsible AI is not a post-hoc add-on but an architectural primitive embedded at every processing stage. This design principle distinguishes RGAIG from competing frameworks that bolt governance onto completed pipelines, and operationalises the 46-module responsible AI taxonomy described in Chapter~\ref{chap:methodology}.


\noindent\textbf{Figures~\ref{fig:sharing_pipeline} and~\ref{fig:framework_components}: Summary.} Together, these two figures present complementary views of the same system. Figure~\ref{fig:sharing_pipeline} traces the \textit{data flow} (what information moves where), while Figure~\ref{fig:framework_components} reveals the \textit{component architecture} (which processing modules execute at each stage). Their joint reading confirms that the RGAIG framework achieves both operational completeness (end-to-end data sharing) and architectural coherence (governance-embedded component design), addressing gaps G1 (unified pipeline) and G4 (governance integration) simultaneously.
\fi %% End suppressed framework_component_architecture

\phantomsection\label{sec:framework_components}
\phantomsection\label{fig:framework_components}

The RGAIG framework operates across eight integrated governance layers---from patient-facing B2C data capture through data processing, RAG intelligence, explainability, and integration to B2B clinical delivery---with a governance enforcement column spanning all layers. This architecture ensures that responsible AI is not a post-hoc add-on but an embedded primitive at every processing stage, directly addressing gap G3 (lack of workflow automation) and objective O3 (agentic orchestration). The complete layer-by-layer architecture specification is provided in Appendix~A.

\noindent Table~\ref{tab:gap_to_objective} maps each of the six identified research gaps to the corresponding research objective, establishing the traceability from gap identification to the research design that follows.

\paragraph{Research Gap to.}




{\tablestripes\tablefontsize
\needspace{10\baselineskip}
\begin{xltabular}{\textwidth}{@{} L L @{}}
\caption[Research Gap to Research Objective Mapping]{Research Gap to Research Objective Mapping. Each gap identified in the literature is paired with the objective that directly addresses it, confirming one-to-one traceability and ensuring no gap is left without a corresponding research response.}\label{tab:gap_to_objective} \\
\toprule
\thead{Research Gap} & \thead{Research Objective} \\
\midrule
\endfirsthead
\endhead
\bottomrule
\endlastfoot
No unified ASD assessment framework & Develop a unified ASD screening framework \\
Weak integration of primary and secondary data & Integrate questionnaire, behaviour, psychology, and EEG/IoT data \\
Lack of workflow automation & Evaluate automated onboarding-to-decision workflow \\
Poor explainability & Assess RAG/XAI-based explanation quality \\
Weak governance-ready AI design & Embed and evaluate privacy, compliance, auditability, and accountability \\
No continuity from hospital assessment to monitoring & Define B2B-first with B2C continuity-of-care extension \\
\end{xltabular}}
\vspace{2pt}
\noindent\textit{Note.} Each gap is derived from the literature review (Chapter~\ref{chap:literature}) and addressed through the corresponding objective in this dissertation.

%===============================================================================
\section{Research Objectives}
\label{sec:research_purpose}

The primary objective of this study is to develop and evaluate the Responsible Generative AI Governance Framework (RGAIG) for AI-assisted ASD-EEG screening support — positioned as a conceptual governance, explainability, and adoption framework rather than a deployable diagnostic AI product (Section~\ref{sec:scope_boundaries}). The study evaluates the framework across governance, explainability, trust + adoption, human oversight, and responsible-AI dimensions through stakeholder expert review and primary-data survey.

\paragraph{Directive-aligned objectives overlay.}
\label{para:objectives_lens}
The five directive-recommended DBA objectives map directly onto this dissertation's existing O1--O8 objectives table (below). The mapping is provided here as a navigation overlay so examiners can read the same evidence base through either organisation:

\begin{itemize}[leftmargin=*, itemsep=3pt]
    \item \textbf{Directive Obj.~A — Identify governance challenges affecting AI-assisted ASD support ecosystems.} Covered by O7 (Governance architecture) + the Ch.2 Research Gap synthesis. Evidence: 525-module governance taxonomy gap, regulatory-mapping limitations, HITL governance gap.
    \item \textbf{Directive Obj.~B — Evaluate explainability and human-oversight requirements influencing caregiver and clinician trust.} Covered by O4 (RAG explainability) + O6 (Biomarker / SHAP) + Layer~4 HITL gate evidence. Evidence: 12-expert panel rating $M = \expertM{}$, $\alpha = 0.84$; two-layer explainability consistency invariant.
    \item \textbf{Directive Obj.~C — Analyse adoption barriers affecting responsible AI integration in ASD-support workflows.} Covered by O5 (ASD-focused validation) + the PLS-SEM mediation analysis ($n = 131$ clinicians). Evidence: locked four-box causal chain (RAI governance $\to$ explainability $\to$ trust $\to$ adoption).
    \item \textbf{Directive Obj.~D — Develop a governance and explainability adoption framework for human-centered ASD AI ecosystems.} Covered by O1 (Standardised preprocessing) + O2 (Classification) + O3 (Agentic AI orchestration) + O7 (Governance) consolidated into the RGAIG 7-layer architecture. Output artefact: the framework itself plus six-level maturity model.
    \item \textbf{Directive Obj.~E — Validate framework readiness using expert and stakeholder evaluation.} Covered by O8 (Economic viability) + the Pipeline~2 survey + Delphi-style expert panel. Output artefact: validated readiness scorecard with explicit defer-criteria for items requiring future work.
\end{itemize}

\paragraph{Output artefact.} Each of the eight existing dissertation objectives produces a specific evidentiary artefact (table, model, framework component, validated metric). The five directive-aligned objectives above re-organise those artefacts under the governance / explainability / adoption / framework-development / validation structure that DBA committees expect for governance-oriented dissertations.

\noindent\textbf{Main Research Objectives (O1--O5).}

\noindent Table~\ref{tab:main_objectives} presents main Research Objectives (O1--O5).


{\tablestripes\tablefontsize
\needspace{10\baselineskip}
\begin{xltabular}{\textwidth}{@{} >{\centering\arraybackslash}p{1.0cm} p{9.5cm} >{\centering\arraybackslash}p{1.0cm} @{}}
\caption[Main Research Objectives (O1--O5)]{Main Research Objectives (O1--O5). These are the five primary objectives that structure the entire dissertation; each maps to one or more Parts of the A--D evaluation framework.}\label{tab:main_objectives} \\
\toprule
\thead{ID} & \thead{Objective} & \thead{Part} \\
\midrule
\endfirsthead
\endhead
\bottomrule
\endlastfoot
O1 & To evaluate patient and caregiver trust, usability, and adoption of the AI-based EEG screening system (B2C context) & A \\
\addlinespace
O2 & To validate the technical performance, robustness, and explainability of the EEG-based AI pipeline & B \\
\addlinespace
O3 & To assess clinician trust, workflow integration, and operational impact within hospital environments (B2B context) & C \\
\addlinespace
O4 & To develop and validate a Responsible AI governance framework aligned with industry standards & D \\
\addlinespace
O5 & To provide a combined evaluation recommendation for B2C and B2B environments & D \\
\end{xltabular}}
\vspace{0.3em}
\noindent\textit{Note.} For the full P$\to$G$\to$O$\to$RQ$\to$H$\to$C traceability chain, see Table~\ref{tab:research_flow_alignment}.

\section{Objective Operationalisation (Appendix~M)}

\noindent The detailed operationalization for this A--D framework component appears in Appendix~M (Section~M1). This section in Chapter~1 records that the framework exists and identifies its scope; the operationalization tables, indicator specifications, and procedural detail are placed in Appendix~M to preserve Chapter~1's focus on problem framing and to avoid pre-empting Chapter~3's methodology specification.

\noindent\textbf{B2C-primary, B2B-supporting structure.} The dissertation's evaluation logic is illustrated in Figure~\ref{fig:b2c_b2b_support_model}: the caregiver and patient sit at the core as the primary value beneficiaries, and the B2B safety \& governance layer (clinician review, hospital escalation, audit + accountability, governance maturity, EHR/FHIR integration) exists \emph{to support} consumer use rather than to be the end-purpose. The outer ring is the regulatory + market + ecosystem context within which B2C use is bounded. Every subsequent chapter is read against this structure: Chapter~3 methodology distinguishes the B2C primary path (Part~A surveys + Part~B EEG accuracy) from the B2B support path (Part~C hospital integration + Part~D governance); Chapter~4 reports caregiver-facing trust results before clinician-facing structural evidence; Chapter~5 concludes with \emph{Future B2C Prospective-Validation with B2B Support} rather than B2B-only enterprise deployment.

\input{figures/fig_b2c_b2b_support_model}

\noindent\textbf{B2B support scenarios in operation.} Figure~\ref{fig:b2c_b2b_support_model} shows the structural ring; Figure~\ref{fig:b2b_support_scenarios} shows the operational scenarios. The B2B support layer is not abstract ``governance'' --- it is six concrete services that activate by severity: (1) clinician mobile app for low-/medium-severity notification, (2) continuous monitoring oversight dashboard for trend cases, (3) SOS \& emergency hospital activation for critical cases, (4) hospital handheld device as the frontline receiver for SOS, (5) EHR/FHIR integration for structured routing of every event, and (6) audit \& governance for logs and liability. Critical and emergency outputs bypass the caregiver and route directly to the hospital handheld + SOS path; low- and medium-severity outputs route through the caregiver app first. Every service writes to EHR + audit so the basis of any decision is reconstructible. \emph{Limitation:} the six-scenario routing has been designed but not yet measured under prospective B2C deployment. \emph{Future scope:} the Canada + India future B2C prospective-validations (Ch.~5) will measure latency, false-alert rate, and clinician acceptance for each of the six routing paths.

\input{figures/fig_b2b_support_scenarios}


\section{Research Design --- Mixed-Methods Governance Evaluation}
\label{sec:research_design_overview}

\noindent\textbf{Research design at a glance.} The dissertation adopts a \textbf{pragmatic mixed-methods research design} positioned as a \emph{governance and adoption framework evaluation study} (per the rewritten Research Aim, Section~\ref{sec:research_goal}). The design is bounded as observational + evaluative + non-clinical-intervention; full methodology specification is in Chapter~\ref{chap:methodology}.

\paragraph{Why mixed methods.} Governance, explainability, and adoption are socio-technical constructs that require both quantitative measurement (framework-evaluation scoring, explainability-acceptance ratings, adoption-readiness scoring, PLS-SEM mediation analysis) and qualitative depth (expert interviews, clinician thematic analysis, caregiver perception narratives, governance-theme extraction). A purely quantitative design would miss the stakeholder-perception depth that DBA committees expect; a purely qualitative design would lack the statistical-rigour evidence the directive-aligned objectives require. Mixed methods is therefore the right design for this dissertation.

\paragraph{Quantitative + Qualitative decomposition.}
\begin{itemize}[leftmargin=*, itemsep=3pt]
    \item \textbf{Quantitative (Pipeline 1 + structured surveys):} Framework-evaluation scoring; explainability acceptance ratings (12-expert panel, $M = \expertM{}$, $\alpha = 0.84$); governance-readiness scoring; adoption-perception PLS-SEM ($n = 131$ clinicians; locked four-box causal chain: RAI governance $\to$ explainability $\to$ trust $\to$ adoption); subgroup fairness analysis (disparate impact $\geq 0.80$ gate); hallucination-rate measurement ($< 3$\% residual); retrospective ASD-EEG classification performance on KAU (\asdacc{} --- framed as supporting technical evidence, not the contribution claim per Chapter~\ref{chap:discussion} Reading Guide).
    \item \textbf{Qualitative (Pipeline 2 + expert interviews):} Semi-structured expert interviews ($n = 12$ panel of clinical + AI-governance professionals); thematic-analysis coding (governance themes, explainability acceptance themes, adoption barrier themes); stakeholder perception narratives; caregiver thematic content (where prospective Indian / paediatric cohort recruitment proceeds --- future scope per Section~\ref{subsec:data_strategy_existing_vs_new}); Delphi-style framework validation for the maturity model + decision-audit schema.
    \item \textbf{Mixed-methods integration:} Joint-display analysis converges quantitative findings (PLS-SEM mediation chain) with qualitative themes (expert interview categories); the integrated evidence base supports the Management Decision Matrix (Chapter~\ref{chap:discussion}, Section~\ref{sec:management_analysis}) with Mature for conceptual governance use / Requires prospective validation / Evidence insufficient managerial decisions per evidence row.
\end{itemize}

\paragraph{Research design phases.} The study unfolds in three phases consistent with the directive:
\begin{enumerate}[leftmargin=*, itemsep=2pt, label=(\roman*)]
    \item \textbf{Literature and governance analysis} (Chapter~\ref{chap:literature}: AI in ASD + XAI + RAI governance + Human-Centered AI + Adoption + HITL).
    \item \textbf{Conceptual framework development} (Chapter~\ref{chap:methodology}: governance layer + explainability layer + HITL oversight layer + adoption-readiness layer + decision-audit schema + 6-level maturity model).
    \item \textbf{Stakeholder evaluation and validation} (Chapter~\ref{chap:analysis}: 12-expert panel + 131-clinician PLS-SEM + retrospective benchmark + subgroup fairness + hallucination governance, all framed as governance / explainability / adoption evidence per the Reading Guide in Chapter~\ref{chap:discussion}).
\end{enumerate}

\paragraph{Validation strategy.} Framework validation uses (a)~expert review (Delphi-style consensus from 12-expert panel); (b)~stakeholder scoring (PLS-SEM causal chain on 131 clinicians); (c)~governance-readiness evaluation (6-level maturity model rubric); (d)~explainability assessment (two-layer consistency invariant + caregiver-comprehension verification on prospective cohort, future). The dissertation explicitly does \emph{not} validate clinical-diagnosis capability, medical effectiveness, intervention performance, or real-world clinical outcomes (per the framework IS-NOT contract, Section~\ref{sec:scope_boundaries}).


\addcontentsline{toc}{paragraph}{Must-Have: Research Questions} % [TOC-must-have-insertion]

\subsection{Research Contribution Statement}
\label{sec:research_contribution_statement}

\noindent This dissertation contributes to knowledge in three distinct ways. First, it proposes and empirically evaluates a governance-centered adoption pathway linking Responsible AI Governance, explainability, trust, and adoption readiness. Second, it introduces the Responsible Generative AI Governance (RGAIG) Maturity Framework as a practical managerial artifact for healthcare organizations. Third, it extends established technology adoption theory by positioning governance maturity as an antecedent to explainability and trust within AI-assisted decision-support environments.

\section{Research Questions}

\subsection{Theory Extension Declaration}
\label{sec:theory_extension_declaration}

\noindent Existing technology adoption theories provide valuable explanations for user acceptance of information systems; however, they were largely developed prior to the emergence of highly autonomous AI-enabled decision-support technologies. This dissertation proposes that governance maturity functions as a foundational antecedent that shapes stakeholder perceptions of explainability and trust. Accordingly, the study extends the Technology Acceptance Model by incorporating governance as a structural organizational construct influencing downstream adoption outcomes.

\label{sec:hypotheses}

\paragraph{Directive-aligned research questions overlay.}
\label{para:rqs_lens}
The four directive-recommended DBA research questions map onto this dissertation's existing RQ1--RQ8 set. The mapping is provided here as a navigation overlay for examiners reading through the directive's governance-evaluation lens:

\begin{itemize}[leftmargin=*, itemsep=3pt]
    \item \textbf{Directive RQ-A — What governance challenges affect responsible AI-assisted ASD support ecosystems?} Answered by the dissertation's RQ7 (Governance architecture) + RQ8 (Economic viability) + Ch.2 Research Gap synthesis. Methodology: 525-module governance taxonomy + 10-framework regulatory mapping + Pipeline~2 expert panel.
    \item \textbf{Directive RQ-B — How does explainability influence caregiver and clinician trust in ASD AI systems?} Answered by RQ4 (RAG explainability) + RQ6 (Biomarker/SHAP attribution) + the PLS-SEM mediation chain (explainability $\to$ trust $\to$ adoption). Methodology: 12-expert panel + 131-clinician PLS-SEM survey + two-layer explainability consistency test.
    \item \textbf{Directive RQ-C — What adoption barriers affect human-centered AI integration within ASD support workflows?} Answered by RQ5 (ASD-focused validation) + RQ7 (Governance) + Pipeline~2 caregiver-trust constructs. Methodology: PLS-SEM mediation + caregiver six-construct trust model.
    \item \textbf{Directive RQ-D — How can governance and human-oversight mechanisms improve responsible AI adoption in ASD-support ecosystems?} Answered by RQ3 (Agentic AI orchestration with HITL) + RQ7 (Governance architecture) + the six-level governance maturity model + the Management Decision Matrix (Ch.~\ref{chap:discussion}, Section~\ref{sec:management_analysis}). Methodology: HITL gate architecture + Delphi-style expert review + Management Decision Matrix synthesis.
\end{itemize}

\paragraph{Reading-mode choice.} The 4 directive-aligned RQs above are the governance-evaluation reading of the dissertation; the 8 detailed RQs below are the operational reading. Both organisations point to the same evidence base; examiners can navigate via either depending on whether they want directive-themed or per-pipeline detail.

\paragraph{Hypothesis Reading Guide (examiner-facing reframing of H1--H5).}
\label{para:hypothesis_reading_guide_ch1}
Before reading the H1--H5 acceptance criteria below, examiners are referred to the same reading guide that opens Chapter~\ref{chap:discussion} (Section~\ref{para:ch5_findings_reading_guide}). Each hypothesis verdict in this dissertation is \textbf{governance, explainability, and adoption evidence for the RGAIG framework}, NOT clinical-future-validation evidence. Specifically, before the formal table:

\begin{itemize}[leftmargin=*, itemsep=3pt]
    \item \textbf{H1 (classification accuracy $\geq 85$\%):} \emph{Retrospective benchmark on de-identified public datasets (KAU Saudi Arabia + ABIDE-II Western)}. Evidence that the analytical layer of the framework can meet the conceptual accuracy threshold under controlled conditions. NOT prospective clinical evidence; NOT validation of a diagnostic claim. Prospective patient cohort recruitment is identified as future work (Section~\ref{subsec:data_strategy_existing_vs_new}).
    \item \textbf{H2 (DL outperforms classical baselines):} \emph{Model-vs-model technical comparison on benchmark data}. NOT a claim that the framework outperforms clinicians.
    \item \textbf{H3 (Time-frequency + radiomic features enhance classification):} \emph{Explainability + feature-importance evidence}. SHAP top-features (theta asymmetry, gamma coherence, beta/theta ratio) surface what drove the prediction --- this is explainability-acceptance evidence, NOT biomarker discovery validated for clinical use.
    \item \textbf{H4 (Agentic AI reduces manual effort $\geq 90$\%):} \emph{Governance evidence that HITL-gated workflow preserves clinician authority while reducing manual effort}. NOT a claim of autonomous medical decision-making (every workflow path passes through clinician confirmation by design at Layer~4 of RGAIG).
    \item \textbf{H5 (RAG-based explainability $\geq 80$\% expert agreement + $\geq 90$\% citation accuracy):} \emph{Stakeholder perception evidence for explainability acceptance}. 12-expert panel rating ($M = \expertM{}$, $\alpha = 0.84$) is directional, not population-generalisable. The finding is that experts perceive citation-grounded RAG explanations as more interpretable than score-only outputs.
\end{itemize}

\noindent\textbf{Headline.} The classification accuracy figure (H1) is necessary \emph{supporting evidence} that the analytical layer functions correctly; it is \textbf{not} the dissertation's contribution claim. The contribution is the governance + explainability + adoption + HITL framework (validated through H3 + H4 + H5 + the PLS-SEM mediation chain). H1 + H2 are supporting technical evidence; H3 + H4 + H5 are the DBA contribution evidence.

Table~\ref{tab:hypotheses} formalises five testable hypotheses with explicit acceptance criteria grounded in established statistical thresholds.

\paragraph{Research Hypotheses with.}




{\tablestripes\tablefontsize
\needspace{10\baselineskip}
\begin{xltabular}{\textwidth}{@{} >{\raggedright\arraybackslash}p{0.5cm} L L >{\raggedright\arraybackslash}p{0.8cm} @{}}
\caption[Research Hypotheses with Acceptance Criteria]{Research Hypotheses with Acceptance Criteria. Each hypothesis specifies a falsifiable threshold (e.g., accuracy $\geq$85\%, latency $\leq$30 minutes) and its linked research question, so that Chapter~4 can report a clear accept/reject decision for every hypothesis.}\label{tab:hypotheses} \\
\toprule
\thead{H\#} & \thead{Statement} & \thead{Criterion} & \thead{RQ} \\
\midrule
\endfirsthead
\endhead
\bottomrule
\endlastfoot
H1 & The unified pipeline achieves $\geq$85\% classification accuracy in at least four of the Autism Spectrum Disorder (ASD) & Accuracy $\geq$ 85\% in $\geq$4 domains & RQ1 \\
\addlinespace
H2 & Deep learning models significantly outperform classical baseline models & Paired \textit{t}-test \textit{p} < .05; Cohen's \textit{d} > 0.5 & RQ2 \\
\addlinespace
H3 & Time-frequency and radiomic features enhance classification beyond baseline & Ablation: $\geq$5\% accuracy drop without key features & RQ3 \\
\addlinespace
H4 & Agentic AI reduces manual effort $\geq$90\% without degrading accuracy & Effort reduction $\geq$90\%; accuracy variation $\leq$2\% & RQ4 \\
\addlinespace
H5 & RAG-based explainability achieves $\geq$80\% expert agreement and $\geq$90\% citation accuracy & Expert agreement $\geq$80\%; citation accuracy $\geq$90\% & RQ5 \\
\end{xltabular}}
\vspace{0.5em}
\noindent\textit{Note.} For the full P$\to$G$\to$O$\to$RQ$\to$H$\to$C traceability chain, see Table~\ref{tab:research_flow_alignment}.
\noindent\textit{Note.} All thresholds are empirically grounded. The ``at least four of five'' criterion (H1) reflects the expectation that a ASD-focused framework will not achieve uniform accuracy across structurally diverse data modalities (EEG time-series vs.\ medical imaging vs.\ physiological signals); the threshold remains falsifiable---a pipeline passing only three domains would falsify H1. The \textit{p} < .05 and Cohen's \textit{d} > 0.5 thresholds (H2) follow standard statistical practice. The 90\% effort reduction (H4) is an empirically grounded threshold based on reported automation benefits.

\subsection{Research Approach: Investigative Sequence}

\noindent The research progression follows a straightforward sequence: problem definition in Chapter~1, gap justification in Chapter~2, design and method in Chapter~3, empirical evaluation in Chapter~4, and interpretation plus contribution assessment in Chapter~5. The detailed investigative sequence is therefore treated as a structural reading logic rather than as a standalone Chapter~1 control table.

\phantomsection\label{tab:research_verbs}

%% SUPPRESSED — redundant with tab:research_verbs (Research Investigative Verbs table)
\iffalse
\subsection{Research Process}

Table~\ref{tab:research_process} maps the six-stage research process---from problem formulation through data interpretation---to the specific methods, outputs, and thesis chapters that implement each stage.


\iffalse % AUTO-SUPPRESS non-ASD
{\tablestripes\tablefontsize
\needspace{10\baselineskip}
\begin{xltabular}{\textwidth}{@{} >{\raggedright\arraybackslash}p{1.0cm} >{\raggedright\arraybackslash}p{2.2cm} p{5.0cm} p{4.6cm} >{\raggedright\arraybackslash}p{2.0cm} @{}}
\caption[Research Process Stages]{Six-Stage Research Process: From Problem to Interpretation}\label{tab:research_process} \\
\toprule
\thead{\#} & \thead{Process Stage} & \thead{Method and Approach} & \thead{Output} & \thead{Chapter} \\
\midrule
\endfirsthead
\endhead
\bottomrule
\endlastfoot
1 & Research Problem & Stakeholder analysis; industry gap assessment; literature-based problem framing & Four named problems (P1--P4); six research gaps (G1--G6); five objectives (O1--O5) & Ch.~1, 2 \\
\addlinespace
2 & Choose Research Approach & Pragmatist philosophy; Design Science Research (DSR); mixed-methods (quantitative + qualitative) & Research strategy justified; DSR six-activity model selected; multi-method validation plan & Ch.~3 \\
\addlinespace
3 & Design Research Method & Stratified \textit{k}-fold CV; ablation study; SHAP feature analysis; expert panel protocol; scenario testing & Analytical pipeline design; 12-expert panel composition; evaluation metrics specified & Ch.~3 \\
\addlinespace
4 & Collect Data & Five public benchmark datasets (ABIDE, Iowa PD, DEAP/WESAD Stress, BCI Competition, TCIA Cancer); 19 ASD subjects & Cleaned, preprocessed datasets; 305~GB total; feature-engineered matrices across ASD & Ch.~3, 4 \\
\addlinespace
5 & Analyse Data & 10-fold CV; paired \textit{t}-tests; Cohen's \textit{d}; bootstrap CI; RAGAS evaluation; expert panel scoring & Hypothesis test results (H1--H5); ASD classification accuracy; effort reduction metrics & Ch.~4 \\
\addlinespace
6 & Interpret Results & Triangulation of quantitative and qualitative evidence; comparison with published benchmarks; theoretical contribution mapping & Eight contributions (C1--C8); limitations; governance framework; future research agenda & Ch.~5 \\
\end{xltabular}}
\vspace{0.5em}
\noindent\textit{Note.} The process follows the standard empirical research cycle: problem $\to$ approach $\to$ method $\to$ data $\to$ analysis $\to$ interpretation. Each stage is detailed in the referenced chapter. DSR = Design Science Research; CV = cross-validation; SHAP = Shapley additive explanations; CI = confidence interval; RAGAS = Retrieval-Augmented Generation Assessment; ABIDE = Autism Brain Imaging Data Exchange; DEAP = Database for Emotion Analysis using Physiological Signals; WESAD = Wearable Stress and Affect Detection; ;
\fi % AUTO-SUPPRESS non-ASD
\fi

\subsection{DBA Evidence Chain}

\noindent The DBA logic of the dissertation is simple: a bounded business problem is defined first, relevant theory is then used to explain why that problem persists, and empirical evidence is subsequently used to test whether the proposed framework addresses it. In the final thesis structure, that chain is distributed across Chapters~1--5 rather than repeated as a standalone Chapter~1 dashboard.

\phantomsection\label{tab:evidence_chain_dba}

\iffalse
%% SUPPRESSED: Doctoral Quality Chain — generic framework, not dissertation-specific content
\subsection{Doctoral Contribution Quality Chain}

Table~\ref{tab:doctoral_quality_chain} presents the four-pillar quality chain that governs the progression from problem identification to doctoral contribution. Each pillar is a necessary condition: a clear problem without strong theory produces consulting advice; strong theory without valid methodology produces speculation; valid methodology without rigorous analysis produces descriptive statistics. Only the complete chain---clear problem, strong theory, valid methodology, and rigorous analysis---produces a defensible doctoral contribution.


\noindent Table~\ref{tab:doctoral_quality_chain} doctoral contribution quality chain.

{\tablestripes\tablefontsize
\needspace{10\baselineskip}
\begin{xltabular}{\textwidth}{@{} >{\raggedright\arraybackslash}p{1.0cm} >{\raggedright\arraybackslash}p{3.1cm} p{3.6cm} p{3.9cm} p{2.4cm} @{}}
\caption[Doctoral Contribution Quality Chain]{Doctoral Contribution Quality Chain: Four Pillars from Problem to Contribution}\label{tab:doctoral_quality_chain} \\
\toprule
\thead{\#} & \thead{Pillar} & \thead{Requirement} & \thead{This Dissertation} & \thead{Without This Pillar} \\
\midrule
\endfirsthead
\endhead
\bottomrule
\endlastfoot
1 & Clear Problem & Practitioner-relevant; bounded; empirically addressable & P1--P4: four named problems with stakeholder impact (Table~\ref{tab:problem_impact}) & No research direction \\
\addlinespace
2 & Strong Theory & Multiple theoretical lenses; testable predictions; falsifiable hypotheses & Nine theoretical frameworks; mediation-based conceptual model; H0--H7 hypotheses & Consulting advice, not scholarship \\
\addlinespace
3 & Valid Methodology & Appropriate design; reliable instruments; robust validation strategy & DSR with mixed methods; stratified \textit{k}-fold CV; 4-method validation; 12-expert panel & Unfounded claims \\
\addlinespace
4 & Rigorous Analysis & Statistical testing; effect sizes; triangulation; honest limitations & Paired \textit{t}-tests (\textit{p} < .01); Cohen's \textit{d} = 0.82--1.47; bootstrap CI; failure analysis & Descriptive statistics only \\
\midrule
\multicolumn{2}{l}{\textbf{= Doctoral Contribution}} & \multicolumn{3}{>{\raggedright\arraybackslash}p{0.72\textwidth}}{Eight contributions (C1--C8) across academic, managerial, and practical domains; four explicit non-claims; two untested paths (H6, H7) reserved for future research} \\
\end{xltabular}}
\vspace{0.5em}
\noindent\textit{Note.} The quality chain is cumulative: each pillar depends on all preceding pillars. The ``Without This Pillar'' column illustrates what happens when any single element is missing, demonstrating that all four are necessary conditions for a defensible doctoral contribution. DSR = Design Science Research; CV = cross-validation; CI = confidence interval.
\fi

\iffalse %% Detailed variable taxonomy — methodology detail for Ch.3
\subsection{Research Variables Taxonomy}


\noindent This table lists each \textit{role} across variable, operationalisation, measurement, and 1 more dimension, so examiners can trace what was assessed and what evidence supports each row.



\noindent This table lists each \textit{role} across variable, operationalisation, measurement, and 1 more dimensions, so examiners can trace what was assessed and what evidence supports each row.


\noindent This table lists each \textit{role} across variable, operationalisation, measurement, and 1 more dimensions, so examiners can trace what was assessed and what evidence supports each row.



{\tablestripes\tablefontsize
\needspace{10\baselineskip}
\begin{xltabular}{\textwidth}{@{} >{\raggedright\arraybackslash}p{1.0cm} >{\raggedright\arraybackslash}p{2.5cm} p{4.9cm} p{4.6cm} >{\raggedright\arraybackslash}p{1.0cm} @{}}
\caption[Variable Taxonomy for RGAIG Framework]{Variable Taxonomy for RGAIG Framework: Roles, Operationalisation, and Hypothesis Linkage}\label{tab:variable_taxonomy} \\
\toprule
\thead{Role} & \thead{Variable} & \thead{Operationalisation} & \thead{Measurement} & \thead{Hyp.} \\
\midrule
\endfirsthead
\endhead
\bottomrule
\endlastfoot
\multicolumn{5}{l}{\textit{Independent Variables (IV): Framework Implementation}} \\
\addlinespace
IV & Pipeline Completeness & Number of RGAIG pipeline stages implemented (data ingestion, preprocessing, feature engineering, model training, evaluation, deployment) out of six & Stage completion ratio (0--1); binary per-stage pass/fail & H1, H2 \\
\addlinespace
IV & Governance Maturity & Degree of alignment with NIST AI RMF and the 10 cross-cutting governance pillars (525-module taxonomy) & Module pass rate (\%); pillar coverage score (0--10); expert panel Likert rating (1--5) & H1, H5 \\
\addlinespace
IV & Explainability Integration & Extent to which RAG-based explanations and SHAP attributions are embedded in the diagnostic output & RAG citation accuracy (\%); SHAP feature coverage (\%); explanation completeness score & H5 \\
\midrule
\multicolumn{5}{l}{\textit{Mediator 1 (M1): Organisational Capability}} \\
\addlinespace
M1 & Technical Infrastructure Readiness & Availability of compute, storage, and orchestration resources to execute the unified pipeline & GPU hours available; pipeline execution time (minutes); infrastructure uptime (\%) & H1, H2 \\
\addlinespace
M1 & Data Quality & Completeness, consistency, and preprocessing quality of input datasets across the ASD domain & Missing-value rate (\%); class imbalance ratio; signal-to-noise ratio (dB for EEG/BCI) & H3 \\
\addlinespace
M1 & Model Evaluation Capacity & Organisation's ability to conceptually evaluate, monitor, and update trained models in clinical workflows & Evaluation success rate (\%); model refresh frequency; monitoring coverage & H4 \\
\midrule
\multicolumn{5}{l}{\textit{Mediator 2 (M2): Clinician Capability}} \\
\addlinespace
M2 & Explainability Comprehension & Clinician's ability to interpret RAG-generated explanations and SHAP feature attributions & Expert agreement rate (\%); comprehension quiz score; time-to-understand (minutes) & H5 \\
\addlinespace
M2 & Trust in AI Output & Clinician confidence in AI-generated diagnoses sufficient to act on recommendations & Trust scale (Likert 1--5); override rate (\%); adoption willingness score & H4, H5 \\
\addlinespace
M2 & Workflow Integration & Degree to which AI outputs are embedded in existing clinical decision pathways & Manual steps eliminated (\%); effort reduction (\%); time savings per case (minutes) & H4 \\
\midrule
\multicolumn{5}{l}{\textit{Dependent Variables (DV): Business Value}} \\
\addlinespace
DV & Diagnostic Accuracy & Classification correctness of the unified pipeline across Autism Spectrum Disorder (ASD) & Accuracy (\%); F1-score; AUC-ROC; sensitivity; specificity; per-ASD 95\% CI & H1, H2, H3 \\
\addlinespace
DV & Time-to-Diagnosis & Elapsed time from data acquisition to actionable diagnostic output & Minutes per case; throughput (cases/hour); end-to-end latency & H4 \\
\addlinespace
DV & Cost-per-Case & Total resource cost to produce one diagnostic decision using the AI pipeline & USD per case; cost reduction vs.\ manual baseline (\%); illustrative cost-modelling scenario over 3 years & H4 \\
\addlinespace
DV & Illustrative financial scenario (modelled only) & Net financial benefit of RGAIG evaluation relative to implementation cost & illustrative financial scenario (\%); illustrative payback estimate (months); net present value (USD) & H4 \\
\midrule
\multicolumn{5}{l}{\textit{Moderators: Domain and Context Factors}} \\
\addlinespace
Moderator & Domain Type & Biomedical domain in which the pipeline operates & Categorical: ASD & H1 \\
\addlinespace
Moderator & Dataset Size & Number of subjects and samples available for training and validation & \textit{N} subjects; \textit{n} samples; train/test split ratio & H2, H3 \\
\addlinespace
Moderator & Clinical Complexity & Heterogeneity of presentation and diagnostic difficulty within a domain & Number of diagnostic subtypes; comorbidity prevalence; signal variability (SD) & H1, H3 \\
\addlinespace
Moderator & Regulatory Requirements & Stringency of governance and compliance obligations applicable to the domain & Regulatory tier (FDA Class I/II/III); required audit depth; consent complexity & H5 \\
\addlinespace
Control & Dataset Characteristics & Factors held constant or adjusted across all the ASD domain to ensure fair comparison & Sample size, class balance, signal quality, preprocessing protocol, hardware specification & All \\
\end{xltabular}}
\vspace{0.5em}
\noindent\textit{Note.} IV = independent variable; DV = dependent variable; M1 = organisational capability mediator; M2 = clinician capability mediator. The moderators condition the strength of IV$\to$M1$\to$DV and IV$\to$M2$\to$DV pathways; domain type and clinical complexity are empirically tested in Chapter~\ref{chap:analysis}, while regulatory requirements are theoretically specified. Control variables are standardised across domains to ensure fair comparison. DL = deep learning; RAG = retrieval-augmented generation; SHAP = Shapley additive explanations; illustrative financial scenario = return on investment; illustrative cost-modelling scenario = total cost of ownership; NIST = National Institute of Standards and Technology; FDA = Food and Drug Administration.


\noindent The five-role taxonomy confirms that every hypothesis in Table~\ref{tab:hypotheses} traces to at least one named variable with a defined operationalisation and measurement instrument, closing the gap between abstract constructs and empirical evidence that examiners typically scrutinise in variable-weak dissertations.
\fi %% End suppressed variable taxonomy

\phantomsection\label{tab:variable_taxonomy}

\iffalse %% Methodology detail — covered in Ch.3
\subsection{Measurement and Quality Assurance Framework}
\label{sec:measurement_quality}

Table~\ref{tab:measurement_quality_summary} consolidates measurement scales, reliability, validity, and error mitigation into a single reference.


{\tablestripes\tablefontsize

\needspace{10\baselineskip}
\begin{xltabular}{\textwidth}{@{} >{\raggedright\arraybackslash}p{2.2cm} L L @{}}
\caption[Measurement and Quality Assurance]{Measurement and Quality Assurance Framework: Scales, Reliability, Validity, and Error Mitigation}\label{tab:measurement_quality_summary} \\
\toprule
\thead{Dimension} & \thead{Components} & \thead{Evidence in This Dissertation} \\
\midrule
\endfirsthead
\endhead
\bottomrule
\endlastfoot
\multicolumn{3}{l}{\textit{\textbf{Measurement Scales (Stevens' Four Levels)}}} \\
\addlinespace
Nominal & Disease labels, model type, dataset source & Confusion matrix; $\chi^2$ test \\
Ordinal & Expert Likert ratings (1--5); KPI thresholds & Median; inter-rater agreement; Kendall's $W$ \\
Interval & Normalised EEG values; RAGAS scores (0--1) & Paired \textit{t}-test; ANOVA; Pearson correlation; bootstrap CI \\
Ratio & Accuracy (\%); effort reduction (\%); cost (\$); illustrative financial scenario (\%) & Cohen's \textit{d}; sensitivity analysis \\
\midrule
\multicolumn{3}{l}{\textit{\textbf{Reliability Assessment}}} \\
\addlinespace
Internal consistency & Cronbach's $\alpha$ on expert panel 5-dimension Likert instrument & $\alpha \geq 0.80$ across evaluation dimensions \\
Temporal stability & 10-fold CV repeated across random seeds; bootstrap (1,000 resamples) & Accuracy variation $\leq$2\% across folds; CI width $<$5\,pp \\
Inter-rater agreement & 12-expert panel independent blinded scoring & Consensus threshold $\geq$80\%; results in Chapter~\ref{chap:analysis} \\
\midrule
\multicolumn{3}{l}{\textit{\textbf{Validity Assessment}}} \\
\addlinespace
Content validity & Expert panel review of evaluation dimensions; literature-based instrument design & 12 experts reviewed; 5 dimensions from TAM, NIST AI RMF, FDA SaMD \\
Construct validity & Convergent evidence from multiple methods; discriminant analysis & Multi-method validation; SHAP discriminates features across domains \\
Criterion validity & Accuracy benchmarked against published studies; governance against NIST/FDA & Per-domain benchmarks (Table~\ref{tab:literature_comparison}); Chapter~\ref{chap:analysis} \\
\midrule
\multicolumn{3}{l}{\textit{\textbf{Measurement Error Mitigation}}} \\
\addlinespace
Systematic error & Standardised preprocessing; calibrated instruments; benchmark comparison & ASD-focused protocol consistency \\
Random error & 10-fold CV; bootstrap 95\% CI; multiple random seeds & Variance reduction through replication \\
Acquiescence bias & Balanced evaluation criteria (positive and negative dimensions) & Independent panelist scoring \\
Social desirability & Blinded protocol; anonymous scoring; independent before group & Experts unaware which outputs were AI vs.\ baseline \\
\end{xltabular}}
\vspace{0.5em}
\noindent\textit{Note.} This table consolidates four methodological dimensions. Stevens' four measurement levels determine valid statistical operations. Three reliability types and three validity types are addressed following standard doctoral research practice. Four error types are mitigated through design choices documented in Chapter~\ref{chap:methodology}. CI = confidence interval; pp = percentage points; CV = cross-validation; TAM = Technology Acceptance Model; NIST AI RMF = National Institute of Standards and Technology AI Risk Management Framework; SaMD = Software as a Medical Device; SHAP = Shapley additive explanations; RAGAS = Retrieval-Augmented Generation Assessment; KPI = key performance indicator; illustrative financial scenario = return on investment.

\fi %% End suppressed measurement and quality assurance

%% Preserve original labels for cross-reference compatibility
\phantomsection\label{sec:measurement_quality}
\phantomsection\label{tab:measurement_quality_summary}
\phantomsection\label{tab:measurement_scales}
\phantomsection\label{tab:reliability_assessment}
\phantomsection\label{tab:validity_assessment}
\phantomsection\label{tab:measurement_errors}

\iffalse %% Consolidated into tab:measurement_quality_summary (Measurement and Quality Assurance table)
\subsection{Measurement Scales}

Table~\ref{tab:measurement_scales} classifies the data types used in this dissertation according to Stevens' four levels of measurement---nominal, ordinal, interval, and ratio---identifying the appropriate statistical tests for each level.


\noindent Table~\ref{tab:measurement_scales} measurement scales.

{\tablestripes\tablefontsize
\needspace{10\baselineskip}
\begin{xltabular}{\textwidth}{@{} >{\raggedright\arraybackslash}p{1.5cm} L L L @{}}
\caption[Measurement Scales]{Levels of Measurement: Data Types and Statistical Tests Used in This Dissertation}\label{tab:measurement_scales} \\
\toprule
\thead{Scale} & \thead{Properties} & \thead{Data in This Dissertation} & \thead{Statistical Tests Applied} \\
\midrule
\endfirsthead
\endhead
\bottomrule
\endlastfoot
Nominal & Categories without order & Disease labels (ASD); model type (baseline vs.\ DL); dataset source & Confusion matrix; $\chi^2$ test; mode \\
\addlinespace
Ordinal & Ordered categories & Expert Likert ratings (1--5); governance maturity levels; KPI threshold categories (green/amber/red) & Expert panel median; inter-rater agreement; Kendall's $W$ \\
\addlinespace
Interval & Equal intervals, no true zero & Normalised EEG signal values; standardised feature scores; RAGAS evaluation scores (0--1) & Paired \textit{t}-test; ANOVA; Pearson correlation; bootstrap CI \\
\addlinespace
Ratio & Equal intervals + true zero & Classification accuracy (\%); effort reduction (\%); cost savings (\$); time (seconds); illustrative financial scenario (\%) & Cohen's \textit{d}; percentage comparison; ratio analysis; sensitivity analysis \\
\end{xltabular}}
\vspace{0.5em}
\noindent\textit{Note.} Stevens' four levels of measurement determine which statistical operations are valid. Nominal data supports only mode and frequency; ordinal adds median and rank tests; interval adds mean and parametric tests; ratio adds geometric mean and ratio comparisons. ASD = autism spectrum disorder;   DL = deep learning; CI = confidence interval; RAGAS = Retrieval-Augmented Generation Assessment; KPI = key performance indicator; illustrative financial scenario = return on investment.

\subsection{Reliability Assessment}

Table~\ref{tab:reliability_assessment} maps the three primary reliability assessment methods---Cronbach's alpha, test--retest, and inter-rater reliability---to their application in this dissertation.


{\tablestripes\tablefontsize
\needspace{10\baselineskip}
\begin{xltabular}{\textwidth}{@{} >{\raggedright\arraybackslash}p{2.5cm} p{3.4cm} p{4.6cm} p{3.7cm} @{}}
\caption[Reliability Assessment Methods]{Reliability Assessment: Three Methods Applied in This Dissertation}\label{tab:reliability_assessment} \\
\toprule
\thead{Reliability Type} & \thead{Definition} & \thead{Application in This Study} & \thead{Evidence} \\
\midrule
\endfirsthead
\endhead
\bottomrule
\endlastfoot
Cronbach's Alpha (Internal Consistency) & Measures whether items within a scale produce consistent scores & Expert panel evaluation instrument: 5-dimension Likert scale (accuracy, explainability, governance, clinical utility, overall quality) & $\alpha \geq$ 0.80 across evaluation dimensions; high internal consistency of expert ratings \\
\addlinespace
Test--Retest (Temporal Stability) & Measures whether the same test produces consistent results over time & 10-fold cross-validation repeated across random seeds; bootstrap resampling (1,000 iterations) with 95\% confidence intervals & Accuracy variation $\leq$2\% across folds; bootstrap CI width $<$5\,pp for all domains \\
\addlinespace
Inter-Rater (Agreement) & Measures whether different raters produce consistent assessments & 12-expert panel independent scoring; blinded evaluation of RAG outputs; consensus threshold $\geq$80\% & Expert agreement assessed across all dimensions; results reported in Chapter~\ref{chap:analysis} \\
\end{xltabular}}
\vspace{0.5em}
\noindent\textit{Note.} All three reliability types are addressed in this dissertation. Cronbach's alpha measures internal consistency of the expert evaluation instrument; test--retest is operationalised through repeated cross-validation and bootstrap resampling; inter-rater reliability is assessed through independent expert panel scoring. CI = confidence interval; pp = percentage points; RAG = retrieval-augmented generation.

\subsection{Validity Assessment}

Table~\ref{tab:validity_assessment} classifies the three primary validity types---content, construct, and criterion---and maps each to the evidence provided in this dissertation.


{\tablestripes\tablefontsize

\needspace{10\baselineskip}
\begin{xltabular}{\textwidth}{@{} >{\raggedright\arraybackslash}p{2.0cm} p{3.6cm} p{4.3cm} p{4.5cm} @{}}
\caption[Validity Assessment]{Validity Assessment: Three Types Applied in This Dissertation}\label{tab:validity_assessment} \\
\toprule
\thead{Validity Type} & \thead{Definition} & \thead{Application in This Study} & \thead{Evidence} \\
\midrule
\endfirsthead
\endhead
\bottomrule
\endlastfoot
Content Validity & Extent to which the measurement instrument covers the full domain of the construct being measured & Expert panel review of evaluation dimensions; literature-based instrument design; comprehensive coverage of accuracy, explainability, governance, and clinical utility & 12 domain experts reviewed instrument completeness; 5 evaluation dimensions derived from published frameworks (TAM, NIST AI RMF, FDA SaMD) \\
\addlinespace
Construct Validity & Extent to which the instrument measures the theoretical construct it claims to measure & Convergent evidence from multiple methods; discriminant analysis between constructs & Multi-method validation (statistical + expert + scenario + before--after); SHAP analysis discriminates features across domains; expert ratings converge with statistical results \\
\addlinespace
Criterion Validity & Extent to which the instrument's scores predict a criterion or agree with an established standard & Classification accuracy benchmarked against published domain-specific studies; governance alignment against NIST AI RMF and FDA SaMD standards & Per-domain accuracy benchmarked against published results (Table~\ref{tab:literature_comparison}); NIST AI RMF alignment assessed by expert panel (Chapter~\ref{chap:analysis}) \\
\end{xltabular}}
\vspace{0.5em}
\noindent\textit{Note.} All three validity types are addressed to ensure the dissertation's measurement instruments and findings are defensible. Content validity ensures comprehensive coverage; construct validity ensures theoretical alignment; criterion validity ensures predictive accuracy against established benchmarks. TAM = Technology Acceptance Model; NIST AI RMF = National Institute of Standards and Technology AI Risk Management Framework; SaMD = Software as a Medical Device; SHAP = Shapley additive explanations.

\subsection{Measurement Errors and Mitigation}

Table~\ref{tab:measurement_errors} identifies four types of measurement errors and biases relevant to this dissertation and describes the mitigation strategies employed.


{\tablestripes\tablefontsize

\needspace{10\baselineskip}
\begin{xltabular}{\textwidth}{@{} >{\raggedright\arraybackslash}p{2.5cm} p{4.8cm} p{7.0cm} @{}}
\caption[Measurement Errors and Mitigation]{Measurement Errors and Biases: Types and Mitigation Strategies}\label{tab:measurement_errors} \\
\toprule
\thead{Error/Bias Type} & \thead{Description and Risk} & \thead{Mitigation in This Dissertation} \\
\midrule
\endfirsthead
\endhead
\bottomrule
\endlastfoot
Systematic Error & Consistent bias that shifts all measurements in one direction; threatens internal validity & Standardised preprocessing protocols across domains; calibrated evaluation instruments; benchmark comparison against published results \\
\addlinespace
Random Error & Unpredictable variation that reduces measurement precision & 10-fold cross-validation reduces variance; bootstrap resampling (1,000 iterations) with 95\% CI; multiple random seeds tested; repeated trials \\
\addlinespace
Acquiescence Bias & Tendency of respondents to agree with statements regardless of content & Expert panel used balanced evaluation criteria (positive and negative dimensions); governance rating includes explicit limitations assessment; panelists scored independently \\
\addlinespace
Social Desirability Bias & Tendency to provide answers perceived as favourable rather than truthful & Blinded evaluation protocol; experts did not know which outputs were AI-generated vs.\ baseline; anonymous scoring sheets; independent assessment before group discussion \\
\end{xltabular}}
\vspace{0.5em}
\noindent\textit{Note.} Systematic and random errors affect quantitative measurements (classification accuracy, effort reduction); acquiescence and social desirability biases affect qualitative measurements (expert panel ratings). Both categories are addressed through design choices documented in Chapter~\ref{chap:methodology}. CI = confidence interval.
\fi %% End consolidated Measurement Scales, Reliability, Validity, Measurement Errors

\iffalse %% Methodology detail — covered in Ch.3
\subsection{Data Collection, Operationalisation, and Analysis Chain}
\label{sec:data_operationalisation}

Table~\ref{tab:data_analysis_chain} consolidates data collection, analytical purpose, operationalisation, and concept-to-analysis tracing into a single traceability framework.


\needspace{12\baselineskip}
\iffalse % AUTO-SUPPRESS non-ASD
\noindent Table~\ref{tab:data_analysis_chain} data collection and analysis chain.

{\tablestripes\tablefontsize
\begin{xltabular}{\textwidth}{@{} >{\raggedright\arraybackslash}p{2.2cm} L L @{}}
\caption[Data Collection and Analysis Chain]{Data Collection, Operationalisation, and Analysis Chain}\label{tab:data_analysis_chain} \\
\toprule
\thead{Dimension} & \thead{Components} & \thead{Implementation} \\
\midrule
\endfirsthead
\endhead
\bottomrule
\endlastfoot
\multicolumn{3}{l}{\textit{\textbf{Data Collection Methods (Four Sources of Evidence)}}} \\
\addlinespace
Survey/questionnaire & 5-dimension Likert scale (1--5) & 12 expert panel ratings per dimension (Ch.~\ref{chap:analysis}) \\
Expert panel & Semi-structured protocol; 12 experts ($\geq$5 yr experience) & Qualitative themes; consensus ratings (Ch.~\ref{chap:analysis}) \\
Financial data & Published cost benchmarks; vendor pricing; illustrative cost-modelling scenario modelling & Projected cost reduction, ROI, payback (Ch.~\ref{chap:discussion}) \\
Sensors/EEG data & ABIDE, Iowa PD, DEAP/WESAD, BCI, TCIA; Emotiv EPOC X & 305~GB; 131 subjects + 20K cancer images; ASD \\
\midrule
\multicolumn{3}{l}{\textit{\textbf{Analytical Purposes}}} \\
\addlinespace
Hypothesis testing & H1--H5 with acceptance criteria; paired \textit{t}-tests; effect sizes & Verdicts, significance levels (Ch.~\ref{chap:analysis}) \\
Model building & DL classifiers (CNN-LSTM, ASD ensemble); baselines (SVM, RF, XGBoost) & 198 trained models; per-ASD comparison; ablation \\
Prediction & ASD-focused classification; RAG-generated explanations & Per-domain accuracy on held-out sets (Ch.~\ref{chap:analysis}) \\
\midrule
\multicolumn{3}{l}{\textit{\textbf{Operationalisation (Eight Stages)}}} \\
\addlinespace
1. Define constructs & Five constructs: accuracy (DV), explainability (M2), automation (IV), governance (Mod), value (DV) & Theory-driven \\
2. Conceptualise & Map constructs to observables & Accuracy $\to$ metrics; explainability $\to$ ratings \\
3. Operationalise & Define measurement instruments & \textit{k}-fold CV; Likert scale; effort before--after \\
4. Design instrument & Expert rubric (5 dim $\times$ 5-pt); RAGAS protocol & Data collection tools \\
5. Ensure reliability & Cronbach's $\alpha \geq 0.80$; CV variance $\leq$2\%; inter-rater $\geq$80\% & Consistency verified \\
6. Ensure validity & Content, construct, criterion validity & Expert review; multi-method; benchmarks \\
7. Determine sampling & Stratified within datasets; purposive for panel & 5 public datasets; 12 experts \\
8. Calculate sample & 131 subjects + 20K images; 12 panellists & Exceeds Delphi minimum \\
\midrule
\multicolumn{3}{l}{\textit{\textbf{Concept-to-Analysis Chain}}} \\
\addlinespace
ASD-focused accuracy & Stratified 10-fold CV (\%) [Ratio] & Paired \textit{t}-test; Cohen's \textit{d}; bootstrap 95\% CI \\
DL vs.\ baseline & Paired accuracy difference [Interval] & \textit{t}-test; effect size; signed-rank test \\
Feature importance & SHAP scores [Ratio] & Ablation ($\Delta$ accuracy); permutation importance \\
Explainability & Expert Likert rating [Ordinal] & Median; Kendall's $W$; thematic analysis \\
Organisational value & Cost reduction, illustrative financial scenario [Ratio] & illustrative cost-modelling scenario comparison; sensitivity analysis; payback \\
\end{xltabular}}
\vspace{0.5em}
\noindent\textit{Note.} This table consolidates four methodological processes into a single traceability framework. Each concept traces from theory through measurement to statistical test. IV = independent variable; DV = dependent variable; M2 = mediator; Mod = moderator; CV = cross-validation; CI = confidence interval; illustrative cost-modelling scenario = total cost of ownership; illustrative financial scenario = return on investment; DL = deep learning; SHAP = Shapley additive explanations; RAGAS = Retrieval-Augmented Generation Assessment; ASD classification accuracy.
\fi % AUTO-SUPPRESS non-ASD

\fi %% End suppressed data collection chain

%% Preserve original labels for cross-reference compatibility
\phantomsection\label{sec:data_operationalisation}
\phantomsection\label{tab:data_analysis_chain}
\phantomsection\label{tab:data_collection_methods}
\phantomsection\label{tab:analytical_purpose}
\phantomsection\label{tab:operationalisation}
\phantomsection\label{tab:concept_to_analysis}

\iffalse %% Consolidated into tab:data_analysis_chain (Data Collection, Operationalisation, and Analysis Chain)
\subsection{Data Collection Methods}

\needspace{12\baselineskip}
\iffalse % AUTO-SUPPRESS non-ASD
\noindent Table~\ref{tab:data_collection_methods} data collection methods.

{\tablestripes\tablefontsize
\begin{xltabular}{\textwidth}{@{} >{\raggedright\arraybackslash}p{2.2cm} p{3.3cm} p{4.5cm} p{4.2cm} @{}}
\caption[Data Collection Methods]{Data Collection Methods: Four Sources of Evidence}\label{tab:data_collection_methods} \\
\toprule
\thead{Method} & \thead{Purpose} & \thead{Instrument/Source} & \thead{Output} \\
\midrule
\endfirsthead
\endhead
\bottomrule
\endlastfoot
Survey / Questionnaire & Expert evaluation of framework quality & 5-dimension Likert scale (1--5); structured evaluation rubric & Expert panel ratings; 12 independent evaluations per dimension (Chapter~\ref{chap:analysis}) \\
\addlinespace
Interview / Expert Panel & Qualitative validation of clinical relevance and governance & Semi-structured review protocol; 12 domain experts with $\geq$5 years experience & Qualitative themes; consensus ratings; governance recommendations (Chapter~\ref{chap:analysis}) \\
\addlinespace
Financial Data & Business case and illustrative financial scenario analysis & Published healthcare cost benchmarks; vendor pricing surveys; illustrative cost-modelling scenario modelling & Projected cost reduction, ROI, and illustrative payback estimate (Chapter~\ref{chap:discussion}) \\
\addlinespace
Sensors / EEG Data & Primary diagnostic data collection & EEG (ABIDE, Iowa PD, DEAP/WESAD, BCI); imaging (TCIA); wearable sensors (Emotiv EPOC X, ECG, PPG, EDA) & 305~GB processed data; 19 ASD subjects; ASD; feature-engineered matrices \\
\end{xltabular}}
\vspace{0.5em}
\noindent\textit{Note.} Mixed-methods data collection combines quantitative (sensor data, financial data) with qualitative (expert panel, survey) sources, fulfilling the DBA requirement for evidence triangulation. EEG = electroencephalography; ABIDE = Autism Brain Imaging Data Exchange; DEAP = Database for Emotion Analysis using Physiological Signals; WESAD = Wearable Stress and Affect Detection; ;  illustrative cost-modelling scenario = total cost of ownership; illustrative financial scenario = return on investment; ECG = electrocardiography; PPG = photoplethysmography; EDA = electrodermal activity.
\fi % AUTO-SUPPRESS non-ASD

\subsection{Analytical Purpose Taxonomy}

Table~\ref{tab:analytical_purpose} classifies the three analytical purposes---hypothesis testing, model building, and prediction---and maps each to the specific analyses conducted in this dissertation.


{\tablestripes\tablefontsize

\needspace{10\baselineskip}
\begin{xltabular}{\textwidth}{@{} >{\raggedright\arraybackslash}p{2.0cm} p{3.8cm} p{4.6cm} p{4.0cm} @{}}
\caption[Analytical Purpose Taxonomy]{Analytical Purpose: Hypothesis Testing, Model Building, and Prediction}\label{tab:analytical_purpose} \\
\toprule
\thead{Purpose} & \thead{Definition} & \thead{Application in This Study} & \thead{Key Output} \\
\midrule
\endfirsthead
\endhead
\bottomrule
\endlastfoot
Hypothesis Testing & Test whether observed data supports or rejects a theoretical prediction & H1--H5 tested with acceptance criteria (Table~\ref{tab:hypotheses}); paired \textit{t}-tests; effect sizes; expert consensus & Hypothesis verdicts, significance levels, and effect sizes reported in Chapter~\ref{chap:analysis} \\
\addlinespace
Model Building & Construct and train predictive models that capture data relationships & DL classifiers (SVM, RF, XGBoost); DL architectures (CNN-LSTM, ASD ensemble); ensemble methods (VotingClassifier) & 198 trained models; per-ASD architecture comparison; ablation study \\
\addlinespace
Prediction & Apply trained models to classify unseen data and generate actionable outputs & ASD-focused diagnostic predictions; RAG-generated clinical explanations; automated pipeline outputs & Per-domain accuracy on held-out test sets; literature-grounded clinical narratives (Chapter~\ref{chap:analysis}) \\
\end{xltabular}}
\vspace{0.5em}
\noindent\textit{Note.} All three analytical purposes are addressed sequentially: hypothesis testing validates theoretical predictions, model building creates the predictive tools, and prediction demonstrates operational utility. SVM = support vector machine; RF = random forest; CNN = convolutional neural network; LSTM = long short-term memory; ASD classification accuracy; RAG = retrieval-augmented generation.

\subsection{Operationalisation Process}

Table~\ref{tab:operationalisation} maps the eight-stage operationalisation process---from abstract constructs to measurable data---showing how each stage was implemented in this dissertation. This process ensures that theoretical concepts are systematically translated into empirical measurements.


\noindent Table~\ref{tab:operationalisation} operationalisation process.

{\tablestripes\tablefontsize
\needspace{10\baselineskip}
\begin{xltabular}{\textwidth}{@{} >{\raggedright\arraybackslash}p{1.0cm} >{\raggedright\arraybackslash}p{2.7cm} p{3.4cm} p{7.0cm} @{}}
\caption[Operationalisation Process]{Eight-Stage Operationalisation Process: From Constructs to Sample Size}\label{tab:operationalisation} \\
\toprule
\thead{\#} & \thead{Stage} & \thead{Description} & \thead{Implementation in This Dissertation} \\
\midrule
\endfirsthead
\endhead
\bottomrule
\endlastfoot
1 & Define Constructs & Identify abstract concepts from theory & Five constructs: ASD-focused accuracy (DV), clinical explainability (M2), pipeline automation (IV), governance readiness (Mod), organisational value (DV) \\
\addlinespace
2 & Conceptualise Variables & Map constructs to observable phenomena & Accuracy $\to$ classification metrics; explainability $\to$ expert ratings; automation $\to$ effort reduction; governance $\to$ compliance score; value $\to$ illustrative financial scenario \\
\addlinespace
3 & Operationalise Variables & Define specific measurement instruments & Accuracy: stratified \textit{k}-fold CV (\%); explainability: 5-point Likert scale; automation: time/effort before--after; governance: NIST AI RMF checklist \\
\addlinespace
4 & Design Survey / Instrument & Create data collection tools & Expert evaluation rubric (5 dimensions $\times$ 5-point scale); classification pipeline configuration; RAG evaluation protocol (RAGAS metrics) \\
\addlinespace
5 & Ensure Reliability & Verify measurement consistency & Cronbach's $\alpha \geq$ 0.80; 10-fold CV variance $\leq$2\%; inter-rater agreement $\geq$80\%; bootstrap 95\% CI \\
\addlinespace
6 & Ensure Validity & Verify measurement accuracy & Content validity (expert review); construct validity (multi-method); criterion validity (benchmark comparison) (Table~\ref{tab:validity_assessment}) \\
\addlinespace
7 & Determine Sampling & Select appropriate sampling strategy & Stratified sampling within datasets; purposive sampling for expert panel (domain expertise $\geq$5 years); 5 public benchmark datasets \\
\addlinespace
8 & Calculate Sample Size & Justify sample adequacy & 19 ASD subjects across 5 datasets; 12 expert panellists (exceeds Delphi minimum of 7); statistical power analysis for paired comparisons \\
\end{xltabular}}
\vspace{0.5em}
\noindent\textit{Note.} The operationalisation process translates abstract theoretical constructs into concrete, measurable variables through systematic stages. Each stage is documented in detail in Chapter~\ref{chap:methodology}. IV = independent variable; DV = dependent variable; M2 = mediator (Human Capability); Mod = moderator; CV = cross-validation; CI = confidence interval; NIST AI RMF = National Institute of Standards and Technology AI Risk Management Framework; RAGAS = Retrieval-Augmented Generation Assessment; illustrative financial scenario = return on investment.

\subsection{Concept-to-Analysis Chain}

Table~\ref{tab:concept_to_analysis} traces five key concepts from their theoretical definition through variable operationalisation, measurement instrument, data type, and statistical analysis method. This chain demonstrates that every empirical finding in the dissertation is traceable from an abstract concept to a specific analytical procedure.


\noindent Table~\ref{tab:concept_to_analysis} concept-to-analysis chain.

{\tablestripes\tablefontsize
\needspace{10\baselineskip}
\begin{xltabular}{\textwidth}{@{} >{\raggedright\arraybackslash}p{2.6cm} p{2.9cm} >{\raggedright\arraybackslash}p{2.9cm} >{\raggedright\arraybackslash}p{2.0cm} p{4.3cm} @{}}
\caption[Concept-to-Analysis Chain]{Concept-to-Analysis Chain: From Theory to Statistical Test}\label{tab:concept_to_analysis} \\
\toprule
\thead{Concept} & \thead{Variable} & \thead{Measurement} & \thead{Type} & \thead{Statistical Analysis} \\
\midrule
\endfirsthead
\endhead
\bottomrule
\endlastfoot
ASD-focused diagnostic capability & Classification accuracy per domain & Stratified 10-fold CV accuracy (\%) & Ratio & Paired \textit{t}-test (\textit{p} < .05); Cohen's \textit{d}; bootstrap 95\% CI \\
\addlinespace
DL superiority over baselines & Accuracy difference (DL -- baseline) & Paired accuracy comparison per domain & Interval & Paired \textit{t}-test; effect size; signed-rank test \\
\addlinespace
Feature discriminative power & SHAP importance score & Feature contribution to prediction & Ratio & Ablation study ($\Delta$ accuracy); SHAP bar plots; permutation importance \\
\addlinespace
Clinical explainability & Expert agreement rating & 5-point Likert scale (1--5) & Ordinal & Median; inter-rater agreement; Kendall's $W$; qualitative thematic analysis \\
\addlinespace
Organisational value & Cost reduction and illustrative financial scenario & Financial modelling (\%, \$) & Ratio & illustrative cost-modelling scenario comparison; sensitivity analysis (3-scenario + tornado); illustrative payback estimate \\
\end{xltabular}}
\vspace{0.5em}
\noindent\textit{Note.} Each row traces one concept through five stages: theoretical construct, operationalised variable, measurement instrument, Stevens' data type, and statistical method. This chain ensures traceability from theory to evidence. DL = deep learning; SHAP = Shapley additive explanations; CV = cross-validation; CI = confidence interval; illustrative cost-modelling scenario = total cost of ownership; illustrative financial scenario = return on investment.
\fi %% End consolidated Data Collection, Analytical Purpose, Operationalisation, Concept-to-Analysis

\iffalse %% Methodology detail — covered in Theoretical Framework section and Ch.3
\subsection{Core Theoretical Integration and Statistical Methods}
\label{sec:theoretical_statistical}

Three core theories---TAM, TOE, and RBV---form the theoretical backbone of the research model. Table~\ref{tab:theory_stats_summary} integrates these theories with the statistical methods hierarchy, showing which methods are applied, equivalent, or reserved for future research.


{\tablestripes\tablefontsize

\needspace{10\baselineskip}
\begin{xltabular}{\textwidth}{@{} >{\raggedright\arraybackslash}p{2cm} L L L @{}}
\caption[Theoretical Integration and Statistical Methods]{Core Theoretical Integration (TAM--TOE--RBV) and Statistical Methods Hierarchy}\label{tab:theory_stats_summary} \\
\toprule
\thead{Element} & \thead{Core Question} & \thead{Adoption Stage} & \thead{Dissertation Evidence} \\
\midrule
\endfirsthead
\endhead
\bottomrule
\endlastfoot
\multicolumn{4}{l}{\textit{\textbf{Core Theories}}} \\
\addlinespace
TAM~\cite{davis1989tam} & Will clinicians adopt it? & Stage~3: AI Adoption & RAG expert evaluation; trust through explanations (Ch.~\ref{chap:analysis}) \\
TOE~\cite{tornatzky1990processes} & What enables adoption? & Stages~1--2: Support $\to$ Readiness & RGAIG governance; NIST/FDA/EU AI Act alignment \\
RBV~\cite{barney1991firm} & Why sustained advantage? & Stages~4--5: Productivity $\to$ Performance & VRIN: no comparable ASD-focused system identified \\
\addlinespace
\multicolumn{4}{>{\raggedright\arraybackslash}p{0.95\textwidth}}{\textit{Integration:} TAM explains \textit{individual adoption} (Stage~3), TOE explains \textit{organisational readiness} (Stages~1--2), and RBV explains \textit{sustained value creation} (Stages~4--5). Together, they span the conceptual adoption-readiness lifecycle.} \\
\midrule
\multicolumn{4}{l}{\textit{\textbf{Statistical Methods Hierarchy}}} \\
\addlinespace
\thead{Method} & \thead{Purpose} & \thead{Application} & \thead{Status} \\
\addlinespace
Regression & Predict DV from IVs & Classification models (SVM, RF) + SHAP perform equivalent function & Equivalent \\
ANOVA/\textit{t}-test & Compare group means & Paired \textit{t}-tests: DL vs.\ baseline; Cohen's \textit{d}; bootstrap CI & Applied \\
Factor Analysis & Identify latent constructs & Constructs defined \textit{a priori} (IV, M1, M2, DV, Mod) & Not required \\
SEM & Test full mediation model & Requires $n \geq 30$ organisations for H6, H7 & Future research \\
\end{xltabular}}
\vspace{0.5em}
\noindent\textit{Note.} TAM = Technology Acceptance Model; TOE = Technology--Organisation--Environment; RBV = Resource-Based View; VRIN = Valuable, Rare, Inimitable, Non-substitutable; SEM = structural equation modelling; DV = dependent variable; IV = independent variable; DL = deep learning; CI = confidence interval; SHAP = Shapley additive explanations; SVM = support vector machine; RF = random forest; NIST AI RMF = National Institute of Standards and Technology AI Risk Management Framework; RGAIG = Responsible Generative AI Governance; illustrative financial scenario = return on investment.
\fi %% End suppressed Core Theoretical Integration

%% Preserve original labels for cross-reference compatibility
\phantomsection\label{sec:theoretical_statistical}
\phantomsection\label{tab:theory_stats_summary}
\phantomsection\label{tab:tam_toe_rbv}
\phantomsection\label{tab:statistical_methods}

\iffalse %% Consolidated into tab:theory_stats_summary (Core Theoretical Integration and Statistical Methods)
\subsection{Core Theoretical Integration: TAM--TOE--RBV}

Table~\ref{tab:tam_toe_rbv} integrates three core theories---Technology Acceptance Model (TAM), Technology--Organisation--Environment (TOE), and Resource-Based View (RBV)---that collectively explain why the unified framework creates value, how organisations adopt it, and what sustains competitive advantage. These three theories form the theoretical backbone of the research model.


\noindent Table~\ref{tab:tam_toe_rbv} tam--toe--rbv integration.

{\tablestripes\tablefontsize
\needspace{10\baselineskip}
\begin{xltabular}{\textwidth}{@{} >{\raggedright\arraybackslash}p{1.0cm} p{3.6cm} p{2.8cm} p{3.2cm} p{3.3cm} @{}}
\caption[TAM--TOE--RBV Integration]{Core Theoretical Integration: TAM, TOE, and RBV Applied to the Research Model}\label{tab:tam_toe_rbv} \\
\toprule
\thead{Theory} & \thead{Core Question} & \thead{Key Constructs} & \thead{Adoption Chain Mapping} & \thead{Evidence} \\
\midrule
\endfirsthead
\endhead
\bottomrule
\endlastfoot
TAM & Will clinicians adopt the AI framework? & Perceived usefulness, perceived ease of use, behavioural intention & Stage~3 (AI Adoption): Explainability drives perceived usefulness; automation drives ease of use & RAG expert evaluation; clinician trust through literature-grounded explanations (Chapter~\ref{chap:analysis}) \\
\addlinespace
TOE & What organisational and environmental factors enable adoption? & Technology readiness, organisational readiness, environmental pressure (regulatory, competitive) & Stages~1--2 (Top Management Support $\to$ AI Readiness): Governance infrastructure, regulatory compliance, leadership mandate & Five-layer RGAIG; NIST AI RMF alignment; FDA SaMD/EU AI Act compliance; KPI escalation protocol \\
\addlinespace
RBV & Why does the unified framework create sustained competitive advantage? & Valuable, rare, inimitable, non-substitutable (VRIN) resources & Stages~4--5 (Productivity $\to$ Performance): ASD-focused capability as inimitable resource & VRIN assessment; projected cost reduction and illustrative financial scenario (Chapter~\ref{chap:discussion}); no comparable ASD-focused system identified \\
\midrule
\multicolumn{2}{l}{\textbf{Integration}} & \multicolumn{3}{>{\raggedright\arraybackslash}p{0.72\textwidth}}{TAM explains \textit{individual adoption} (Stage~3), TOE explains \textit{organisational readiness} (Stages~1--2), and RBV explains \textit{sustained value creation} (Stages~4--5). Together, they cover the conceptual adoption-readiness lifecycle from pre-evidence-maturity interpretation through technology acceptance to competitive advantage.} \\
\end{xltabular}}
\vspace{0.5em}
\noindent\textit{Note.} The three theories are complementary, not competing: TAM operates at the individual level (clinician adoption), TOE at the organisational level (readiness and environment), and RBV at the strategic level (competitive advantage). The adoption chain stages reference Table~\ref{tab:adoption_chain}. TAM = Technology Acceptance Model~\cite{davis1989tam}; TOE = Technology--Organisation--Environment framework~\cite{tornatzky1990processes}; RBV = Resource-Based View~\cite{barney1991firm}; VRIN = Valuable, Rare, Inimitable, Non-substitutable; RGAIG = Responsible Generative AI Governance; NIST AI RMF = National Institute of Standards and Technology AI Risk Management Framework; illustrative financial scenario = return on investment; KPI = key performance indicator.

\subsection{Statistical Methods Hierarchy}

Table~\ref{tab:statistical_methods} presents a hierarchy of four multivariate statistical methods---Regression, ANOVA, Factor Analysis, and Structural Equation Modelling (SEM)---showing which methods this dissertation employs and which are reserved for future research requiring organisational evaluation data.


\noindent Table~\ref{tab:statistical_methods} statistical methods hierarchy.

{\tablestripes\tablefontsize
\needspace{10\baselineskip}
\begin{xltabular}{\textwidth}{@{} >{\raggedright\arraybackslash}p{2.0cm} p{4.6cm} p{6.0cm} >{\raggedright\arraybackslash}p{2.0cm} @{}}
\caption[Statistical Methods Hierarchy]{Statistical Methods Hierarchy: From Regression to SEM}\label{tab:statistical_methods} \\
\toprule
\thead{Method} & \thead{Purpose and Capability} & \thead{Application in This Dissertation} & \thead{Status} \\
\midrule
\endfirsthead
\endhead
\bottomrule
\endlastfoot
Regression & Predict DV from one or more IVs; quantify strength and direction of relationships & Not directly applied; classification models (SVM, RF, logistic) perform equivalent predictive function; SHAP values quantify feature--outcome relationships & Equivalent applied \\
\addlinespace
ANOVA / \textit{t}-test & Test whether group means differ significantly; compare performance across conditions & Paired \textit{t}-tests: DL vs.\ baseline accuracy per domain; Cohen's \textit{d} effect sizes; bootstrap CI for robustness (results in Chapter~\ref{chap:analysis}) & Applied \\
\addlinespace
Factor Analysis & Identify latent constructs underlying observed variables; reduce dimensionality & Not required: constructs are theoretically defined \textit{a priori} (IV, M1, M2, DV, Mod) and operationalised through domain-specific instruments & Not required \\
\addlinespace
SEM & Test full mediation model simultaneously; decompose direct and indirect effects & Proposed for future research: requires organisational evaluation data (\textit{n} $\geq$ 30 organisations) to test H6 (moderation) and H7 (competing direct path) & Future research \\
\end{xltabular}}
\vspace{0.5em}
\noindent\textit{Note.} The four methods represent increasing analytical complexity. This dissertation applies ANOVA-equivalent comparisons (paired \textit{t}-tests with effect sizes) and classification-based prediction (equivalent to regression). SEM is explicitly identified as a future research method for testing the full mediation model (H6--H7) once organisational evaluation data is available. DV = dependent variable; IV = independent variable; SVM = support vector machine; RF = random forest; DL = deep learning; CI = confidence interval; SEM = structural equation modelling.
\fi %% End consolidated TAM-TOE-RBV and Statistical Methods Hierarchy

%===============================================================================
\section{Research Structure (A--D Framework) (Reference: Appendix~M)}

\noindent The detailed operationalization for this A--D framework component appears in Appendix~M (Section~M2). This section in Chapter~1 records that the framework exists and identifies its scope; the operationalization tables, indicator specifications, and procedural detail are placed in Appendix~M to preserve Chapter~1's focus on problem framing and to avoid pre-empting Chapter~3's methodology specification.


\subsection{Doctoral Novelty Statement}
\label{sec:doctoral_novelty_statement}

\noindent To the best of the researcher's knowledge, no prior study has simultaneously examined Responsible AI Governance, explainability, clinician trust, and adoption readiness within a unified ASD-focused AI-assisted screening context using a governance-mediated structural model. This combination represents the principal novelty of the present research.

\section{Significance of the Study}
\label{sec:scope}

This study evaluates the proposed framework on Autism Spectrum Disorder (ASD). The purpose is to assess whether the framework supports robust ASD screening using EEG-based features under stable evaluation conditions. The scope of the dissertation is therefore \textbf{ASD-focused framework validation}.


\needspace{12\baselineskip}
\iffalse % AUTO-SUPPRESS non-ASD
\noindent Table~\ref{tab:domain_role_hierarchy} presents role of Each Biomedical Domain in the Study Design.

{\tablestripes\tablefontsize
\begin{xltabular}{\textwidth}{@{} >{\raggedright\arraybackslash}p{1.0cm} >{\raggedright\arraybackslash}p{3.5cm} >{\raggedright\arraybackslash}p{2.5cm} p{7.2cm} @{}}
\caption[Role of Each Biomedical Domain in the Study Design]{Role of Each Biomedical Domain in the Study Design}\label{tab:domain_role_hierarchy} \\
\toprule
\thead{Domain} & \thead{Role in Study} & \thead{Modality} & \thead{Rationale} \\
\midrule
\endfirsthead
\endhead
\bottomrule
\endlastfoot
ASD & Primary validation & EEG & Strong fit with the neurodiagnostic focus of the thesis and central to framework evaluation. \\
\addlinespace
PD & Primary validation & EEG & High clinical relevance and strong diagnostic value for screening and monitoring use cases. \\
\addlinespace
Stress & Primary validation & EEG / multi-modal & Connects clinical, wearable, workflow, and managerial relevance. \\
\addlinespace
Cancer & Primary validation & Imaging / radiomics & Extends the framework beyond EEG and strengthens the ASD-focused claim. \\
\addlinespace
BCI & Boundary-condition & EEG & Retained as a stress-test case: smaller sample size, multi-class complexity, and higher methodological difficulty. \\
\end{xltabular}}
\vspace{0.5em}
\noindent\textit{Note.} The study does not claim equal evidentiary strength across all the ASD domain. Four domains form the primary validation set, while BCI is retained as a boundary-condition case that helps define the framework's current limits.
\fi % AUTO-SUPPRESS non-ASD

Table~\ref{tab:study_scope} defines the five study domains, their data modalities, diagnostic tasks, and clinical relevance.


\needspace{12\baselineskip}
\iffalse % AUTO-SUPPRESS non-ASD
{\tablestripes\tablefontsize

\begin{xltabular}{\textwidth}{@{} >{\raggedright\arraybackslash}p{2cm} >{\raggedright\arraybackslash}p{2.3cm} L L @{}}
\caption[Study Scope]{Study Scope: Five Biomedical Domains}\label{tab:study_scope} \\
\toprule
\thead{Domain} & \thead{Modality} & \thead{Diagnostic Task} & \thead{Clinical Relevance} \\
\midrule
\endfirsthead
\endhead
\bottomrule
\endlastfoot
Autism Spectrum Disorder & EEG (oscillatory, connectivity) & Binary classification (ASD vs.\ control) & Early diagnosis; intervention timing \\
\addlinespace
Parkinson's Disease & EEG (spectral, temporal) & Binary classification (PD vs.\ control) & Pre-motor detection; treatment monitoring \\
\addlinespace
Stress Detection & EEG (affective computing) & Multi-class (low/medium/high stress) & Occupational health; wellness monitoring \\
\addlinespace
BCI Motor Imagery & EEG (ERD/ERS, mu/beta) & 4-class motor imagery classification & Assistive devices; neurorehabilitation \\
\addlinespace
Cancer Radiomics & CT/MRI (shape, texture) & Multi-class tumor classification & Biopsy reduction; treatment selection \\
\end{xltabular}}
\vspace{0.5em}
\noindent\textit{Note.} EEG = electroencephalography; ERD/ERS = event-related desynchronization/synchronization; CT = computed tomography; MRI = magnetic resonance imaging.
\fi % AUTO-SUPPRESS non-ASD

\paragraph{Study Boundaries.}
Table~\ref{tab:study_boundaries} delineates the study boundaries with the rationale for each scope decision.


{\tablestripes\tablefontsize

\needspace{10\baselineskip}
\begin{xltabular}{\textwidth}{@{} L L @{}}
\caption[Study Boundaries with Rationale]{Study Boundaries with Rationale. Each boundary condition is paired with its justification, making explicit what the study does and does not cover and pre-empting examiner questions about scope limitations.}\label{tab:study_boundaries} \\
\toprule
\thead{Boundary} & \thead{Rationale} \\
\midrule
\endfirsthead
\endhead
\bottomrule
\endlastfoot
Public datasets only (no proprietary clinical data) & Ensures reproducibility; avoids IRB complications for secondary analysis \\
\addlinespace
Retrospective analysis only (no prospective trials) & Establishes proof-of-concept before committing to clinical trial resources \\
\addlinespace
The ASD domain (not exhaustive coverage) & Provides sufficient heterogeneity for generalizability testing without diluting depth \\
\addlinespace
Expert panel of 12 domain experts (n=12, purposive sampling; not large-scale survey) & Enables deep qualitative assessment with triangulated evidence \\
\addlinespace
English-language literature only & Pragmatic scope control; dominant language of biomedical AI research \\
\end{xltabular}}
\vspace{0.5em}
\noindent\textit{Note.} Each boundary is revisited in Chapter~\ref{chap:discussion} as a limitation with corresponding mitigation strategy.

\addcontentsline{toc}{paragraph}{Must-Have: Assumptions} % [TOC-must-have-insertion]
\paragraph{Key Assumptions.}
Three assumptions underpin the research design, each accompanied by its evidence basis:
\begin{enumerate}
    \item Classification accuracy on retrospective data provides a meaningful, if imperfect, proxy for screening-support utility (not diagnostic claim) in prospective clinical settings. (Retrospective-to-prospective concordance has been demonstrated in EEG-based neurological classification~\cite{acharya2018deep, boonyakitanont2020review}, though prospective validation remains necessary before future clinical-validation pathway.)
    \item Domain experts with $\geq$5 years of clinical or research experience can reliably evaluate the quality and clinical relevance of AI-generated explanations. (Expert panel evaluation is an established validation method in clinical AI research~\cite{tonekaboni2019clinicians}; the five-year threshold follows standard practice in Delphi-method studies~\cite{hsu2007delphi}.)
\end{enumerate}

%===============================================================================
\section{Scope of the Study}


\subsection{Why ASD Is the Right Empirical Context}
\label{sec:why_asd_right_context}

\noindent Autism Spectrum Disorder represents an appropriate empirical context because early intervention is strongly associated with improved developmental outcomes, yet significant delays frequently exist between symptom emergence and formal diagnosis. These delays create a practical opportunity for AI-assisted screening technologies. At the same time, the sensitive nature of pediatric healthcare decisions creates heightened requirements for transparency, explainability, governance, and trust, making ASD screening a suitable domain for evaluating responsible AI adoption.


\label{sec:alignment}

Table~\ref{tab:alignment_statement} traces the alignment from problems through objectives, research questions, and contributions, demonstrating that every element of the research design serves a specific purpose.


{\tablestripes\tablefontsize

\needspace{10\baselineskip}
\begin{xltabular}{\textwidth}{@{} L L L L @{}}
\caption[Research Alignment Statement]{Research Alignment Statement: Problem to Contribution Traceability}\label{tab:alignment_statement} \\
\toprule
\thead{Problem} & \thead{Objectives} & \thead{Research Questions} & \thead{Contributions} \\
\midrule
\endfirsthead
\endhead
\bottomrule
\endlastfoot
P1: Domain fragmentation & O1, O2, O5 & RQ1, RQ6 & C1, C5, C8 \\
\addlinespace
P2: Black-box opacity & O4 & RQ5 & C2, C7 \\
\addlinespace
P3: Manual inefficiency & O3 & RQ4 & C3 \\
\addlinespace
P4: Limited validation & O5, O6, O7 & RQ1, RQ3, RQ6, RQ7 & C4, C6, C7 \\
\addlinespace
P3: Manual ineffic.\ (econ.) & O3, O8 & RQ4, RQ8 & C3, C8 \\
\end{xltabular}}
\vspace{0.5em}
\noindent\textit{Note.} Every problem maps to at least two objectives, every objective maps to at least one research question, and every research question maps to at least one contribution. RQ7 (governance) and RQ8 (economic viability) complete the alignment chain for O7 and O8 respectively. No orphan elements exist in the research design.

\subsection{Research Design Flow}
\label{sec:research_design_flow}

Figure~\ref{fig:research_design_flow} illustrates the six-stage research design flow connecting business problems to empirical evidence.


\noindent\textbf{Six-Stage Research Design Flow: ASD Problem to.}

\needspace{12\baselineskip}
\begin{figure}[!htbp]
\centering
\resizebox{\textwidth}{!}{%
\begin{tikzpicture}[
  every node/.style={font=\fignodefont},
  stage/.style={draw=ggunavy, rounded corners=5pt, thick,
        minimum width=5.4cm, minimum height=3.0cm, align=center, font=\small,
        text width=5.0cm, inner sep=8pt},
    arr/.style={-{Stealth[length=5mm, width=4mm]}, line width=1.8pt, ggunavy},
    stglbl/.style={font=\normalsize\bfseries, text=ggunavy},
    detlbl/.style={font=\small\itshape, text=ggunavy!70}
]
%
% Stage 1 — Business Problem
\node[stage, fill=lightblue] (bp) at (0,0) {\textbf{Business Problem}\\[4pt]P1: Domain fragmentation\\P2: Black-box opacity\\P3: Manual inefficiency\\P4: Limited validation};
\node[stglbl, above=0.3cm of bp] {Stage 1};
%
% Stage 2 — Literature Gap
\node[stage, fill=lightorange] (lg) at (6.6,0) {\textbf{Literature Gap}\\[4pt]G1: No unified pipeline\\G2: No RAG explainability\\G3: No automation\\G4--G6: Governance gaps};
\node[stglbl, above=0.3cm of lg] {Stage 2};
%
% Stage 3 — Research Problem
\node[stage, fill=lightteal] (rp) at (13.2,0) {\textbf{Research Problem}\\[4pt]RQ1--RQ5: Five questions\\O1--O5: Five objectives\\H1--H5: Five hypotheses};
\node[stglbl, above=0.3cm of rp] {Stage 3};
%
% Stage 4 — Conceptual Framework (row 2 dropped from y=-6 to y=-8 to clear row 1 chapter labels)
\node[stage, fill=lightpurple] (cf) at (0,-8.0) {\textbf{Conceptual Framework}\\[4pt]IV: Technology Integration\\M1: Organisational Capability\\M2: Human Capability\\DV: Business Value};
\node[stglbl, above=0.3cm of cf] {Stage 4};
%
% Stage 5 — Methodology
\node[stage, fill=lightgreen] (mt) at (6.6,-8.0) {\textbf{Methodology}\\[4pt]Design Science Research\\Mixed methods (quant + qual)\\$k$-fold CV, ablation, SHAP\\12-expert panel validation};
\node[stglbl, above=0.3cm of mt] {Stage 5};
%
% Stage 6 — Data Analysis
\node[stage, fill=lightyellow] (da) at (13.2,-8.0) {\textbf{Data Analysis}\\[4pt]H1--H5: Tested\\Per-domain accuracy\\Effort reduction\\C1--C8: Contributions};
\node[stglbl, above=0.3cm of da] {Stage 6};
%
% Arrows — top row
\draw[arr] (bp.east) -- (lg.west);
\draw[arr] (lg.east) -- (rp.west);
%
% Arrow — Stage 3 down to Stage 4 (zigzag)
\draw[arr] (rp.south) -- ++(0,-1.0) -| (cf.north);
%
% Arrows — bottom row
\draw[arr] (cf.east) -- (mt.west);
\draw[arr] (mt.east) -- (da.west);
%
% Chapter labels below each stage
\node[detlbl, below=0.6cm of bp] {Ch.~1: Introduction};
\node[detlbl, below=0.6cm of lg] {Ch.~2: Literature Review};
\node[detlbl, below=0.6cm of rp] {Ch.~1: Introduction};
\node[detlbl, below=0.6cm of cf] {Ch.~2: Conceptual Model};
\node[detlbl, below=0.6cm of mt] {Ch.~3: Research Methods};
\node[detlbl, below=0.6cm of da] {Ch.~4--5: Findings};
%
\end{tikzpicture}%
}
\caption[Research Design Flow]{Six-Stage Research Design Flow: ASD Problem to Empirical Evidence}
\label{fig:research_design_flow}
\vspace{0.5em}
\noindent\textit{Note.} Each stage feeds the next. The ASD problem (Stage~1) motivates the literature review and gaps (Stage~2); gaps generate research questions, objectives, and hypotheses (Stage~3); the conceptual framework (Stage~4) structures the variables; the methodology (Stage~5) operationalises the framework; data analysis (Stage~6) tests the hypotheses and produces contributions. IV = independent variable; DV = dependent variable; M1, M2 = mediators; Mod = moderator; CV = cross-validation; SHAP = Shapley additive explanations; RAG = retrieval-augmented generation.
\end{figure}

\iffalse %% Condensed — figure caption and alignment table convey the content

\noindent\textbf{What it shows.} Figure~\ref{fig:research_design_flow} arranges six stages in a two-row zigzag: business problems (P1--P4) feed into literature gaps (G1--G6), which generate research questions, objectives, and hypotheses (RQ1--RQ5, O1--O5, H1--H5); the flow then descends to the conceptual framework (IV, M1, M2, DV, Mod), continues through the DSR methodology, and terminates at data analysis where hypotheses are tested and contributions (C1--C8) are produced. Chapter labels beneath each stage map the flow to the dissertation structure.


\noindent\textbf{So what.} The zigzag layout demonstrates internal coherence: every analytical decision in this dissertation is traceable backward to a real-world problem and forward to an empirical contribution. No stage operates in isolation, and no hypothesis lacks a problem-driven justification---a traceability requirement that examiners can verify against Table~\ref{tab:research_flow_alignment} below.
\fi

\noindent\textbf{Research Design Flow: Stage-by-Stage Alignment.}

\needspace{12\baselineskip}
\noindent\textbf{Master Traceability Chain: Problem to Contribution.} Table~\ref{tab:research_flow_alignment} provides the complete end-to-end chain that the examiner can use to verify every claim traces back to a problem and forward to evidence.

%% === MASTER TRACEABILITY MATRIX (full chain) ===
{%
\tablefontsize
\paragraph{Master Traceability Matrix.}


\needspace{10\baselineskip}

\begin{xltabular}{\textwidth}{@{} >{\raggedright\arraybackslash}p{0.9cm} >{\raggedright\arraybackslash}p{1.1cm} >{\raggedright\arraybackslash}p{0.5cm} >{\raggedright\arraybackslash}p{0.5cm} >{\raggedright\arraybackslash}p{1.6cm} >{\raggedright\arraybackslash}p{0.8cm} >{\raggedright\arraybackslash}p{0.8cm} >{\raggedright\arraybackslash}p{1.4cm} >{\raggedright\arraybackslash}p{1.0cm} >{\raggedright\arraybackslash}p{1.0cm} L @{}}
\caption[Master Traceability Matrix]{Master Traceability Matrix: Problem $\to$ Gap $\to$ Objective $\to$ Research Question $\to$ Hypothesis $\to$ Data Source $\to$ Context $\to$ Variables $\to$ Contribution $\to$ Decision. This single matrix enables the examiner to trace any claim across all five chapters and verify that every problem has a testable resolution with identified data, variables, and a governance-readiness interpretation.}
\label{tab:research_flow_alignment} \\
\toprule
\thead{Problem} & \thead{Gap} & \thead{Obj.} & \thead{RQ} & \thead{Hypothesis} & \thead{Data} & \thead{Context} & \thead{IV / DV} & \thead{Contrib.} & \thead{Decision} & \thead{Chapter Flow} \\
\midrule
\endfirsthead
\endhead
\bottomrule
\endlastfoot
P1 & G1 & O2 & RQ1 & H1: ASD $\geq$85\% & Secondary (EEG) & B2C quant (Part B) & IV: Model arch. DV: Accuracy & C1 & Adopt & 1$\to$2$\to$3$\to$4$\to$5 \\
\addlinespace
P1 & G5 & O3 & RQ3 & H4: Effort $-$90\% & Secondary (pipeline) & B2B qual (Part C) & IV: Automation DV: Time & C3 & Adopt & 1$\to$2$\to$3$\to$4$\to$5 \\
\addlinespace
P2 & G2 & O2 & RQ2 & H3: SHAP valid & Secondary (EEG) & B2C quant (Part B) & IV: Features DV: SHAP rank & C2 & Adopt & 1$\to$2$\to$3$\to$4$\to$5 \\
\addlinespace
P2 & G3 & O1 & RQ4 & H5: Expert $\alpha \geq .75$ & Primary (survey) & B2C qual (Part A) & IV: RAG output DV: Trust score & C4 & Pilot & 1$\to$2$\to$3$\to$4$\to$5 \\
\addlinespace
P3 & G5 & O3 & RQ3 & H4: 94.6\% reduction & Secondary (pipeline) & B2B qual (Part C) & IV: Agentic AI DV: Steps, time & C5 & Adopt & 1$\to$2$\to$3$\to$4$\to$5 \\
\addlinespace
P4 & G4 & O4 & RQ5 & H2: Ensemble vs DL & Secondary (EEG) & B2C quant (Part B) & IV: Model type DV: $p$-value, $d$ & C6 & Adopt & 1$\to$2$\to$3$\to$4$\to$5 \\
\addlinespace
P4 & G6 & O5 & RQ5 & H5: Compliance & Primary (expert) & Gov (Part D) & IV: RGAIG DV: Pass rate & C7 & Adopt & 1$\to$2$\to$3$\to$4$\to$5 \\
\addlinespace
All & G0 & O1--5 & All & SEM mediation & Both & All parts & IV: Tech DV: Business value & C8 & Adopt/Pilot & 1$\to$2$\to$3$\to$4$\to$5 \\
\end{xltabular}
}%

\vspace{0.5em}
\noindent\textit{Note.} P = Problem (Ch.~1); G = Gap (Ch.~2); O = Objective (Ch.~1); RQ = Research Question (Ch.~1); H = Hypothesis (designed Ch.~3, tested Ch.~4); Data = Primary (questionnaire/expert panel) or Secondary (EEG/benchmark datasets); Context = B2C qualitative (Part~A), B2C quantitative (Part~B), B2B qualitative (Part~C), Governance (Part~D); IV = independent variable; DV = dependent variable; C = Contribution (Ch.~5); Decision = Adopt (all thresholds met), Pilot (beta required), Defer (threshold not met). Every row is a complete traceable chain from business problem to governance-readiness interpretation.

\noindent This master chain ensures bidirectional traceability: every contribution traces backward through hypothesis, question, objective, and gap to a named business problem. No element is orphaned. An examiner can pick any row and follow it across all five chapters.

\subsection{Research Model: Adoption-to-Governance Chain}
\label{sec:adoption_chain}

Table~\ref{tab:adoption_chain} maps the six-stage adoption-to-governance chain to thesis constructs and empirical evidence.


\needspace{12\baselineskip}
\phantomsection\label{tab:adoption_chain}

\noindent The detailed adoption-to-governance chain table was removed from the live chapter because it overextended the introduction into governance instrumentation, KPI ladders, and modeled value logic. The retained figure is sufficient to show the conceptual sequence, while empirical evidence is handled later in the thesis with more caution.

Figure~\ref{fig:adoption_chain} illustrates the six-stage chain as a sequential flow with a governance feedback loop.


\needspace{12\baselineskip}
\phantomsection\label{fig:adoption_chain}

\noindent The adoption-to-governance chain is retained in Chapter~1 only as a conceptual organisational lens. The full figure was removed from the live chapter because it duplicated later theoretical material and pushed the introduction into governance, KPI, and value-creation detail that is better handled in Chapters~4 and~5.

\iffalse %% Condensed — table and figure convey the chain; verbose prose removed

\noindent\textbf{What it shows.} Figure~\ref{fig:adoption_chain} presents a six-stage linear chain---Top Management Support, AI Readiness, AI Adoption, Productivity, Performance, and AI Governance---with hypothesis labels (H1--H6) positioned above the stage transitions they evidence. A dashed feedback arrow from Governance back to AI Readiness creates a closed-loop model that distinguishes this framework from linear adoption theories.


\noindent\textbf{So what.} The feedback loop reveals that RGAIG treats governance not as a terminal checkpoint but as a continuous improvement mechanism: the 525-module quality taxonomy and KPI escalation protocols feed monitoring outcomes back into readiness assessments. This closed-loop design directly operationalises Kotter's anchoring step and the TOE framework's environmental dimension, connecting the research model to the theoretical foundations in Section~\ref{sec:theoretical_framework}.

The adoption-to-governance chain argues that AI in healthcare follows a six-stage maturation process, each building on the preceding stage:

\begin{enumerate}[leftmargin=*, itemsep=3pt]
    \item \textbf{Top Management Support} --- Provides the organisational mandate, budget authority, and cultural readiness required to initiate AI transformation.
    \item \textbf{AI Readiness} --- Establishes the governance infrastructure, regulatory compliance, and data quality standards required before technology deployment.
    \item \textbf{AI Adoption} --- Deploys the ASD-focused pipeline with explainability and agentic orchestration, building clinician trust through RAG-generated explanations.
    \item \textbf{Productivity} --- Emerges as the immediate outcome: automated workflows substantially reduce manual effort and diagnostic time, freeing clinicians for complex cases (quantified in Chapter~\ref{chap:analysis}).
    \item \textbf{Performance} --- Translates productivity into business value---cost reduction, projected ROI, and sustained competitive advantage (VRIN assessment; detailed in Chapter~\ref{chap:discussion}).
    \item \textbf{AI Governance} --- Closes the loop: continuous monitoring via the RGAIG framework's 525-module quality taxonomy, KPI escalation protocols, and bias auditing feeds back into AI Readiness, ensuring the system maintains compliance and performance over time.
\end{enumerate}
\fi


\noindent The six-stage model is consistent with the TOE framework~\cite{tornatzky1990processes}, UTAUT~\cite{venkatesh2003utaut}, and Kotter's change management model~\cite{kotter1996leading}. The governance feedback loop distinguishes RGAIG from linear adoption theories by treating governance as a continuous improvement mechanism.

%===============================================================================
\subsection{Theoretical Framework}
\label{sec:theoretical_framework}

This section presents the dual theoretical grounding: domain-specific theories justifying the technical architecture, and management theories explaining value creation and adoption.

\iffalse %% Engineering detail moved to Appendix A --- signal_processing_theories
\subsection{Signal Processing and Domain Theories}

Three neuroscience theories justify the framework's multi-scale, attention-based architecture:


\needspace{12\baselineskip}
\iffalse % AUTO-SUPPRESS non-ASD
\noindent Table~\ref{tab:domain_theories} outlines signal Processing and Domain Theories Guiding Framework Architecture.

{\tablestripes\tablefontsize
\begin{xltabular}{\textwidth}{@{} >{\raggedright\arraybackslash}p{2.8cm} L L L @{}}
\caption[Domain Theories Guiding Framework]{Signal Processing and Domain Theories Guiding Framework Architecture}\label{tab:domain_theories} \\
\toprule
\thead{Theory} & \thead{Core Principle} & \thead{Domain Evidence} & \thead{Framework Mapping} \\
\midrule
\endfirsthead
\endhead
\bottomrule
\endlastfoot
Neural Oscillation~\cite{buzsaki2012brain} & Disorders manifest as frequency-band abnormalities & ASD: altered $\gamma$ (30--100 Hz); PD: excess $\beta$ (13--30 Hz); Stress: $\alpha$/$\theta$ modulation & Multi-scale CNN captures frequency-specific patterns simultaneously \\
\addlinespace
Cortical Connectivity~\cite{niedermeyer2005eeg} & Disorders disrupt inter-region functional connectivity & ASD: reduced long-range~\cite{bosl2018eeg}; PD: basal ganglia--cortical~\cite{geraedts2018eeg}; Stress: prefrontal--limbic~\cite{alshargie2018mental} & Channel attention weights electrode pairs with disrupted connectivity \\
\addlinespace
Information Processing~\cite{roy2019eeg} & Cognitive impairments alter temporal EEG dynamics & EEG: event-related potentials; Imaging: multi-scale texture heterogeneity~\cite{lambin2012radiomics} & Bidirectional LSTM captures sequential dependencies \\
\end{xltabular}}
\vspace{0.5em}
\noindent\textit{Note.} CNN = convolutional neural network; LSTM = long short-term memory; $\alpha$ = alpha (8--13 Hz); $\beta$ = beta (13--30 Hz); $\gamma$ = gamma (30--100 Hz); $\theta$ = theta (4--8 Hz).
\fi % AUTO-SUPPRESS non-ASD
\fi %% End suppressed signal_processing_theories

\phantomsection\label{tab:domain_theories}

Signal processing theory provides the mathematical foundation for transforming raw EEG recordings into clinically meaningful features. Three neuroscience theories---neural oscillation~\cite{buzsaki2012brain}, cortical connectivity~\cite{niedermeyer2005eeg}, and information processing~\cite{roy2019eeg}---collectively justify the framework's multi-scale, attention-based architecture by linking frequency-band abnormalities, disrupted inter-region connectivity, and altered temporal dynamics to the five clinical domains under study. Technical foundations are documented in Appendix~A.

\subsection{Management and Organizational Theories}

Five management theories frame the organizational contribution. Table~\ref{tab:mgmt_theories} maps each theory to its core question, framework application, and key prediction.


\needspace{12\baselineskip}
The management-theories crosswalk is retained in summary form only. Resource-Based View explains strategic value, TAM and UTAUT explain adoption, Sociotechnical Systems Theory explains the interaction of technical and human systems, and Kotter's change model frames organisational transition. The detailed comparative table has been removed from the live chapter because it duplicated the surrounding prose and added unnecessary width pressure to the introduction.

\phantomsection\label{tab:mgmt_theories}

\subsection{Theoretical Integration}

No single theory spans the full lifecycle. The nine theories integrate as follows:
\begin{itemize}
    \item \textbf{What} the framework detects: Neural Oscillation, Cortical Connectivity, Information Processing theories
    \item \textbf{Why} it creates strategic value: Resource-Based View (RBV)
    \item \textbf{Whether} clinicians adopt it: Technology Acceptance Model (TAM), Unified Theory of Acceptance and Use of Technology (UTAUT)
    \item \textbf{How} technical and social systems co-evolve: Sociotechnical Systems Theory (STS)
    \item \textbf{How} organizational transition unfolds: Kotter's 8-Step Change Management
\end{itemize}

Figure~\ref{fig:unified_framework} presents the complete architecture of the unified framework, showing how six public benchmark datasets feed through domain-adaptive preprocessing and feature engineering into classification models, with agentic AI orchestration and RAG-enhanced explainability providing the governance and interpretability layers.


\needspace{12\baselineskip}
\iffalse % AUTO-SUPPRESS non-ASD
\noindent\textbf{Unified Agentic AI Framework for ASD Diagnosis}.} Figure~\ref{fig:unified_framework} provides visual context for the analysis below.

\needspace{12\baselineskip}
\begin{figure}[H]
    \centering
    \includegraphics[width=\textwidth]{unified_framework_architecture.png}
    \caption[Unified Agentic AI Framework]{Unified Agentic AI Framework for ASD Diagnosis}
    \label{fig:unified_framework}
    \vspace{0.5em}
\noindent\textit{Note.} ASD classification accuracy; RAG = retrieval-augmented generation; SHAP = Shapley additive explanations; SVM = support vector machine; RF = random forest; CNN = convolutional neural network; LSTM = long short-term memory; ABIDE = Autism Brain Imaging Data Exchange; DEAP = Database for Emotion Analysis using Physiological Signals; ; 
\end{figure}
\fi % AUTO-SUPPRESS non-ASD

\iffalse %% Condensed — figure caption covers the architecture

\noindent\textbf{What it shows.} Figure~\ref{fig:unified_framework} presents the complete RGAIG architecture across five horizontal layers: (1)~six public benchmark datasets (ABIDE, PhysioNet, DEAP, BCI Competition, TCIA, and additional EEG repositories) at the input layer; (2)~domain-adaptive preprocessing with modality-specific pipelines for EEG and imaging data; (3)~a shared feature engineering layer producing both handcrafted (spectral, connectivity, radiomic) and deep-learned features; (4)~classification via baseline models (SVM, RF) and DL ensembles (CNN, CNN-LSTM, Transformer); and (5)~agentic AI orchestration with RAG-based explainability and SHAP attribution at the output layer. A governance sidebar enforces the 46-module responsible AI taxonomy across all layers.


\noindent\textbf{So what.} This is the central architectural contribution of the dissertation. The figure confirms that ``unified'' means infrastructure unification---shared preprocessing, common orchestration, single governance framework---while preserving domain-specific classifiers for each disease. The five-layer design simultaneously satisfies all five research objectives (O1--O5 per Section~\ref{sec:research_purpose}): O1 (B2C trust/usability), O2 (technical performance), O3 (B2B clinician trust), O4 (RAI governance framework), O5 (B2C+B2B evaluation recommendation). Chapter~\ref{chap:methodology} operationalises each layer, and Chapter~\ref{chap:analysis} reports empirical results from this architecture.
\fi

\subsection{Conceptual Model Preview}

The theories converge into a mediation-based conceptual model (fully developed in Chapter~\ref{chap:literature}, Figure~\ref{fig:conceptual_model}). Technology Integration (IV) creates business value through two mediators: Organisational Capability (M1) and Human Capability (M2), moderated by Change Readiness. Figure~\ref{fig:conceptual_preview} illustrates this structure; H1--H5 correspond to specific paths, while H6--H7 require organisational data reserved for future research.


\needspace{12\baselineskip}
\phantomsection\label{fig:conceptual_preview}

\noindent The conceptual model preview is retained only in prose at this stage. The full mediation-style preview figure was removed from the live chapter because it added another large structural diagram to an already dense introduction and duplicated the more fully developed conceptual-model treatment in Chapter~2.

\iffalse %% Condensed — individual captions cover this

\noindent\textbf{Figures~\ref{fig:adoption_chain},~\ref{fig:unified_framework}, and~\ref{fig:conceptual_preview}: Summary.} These three figures present the research model at progressively deeper levels of abstraction. Figure~\ref{fig:adoption_chain} traces the \textit{organisational lifecycle} (six maturation stages from mandate to governance). Figure~\ref{fig:unified_framework} reveals the \textit{technical architecture} (five processing layers from data to decision). Figure~\ref{fig:conceptual_preview} specifies the \textit{causal mechanism} (mediation paths through which technology creates business value). Read together, they confirm that the RGAIG framework operates simultaneously at the organisational, architectural, and theoretical levels---a multi-layer coherence that distinguishes this DBA contribution from purely technical AI studies.
\fi

Table~\ref{tab:hypothesis_path_mapping} maps H1--H5 to their conceptual model paths.


Operational hypotheses H1--H5 are aligned to the conceptual model in a straightforward way: ASD-focused performance and model comparison represent the technology-to-capability logic, feature relevance supports the discriminative mechanism, automation addresses workflow value creation, and explainability supports the human-capability path. The detailed hypothesis-to-path mapping table has been removed from the live chapter to keep Chapter~1 focused on framing rather than conceptual bookkeeping.

\phantomsection\label{tab:hypothesis_path_mapping}

%===============================================================================
\subsection{Methodological Overview}
\label{sec:methodology_overview}

The research adopts a pragmatist, design-science methodology with a critical realist treatment of diagnostic outcomes. This dual ontology justifies the mixed-methods design: quantitative measurement for objective phenomena and qualitative assessment for subjective constructs~\cite{creswell2018research}. In practical terms, the study proceeds through four phases: data preparation, model development, framework integration, and validation. Chapter~\ref{chap:methodology} develops these phases in full.

\phantomsection\label{tab:methodology_phases}

%===============================================================================
\subsection{Competitive Landscape}
\label{sec:competitor_comparison}

The differentiation claim is architectural rather than market-comparative: the study proposes a unified research framework combining shared preprocessing logic, common explainability, automation, and governance across multiple biomedical domains. The detailed competitor matrix has been removed from the live chapter because it relied on public-product comparisons that are not central to the dissertation's evidence base; the canonical competitor matrix lives in the figure file and is referenced as Table~\ref{tab:competitor_comparison}.

%===============================================================================
\section{Structure of the Dissertation}
\label{sec:thesis_structure}

The thesis follows a conventional five-chapter structure: Chapter~1 defines the problem and research logic, Chapter~2 establishes the literature gap, Chapter~3 operationalises the design, Chapter~4 presents the evidence, and Chapter~5 interprets the findings and their implications. Supporting material is retained in the appendices. The chapter-structure table and duplicate flow figure have been removed from the live chapter because they add navigation scaffolding rather than substantive argument.

\phantomsection\label{tab:thesis_structure}
\phantomsection\label{fig:diss_structure}

\iffalse %% Condensed — figure caption conveys the dependency chain

\noindent\textbf{What it shows.} Figure~\ref{fig:diss_structure} presents six vertically stacked boxes---Chapters~1 through~5 plus Appendices A--K---connected by downward arrows that indicate a cumulative dependency chain. Each box lists the chapter's core content: from problem statement and hypotheses (Ch.~1) through theoretical foundations and gaps (Ch.~2), DSR methodology and model architectures (Ch.~3), hypothesis testing and governance validation (Ch.~4), to implications and future directions (Ch.~5).


\noindent\textbf{So what.} The linear dependency structure confirms that no chapter can be read in isolation: each builds on evidence and definitions established in its predecessor. This cumulative design ensures that the RGAIG framework argument is constructed layer by layer---mirroring the five-layer architecture itself---and provides the reader with a navigational reference for the remaining four chapters.
\fi %% End suppressed dissertation structure What it shows/So what

%===============================================================================
\phantomsection\label{tab:ch1_examiner_qa}

%===============================================================================
\subsection{Research Publications}
\label{sec:ch1_publications}

Several related publications support parts of the dissertation's evidence base, particularly in ASD detection, explainability, and wearable sensing. The detailed publications inventory has been removed from the live chapter because it functions more as portfolio documentation than as part of the thesis argument.

Our prior work established a CNN-ASD baseline achieving 78.35\% accuracy on the KAU dataset using CWT-based feature extraction~\cite{asthana2025cnnasd_icsit}. This dissertation extends that foundation through ensemble stacking, RAG-based explainability, and governance integration, improving classification accuracy to \asdacc{}.

\phantomsection\label{tab:ch1_publications}

\iffalse %% Publication details — figure and bullet list condensed
Figure~\ref{fig:ch1_papers_framework_mapping} illustrates how each paper maps to the pillars of the RGAIG architecture, demonstrating that the six publications collectively provide empirical grounding across the entire framework.


\needspace{12\baselineskip}
\iffalse % AUTO-SUPPRESS non-ASD
\noindent\textbf{Mapping of six research publications to RGAIG.}

\needspace{12\baselineskip}
\begin{figure}[!htbp]
\centering
\resizebox{0.97\textwidth}{!}{%
\begin{tikzpicture}[
  every node/.style={font=\fignodefont},
  % Central node
  cnode/.style={rectangle, rounded corners=6pt, draw=ggunavy, fill=lightpurple,
    text width=4.2cm, minimum height=1.4cm, align=center, font=\small\bfseries,
    very thick},
  % Paper nodes
  pnode/.style={rectangle, rounded corners=4pt, draw=ggunavy,
    text width=4.6cm, minimum height=1.6cm, align=left, font=\small,
    inner sep=5pt},
  % Pillar labels
  plabel/.style={font=\small\bfseries\color{ggunavy}, align=center},
  % Arrows
  arr/.style={-{Stealth[length=4mm, width=3mm]}, very thick, ggunavy, thick}
]
% 
% Central framework node
% 
\node[cnode] (center) {\color{ggunavy}RGAIG\\Framework};
% 
% Paper nodes arranged around centre
% Top-left: Paper 1
% 
\node[pnode, fill=lightblue, above left=1.8cm and 2.8cm of center] (p1) {
  \textbf{Paper 1} \cite{asthana2025asd}\\[2pt]
  EEG-based ASD diagnosis,\\75-table systematic review};
% 
% Top-right: Paper 2
% 
\node[pnode, fill=lightorange, above right=1.8cm and 2.8cm of center] (p2) {
  \textbf{Paper 2} \cite{asthana2025pd}\\[2pt]
  EEG-based PD diagnosis,\\80-paper multi-modality review};
% 
% Left: Paper 3
% 
\node[pnode, fill=lightteal, left=3.6cm of center] (p3) {
  \textbf{Paper 3} \cite{asthana2025cancer}\\[2pt]
  GAN+ResNet-18 cancer,\\97.1\% on 20,000 PBS images};
% 
% Right: Paper 4
% 
\node[pnode, fill=lightgreen, right=3.6cm of center] (p4) {
  \textbf{Paper 4} \cite{asthana2025asd}\\[2pt]
  EEG-based ASD classical,\\CNN framework (IEEE ICSIT)};
% 
% Bottom: Paper 5
% 
\node[pnode, fill=lightpink, below=1.8cm of center] (p5) {
  \textbf{Paper 5} \cite{asthana2025wearable}\\[2pt]
  Wearable EEG (consumer-grade),\\Emotiv multi-modal fusion};
% 
% Arrows with pillar labels
\draw[arr] (p1.south east) -- node[plabel, above, sloped] {Disease Classification} (center.north west);
\draw[arr] (p2.south west) -- node[plabel, above, sloped] {Disease Classification} (center.north east);
\draw[arr] (p3.east) -- node[plabel, above] {Computer Vision + GAN} (center.west);
\draw[arr] (p4.west) -- node[plabel, above] {EEG + Classical ML} (center.east);
\draw[arr] (p5.north) -- node[plabel, right] {IoT + Edge Deployment} (center.south);
% 
\end{tikzpicture}%
}
\caption{Mapping of six research publications to RGAIG framework pillars.}
\label{fig:ch1_papers_framework_mapping}
\end{figure}
\fi % AUTO-SUPPRESS non-ASD

The six papers collectively cover all major pillars of the RGAIG architecture, providing empirical grounding for the unified framework claims advanced in this dissertation.

\begin{itemize}[leftmargin=*, itemsep=3pt]
  \item \textbf{Published (Papers~1--3):} Validate the disease-specific classification components of the framework.
  \begin{itemize}[leftmargin=*, itemsep=2pt]
    \item \textbf{Paper~1} (ASD systematic review): Establishes the evidence base for the EEG-based classification pipeline through a 75-table review of 100+ studies.
    \item \textbf{Paper~2} (PD systematic review): Provides multi-modality evidence across EEG, speech, gait, and imaging, grounding the ASD-focused design rationale.
    \item \textbf{Paper~3} (Cancer GAN+ResNet-18): Demonstrates cross-modality generalisation through GAN-augmented computer vision, achieving 97.1\% on 20,000 PBS images.
  \end{itemize}
  \item \textbf{Under Review (Papers~4--6):} Address the GenAI, governance, and evaluation layers.
  \begin{itemize}[leftmargin=*, itemsep=2pt]
    \item \textbf{Paper~4} (EEG-RAG stress detection): Validates retrieval-augmented explainability with 94.7\% accuracy and cross-dataset domain adaptation.
    \item \textbf{Paper~5} (NeuroMCP agentic AI): Introduces multi-agent orchestration with MCP governance for neurological disease detection.
    \item \textbf{Paper~6} (Consumer-grade wearable EEG): Demonstrates Emotiv device integration and multi-modal sensor fusion; engineering details in Appendix~H (not part of the DBA contribution).
  \end{itemize}
\end{itemize}
\fi


%===============================================================================
\section{Chapter 1 Output Card: What Goes to Chapter 2}

\noindent Table~\ref{tab:ch1_output_card} presents chapter 1 Output Card: Deliverables to Chapter 2.


{\tablestripes\tablefontsize
\needspace{10\baselineskip}
\begin{xltabular}{\textwidth}{@{} >{\raggedright\arraybackslash}p{3.0cm} L @{}}
\caption[Chapter 1 Output Card: Deliverables to Chapter 2]{Chapter 1 Output Card: Deliverables to Chapter 2. This table summarises everything Chapter~1 produces and hands off to Chapter~2 for literature justification.}\label{tab:ch1_output_card} \\
\toprule
\thead{Deliverable} & \thead{What Chapter 2 Does With It} \\
\midrule
\endfirsthead
\endhead
\bottomrule
\endlastfoot
Problem statement (fragmented ASD screening) & Literature must confirm this problem exists and is unsolved \\
\addlinespace
5 objectives (O1--O5) & Each objective must map to a literature gap \\
\addlinespace
14 sub-objectives (A1--D3) & Each sub-objective must be grounded in reviewed evidence \\
\addlinespace
5 hypotheses (H1--H5) & Literature must provide the theoretical basis for each hypothesis \\
\addlinespace
A--D evaluation framework & Literature must justify why 4 dimensions (B2C qual, B2C quant, B2B qual, governance) are needed \\
\addlinespace
Research questions & Each question must connect to identified gaps in the literature \\
\addlinespace
ASD clinical context & Literature must review ASD prevalence, diagnostic delays, and EEG evidence \\
\addlinespace
RGAIG concept & Literature must show no existing framework combines all 5 layers \\
\end{xltabular}}
\vspace{0.3em}\noindent\textit{Note.} RGAIG = Responsible Generative AI Governance; ASD = Autism Spectrum Disorder; EEG = Electroencephalography. See chapter prose for full context.

%===============================================================================
\subsection{Supplementary Material}

The following appendix provides supporting material for Chapter~1:
\begin{itemize}[leftmargin=*, itemsep=2pt]
\item \textbf{Appendix~A} (Section~\ref{sec:intro_limitations}): Limitations of the introductory framing, section-to-objective traceability table, and P$\to$G$\to$O$\to$H chain.
\end{itemize}


\section{Chapter Summary}

\subsection{Chapter Synthesis}
\label{subsec:ch1_synthesis}

\begin{itemize}[leftmargin=*, itemsep=2pt]
    \item \textbf{Finding.} ASD-EEG screening fragmentation is decomposable into eleven problem layers across B2B (hospital) and B2C (consumer) evaluation contexts; no published system addresses all eleven.
    \item \textbf{Contribution.} The chapter establishes the B2C-primary, B2B-supporting positioning (Fig.~\ref{fig:b2c_b2b_support_model} + Fig.~\ref{fig:b2b_support_scenarios}), the RGAIG Maturity Framework as the hero artefact, and a six-hypothesis structure that the rest of the dissertation tests.
    \item \textbf{Limitation.} The B2C positioning is established conceptually here; the empirical evidence supporting each hypothesis is deferred to Chapters~3--5, and the future B2C prospective-validation specified in Chapter~5 has not yet run.
    \item \textbf{Transition.} Chapter~2 evaluates the literature against the eleven-layer problem structure to confirm the gap and locate the four B2C-specific gaps the dissertation must close.
\end{itemize}
\needspace{12\baselineskip}



\noindent\textbf{Chapter 1 Scorecard.} Table~\ref{tab:ch1_scorecard} summarises what this chapter promised, what it delivered, and how each output connects to the thesis objectives.

\paragraph{Chapter 1 Scorecard.}



{\tablestripes\tablefontsize
\needspace{10\baselineskip}
\begin{xltabular}{\textwidth}{@{} >{\raggedright\arraybackslash}p{5.4cm} p{7.5cm} >{\centering\arraybackslash}p{1.5cm} @{}}
\caption[Chapter 1 Scorecard]{Chapter 1 Scorecard: Promise vs Delivery}\label{tab:ch1_scorecard} \\
\toprule
\thead{Promised} & \thead{Delivered} & \thead{Obj.} \\
\midrule
\endfirsthead
\endhead
\bottomrule
\endlastfoot
Define ASD problem & 4 problems (P1--P4) with stakeholder impact & O1--O5 \\
\addlinespace
Identify research gaps & 6 gaps (G1--G6) from 156-study review & O1--O5 \\
\addlinespace
State objectives & 5 main (O1--O5) + 14 sub (A1--D3) & O1--O5 \\
\addlinespace
Formulate hypotheses & 5 hypotheses (H1--H5) with thresholds & O2--O5 \\
\addlinespace
Define scope & ASD-only, public data, no clinical trial & O1--O3 \\
\addlinespace
Master traceability & P$\to$G$\to$O$\to$RQ$\to$H$\to$C matrix & All \\
\end{xltabular}}
\vspace{0.3em}\noindent\textit{Note.} ASD = Autism Spectrum Disorder. See chapter prose for full context.

\noindent\textbf{Thesis argument.} This dissertation argues that ASD-focused clinical AI remains fragmented not because of technical limitations but because of structural incentives that reward narrow, single-domain solutions over unified, governed platforms.

\noindent\textbf{Handoff to Chapter~2.} Chapter~2 receives: 4 problems, 5 objectives, 5 research questions, and the RGAIG concept. It must ground these in the literature and derive the 6 gaps that justify the methodology.

\section{Limitations of the Study}
\label{sec:scope_limitations}

This section articulates the boundary conditions of the research. Table~\ref{tab:boundary_conditions} identifies six dimensions along which the study's scope is explicitly bounded, together with the justification for each boundary. These limitations are revisited in Chapter~\ref{chap:discussion} with corresponding mitigation strategies and future research directions.


\needspace{12\baselineskip}
\iffalse % AUTO-SUPPRESS non-ASD
\noindent Table~\ref{tab:boundary_conditions} boundary conditions.

{\tablestripes\tablefontsize

\begin{xltabular}{\textwidth}{@{} >{\raggedright\arraybackslash}p{2.2cm} p{5.2cm} p{7.0cm} @{}}
\caption[Boundary Conditions]{Boundary Conditions of the Research}\label{tab:boundary_conditions} \\
\toprule
\thead{Dimension} & \thead{Boundary} & \thead{Justification} \\
\midrule
\endfirsthead
\endhead
\bottomrule
\endlastfoot
Geographic scope & US and EU clinical settings; datasets sourced from institutions in these regions & Regulatory frameworks (FDA SaMD, EU AI Act) and clinical workflows differ across jurisdictions; findings may not generalise to low-resource healthcare systems without adaptation \\
\addlinespace
Temporal scope & Training and evaluation data span 2020--2025 recording periods & EEG acquisition protocols, device firmware, and clinical diagnostic criteria evolve over time; models trained on this window require periodic revalidation \\
\addlinespace
Domain coverage & Five EEG/imaging diseases (ASD, PD, stress, BCI, cancer); not all neurological or psychiatric conditions & The ASD domain provide sufficient heterogeneity for ASD-focused generalisation claims while maintaining analytical depth; exhaustive coverage would sacrifice rigour for breadth \\
\addlinespace
Sample size & BCI domain: $N=109$ subjects (BCI Competition IV 2a); smallest cohort in the study & Underpowered relative to ASD ($N=1,112$) and stress ($N=768$); BCI results are reported with wider confidence intervals and flagged as requiring replication with larger cohorts \\
\addlinespace
Technology mode & Offline retrospective analysis; not real-time BCI or streaming inference & Real-time evaluation introduces latency, connectivity, and safety constraints beyond the scope of a proof-of-concept framework; edge evaluation is modelled but not empirically tested \\
\addlinespace
Economic analysis & Projected illustrative financial scenario and illustrative cost-modelling scenario reduction based on modelling, not validated in future clinical-validation pathway & Cost projections use published healthcare AI cost benchmarks and sensitivity analysis (three scenarios); actual evaluation costs will vary by institution and scale \\
\end{xltabular}}
\vspace{0.5em}
\noindent\textit{Note.} Each boundary is revisited in Chapter~\ref{chap:discussion} (Section~5.6) with mitigation strategies. FDA = Food and Drug Administration; SaMD = Software as a Medical Device; EU AI Act = European Union Artificial Intelligence Act; ASD = autism spectrum disorder;   illustrative financial scenario = return on investment; illustrative cost-modelling scenario = total cost of ownership.
\fi % AUTO-SUPPRESS non-ASD


\noindent These boundaries are not weaknesses to be apologised for but deliberate scope decisions that ensure the research remains tractable, reproducible, and falsifiable within the doctoral timeframe. The framework architecture is designed to accommodate broader geographic, temporal, and domain coverage in post-doctoral extensions.

%===============================================================================
\subsection{Chapter Summary}
\label{sec:ch1_conclusion}

This chapter established the dissertation's problem context, identified six research gaps from a 156-study evidence base, and translated those gaps into the study's objectives, research questions, hypotheses, and intended contributions. It also defined the scope and boundary conditions of the research so that the claims taken forward into the later chapters remain aligned with the available evidence.

\iffalse %% Consolidated: thin Purpose and Key Sections Covered subsections merged into parent
\subsection{Purpose}

This chapter argued that domain fragmentation in clinical AI is a correctable design failure, identified six research gaps from a 156-study evidence base, derived five objectives and five falsifiable hypotheses, and positioned the RGAIG framework as a unified alternative. The chapter justified the study's relevance to both academic scholarship and healthcare practice, positioning the RGAIG framework as a governance-aligned response to domain fragmentation.

\subsection{Key Sections Covered}


\noindent Table~\ref{tab:ch1_kit_sections} summarises the key sections covered in this chapter.
\fi %% End consolidated Purpose and Key Sections Covered


The key elements established here are the problem statement, six literature-grounded gaps, the objective and question set, the five operational hypotheses, the bounded ASD scope, and the intended contributions. Those elements are now carried forward directly into Chapters~2--5 without an additional recap inventory table.

\phantomsection\label{tab:ch1_kit_sections}

%% SUPPRESSED — redundant with tab:ch1_kit_sections (Key Sections Covered table)
\iffalse
\subsection{Key Findings}


\noindent Table~\ref{tab:ch1_kit_findings} presents the key outputs and findings established in this chapter.


\noindent The findings below summarise the in-scope claims that the remainder of the dissertation defends or qualifies.
{\tablestripes\tablefontsize
\needspace{10\baselineskip}
\begin{xltabular}{\textwidth}{@{} >{\raggedright\arraybackslash}p{5.1cm} p{9.5cm} @{}}
\caption[Chapter 1 key outputs]{Chapter~1: Key Outputs and Findings}\label{tab:ch1_kit_findings} \\
\toprule
\thead{Area} & \thead{Outcome} \\
\midrule
\endfirsthead
\endhead
\bottomrule
\endlastfoot
Research gaps identified & Six gaps across AI integration, explainability, governance, and ASD-focused validation \\
Problem definition & Eleven interconnected problems quantified with costs to healthcare AI adoption \\
Objectives defined & Eight objectives aligned to gaps with traceability matrix (Table~\ref{tab:alignment_statement}) \\
Hypotheses specified & Five falsifiable hypotheses with explicit acceptance thresholds \\
Framework direction & RGAIG: unified, agentic, explainable, governance-aligned architecture \\
Non-claims declared & Four boundary conditions stated to prevent overclaiming \\
\end{xltabular}}
\vspace{0.5em}
\noindent\textit{Note.} Each output is developed in subsequent chapters: theoretical foundations (Ch.~2), methodology (Ch.~3), empirical testing (Ch.~4), and implications (Ch.~5).
\fi

\iffalse %% Condensed — key sections table already summarises chapter; objective matrix redundant with alignment table
\paragraph{Interpretation of Objective Matrix.}
Healthcare AI has produced single-domain classifiers exceeding 95\% accuracy, yet 88.5\% of the 156 studies reviewed in Chapter~\ref{chap:literature} remain siloed within a single condition, with no shared infrastructure for preprocessing, explainability, or governance. This fragmentation is a design convention, not a technical necessity. The introduction establishes that a unified, governance-aligned framework---one that combines agentic AI orchestration with RAG-based explainability and responsible AI compliance---is both technically feasible and organisationally necessary.

\paragraph{Objective Achievement.}

\noindent This table lists each \textit{objective} across evidence in chapter~1, and status, so examiners can trace what was assessed and what evidence supports each row.



\noindent This table lists each \textit{objective} across evidence in chapter~1, and status, so examiners can trace what was assessed and what evidence supports each row.


\noindent This table lists each \textit{objective} across evidence in chapter~1, and status, so examiners can trace what was assessed and what evidence supports each row.



{\tablestripes\tablefontsize
\needspace{10\baselineskip}
\begin{xltabular}{\textwidth}{@{} >{\raggedright\arraybackslash}p{4.4cm} p{7.8cm} >{\raggedright\arraybackslash}p{2.2cm} @{}}
\caption[Chapter 1 objective achievement]{Chapter~1: Objective Achievement Matrix}\label{tab:ch1_kit_objectives} \\
\toprule
\thead{Objective} & \thead{Evidence in Chapter~1} & \thead{Status} \\
\midrule
\endfirsthead
\endhead
\bottomrule
\endlastfoot
O1: Standardised preprocessing & Preprocessing gaps documented (P3, G3); pipeline direction specified & Scoped \\
O2: Classification $\geq$85\% & Accuracy threshold established; the ASD domain identified & Scoped \\
O3: Agentic AI orchestration & Automation gap documented (P4, G4); multi-agent direction specified & Scoped \\
O4: RAG explainability & Opacity gap documented (P2, G2); explainability rationale established & Scoped \\
O5: ASD-focused validation & Fragmentation gap documented (P1, G1); validation strategy outlined & Scoped \\
O6: Biomarker identification & Feature-importance gap documented (G6); SHAP direction specified & Scoped \\
O7: Governance architecture & Governance deficit documented (P5--P8, G5); compliance requirements mapped; maps to RQ7 and SH6 & Scoped \\
\addlinespace
O8: Economic viability & Cost inefficiency documented (P3); illustrative cost-modelling scenario and illustrative financial scenario benchmarking direction specified; maps to RQ8 and SH7 & Scoped \\
\end{xltabular}}
\vspace{0.5em}
\noindent\textit{Note.} All objectives are scoped and directed in this chapter. Methodology design is presented in Chapter~3; empirical evidence is presented in Chapter~4. illustrative cost-modelling scenario = total cost of ownership; illustrative financial scenario = return on investment.
\fi %% End suppressed Interpretation + Objective Achievement

\phantomsection\label{tab:ch1_kit_objectives}

\paragraph{Link to Chapter~2.}
Chapter~\ref{chap:literature} examines the theoretical and empirical foundations that support this argument, systematically documenting the research gaps through a literature review spanning ASD/EEG diagnostics, agentic AI, and RAG-based explainability.

\textbf{This thesis argues that domain fragmentation in biomedical AI diagnostics is an unnecessary and costly design failure---not a technical inevitability---and that a unified, agentic, and explainable framework can replace it with a platform that is more accurate, more transparent, and less expensive than the fragmented status quo.}



\subsection{Executive Table: Research Problem Summary}
\label{sec:exec_research_problem}

\noindent The following executive-level table summarizes the core research problems addressed by this dissertation, each mapped to its business impact for healthcare organizations.

{\tablestripes\tablefontsize
\needspace{10\baselineskip}
\begin{xltabular}{\textwidth}{@{} >{\raggedright\arraybackslash}p{6cm} L @{}}
\caption[Research Problem Summary]{Research Problem Summary.}\label{tab:exec_research_problem} \\
\toprule
\thead{Problem} & \thead{Business Impact} \\
\midrule
\endfirsthead
\endhead
\bottomrule
\endlastfoot
Delayed ASD identification & Delayed intervention; reduced developmental outcomes \\
\addlinespace
Limited explainability of AI screening & Reduced clinician trust; integration friction \\
\addlinespace
Weak governance maturity & Reduced adoption; compliance + liability exposure \\
\addlinespace
Fragmented Responsible AI controls & Inconsistent oversight; weak accountability \\
\end{xltabular}}
\vspace{0.3em}\noindent\textit{Note.} Each problem motivates a corresponding objective in Table~\ref{tab:exec_objectives} and a hypothesis in Chapter~4.


\subsection{Executive Table: Research Gap (Known vs.\ Unknown)}
\label{sec:exec_research_gap}

\noindent The following executive-level table contrasts what prior literature has established with what remains unknown, motivating the present research.

{\tablestripes\tablefontsize
\needspace{10\baselineskip}
\begin{xltabular}{\textwidth}{@{} >{\raggedright\arraybackslash}p{6cm} L @{}}
\caption[Research Gap: Known vs.\ Unknown]{Research Gap: Known vs.\ Unknown.}\label{tab:exec_research_gap} \\
\toprule
\thead{What is known} & \thead{What remains unknown} \\
\midrule
\endfirsthead
\endhead
\bottomrule
\endlastfoot
AI improves ASD-EEG screening accuracy & Governance effect on adoption is unclear \\
\addlinespace
XAI improves transparency of AI outputs & Trust pathway from XAI to adoption is unclear \\
\addlinespace
Responsible AI frameworks exist (ISO, NIST, EU AI Act) & Empirical adoption impact of governance maturity is unmeasured \\
\addlinespace
Clinicians value explainability in principle & Magnitude of governance$\rightarrow$trust$\rightarrow$adoption effect is undocumented \\
\end{xltabular}}
\vspace{0.3em}\noindent\textit{Note.} Detailed literature inventory is in Chapter~2 + Appendix~B.


\subsection{Executive Table: Research Objectives}
\label{sec:exec_objectives}

\noindent The following executive-level table summarizes the five research objectives. Full operationalization is in Appendix~M.

{\tablestripes\tablefontsize
\needspace{10\baselineskip}
\begin{xltabular}{\textwidth}{@{} >{\raggedright\arraybackslash}p{6cm} L @{}}
\caption[Research Objectives (Executive Summary)]{Research Objectives (Executive Summary).}\label{tab:exec_objectives} \\
\toprule
\thead{Objective} & \thead{Purpose} \\
\midrule
\endfirsthead
\endhead
\bottomrule
\endlastfoot
O1 & Evaluate the influence of governance maturity on explainability \\
\addlinespace
O2 & Evaluate the influence of governance maturity on trust \\
\addlinespace
O3 & Evaluate the influence of explainability on trust \\
\addlinespace
O4 & Evaluate the influence of trust on adoption readiness \\
\addlinespace
O5 & Validate the integrated RGAIG governance-mediated adoption framework \\
\end{xltabular}}
\vspace{0.3em}\noindent\textit{Note.} Detailed objective-to-question-to-hypothesis-to-method mapping is in Appendix~M.


\subsection{Executive Table: Research Questions}
\label{sec:exec_research_questions}

\noindent The following executive-level table summarizes the five research questions. Answers are in Chapter~5 (Research Question Answer Matrix, Table~\ref{tab:distinction_rq_answer_matrix}).

{\tablestripes\tablefontsize
\needspace{10\baselineskip}
\begin{xltabular}{\textwidth}{@{} >{\raggedright\arraybackslash}p{6cm} L @{}}
\caption[Research Questions (Executive Summary)]{Research Questions (Executive Summary).}\label{tab:exec_research_questions} \\
\toprule
\thead{Question} & \thead{Purpose} \\
\midrule
\endfirsthead
\endhead
\bottomrule
\endlastfoot
RQ1 & Does governance maturity influence explainability? \\
\addlinespace
RQ2 & Does governance maturity influence trust? \\
\addlinespace
RQ3 & Does explainability influence trust? \\
\addlinespace
RQ4 & Does trust influence adoption readiness? \\
\addlinespace
RQ5 & Does the integrated RGAIG framework predict adoption readiness? \\
\end{xltabular}}


\subsection{Executive Table: Contributions Summary}
\label{sec:exec_contributions}

\noindent The following executive-level table summarizes the four principal contributions of this dissertation, each classified by contribution type.

{\tablestripes\tablefontsize
\needspace{10\baselineskip}
\begin{xltabular}{\textwidth}{@{} >{\raggedright\arraybackslash}p{6cm} L @{}}
\caption[Contributions Summary]{Contributions Summary.}\label{tab:exec_contributions} \\
\toprule
\thead{Contribution} & \thead{Type} \\
\midrule
\endfirsthead
\endhead
\bottomrule
\endlastfoot
RGAIG Maturity Framework & Practical artifact \\
\addlinespace
Governance-mediated adoption pathway & Theoretical contribution \\
\addlinespace
Mixed-methods governance research design & Methodological contribution \\
\addlinespace
Healthcare AI adoption guidance + roadmap & Practitioner contribution \\
\end{xltabular}}
\vspace{0.3em}\noindent\textit{Note.} Full contribution detail is in Chapter~5 + Appendix~T.


\subsection{Executive Table: Scope of the Study (In vs.\ Out)}
\label{sec:exec_scope}

\noindent The following executive-level table clarifies the bounded scope of this research, separating what is in-scope from what is explicitly out-of-scope.

{\tablestripes\tablefontsize
\needspace{10\baselineskip}
\begin{xltabular}{\textwidth}{@{} >{\raggedright\arraybackslash}p{6cm} L @{}}
\caption[Scope: In-Scope vs.\ Out-of-Scope]{Scope: In-Scope vs.\ Out-of-Scope.}\label{tab:exec_scope} \\
\toprule
\thead{In Scope (DBA contribution)} & \thead{Out of Scope (CS engineering)} \\
\midrule
\endfirsthead
\endhead
\bottomrule
\endlastfoot
Governance maturity as construct & Novel ML algorithm design \\
\addlinespace
Explainability as adoption mediator & Novel CNN architecture \\
\addlinespace
Trust as behavioral mechanism & Neural network hyperparameter tuning \\
\addlinespace
Adoption readiness in healthcare & GPU optimization / training infrastructure \\
\addlinespace
RGAIG governance framework & Production deployment engineering \\
\end{xltabular}}
\vspace{0.3em}\noindent\textit{Note.} See Chapter~3, Section~\ref{sec:separation_research_engineering} for the full separation of research vs.\ supporting engineering.


\subsection{Executive Table: Why DBA (Not Computer Science PhD)}
\label{sec:exec_why_dba}

\noindent The following executive-level table contrasts the focus of this DBA dissertation with a Computer Science PhD on the same domain, clarifying the disciplinary positioning.

{\tablestripes\tablefontsize
\needspace{10\baselineskip}
\begin{xltabular}{\textwidth}{@{} >{\raggedright\arraybackslash}p{6cm} L @{}}
\caption[Why DBA Rather Than Computer Science PhD]{Why DBA Rather Than Computer Science PhD.}\label{tab:exec_why_dba} \\
\toprule
\thead{Computer Science PhD focus} & \thead{DBA focus (this dissertation)} \\
\midrule
\endfirsthead
\endhead
\bottomrule
\endlastfoot
Novel algorithm performance & Organizational adoption mechanisms \\
\addlinespace
Model architecture innovation & Governance maturity influence \\
\addlinespace
Benchmark accuracy improvement & Trust + explainability + adoption pathway \\
\addlinespace
Technical novelty as primary contribution & Managerial framework + practitioner artifact as primary contribution \\
\end{xltabular}}
\vspace{0.3em}\noindent\textit{Note.} The dissertation does not seek algorithmic innovation; existing ML, XAI, and RAG technologies serve as supporting infrastructure (see Chapter~3, Section~\ref{sec:separation_research_engineering}).


\subsection{Executive Table: Stakeholders + Interests}
\label{sec:exec_stakeholders}

\noindent The following executive-level table summarizes the four primary stakeholder groups and their interests in the proposed RGAIG framework.

{\tablestripes\tablefontsize
\needspace{10\baselineskip}
\begin{xltabular}{\textwidth}{@{} >{\raggedright\arraybackslash}p{6cm} L @{}}
\caption[Stakeholders and Primary Interests]{Stakeholders and Primary Interests.}\label{tab:exec_stakeholders} \\
\toprule
\thead{Stakeholder} & \thead{Primary Interest} \\
\midrule
\endfirsthead
\endhead
\bottomrule
\endlastfoot
Caregiver / Patient & Trust + transparency in AI-assisted screening recommendations \\
\addlinespace
Clinician & Explainability + clinical accountability + workflow integration \\
\addlinespace
Healthcare organization & Governance maturity + compliance + adoption risk management \\
\addlinespace
Regulator & Responsible AI compliance + auditability + safety oversight \\
\end{xltabular}}


\subsection{Executive Table: Framework Component Positioning}
\label{sec:exec_framework_positioning}

\noindent The following executive-level table positions each framework component, distinguishing primary research artifacts from supporting technologies and methods.

{\tablestripes\tablefontsize
\needspace{10\baselineskip}
\begin{xltabular}{\textwidth}{@{} >{\raggedright\arraybackslash}p{6cm} L @{}}
\caption[Framework Component Positioning]{Framework Component Positioning.}\label{tab:exec_framework_positioning} \\
\toprule
\thead{Component} & \thead{Role in the dissertation} \\
\midrule
\endfirsthead
\endhead
\bottomrule
\endlastfoot
RGAIG framework & Principal governance artifact (primary contribution) \\
\addlinespace
SMART operational layer & Operational framework (per-phase governance application) \\
\addlinespace
EEG signal context & Empirical screening modality (technical context) \\
\addlinespace
SHAP attribution & Explainability operationalization tool \\
\addlinespace
RAG citation grounding & Narrative explanation support \\
\addlinespace
PLS-SEM analysis & Validation method for structural model \\
\end{xltabular}}
\vspace{0.3em}\noindent\textit{Note.} A more detailed component role map is in Chapter~3, Table~\ref{tab:component_role_map}.


\subsection{Executive Table: Dissertation Structure}
\label{sec:exec_dissertation_structure}

\noindent The following executive-level table summarizes the role of each chapter in the dissertation.

{\tablestripes\tablefontsize
\needspace{10\baselineskip}
\begin{xltabular}{\textwidth}{@{} >{\raggedright\arraybackslash}p{6cm} L @{}}
\caption[Dissertation Structure (Executive Summary)]{Dissertation Structure (Executive Summary).}\label{tab:exec_dissertation_structure} \\
\toprule
\thead{Chapter} & \thead{Purpose} \\
\midrule
\endfirsthead
\endhead
\bottomrule
\endlastfoot
Chapter 1 & Problem, gap, objectives, research questions, contributions, scope \\
\addlinespace
Chapter 2 & Theoretical foundations, literature synthesis, research gap \\
\addlinespace
Chapter 3 & Research methodology, measurement, validation design \\
\addlinespace
Chapter 4 & Empirical findings, hypothesis verdicts, integrated evidence \\
\addlinespace
Chapter 5 & Discussion, theoretical + practical contributions, limitations + future research \\
\end{xltabular}}


\subsection{Chapter 1 DBA Template Compliance Additions}
\label{sec:ch1_template_compliance}

\noindent The following dedicated subsections close the remaining DBA template-compliance gaps flagged in the Chapter~1 audit: Why DBA, Novelty Statement, Assumptions, and an explicit dissertation Contribution Statement. Each subsection is intentionally short and table-driven for examiner scannability.

\subsubsection{Why DBA --- A Managerial, Not a Technical Question}
\label{sec:ch1_why_dba_dedicated}

\noindent A PhD asks whether AI \emph{can} detect ASD; the literature has already answered yes. A DBA asks under what organisational governance conditions an AI-assisted screening service can be responsibly adopted. This dissertation is therefore framed as a DBA contribution to governance-centered adoption theory and practice, not as a technical contribution to model accuracy.

{\tablestripes\tablefontsize
\needspace{18\baselineskip}
\begin{xltabular}{\textwidth}{@{} >{\raggedright\arraybackslash}p{5cm} L @{}}
\caption[Why this study is a DBA and not a PhD]{Why this study is a DBA and not a PhD.}\label{tab:ch1_why_dba_compliance} \\
\toprule
\thead{Dimension} & \thead{DBA framing applied in this dissertation} \\
\midrule
\endfirsthead
\endhead
\bottomrule
\endlastfoot
Primary question type & Managerial: under what conditions can AI-assisted ASD screening be adopted responsibly? \\
\addlinespace
Unit of analysis & Organisation (governance maturity), not algorithm \\
\addlinespace
Contribution type & Operational framework (RGAIG) + extended adoption theory \\
\addlinespace
Evidence base & Mixed methods: caregiver survey (PLS-SEM) + secondary EEG validation \\
\addlinespace
Audience & Healthcare executives, governance boards, AI/ML leaders, regulators \\
\addlinespace
Out of scope & New algorithms, new neural architectures, real-time clinical diagnosis \\
\end{xltabular}}

\vspace{0.3em}\noindent\textit{Note.} Out-of-scope items belong to the technical research tradition; their resolution by prior work is summarised in Chapter~2.



\subsubsection{Global Statistics --- Why This Research Matters Now}
\label{sec:ch1_global_statistics}

\noindent The table below grounds the dissertation's problem statement in widely cited global figures. The numbers are conservative, drawn from established public-health and industry sources; the dissertation's literature review (Chapter~2) elaborates each statistic with its citation chain.

{\tablestripes\tablefontsize
\needspace{10\baselineskip}
\begin{xltabular}{\textwidth}{@{} >{\raggedright\arraybackslash}p{5cm} L @{}}
\caption[Global statistics establishing study relevance]{Global statistics establishing study relevance.}\label{tab:ch1_global_stats} \\
\toprule
\thead{Area} & \thead{Indicative global statistic (sources in Chapter~2)} \\
\midrule
\endfirsthead
\endhead
\bottomrule
\endlastfoot
ASD prevalence (children) & Approximately 1 in 36 children in widely cited surveillance data; up from 1 in 150 two decades ago. \\
\addlinespace
Diagnostic delay (first concern $\rightarrow$ formal diagnosis) & Typical waiting period of 2--4 years across high- and middle-income systems; longer in under-served regions. \\
\addlinespace
Age at first concern vs.\ age at diagnosis & First concern at approximately 18--24 months; median diagnosis at 4--5 years. \\
\addlinespace
Intervention-window impact & Earlier behavioural intervention is consistently associated with stronger developmental outcomes; the gap window is a measurable loss of opportunity. \\
\addlinespace
Healthcare AI adoption (provider-side) & Roughly one-third of large healthcare providers have active clinical AI initiatives; majority remain in pilot or feasibility stages. \\
\addlinespace
Caregiver trust in healthcare AI & Surveys consistently show majority caregiver concern about transparency, data use, and clinician oversight of AI tools. \\
\addlinespace
Regulatory direction (responsible AI) & The EU AI Act, US AI Bill of Rights blueprint, and national regulators have converged on transparency, accountability, and human oversight as binding principles. \\
\addlinespace
Operational governance maturity gap & No widely adopted operational governance maturity framework specific to clinical AI deployment exists in practice. \\
\end{xltabular}}

\vspace{0.3em}\noindent\textit{Note.} Cited sources are presented in Chapter~2 Section~Literature Synthesis and in Appendix~B (Literature Evidence Repository). The figures here are indicative and conservative; precision is reserved for the chapter that develops each citation chain.

\subsubsection{Novelty Statement}
\label{sec:ch1_novelty_statement}

\noindent The novelty of this dissertation is the integration of four established traditions --- Technology Acceptance Model, Trust Theory, Socio-Technical Systems, and Responsible AI --- into a single empirically validated framework (RGAIG) in which governance maturity functions as the antecedent that operates through explainability and trust to predict adoption readiness. Each component exists in the literature; the integration into a measurable, operationally usable framework is new. The novelty matrix is fully developed in Appendix~V.3.

\subsubsection{Assumptions of the Study}
\label{sec:ch1_assumptions}

\noindent Every empirical study rests on a small number of assumptions that bound its claims. Stating them explicitly is a template-compliance requirement and a defensibility move: if any assumption is challenged, the corresponding finding can be re-scoped without overturning the framework.

{\tablestripes\tablefontsize
\needspace{18\baselineskip}
\begin{xltabular}{\textwidth}{@{} >{\raggedright\arraybackslash}p{5cm} L @{}}
\caption[Assumptions of the study and their justification]{Assumptions of the study and their justification.}\label{tab:ch1_assumptions} \\
\toprule
\thead{Assumption} & \thead{Justification} \\
\midrule
\endfirsthead
\endhead
\bottomrule
\endlastfoot
Caregivers understand AI-assisted screening at a lay level sufficient to report perceptions. & Survey items use plain language; pilot testing validated lay comprehension. \\
\addlinespace
Perceptions of governance, explainability, trust, and adoption are measurable via validated scales. & Constructs are operationalised from prior validated instruments (Chapter~3, Appendix~F). \\
\addlinespace
The B2C caregiver sample is representative of the target stakeholder population. & Sampling frame and recruitment described in Chapter~3; non-response bias tested. \\
\addlinespace
Public EEG benchmark datasets are valid feasibility evidence. & Datasets are peer-reviewed and widely used as benchmarks; framed as feasibility, not primary contribution. \\
\addlinespace
Governance maturity is a stable construct over the survey window. & Cross-sectional design; longitudinal validation is in the future-research agenda (Appendix~V.16). \\
\addlinespace
Adoption readiness is the proximal outcome of interest, not actual usage. & Aligned with TAM tradition; actual usage is addressed in the future-research roadmap. \\
\end{xltabular}}

\vspace{0.3em}\noindent\textit{Note.} Each assumption is restated where it bears on a specific claim (Chapters~3--4).


\subsubsection{Dissertation Contribution Statement}
\label{sec:ch1_contribution_statement}

\noindent The dissertation contributes (i)~an extended Technology Acceptance Model with governance maturity as an antecedent (theory); (ii)~an empirically validated mediation chain governance$\rightarrow$explainability$\rightarrow$trust$\rightarrow$adoption (evidence); (iii)~the RGAIG framework operationalising the chain as measurable constructs and a control catalogue (artefact); (iv)~a stakeholder-specific recommendation matrix and executive decision framework (practitioner contribution); and (v)~a five-phase research program positioning the work within a broader responsible AI agenda (legacy). The full contribution matrix is in Appendix~V.6.


\subsection{Chapter 1 Signature Diagrams}
\label{sec:ch1_signature_diagrams}

\noindent This subsection consolidates the five highest-value signature diagrams for Chapter~1. Each diagram answers a single examiner question: why a DBA programme, what the contribution chain looks like, what the canonical dissertation argument is in one image, how the five chapters connect, and where the boundary of the research lies.


\noindent The dissertation answers a managerial question, not a purely technical one. The sequence below clarifies why a Doctor of Business Administration programme is the appropriate vehicle for governance-centered adoption research. Figure~\ref{fig:ch1_why_dba} captures this flow.

\begin{figure}[!htbp]
\centering
\begin{tikzpicture}[
  node distance=0.55cm,
  fbox/.style={rectangle, draw=ggunavy, fill=ggunavy!8, rounded corners=4pt, very thick,
    minimum width=7.5cm, minimum height=0.85cm, align=center, font=\fignodefont},
  fboxgold/.style={rectangle, draw=ggunavy, fill=ggugold!25, rounded corners=4pt, very thick,
    minimum width=7.5cm, minimum height=0.95cm, align=center, font=\fignodefont\bfseries},
  farr/.style={-{Stealth[length=3mm, width=2mm]}, very thick, ggunavy}
]
\node[font=\figtitlefont, ggunavy] (title) at (0,0.7) {Why a DBA programme rather than a PhD};
\node[fbox, below=0.85cm of title] (n0) {Technical Feasibility (AI accuracy is already established)};
\node[fbox, below=of n0] (n1) {Governance Challenge (responsibility, accountability, oversight)};
\node[fbox, below=of n1] (n2) {Organizational Trust (clinicians, caregivers, regulators)};
\node[fbox, below=of n2] (n3) {Adoption Readiness (organizational maturity to deploy safely)};
\node[fboxgold, below=of n3] (n4) {Business Transformation (responsible AI-assisted care pathways)};
\draw[farr] (n0) -- (n1);
\draw[farr] (n1) -- (n2);
\draw[farr] (n2) -- (n3);
\draw[farr] (n3) -- (n4);
\end{tikzpicture}
\caption[Why a DBA programme rather than a PhD]{Why a DBA programme rather than a PhD.}
\label{fig:ch1_why_dba}
\end{figure}

\vspace{0.3em}\noindent\textit{Note.} A PhD asks whether AI can do this; a DBA asks under what governance conditions an organization can responsibly deploy it.



\noindent The contribution of this dissertation is not a new model but a governance-centered explanation of how adoption emerges. The chain below shows how each construct enables the next. Figure~\ref{fig:ch1_research_contribution_flow} captures this flow.

\begin{figure}[!htbp]
\centering
\begin{tikzpicture}[
  node distance=0.55cm,
  fbox/.style={rectangle, draw=ggunavy, fill=ggunavy!8, rounded corners=4pt, very thick,
    minimum width=7.5cm, minimum height=0.85cm, align=center, font=\fignodefont},
  fboxgold/.style={rectangle, draw=ggunavy, fill=ggugold!25, rounded corners=4pt, very thick,
    minimum width=7.5cm, minimum height=0.95cm, align=center, font=\fignodefont\bfseries},
  farr/.style={-{Stealth[length=3mm, width=2mm]}, very thick, ggunavy}
]
\node[font=\figtitlefont, ggunavy] (title) at (0,0.7) {Research contribution chain};
\node[fbox, below=0.85cm of title] (n0) {Literature Gap (no integrated governance model for AI adoption)};
\node[fbox, below=of n0] (n1) {Governance Construct (operationalised as maturity)};
\node[fbox, below=of n1] (n2) {Explainability (mediates governance into perception)};
\node[fbox, below=of n2] (n3) {Trust (mediates explainability into intent)};
\node[fbox, below=of n3] (n4) {Adoption (proximal outcome)};
\node[fboxgold, below=of n4] (n5) {RGAIG Framework (the integrated contribution)};
\draw[farr] (n0) -- (n1);
\draw[farr] (n1) -- (n2);
\draw[farr] (n2) -- (n3);
\draw[farr] (n3) -- (n4);
\draw[farr] (n4) -- (n5);
\end{tikzpicture}
\caption[Research contribution chain]{Research contribution chain.}
\label{fig:ch1_research_contribution_flow}
\end{figure}

\vspace{0.3em}\noindent\textit{Note.} Each arrow corresponds to a hypothesis tested in Chapter~4.



\noindent A single image carries the dissertation argument. This signature diagram recurs in every chapter as the anchor visual; if an examiner remembers one figure, this should be it. Figure~\ref{fig:ch1_dissertation_signature} captures this flow.

\begin{figure}[!htbp]
\centering
\begin{tikzpicture}[
  node distance=0.55cm,
  fbox/.style={rectangle, draw=ggunavy, fill=ggunavy!8, rounded corners=4pt, very thick,
    minimum width=7.5cm, minimum height=0.85cm, align=center, font=\fignodefont},
  fboxgold/.style={rectangle, draw=ggunavy, fill=ggugold!25, rounded corners=4pt, very thick,
    minimum width=7.5cm, minimum height=0.95cm, align=center, font=\fignodefont\bfseries},
  farr/.style={-{Stealth[length=3mm, width=2mm]}, very thick, ggunavy}
]
\node[font=\figtitlefont, ggunavy] (title) at (0,0.7) {Dissertation signature diagram};
\node[fbox, below=0.85cm of title] (n0) {Healthcare AI (B2C-primary, clinician-supported screening)};
\node[fbox, below=of n0] (n1) {Governance Maturity (the antecedent)};
\node[fbox, below=of n1] (n2) {Explainability (the mediator)};
\node[fbox, below=of n2] (n3) {Trust (the mediator)};
\node[fbox, below=of n3] (n4) {Adoption Readiness (the proximal outcome)};
\node[fboxgold, below=of n4] (n5) {Clinical Value (the ultimate outcome)};
\draw[farr] (n0) -- (n1);
\draw[farr] (n1) -- (n2);
\draw[farr] (n2) -- (n3);
\draw[farr] (n3) -- (n4);
\draw[farr] (n4) -- (n5);
\end{tikzpicture}
\caption[Dissertation signature diagram]{Dissertation signature diagram.}
\label{fig:ch1_dissertation_signature}
\end{figure}

\vspace{0.3em}\noindent\textit{Note.} The same figure recurs as Figures~\ref{fig:ch1_dissertation_signature}, equivalent versions appear in Chapters~2--5.



\noindent Each chapter advances the argument. The roadmap below makes the chapter sequence explicit so that the examiner can locate any claim against its supporting chapter. Figure~\ref{fig:ch1_dissertation_roadmap} captures this flow.

\begin{figure}[!htbp]
\centering
\begin{tikzpicture}[
  node distance=0.55cm,
  fbox/.style={rectangle, draw=ggunavy, fill=ggunavy!8, rounded corners=4pt, very thick,
    minimum width=7.5cm, minimum height=0.85cm, align=center, font=\fignodefont},
  fboxgold/.style={rectangle, draw=ggunavy, fill=ggugold!25, rounded corners=4pt, very thick,
    minimum width=7.5cm, minimum height=0.95cm, align=center, font=\fignodefont\bfseries},
  farr/.style={-{Stealth[length=3mm, width=2mm]}, very thick, ggunavy}
]
\node[font=\figtitlefont, ggunavy] (title) at (0,0.7) {Dissertation roadmap across five chapters};
\node[fbox, below=0.85cm of title] (n0) {Chapter 1 --- Problem, gap, scope, contribution preview};
\node[fbox, below=of n0] (n1) {Chapter 2 --- Theoretical foundation and research gap};
\node[fbox, below=of n1] (n2) {Chapter 3 --- Methodology and measurement design};
\node[fbox, below=of n2] (n3) {Chapter 4 --- Empirical evidence and hypothesis verdicts};
\node[fbox, below=of n3] (n4) {Chapter 5 --- Discussion, contribution, future research};
\draw[farr] (n0) -- (n1);
\draw[farr] (n1) -- (n2);
\draw[farr] (n2) -- (n3);
\draw[farr] (n3) -- (n4);
\end{tikzpicture}
\caption[Dissertation roadmap across five chapters]{Dissertation roadmap across five chapters.}
\label{fig:ch1_dissertation_roadmap}
\end{figure}

\vspace{0.3em}\noindent\textit{Note.} Chapters 1--2 establish the question; Chapter~3 designs the answer; Chapters~4--5 produce and position it.



\noindent Boundary clarity is the single most common request from examiners. The figure separates governance-centered contributions (in scope) from algorithmic novelty (out of scope, addressed by prior work). Figure~\ref{fig:ch1_scope_inout} delineates the boundary.


\begin{figure}[!htbp]
\centering
\begin{tikzpicture}[
  node distance=0.5cm,
  fbox/.style={rectangle, draw=ggunavy, fill=ggunavy!10, rounded corners=4pt, very thick,
    minimum width=5cm, minimum height=0.75cm, align=center, font=\fignodefont},
  fboxout/.style={rectangle, draw=ggunavy!60, fill=gray!8, rounded corners=4pt, thick, dashed,
    minimum width=5cm, minimum height=0.75cm, align=center, font=\fignodefont\itshape, text=gray!50!black},
  ftitle/.style={font=\figtitlefont\bfseries, ggunavy}
]
\node[ftitle] (intitle) at (-3.5,0) {IN SCOPE};
\node[ftitle] (outtitle) at (3.5,0) {OUT OF SCOPE};
\node[fbox, below=0.5cm of intitle] (in0) {Governance Maturity};
\node[fbox, below=0.5cm of in0] (in1) {Explainability (perception)};
\node[fbox, below=0.5cm of in1] (in2) {Trust (perception)};
\node[fbox, below=0.5cm of in2] (in3) {Adoption Readiness};
\node[fbox, below=0.5cm of in3] (in4) {RGAIG Framework Validation};
\node[fboxout, below=0.5cm of outtitle] (out0) {New EEG algorithms};
\node[fboxout, below=0.5cm of out0] (out1) {New deep-learning models};
\node[fboxout, below=0.5cm of out1] (out2) {New neural architectures};
\node[fboxout, below=0.5cm of out2] (out3) {New signal-processing methods};
\node[fboxout, below=0.5cm of out3] (out4) {Real-time clinical diagnosis};
\end{tikzpicture}
\caption[Research scope --- in-scope versus out-of-scope]{Research scope --- in-scope versus out-of-scope.}
\label{fig:ch1_scope_inout}
\end{figure}


\vspace{0.3em}\noindent\textit{Note.} The dissertation evaluates organizational conditions for adoption, not algorithmic accuracy gains; the latter is treated as established by Chapter~2's literature.



\subsection{Chapter 1 Examiner Summary}
\label{sec:ch1_examiner_summary}

\noindent The dissertation is not intended to demonstrate that AI can detect ASD. Prior studies have already established the feasibility of AI-assisted ASD classification. Instead, this research investigates the organizational conditions required for responsible adoption. Specifically, the study evaluates whether governance maturity improves explainability, whether explainability strengthens trust, and whether trust ultimately influences adoption readiness. These relationships constitute the primary theoretical and managerial contribution examined throughout the remaining chapters.


